% 本文件由 Word 文档转换生成，可直接由 main.tex 使用
% % 本文件由 Word 文档转换生成，可直接由 main.tex 使用
% % 本文件由 Word 文档转换生成，可直接由 main.tex 使用
% % 本文件由 Word 文档转换生成，可直接由 main.tex 使用
% \include{chapter_materials}
% 文件中不加载任何宏包，所有样式均调用项目现有的 package.tex。
\part*{微模块5\quad 人工智能在材料科学中的应用}
\addcontentsline{toc}{part}{微模块5\texorpdfstring{\quad}{ }人工智能在材料科学中的应用}

\chapter{材料认知与性能预测}
\label{chap:materials-cognition-prediction}


理解材料性能的来源，是材料科学研究的核心任务。材料的力学、热学、电学、磁学和光学性能，往往不是由单一因素决定的，而是受到成分组成、微观结构和制备工艺的共同影响。随着研究对象从简单合金、陶瓷和半导体扩展到高熵合金、复合材料、低维材料和复杂功能材料，材料候选空间迅速增大，单纯依靠实验试错已经难以系统揭示其中的规律。在这一背景下，人工智能方法的价值，主要体现在对分散实验数据、计算数据和文献知识的组织与学习上，使模型能够从中提取“成分--结构--工艺--性能”之间的关联，从而辅助材料性能预测和后续实验设计。

本章围绕“材料认知与性能预测”展开，重点解决四个问题。第一节说明材料人工智能为什么需要可靠数据基础，以及实验、计算、文献和数据库数据应如何整理。第二节讨论材料成分和晶体结构如何转化为描述符或图表示，使模型能够读取材料信息。第三节介绍描述符模型、图神经网络和机器学习原子间势（下文简称“机器学习势”）如何服务硬度预测、性质预测和熔点模拟。第四节进一步说明模型解释与物理知识嵌入的作用，帮助研究者判断模型结果是否符合材料机理和实验事实。通过本章学习，研究者需要建立一个基本认识：人工智能并不是替代材料理论和实验验证，而是与数据、计算和材料机理共同构成辅助材料认知的新工具。




\clearpage
\section{范式变革与数据基础}


\begin{sectionkeypoints}
\item 材料人工智能中的数据并不只是单个性能数值，而是包含材料组成、结构信息、制备工艺、测试条件、数据来源和目标性能的完整记录。
\item 材料数据主要来自实验、计算、文献和数据库四类渠道。不同来源的数据各有优势，也各有局限，使用时需要结合具体研究任务进行判断。
\item 材料数据库为模型训练、性能预测和候选筛选提供了重要的数据基础，但数据库中的数据不宜直接等同于最终模型输入，还需要经过字段筛选、来源核验和格式整理。
\item 在正式建模之前，应先整理原始数据清单，明确记录材料是什么、结构是什么、性能数值是多少、数据来自哪里，以及采用了什么计算方法或实验条件。
\item AI 可辅助解释数据库字段、设计表格框架和检查数据规范性，但不宜替代数据库原始记录、实验记录和论文原文核验。材料数据整理的核心要求是清楚、可追溯、可检查。
\end{sectionkeypoints}

\subsection{实验试错与数据驱动}

材料科学的核心任务之一，是理解和建立材料成分、结构、工艺与性能之间的关联。研究者希望理解：为什么某种成分的合金具有更高强度？为什么某种陶瓷能够承受更高温度？为什么某种半导体材料具有合适的带隙并能用于发光器件？这些问题表面上属于不同材料体系，但本质上都指向同一个核心：材料的内部组成与结构如何决定其外部性能。


在传统材料研发中，研究者通常依靠已有理论、文献经验和实验试错来提出材料方案，再通过制备、表征和性能测试不断调整成分与工艺。例如，为了提高合金强度，可尝试加入不同合金元素；为了改善陶瓷耐高温性能，可调节烧结温度、晶粒尺寸或相组成；为了优化光电材料性能，可改变材料的元素组成和晶体结构。这种方式直接可靠，也符合材料实验研究的基本规律。它的优势在于结论来自真实样品和真实测试条件，缺点是往往需要大量重复实验，研发周期较长，成本也较高。


当材料体系从简单单组元或二元材料扩展到多组元合金、复合材料、异质结构、低维材料和复杂光电材料时，实验试错面对的组合空间会迅速扩大。以高熵合金为例，四种、五种甚至更多主要元素的组合、比例和热处理工艺都可能改变组织结构和性能；同时，强度、硬度、韧性、耐腐蚀性、导电性或发光性能又受到晶体结构、缺陷、晶粒尺寸、相组成、界面状态和服役环境共同影响。传统经验仍然重要，但单纯依靠人工逐一尝试，已经难以在有限时间内完成系统筛选。


在这一背景下，材料研发逐渐从单纯依赖实验试错，转向实验、计算、数据和人工智能相结合的新方式。实验能够提供真实可靠的性能结果，计算能够在实验之前预测部分结构和能量信息，数据库能够积累已有材料的成分、结构和性能数据，而人工智能则能够从这些数据中学习复杂的材料规律。换言之，人工智能并不是脱离实验和理论单独发挥作用，而是建立在已有材料知识、实验数据和计算数据基础上的辅助工具。


数据驱动并不是否定实验和理论，而是把已有实验数据、计算数据、文献数据和数据库资源转化为可分析、可学习和可预测的信息。人工智能模型可从大量材料数据中学习成分、结构与性能之间的关系，并用于预测未知材料的性质。这样，研究者不再需要在庞大的候选空间中完全盲目试错，而可借助数据和模型缩小研究范围，把实验资源集中到更有希望的材料体系上。


\subsubsection*{\textbf{1. 人工智能用于材料科学的基本流程}}

在人工智能用于材料科学的过程中，研究者通常需要依次解决几个关键环节。第一个环节是材料数据整理。需要收集和规范实验数据、计算数据、文献数据和数据库资源，并记录材料的成分、结构、工艺、性能以及数据来源等信息。只有数据来源清楚、含义明确，后续模型学习才具有可靠基础。

第二个环节是材料表示。机器学习模型不宜直接理解“合金”“晶体结构”或“热处理工艺”等材料概念，因此需要把材料转化为机器可处理的形式。例如，可用描述符表示化学成分和结构信息，也可用图结构表示原子之间的连接关系。

第三个环节是模型预测。在材料被合理表示之后，可建立机器学习模型，预测材料的形成能、带隙、弹性模量、熔点、硬度等性能。对于涉及原子运动和微观演化的问题，还可进一步构建机器学习势（利用机器学习模型学习原子结构与能量、力之间关系的势函数），用于模拟扩散、相变和熔化等原子尺度行为。

最后，模型预测结果还需要回到材料科学问题本身。通过可解释性方法，可分析模型为什么做出某种预测，识别哪些成分、结构或工艺因素对性能影响更大，从而把模型结果进一步转化为可理解的材料规律。


在材料研究中还必须明确，人工智能不宜替代材料专业知识、实验验证和物理判断。人工智能的作用更像是一种增强手段：它可帮助研究者整理数据、发现规律、提出候选、预测性能和辅助决策，但最终结论仍然需要通过材料理论分析和实验结果进行验证。尤其是在材料科学中，模型给出的预测结果并不等同于真实材料性能。材料是否能够被实际制备出来，结构是否稳定，以及能否在实际服役环境中长期可靠工作，仍然需要进一步验证。


围绕材料性能预测和材料设计两类问题，人工智能用于材料科学可概括为两条基本主线。第一条是正向预测，即在材料成分、结构或实验条件已经给定的情况下，预测材料可能具有的性能和行为。例如，已知一种晶体材料的结构，可预测它的形成能、带隙、弹性模量或熔点。第二条是逆向设计，即先给出目标性能，再反过来寻找可能满足要求的材料。从概念上看，正向预测关注已知材料可能表现出什么性能，逆向设计关注为了获得目标性能应当寻找什么材料。


本章主要讨论第一条主线，即材料科学中的正向预测。本章将从材料数据出发，进一步讨论材料如何被表示成机器可理解的形式，机器学习如何预测材料性能，机器学习势如何模拟原子尺度行为，以及可解释人工智能如何帮助研究者从模型结果中提炼材料规律。


\subsection{数据来源与材料数据库}

人工智能要预测材料性能，首先需要数据。对于材料科学而言，数据并不只是一个性能数值，而是对材料组成、制备过程、测试条件和性能结果的完整记录。一个可靠的材料数据条目，至少应该说明材料的化学组成、结构信息、制备工艺、测试条件、数据来源和目标性能。缺少这些基本信息，模型就难以建立清晰的输入与输出对应关系，所得预测结果也难以作为可靠依据。


材料数据主要来自实验、计算、文献和数据库四类渠道。


第一类是实验数据。实验数据来自材料制备、显微表征和性能测试过程，例如硬度、拉伸强度、腐蚀电流密度、熔点、热导率、带隙、发光波长等。实验数据最接近真实材料和真实服役环境，但获取成本较高，也容易受到制备工艺、测试条件和仪器参数的影响。例如，同一种合金在不同热处理温度下可能得到不同硬度，同一种涂层在不同腐蚀介质中可能表现出不同耐蚀性能。因此，实验数据不应只记录结果，还要尽量记录测试条件。


第二类是计算数据。计算数据主要来自第一性原理计算、分子动力学模拟、相场模拟等方法。例如，第一性原理计算可得到形成能、能带结构、弹性常数、态密度、晶体结构稳定性等信息；分子动力学可模拟扩散、相变和高温行为。计算数据的优点是容易标准化，也适合大规模生成；但计算数据依赖具体计算方法和参数，不宜简单等同于实验结果。


第三类是文献数据。大量材料知识存在于论文、专利、实验报告和技术资料中。例如，一篇论文可能同时记录材料成分、制备工艺、热处理条件、显微组织和性能结果。文献数据的优点是信息丰富，缺点是格式不统一，很多信息隐藏在正文、表格、图片或补充材料中，需要人工整理或借助文本处理工具提取。


第四类是材料数据库。材料数据库可理解为经过整理的材料信息资源库，它把大量材料的化学组成、晶体结构、计算性质、实验性质和文献信息集中到一个可检索的平台中。对于材料人工智能而言，数据库是重要数据来源，因为它可为模型训练、性能预测和材料筛选提供相对规范的数据基础。


常见材料数据库包括 Materials Project、OQMD\footnote{Open Quantum Materials Database，开放量子材料数据库。}、AFLOW\footnote{Automatic Flow Framework for Materials Discovery，自动化材料发现计算框架。}、NOMAD\footnote{Novel Materials Discovery，材料数据管理与共享平台。}、ICSD\footnote{Inorganic Crystal Structure Database，无机晶体结构数据库。}、JARVIS\footnote{Joint Automated Repository for Various Integrated Simulations，材料计算数据库与工具平台。} 等。不同数据库的侧重点并不完全相同：有的偏向无机晶体结构和稳定性数据，有的偏向大规模热力学计算结果，有的强调计算数据管理和共享，还有的主要提供实验晶体结构信息。教材中列举数据库名称，目的不是要求逐一掌握平台功能，而是说明不同数据来源对应不同研究任务。

对于入门阶段的研究者而言，不必一开始深入掌握每个数据库的所有功能，但需要知道不同数据库主要提供哪类材料数据，以及这些数据适合用于哪些研究任务。只有明确数据来源和适用范围，才能更合理地将数据库用于材料性能预测、候选筛选和模型训练。


使用材料数据库时，首先要明确自己的研究任务。不同任务需要的数据字段不同。例如，如果任务是预测无机光电材料的带隙，就需要关注化学式、晶体结构、空间群、带隙、形成能、数据来源和计算方法等字段；如果任务是预测合金硬度，就需要关注成分、相结构、制备工艺、热处理条件、硬度值和测试方法等字段；如果任务是筛选超高温陶瓷，就需要关注化学组成、晶体结构、熔点、密度、热稳定性和抗氧化相关数据。


因此，数据库只是提供了材料信息的来源，真正进入模型训练之前，还需要把这些信息整理成结构清楚、字段明确、来源可追溯的原始数据清单。下面以无机光电材料带隙数据整理为例，说明从数据库原始记录到可建模数据清单的具体做法。

\subsubsection*{\textbf{1. 案例示范：无机光电材料带隙数据整理}}

以“无机光电材料带隙预测”为例，数据准备的第一步不是训练模型，而是从数据库或文献中整理出一张原始数据清单。这个清单应该回答几个基本问题：材料是什么？结构是什么？带隙是多少？这个带隙来自哪里？是实验值还是计算值？如果是计算值，采用了什么计算方法？如果是实验值，是否记录了测试条件？


一张初步整理后的原始数据清单可按表 \ref{tab:raw-bandgap-data} 设计。表中的 eV\footnote{Electron volt，电子伏特，常用于表示电子能量和带隙大小。} 是带隙常用单位。


\begin{formalTable}[htbp]
\caption{原始带隙数据清单示例}
\label{tab:raw-bandgap-data}
\scriptsize
\setlength{\tabcolsep}{2.8pt}
\begin{tabularx}{\textwidth}{@{}YYYYYYYYY@{}}
\toprule
\textbf{材料编号} & \textbf{化学式} & \textbf{晶体结构} & \textbf{空间群} & \textbf{带隙/eV} & \textbf{数据类型} & \textbf{数据来源} & \textbf{方法或条件} & \textbf{备注} \\
\midrule
001 & GaN & 纤锌矿 & P6\textsubscript{3}mc & 3.40 & 实验 & 文献 & 室温测试 & 直接带隙 \\
002 & ZnO & 纤锌矿 & P6\textsubscript{3}mc & 3.37 & 实验 & 文献 & 室温测试 & 直接带隙 \\
003 & SrTiO\textsubscript{3} & 钙钛矿 & Pm-3m & 3.20 & 数据库 & 数据库记录 & 需注明方法 & 需核对来源 \\
\bottomrule
\end{tabularx}
\end{formalTable}

表 \ref{tab:raw-bandgap-data} 的作用是保存原始材料信息，而不是直接代表最终模型输入。它的重点在于把分散的信息整理成清楚、可追溯、可检查的记录。后续是否把化学式转化为成分特征，是否把晶体结构转化为结构描述符，属于下一节要讨论的材料表示问题。


在整理这类原始数据清单时，可借助 AI 辅助完成部分工作。此处的重点不是让 AI 代替研究者判断数据是否可靠，而是让 AI 帮助理解字段含义、设计表格框架和检查数据规范性。以下以无机光电材料带隙数据整理为例给出提示词示范。

\begin{samepage}
\begin{userbox}
我正在整理无机光电材料的带隙数据，准备用于后续机器学习建模。数据可能来自 Materials Project、JARVIS、文献或实验记录。请帮我设计一张原始数据清单的表头，并说明每一列的作用。

要求如下：
\begin{enumerate}
\item 表格要保留材料编号、化学式、晶体结构、空间群、带隙数值、带隙单位、数据类型、数据来源、计算方法或实验条件等信息；
\item 区分实验值和计算值；
\item 提醒哪些字段是后续建模前必须核验的；
\item 不要虚构具体材料数据，只需要给出表头设计和填写说明。
\end{enumerate}
\end{userbox}

\begin{aibox}
可将原始数据清单设计为以下字段：

\begin{itemize}
\item 材料编号：用于唯一标识每一条材料记录，便于后续追踪和检查。
\item 化学式：记录材料的元素组成，例如 GaN、ZnO、SrTiO$_3$ 等。
\item \textbf{晶体结构：}记录材料所属结构类型，例如纤锌矿、钙钛矿等。
\item \textbf{空间群：}记录晶体结构的对称性信息。
\item 带隙数值：记录材料的带隙大小。
\item 带隙单位：通常为 eV，避免不同来源数据混用时产生单位错误。
\item \textbf{数据类型：}说明该带隙是实验值还是计算值。
\item \textbf{数据来源：}记录数据来自数据库、论文、实验报告还是其他资料。
\item \textbf{方法或条件：}如果是计算值，应注明计算方法；如果是实验值，应注明测试温度、测试方法等条件。
\item \textbf{备注：}记录是否为直接带隙、是否需要进一步核验等补充信息。
\end{itemize}

在正式建模前，尤其需要核验带隙数值、数据类型、数据来源和方法或条件等字段。同一种材料在不同晶体结构、不同计算方法或不同测试条件下可能具有不同带隙，因此这些信息不应省略。
\end{aibox}
\end{samepage}

上面的提示词示范说明，AI 可帮助研究者快速形成数据表框架，但 AI 输出并不宜直接作为最终数据结果。研究者还需要回到数据库原始页面或论文原文，核对具体数值、计算方法、实验条件和材料结构是否准确。

因此，在材料数据整理阶段，AI 更适合充当数据整理助手，而不是数据来源本身。它可帮助设计表格、解释字段和生成检查规则，但具体数据是否准确，仍然需要根据数据库原始记录、实验记录或论文原文进行确认。尤其在材料科学中，同一种化学式可能对应不同晶体结构，不同计算方法也可能得到不同带隙值，因此数据来源、计算方法和实验条件必须保留。


为了提高数据质量，原始数据整理时应注意几个问题。第一，要统一单位。例如带隙通常用 eV，硬度可能用 HV\footnote{Vickers Hardness，维氏硬度单位。} 或 GPa\footnote{Gigapascal，吉帕，常用于表示压力、应力或弹性模量。}，能量可能用 eV/atom 或 kJ/mol。第二，要区分实验数据和计算数据。二者可同时记录，但不能在没有说明的情况下混用。第三，要保留数据来源。每一条数据宜能追溯到数据库编号、论文或实验记录。第四，要检查缺失值和重复值。同一个材料如果重复出现，需要判断它们是否对应不同结构、不同测试条件或不同计算方法。第五，要记录必要条件。对于实验数据，应尽量记录温度、压力、测试环境和样品状态；对于计算数据，应记录计算方法和参数来源。

\subsubsection*{\textbf{小结}}

本节说明材料研发从实验试错走向数据、计算和实验协同的基本背景。
材料数据应包含成分、结构、工艺、性能、来源和测试条件，而不只是单一性能数值。
数据库和 AI 可提高整理效率，但数据核验、条件记录和材料判断仍需由研究者完成。

\subsection*{习题}

\begin{enumerate}[leftmargin=*]
\item 材料数据为什么不能只记录性能值？还应补充哪些字段？
\item 比较四类材料数据来源，并说明选择依据。
\end{enumerate}

\clearpage
\section{特征构建与结构表示}


\begin{sectionkeypoints}
\item 材料特征构建的目的，是把化学式、晶体结构、空间群和性能数值等原始记录转化为机器学习模型能够读取的数值特征。
\item 成分描述符可由元素组成和元素属性计算得到，例如平均原子序数、平均电负性、电负性差和平均原子半径等。
\item 结构描述符用于补充原子排列信息，常见字段包括空间群编号、晶胞体积、密度和晶格常数等。
\item 图表示把晶体看成原子节点和原子间边组成的结构，可更直接地保留局域配位、原子间距离和周期性信息。
\item 描述符和图表示都需要材料判断。特征不是越多越好，应避免目标泄漏，并检查所选特征是否具有明确的材料意义。
\end{sectionkeypoints}

\subsection{成分与结构描述符}

上一节介绍了材料数据的来源和材料数据库的作用。通过实验、计算、文献或数据库，研究者可获得材料的化学式、晶体结构、空间群、性能数值和数据来源等信息。但是，这些信息只是原始材料记录，还不宜直接被大多数机器学习模型使用。机器学习模型真正能够处理的，通常是一组数值特征。因此，在进行性能预测之前，需要把材料信息进一步转化为模型可读取的数值表达。


这一过程称为材料的特征化，也常称为描述符构建。描述符指用数字表达材料中可能影响性能的关键信息。描述符不是随意选择的一组数字，而应该尽量反映材料的成分、结构和物理化学特征。例如，对于光电材料带隙预测，元素电负性、原子半径、晶体结构、空间群和晶胞参数等信息都可能与电子结构有关，因此可作为候选描述符。


\subsubsection*{\textbf{1. 案例示范：无机光电材料特征构建}}

本节围绕一个任务展开：无机光电材料带隙预测。假设研究者已经从材料数据库或文献中整理出表 \ref{tab:bandgap-raw-records} 所示的原始数据清单。


\begin{formalTable}[htbp]
\caption{无机光电材料原始记录示例}
\label{tab:bandgap-raw-records}
\footnotesize
\begin{tabularx}{\textwidth}{@{}YYYYY@{}}
\toprule
\textbf{材料} & \textbf{化学式} & \textbf{晶体结构} & \textbf{空间群} & \textbf{带隙/eV} \\
\midrule
氮化镓 & GaN & 纤锌矿 & P6\textsubscript{3}mc & 3.40 \\
氧化锌 & ZnO & 纤锌矿 & P6\textsubscript{3}mc & 3.37 \\
钛酸锶 & SrTiO\textsubscript{3} & 钙钛矿 & Pm-3m & 3.20 \\
\bottomrule
\end{tabularx}
\end{formalTable}

这个表格适合人阅读，但模型不宜直接理解“GaN”“纤锌矿”或“P6\textsubscript{3}mc”这些符号和文字。因此，第一步是把化学式转化为元素组成。


以 GaN 为例，它由 Ga 和 N 两种元素组成，原子比例为 1:1。为了便于模型处理，可把这个比例归一化为原子分数，即 Ga 为 0.50，N 为 0.50。ZnO 同理，Zn 为 0.50，O 为 0.50。SrTiO\textsubscript{3} 中 Sr、Ti 和 O 的比例为 1:1:3，归一化后分别为 0.20、0.20 和 0.60。


经过这一步，原始化学式可转化为表 \ref{tab:formula-composition-example} 所示形式。


\begin{formalTable}[htbp]
\caption{化学式到元素组成的转换示例}
\label{tab:formula-composition-example}
\begin{tabularx}{\textwidth}{@{}YYY@{}}
\toprule
\textbf{化学式} & \textbf{元素组成} & \textbf{原子分数} \\
\midrule
GaN & Ga, N & Ga: 0.50；N: 0.50 \\
ZnO & Zn, O & Zn: 0.50；O: 0.50 \\
SrTiO\textsubscript{3} & Sr, Ti, O & Sr: 0.20；Ti: 0.20；O: 0.60 \\
\bottomrule
\end{tabularx}
\end{formalTable}

这一步表面上简单，但当数据表中有几百或几千个化学式时，人工逐个拆解就很低效。实际建模时，通常需要用程序批量解析化学式，并检查程序是否能够正确处理下标、括号、掺杂和复杂化学式。只有化学式解析正确，后续成分描述符才有意义。


得到元素组成后，第二步是引入元素属性。机器学习模型并不知道 Ga、N、Zn、O、Sr、Ti 这些元素各自有什么物理化学特征，因此需要查找或建立元素属性表。这个属性表可包括原子序数、原子半径、电负性、价电子数、元素熔点等信息。


例如，GaN 的两个组成元素是 Ga 和 N。如果需要计算平均电负性，就需要知道 Ga 和 N 各自的电负性，再按照原子分数加权平均。一般来说，若第 i 种元素的原子分数为 xi，对应元素属性为 Pi，则该材料的平均属性可表示为：


\begin{equation*}\text{平均属性}=\sum_i x_i P_i\end{equation*}

这个公式可用于计算平均原子序数、平均电负性、平均原子半径、平均价电子数等特征。除了平均值，还可计算最大值、最小值、极差和方差。例如，电负性差可用组成元素中最大电负性与最小电负性的差值表示，用来描述元素之间化学性质差异。


对于带隙预测来说，电负性和电负性差具有直观意义。带隙与材料的电子结构和成键特征有关，而元素之间电负性差异会影响电子分布和键合方式。因此，平均电负性、电负性差、平均原子半径等都可作为初步描述符。当然，这些描述符并不能完全决定带隙，但它们为模型提供了可学习的材料信息。


经过这一步，数据表可进一步转化为表 \ref{tab:composition-descriptor-example} 所示的成分描述符表。


\begin{formalTable}[htbp]
\caption{成分描述符表构建示例}
\label{tab:composition-descriptor-example}
\footnotesize
\setlength{\tabcolsep}{2.8pt}
\begin{tabularx}{\textwidth}{@{}YYYYYY@{}}
\toprule
\textbf{化学式} & \textbf{平均原子序数} & \textbf{平均电负性} & \textbf{电负性差} & \textbf{平均原子半径} & \textbf{带隙/eV} \\
\midrule
GaN & 数值 & 数值 & 数值 & 数值 & 3.40 \\
ZnO & 数值 & 数值 & 数值 & 数值 & 3.37 \\
SrTiO\textsubscript{3} & 数值 & 数值 & 数值 & 数值 & 3.20 \\
\bottomrule
\end{tabularx}
\end{formalTable}

此处的“数值”需要根据元素属性表计算得到。实际操作时，可先准备一个元素属性 CSV\footnote{Comma-Separated Values，逗号分隔值文件。} 文件，其中包含元素符号、原子序数、电负性、原子半径等字段。然后将化学式解析结果与元素属性表关联起来，自动计算每个材料的平均值和差值。这样可减少重复计算，但计算结果仍然需要通过抽样检查来确认。


第三步，是加入结构描述符。仅有成分描述符还不够，因为材料性能不仅取决于由哪些元素组成，还取决于这些原子如何排列。GaN 和 ZnO 都是纤锌矿结构，而 SrTiO\textsubscript{3} 是钙钛矿结构。不同晶体结构对应不同的原子排列、局域环境和电子能带特征，因此结构信息对于带隙预测也很重要。


对于入门阶段的研究者，可先使用比较容易获得的结构描述符，例如空间群编号、晶胞体积、密度、晶格常数等。以空间群为例，P6\textsubscript{3}mc 可对应空间群编号 186，Pm-3m 可对应空间群编号 221。这样，文字形式的空间群就可转化为数值字段。晶胞体积和密度则可从晶体结构文件或数据库记录中获得。


加入结构描述符后，表格可进一步扩展为表 \ref{tab:structure-descriptor-example}。


\begin{formalTable}[htbp]
\caption{成分与结构描述符表构建示例}
\label{tab:structure-descriptor-example}
\scriptsize
\setlength{\tabcolsep}{2.8pt}
\begin{tabularx}{\textwidth}{@{}YYYYYYYY@{}}
\toprule
\textbf{化学式} & \textbf{平均原子序数} & \textbf{平均电负性} & \textbf{电负性差} & \textbf{平均原子半径} & \textbf{空间群编号} & \textbf{晶胞体积} & \textbf{带隙/eV} \\
\midrule
GaN & 数值 & 数值 & 数值 & 数值 & 186 & 数值 & 3.40 \\
ZnO & 数值 & 数值 & 数值 & 数值 & 186 & 数值 & 3.37 \\
SrTiO\textsubscript{3} & 数值 & 数值 & 数值 & 数值 & 221 & 数值 & 3.20 \\
\bottomrule
\end{tabularx}
\end{formalTable}

在这个表中，平均原子序数、平均电负性、电负性差、平均原子半径、空间群编号和晶胞体积等列就是模型输入特征，带隙/eV 是模型要预测的目标值。这样，原始的材料记录就被转化为了机器学习可处理的特征表。


该流程可概括为下列步骤，并可结合图 \ref{fig:materials-data-feature} 理解：


\begin{lightquote}
\centering 原始材料记录 → 化学式解析 → 元素属性计算 → 结构字段数值化 → 特征表
\end{lightquote}

\begin{figure}[H]
\centering
\BookFigureImage{images_11/1.png}
\caption{材料数据整理与特征构建示意图}
\label{fig:materials-data-feature}
\end{figure}
\FloatBarrier

在这一环节，如果借助 AI 辅助处理，应像数据整理阶段一样把任务、已有字段、期望输出和限制条件说明清楚。可按下面方式提出任务。

\begin{samepage}
\begin{userbox}
我已经整理了一张无机光电材料带隙数据表，字段包括化学式、晶体结构、空间群、带隙/eV。现在准备用于机器学习建模，请帮我设计从原始表格到特征表的处理流程。

要求如下：
\begin{enumerate}
\item 说明如何把化学式解析为元素组成和原子分数；
\item 说明如何利用元素属性表计算平均原子序数、平均电负性、电负性差和平均原子半径；
\item 指出晶体结构、空间群等文字字段哪些需要进一步数值化；
\item 提醒哪些字段不宜直接作为模型输入，以及如何避免目标泄漏；
\item 不要虚构具体数值，只给出处理步骤和检查要点。
\end{enumerate}
\end{userbox}

\begin{aibox}
可把处理流程分为五步：

\begin{enumerate}
\item 化学式解析：将 GaN、ZnO、SrTiO\textsubscript{3} 等化学式拆分为元素种类和原子比例，再归一化为原子分数。
\item \textbf{元素属性匹配：}准备元素属性表，至少包含元素符号、原子序数、电负性、原子半径等字段，并与解析后的元素组成对应。
\item \textbf{成分描述符计算：}按原子分数加权计算平均原子序数、平均电负性、平均原子半径，并计算电负性差等差异类特征。
\item \textbf{结构字段转换：}晶体结构名称和空间群符号不宜直接输入多数模型，应转换为空间群编号、晶胞体积、密度或其他可计算结构描述符。
\item \textbf{建模前检查：}带隙/eV 是预测目标，不能作为输入特征；发光颜色等由带隙直接决定的字段也要谨慎使用，避免目标泄漏。
\end{enumerate}

最终得到的特征表应做到字段清楚、单位一致、输入特征和预测目标分离，并保留原始数据来源，便于后续核验。
\end{aibox}
\end{samepage}

在实践中，描述符不是越多越好。如果数据样本较少，而输入特征过多，模型可能会过拟合。对于入门任务，可先选择少量具有明确材料意义、且容易从数据库获得的描述符，例如平均原子序数、平均电负性、电负性差、平均原子半径、空间群编号和晶胞体积。等研究者理解基本流程后，再逐步加入更复杂的局域环境描述符。


还要避免目标泄漏。目标泄漏指的是把与预测目标直接相关、甚至由目标结果推导出来的信息放入输入特征中。例如，如果目标是预测带隙，就不应把“发光颜色”作为输入特征，因为发光颜色本身往往已经反映了带隙大小。这样的模型即使看起来预测准确，也没有真正学习材料成分和结构对带隙的影响。


\subsection{图表示与物理对称性}

上一节介绍了成分描述符和结构描述符。通过这些描述符，研究者可将化学式、元素属性、空间群和晶胞体积等信息转化为数值特征，用于机器学习模型预测材料性能。这种方法直观、容易理解，也适合入门学习。但它仍然有一个明显问题：很多描述符是人工设计的，模型看到的是人工预先整理好的特征，而不是材料结构本身。


成分描述符依赖人工设计，适合快速形成特征表，但难以完整保留原子排列、局域配位和相互作用距离等结构信息。带隙、形成能、弹性性质、离子扩散能力等，都与原子之间的连接关系、距离和局域环境密切相关。如果模型只使用平均电负性、平均原子半径、空间群编号等手工描述符，可能会丢失一部分精细结构信息。


图表示为这一问题提供了另一种思路：把材料看作原子节点与原子间边组成的图结构。在这张图中，原子是节点，原子之间的邻近关系或相互作用是边。这样，模型就不再只是读取一组人工设计的表格特征，而是可直接从材料的原子结构中学习信息。


继续以无机光电材料带隙预测为例，假设研究者要预测 GaN 的带隙。上一节中，研究者可使用平均电负性、电负性差、平均原子半径、空间群编号等描述符表示 GaN。但如果采用图表示，问题就变为：Ga 原子和 N 原子在晶体中如何排列，它们之间的距离是多少，每个原子周围有多少邻居。


在晶体图中，节点通常代表原子。对于 GaN 来说，图中的节点可以是 Ga 原子和 N 原子。每个节点需要带有节点特征，用来描述这个原子的基本信息。例如，一个 Ga 节点可包含 Ga 的原子序数、元素类型、电负性、原子半径、价电子数等信息；一个 N 节点也可包含 N 的对应元素属性。这样，模型就知道每个节点代表什么元素。


边通常代表原子之间的邻近关系。晶体中原子不是孤立存在的，材料性能往往取决于原子周围的局域环境。因此，需要根据一定规则判断哪些原子之间建立边。最常见的方法是设置一个截断半径：如果两个原子之间的距离小于某个阈值，就认为它们在图中相连。例如，可规定距离小于 4 \AA{} 的原子之间建立边。这样，Ga 原子周围一定距离内的 N 原子会与它相连，形成局域原子环境。


每条边也可带有边特征。边特征可包括原子间距离、相对位置、键类型或距离展开后的数值表示。对于带隙预测来说，原子间距离和局域配位环境会影响轨道重叠、成键方式和电子能带结构。因此，图表示比简单成分描述符更接近材料真实结构。


\subsubsection*{\textbf{1. 案例示范：GaN 晶体图构建}}

以 GaN 为例，构建晶体图可按照以下步骤进行。


第一步，获得 GaN 的晶体结构文件。这个文件可来自材料数据库，常见格式包括 CIF\footnote{Crystallographic Information File，晶体学信息文件。} 文件或其他晶体结构格式。结构文件中记录了晶格参数、原子种类和原子坐标。对于图表示来说，晶体结构文件是关键输入，因为它提供了原子在空间中的排列信息。


第二步，读取原子信息。程序需要从结构文件中读取每个原子的元素类型和坐标。例如，一个晶胞中可能包含 Ga 原子和 N 原子，每个原子都有对应的空间位置。此时，每个原子就可被看作图中的一个节点。


第三步，为节点添加特征。对于每个原子节点，可用元素属性作为节点特征。例如，Ga 节点可用原子序数、电负性、原子半径、价电子数等表示；N 节点也用同样类型的属性表示。这样，不同元素虽然都是“节点”，但节点特征不同，模型可区分它们的化学性质。


第四步，根据原子间距离建立边。设定一个截断半径，例如 4 \AA{}。程序计算所有原子之间的距离，如果两个原子距离小于截断半径，就在它们之间建立一条边。由于晶体具有周期性，计算距离时还要考虑相邻晶胞中的原子，而不是只看一个孤立晶胞。否则，边界附近原子的邻居关系会被错误截断。


第五步，为边添加特征。最简单的边特征是两个原子之间的距离。例如，Ga--N 距离可作为边特征输入模型。更复杂的方法会把距离展开成一组函数值，使模型更容易学习不同距离对材料性质的影响。


第六步，将整张图与目标性能对应起来。对于带隙预测任务，GaN 的晶体图作为模型输入，GaN 的带隙数值作为模型输出。模型学习的目标，就是从晶体图中预测带隙。换言之，模型不再只学习“平均电负性--带隙”的关系，而是学习“原子种类、局域连接、原子距离和晶体结构--带隙”的关系。


用表格形式表示，图数据可概括为表 \ref{tab:gan-crystal-graph-data-components} 所示内容。


\begin{formalTable}[htbp]
\caption{GaN 晶体图数据组成示例}
\label{tab:gan-crystal-graph-data-components}
\begin{tabularx}{\textwidth}{@{}YYY@{}}
\toprule
\textbf{图的组成} & \textbf{在 GaN 中对应什么} & \textbf{作用} \\
\midrule
节点 & Ga 原子、N 原子 & 表示晶体中的原子 \\
节点特征 & 原子序数、电负性、原子半径等 & 描述每个原子的元素属性 \\
边 & 截断半径内相邻原子之间的连接 & 表示局域原子关系 \\
边特征 & 原子间距离、相对位置等 & 描述原子之间的空间关系 \\
图标签 & GaN 的带隙值 & 模型需要预测的目标性能 \\
\bottomrule
\end{tabularx}
\end{formalTable}

晶体图构建涉及结构文件读取、周期性边界条件和邻接关系检查，适合借助 AI 辅助梳理代码流程。但提示词不应只写“构建晶体图”，而要明确输入文件、截断半径、输出内容和检查要求。

\begin{samepage}
\begin{userbox}
我有一个 GaN 的 CIF 晶体结构文件，准备用它构建用于带隙预测的晶体图。请帮我设计一段 Python 处理流程，用 pymatgen 读取结构，并根据 4 \AA{} 截断半径建立原子邻接关系。

要求如下：
\begin{enumerate}
\item 输出每个原子的元素种类和坐标；
\item 将原子作为节点，并说明可使用哪些节点特征；
\item 根据 4 \AA{} 截断半径生成边列表，并说明是否考虑周期性边界条件；
\item 输出节点数量、边数量和部分邻居关系，便于检查图是否合理；
\item 不要假设 CIF 文件一定正确，提醒需要核对哪些结构信息。
\end{enumerate}
\end{userbox}

\begin{aibox}
可按以下流程处理：

\begin{enumerate}
\item 用 pymatgen 读取 GaN.cif，先检查晶格参数、原子种类、原子坐标和空间群信息是否完整。
\item 将每个原子作为图中的一个节点。节点特征可从元素属性中选取，例如原子序数、电负性、原子半径和价电子数。
\item 使用结构对象的邻居搜索功能，在 4 \AA{} 截断半径内寻找相邻原子。对于晶体材料，应使用考虑周期性边界条件的方法，而不是只计算单个晶胞内部距离。
\item 将邻居关系保存为边列表，边特征可先使用原子间距离。后续如果用于图神经网络，还可把距离展开为多个函数值。
\item 检查节点数、边数和每类原子的邻居数量。如果某些原子没有邻居，或者边数量明显异常，应重新检查截断半径、CIF 文件和周期性设置。
\end{enumerate}

这段流程只能帮助完成结构到图的转换。截断半径是否合理、节点特征是否足够、图结构是否符合 GaN 的晶体环境，仍然需要结合晶体结构知识判断。
\end{aibox}
\end{samepage}

图表示之所以适合材料结构，还因为它能够更自然地满足材料中的一些物理要求。材料性质不应该因为坐标系移动、整体旋转或原子编号改变而发生变化。例如，同一个 GaN 晶体，如果研究者把整个晶体平移一段距离，它的带隙不应该改变；如果研究者重新编号原子，材料性质也不应该改变；如果只是改变观察角度，材料本身也没有发生变化。


这些要求可概括为物理对称性。对于材料机器学习，常见的物理对称性包括平移不变性、旋转不变性和置换不变性。


平移不变性是指材料整体移动后，性质不应改变。换言之，原子坐标整体加上同一个位移，材料的带隙、形成能等性质不应变化。因此，模型不应只依赖原子的绝对坐标，而应该更多使用原子之间的相对位置、距离或局域环境。


旋转不变性是指材料整体旋转后，性质不应改变。对于带隙、形成能、总能量等标量性质，材料旋转一个角度后，预测值应该保持不变。因此，很多模型会使用原子间距离这类旋转不变的量作为边特征。


置换不变性是指相同类型原子的编号顺序改变后，性质不应改变。在晶体结构文件中，原子编号只是人为排列顺序，不代表真实物理差异。例如，同一个晶体中两个等价 Ga 原子交换编号后，化学组成、局域环境和带隙都没有改变，模型输出也不应改变。图神经网络通过对邻居信息进行聚合，可在一定程度上满足这种要求。


除了不变性，还有一个更进一步的概念叫等变性。从概念上看，不变性要求某些输出在变换后保持不变，而等变性要求输出随着输入变换而按照物理规律相应变化。例如，能量是标量，材料整体旋转后能量不变；力是矢量，如果结构旋转，力的方向也应该随之旋转。可以把它理解为：同一个原子环境换了观察方向，能量数值不变，但每个原子受力的方向要跟着坐标系一起转动。对于带隙预测这类标量任务，理解不变性已经足够；对于后续机器学习势、力预测和分子动力学模拟，等变性会更加重要。


因此，图表示不仅是一种数据格式，更是一种更符合材料物理特征的表示方式。它把材料看作原子节点和相互作用边组成的结构，同时尽量尊重材料中的平移、旋转和置换对称性。相比只使用手工描述符，图表示能够保留更多局域结构信息，也为后续图神经网络模型奠定基础。


图表示并不意味着完全不需要人工判断。构建晶体图时，截断半径、节点特征、边特征表示和周期性边界条件处理都会影响模型结果。即使借助工具生成代码并检查图结构，上述关键设置是否合理，仍然需要结合材料结构和预测任务来判断。

\subsubsection*{\textbf{小结}}

本节说明材料信息需要转化为描述符、结构字段或晶体图，才能进入机器学习模型。
成分和结构描述符适合表格建模，图表示则更适合保留原子连接、距离和局域环境。
特征构建必须避免目标泄漏，并结合物理对称性、截断半径和边特征检查输入是否合理。

\subsection*{习题}

\begin{enumerate}[leftmargin=*]
\item 比较成分描述符和结构描述符的作用及适用场景。
\item 图表示如何保留比成分描述符更多的结构信息？
\item 成分描述符和图表示各有什么优势与局限？请结合带隙预测举例说明。
\end{enumerate}

\clearpage
\section{性能预测与尺度模拟}


\begin{sectionkeypoints}
\item 材料性能预测的核心，是让模型学习材料特征与目标性能之间的对应关系。硬度、带隙、形成能、弹性模量和熔点都可成为预测目标。
\item 描述符模型适合表格数据和入门任务，关键步骤包括整理特征表、划分训练测试集、训练模型、评价误差和检查特征重要性。
\item 图神经网络适合处理晶体结构数据，能够从原子节点、原子间边和局域环境中学习材料性质。
\item 机器学习势学习的是原子结构与能量、力之间的关系，可连接第一性原理计算和分子动力学模拟。
\item 模型预测和尺度模拟都必须进行可靠性检查。训练集外推、过拟合、数据泄漏、结构文件错误和模拟参数不合理都会影响结论。
\end{sectionkeypoints}

\subsection{描述符模型与硬度预测}
\label{subsec:descriptor-hardness-prediction}

有了材料描述符特征表之后，下一步是建立特征与性能之间的映射关系。也就是说，模型需要根据成分、结构、工艺或显微组织等输入特征，预测硬度、带隙、形成能、熔点等目标性能。


材料性能预测的基本思想，是让模型学习“材料特征”和“目标性能”之间的对应关系。材料特征可来自成分、结构、工艺和显微组织，目标性能可以是硬度、强度、带隙、形成能、熔点或耐腐蚀性能等。对于传统机器学习方法而言，材料首先要被转化成一张特征表，每一行代表一个材料样品，每一列代表一个可计算的特征，最后一列通常是需要预测的性能。


\subsubsection*{\textbf{1. 案例示范：高熵合金涂层硬度预测}}

本节只围绕一个具体任务展开：高熵合金涂层硬度预测。假设研究者希望根据高熵合金涂层的成分和组织信息，预测涂层硬度。这个例子适合说明传统机器学习预测流程，因为硬度是材料实验中常见的数值型性能，既受成分影响，也受相结构、晶粒尺寸和硬质相等组织因素影响。


对于一个高熵合金涂层样品，原始实验记录可能包括名义成分、制备工艺、XRD\footnote{X-ray Diffraction，X 射线衍射。} 相分析、显微组织观察和显微硬度测试结果。经过前面的描述符构建过程后，这些信息可整理为表 \ref{tab:hea-hardness-features} 所示形式。


\begin{formalTable}[htbp]
\caption{高熵合金涂层硬度预测特征表示示例}
\label{tab:hea-hardness-features}
\scriptsize
\setlength{\tabcolsep}{2.8pt}
\begin{tabularx}{\textwidth}{@{}YYYYYYYY@{}}
\toprule
\textbf{样品编号} & \textbf{原子半径差} & \textbf{平均电负性} & \textbf{电负性差} & \textbf{价电子浓度} & \textbf{是否含硬质相} & \textbf{平均晶粒尺寸/\(\mu\)m} & \textbf{硬度/HV} \\
\midrule
S1 & 数值 & 数值 & 数值 & 数值 & 0 & 数值 & 数值 \\
S2 & 数值 & 数值 & 数值 & 数值 & 1 & 数值 & 数值 \\
S3 & 数值 & 数值 & 数值 & 数值 & 1 & 数值 & 数值 \\
S4 & 数值 & 数值 & 数值 & 数值 & 1 & 数值 & 数值 \\
\bottomrule
\end{tabularx}
\end{formalTable}

表 \ref{tab:hea-hardness-features} 中，“原子半径差、平均电负性、电负性差、价电子浓度、是否含硬质相、平均晶粒尺寸”等列是模型输入特征，“硬度/HV”是模型需要预测的目标值。换成机器学习语言，就是把输入特征组成矩阵 X，把硬度组成目标向量 y。模型训练的目的，就是从 X 中学习到 y 的变化规律。


在这个任务中，硬度是一个连续数值，因此它属于回归任务。回归任务的目标不是判断材料属于哪一类，而是预测一个具体数值。例如，模型输入某个涂层样品的成分和组织特征后，输出预测硬度是多少 HV。与之相对，如果任务是判断材料是否耐腐蚀、是否稳定、是否属于某种相结构，那就更接近分类任务。


为了让这个过程更加具体，以下给出一个精简的机器学习代码框架。假设研究者已经整理好了一个高熵合金涂层硬度数据表，文件名为 hea\_\allowbreak{}coating\_\allowbreak{}hardness.csv。表格中包含若干材料描述符和目标性能“硬度/HV”。代码的目标是读取数据、划分训练集和测试集、训练随机森林回归模型，并评价预测误差。代码采用平均绝对误差（Mean Absolute Error，MAE）、均方根误差（Root Mean Squared Error，RMSE）和决定系数（Coefficient of Determination，\(R^2\)）评价预测误差与拟合程度。


\begin{lstlisting}[language=Python,escapeinside={(*@}{@*)}]
# (*@\textnormal{1. 导入常用库}@*)
import pandas as pd
from sklearn.model_selection import train_test_split
from sklearn.ensemble import RandomForestRegressor
from sklearn.metrics import mean_absolute_error, mean_squared_error, r2_score
# (*@\textnormal{2. 读取高熵合金涂层硬度数据表}@*)
data = pd.read_csv("hea_coating_hardness.csv")
# (*@\textnormal{3. 指定输入特征和预测目标}@*)
# (*@\textnormal{这些特征来自前面已经构建好的材料描述符}@*)
feature_cols = [
    "atomic_radius_diff",             # (*@\textnormal{原子半径差}@*)
    "mean_electronegativity",         # (*@\textnormal{平均电负性}@*)
    "electronegativity_diff",         # (*@\textnormal{电负性差}@*)
    "valence_electron_concentration", # (*@\textnormal{价电子浓度}@*)
    "has_hard_phase",                 # (*@\textnormal{是否含硬质相：0 表示否，1 表示是}@*)
    "grain_size"                      # (*@\textnormal{平均晶粒尺寸}@*)
]
X = data[feature_cols]
y = data["hardness_HV"]
# (*@\textnormal{4. 划分训练集和测试集}@*)
X_train, X_test, y_train, y_test = train_test_split(
    X,
    y,
    test_size=0.2,
    random_state=42
)
# (*@\textnormal{5. 建立并训练随机森林回归模型}@*)
model = RandomForestRegressor(
    n_estimators=200,
    random_state=42
)
model.fit(X_train, y_train)
# (*@\textnormal{6. 在测试集上进行硬度预测}@*)
y_pred = model.predict(X_test)
# (*@\textnormal{7. 评价模型预测效果}@*)
mae = mean_absolute_error(y_test, y_pred)
rmse = mean_squared_error(y_test, y_pred) ** 0.5
r2 = r2_score(y_test, y_pred)
print("MAE:", mae)
print("RMSE:", rmse)
print("R2:", r2)
# (*@\textnormal{8. 输出特征重要性}@*)
importance = pd.DataFrame({
    "feature": feature_cols,
    "importance": model.feature_importances_
}).sort_values(by="importance", ascending=False)
print(importance)
\end{lstlisting}

这段代码对应了材料性能预测的基本流程。首先，pd.read\_\allowbreak{}csv() 用于读取已经整理好的材料数据表。此处的数据表不是原始实验记录，而是已经经过描述符构建后的特征表。换言之，代码中的输入不是“FeCoCrNiTi”这样的化学式，而是“原子半径差、平均电负性、电负性差、价电子浓度、是否含硬质相、平均晶粒尺寸”等数值特征。


其中，X = data[feature\_\allowbreak{}cols] 表示模型输入，y = data["hardness\_\allowbreak{}HV"] 表示模型输出。对于本例来说，X 是涂层样品的成分和组织特征，y 是实验测得的硬度值。模型训练时要学习的，就是这些特征与硬度之间的关系。


train\_\allowbreak{}test\_\allowbreak{}split() 用于划分训练集和测试集。训练集用于让模型学习规律，测试集用于检验模型对新样品的预测能力。此处设置 test\_\allowbreak{}size=0.2，表示 80\% 的样品用于训练，20\% 的样品用于测试。对于材料数据来说，样本数量往往不大，因此测试集必须严格独立，不能参与模型训练。否则，模型评估结果会偏高，不能代表真实预测能力。


本例使用的模型是随机森林回归模型。随机森林由多棵决策树组成，能够处理非线性关系，对小规模表格数据比较友好，也可输出特征重要性。对于高熵合金涂层硬度预测来说，成分描述符、硬质相和晶粒尺寸与硬度之间通常不是简单线性关系，因此随机森林适合作为入门模型。此处选择随机森林并不意味着它一定最优，而是因为它便于研究者理解完整建模流程。


模型训练完成后，model.predict(X\_\allowbreak{}test) 会对测试集样品进行硬度预测。此时模型只看到测试样品的输入特征，不知道真实硬度。随后，研究者用真实硬度和预测硬度进行比较，计算 MAE、RMSE 和 R\textsuperscript{2}。


MAE 是平均绝对误差，表示模型预测值与真实值平均相差多少。例如，如果 MAE 为 30 HV，可理解为模型对测试样品的硬度平均预测误差约为 30 HV。RMSE 是均方根误差，对较大的预测误差更加敏感。如果少数样品预测偏差很大，RMSE 会明显增大。R\textsuperscript{2} 是决定系数，用来衡量预测值与真实值变化趋势的一致程度，越接近 1，说明模型预测效果越好。


对于材料研究来说，仅有误差指标还不够。宜进一步绘制“真实硬度--预测硬度”散点图。如果模型预测较好，散点应大致分布在 y = x 参考线附近。图中的坐标轴、单位、图例和参考线都要清楚，否则模型结果不便于与实验报告或论文表达衔接。


除了预测误差，随机森林模型还可输出特征重要性。代码中的 model.feature\_\allowbreak{}importances\_\allowbreak{} 可显示不同输入特征对硬度预测的贡献大小。如果结果显示“是否含硬质相”“平均晶粒尺寸”“原子半径差”等特征重要性较高，就可结合材料强化机制进行分析。例如，硬质相可能带来第二相强化，晶粒细化可能带来细晶强化，原子半径差可能与固溶强化和晶格畸变有关。


在代码实现阶段，AI 更适合承担“字段核对、代码改写、报错解释和绘图补充”等辅助工作，而不是替研究者判断哪些特征一定正确。可用下面的提示词约束任务边界。

\begin{samepage}
\begin{userbox}
我正在用随机森林预测高熵合金涂层硬度，CSV 表头包括原子半径差、平均电负性、电负性差、价电子浓度、是否含硬质相、晶粒尺寸和硬度/HV。请帮我检查这段建模流程是否清楚，并补充测试集真实硬度与预测硬度散点图的绘图步骤。

要求如下：
\begin{enumerate}
\item 判断哪些列可作为输入特征，哪一列是预测目标；
\item 提醒哪些字段需要确认数据来源或实验条件；
\item 说明如果出现列名不存在、缺失值或字符串字段无法输入模型，应如何检查；
\item 补充绘制真实值--预测值散点图和 y=x 参考线的思路；
\item 不要替我判断模型一定可靠，只给出代码检查和结果呈现建议。
\end{enumerate}
\end{userbox}

\begin{aibox}
表 \ref{tab:hea-hardness-features} 中，原子半径差、平均电负性、电负性差、价电子浓度、是否含硬质相和晶粒尺寸可作为候选输入特征，硬度/HV 是预测目标，不应同时放入输入特征。

建模前需要重点检查三类问题：
\begin{enumerate}
\item \textbf{字段是否存在：}代码中指定的输入特征名称必须与 CSV 表头完全一致，否则会出现列名不存在的错误。
\item \textbf{数据是否可输入模型：}输入特征应为数值型；如果存在空值、文字标签或单位混用，需要先清洗或编码。
\item \textbf{材料含义是否可靠：}是否含硬质相应来自 XRD、SEM\footnote{Scanning Electron Microscopy，扫描电子显微镜。} 或其他表征结果，晶粒尺寸应说明统计方法，硬度值应记录测试载荷和保载时间。
\end{enumerate}

模型评价时，除 MAE、RMSE 和 R\textsuperscript{2} 外，可绘制测试集真实硬度与预测硬度散点图，并添加 y=x 参考线。如果散点明显偏离参考线，说明模型在部分样品上的预测误差较大，需要回到数据质量、样本数量和特征设计中继续检查。
\end{aibox}
\end{samepage}

在实际建模时，还应特别注意过拟合。材料数据往往样本少、特征多，如果只有十几个样品却加入几十个描述符，模型可能在训练集上表现很好，但在测试集上预测很差。这说明模型可能只是记住了已有样品，而没有真正学到可推广的规律。为了降低过拟合风险，可减少无关特征，采用交叉验证，增加样本数量，或者选择更简单、更稳定的模型。是否属于过拟合，仍然要结合训练误差、测试误差、数据规模和材料任务共同判断。


数据泄漏同样需要警惕。预测硬度时，不应把由硬度结果直接推导出的字段作为输入特征。例如，“是否高硬度样品”“硬度等级”这类字段不能作为输入，因为它们本身已经包含了目标信息。否则模型看起来预测很准，实际上并没有从成分和组织中学习硬度规律。


\subsection{图神经网络与性质预测}

上一节以高熵合金涂层硬度预测为例，介绍了传统机器学习的基本流程。那种方法的特点是：研究者首先根据材料知识人工构造描述符，例如原子半径差、电负性差、价电子浓度、是否含硬质相和晶粒尺寸，然后再把这些描述符输入随机森林等模型中预测硬度。这种方法直观、容易操作，也适合材料数据量较少的入门任务。


但是，传统描述符方法也有局限。材料性能不仅取决于几个平均特征，还与原子在晶体中的排列方式、局域配位环境、原子间距离和结构稳定性有关。尤其对于高熵合金涂层这类复杂体系，材料中可能同时存在 FCC\footnote{Face-Centered Cubic，面心立方结构。} 固溶体、BCC\footnote{Body-Centered Cubic，体心立方结构。} 相、碳化物或硼化物等不同结构。仅用“是否含硬质相”“平均晶粒尺寸”这样的表格特征，很难完整表达原子尺度结构对性能的影响。


图神经网络提供了另一种建模思路。它不再把材料完全压缩成一行人工描述符，而是把材料看作由原子节点和原子间关系组成的图结构。在晶体图中，每个原子是一个节点，原子之间的邻近关系是边，节点和边都可带有特征。模型通过在图上进行信息传递，学习局域原子环境与材料性质之间的关系。


这里要限定问题边界：高熵合金涂层是一个多尺度、多相组织体系，真实涂层中包含晶粒、晶界、第二相、孔隙、界面和残余应力等复杂因素。图神经网络不宜简单地把一张 SEM 图或一个涂层样品直接转化成完整晶体图，然后准确预测整体硬度。更合理的入门方式是：选取涂层中的代表性晶体相或候选晶体结构，例如 FCC 固溶体、BCC 固溶体或 TiC 等硬质相，把它们转化为晶体图，学习结构与形成能、弹性模量或稳定性等性质之间的关系。然后再把这些原子尺度信息作为理解涂层硬度和强化机制的补充。


这里的重点不是用图神经网络直接预测整个涂层硬度，而是说明如何把涂层中的代表性晶体结构转化为图，并用图神经网络预测与性能相关的材料性质。


假设研究者研究的是一种 FeCoCrNiTi 高熵合金涂层。实验表明，涂层中可能存在 FCC 固溶体、BCC 相和 TiC 等硬质相。上一节中，研究者用人工描述符预测涂层硬度；现在研究者从更微观的角度思考：如果想理解某个相为什么可能更稳定、为什么某种硬质相会提高硬度，就需要进一步关注原子尺度结构。图神经网络正适合处理这类晶体结构信息。


对于一个代表性晶体结构，原始数据通常来自 CIF 文件或数据库结构文件。一个结构文件中记录了晶格参数、原子种类和原子坐标。图神经网络的第一步，就是把这个晶体结构转化为晶体图。


\subsubsection*{\textbf{1. 案例示范：TiC 晶体图性质预测}}

以 TiC 硬质相为例，它在高熵合金涂层中常被视为提高硬度的重要相之一。构建 TiC 晶体图时，可按照以下步骤进行。


第一步，读取晶体结构文件。结构文件中包含 Ti 和 C 原子的种类、坐标以及晶格参数。对于模型来说，这些信息是图结构的基础。


第二步，把原子转化为节点。TiC 中的 Ti 原子和 C 原子都可看作图中的节点。每个节点需要包含节点特征，例如原子序数、电负性、原子半径、价电子数等。这样，模型能够区分 Ti 节点和 C 节点，并理解它们具有不同化学属性。


第三步，根据原子间距离建立边。在晶体中，一个原子的性质受到周围邻居原子的影响。因此，可设置一个截断半径，例如 5 \AA{}。如果两个原子之间的距离小于这个截断半径，就在它们之间建立一条边。这样，模型就可知道哪些原子彼此相邻。


第四步，为边添加特征。最常见的边特征是原子间距离。距离越近，原子之间相互作用通常越强；距离不同，键合环境也不同。因此，边特征可帮助模型学习局域结构对性质的影响。


第五步，把整个晶体图与目标性质对应起来。如果目标是预测形成能，那么每一个晶体图对应一个形成能数值；如果目标是预测弹性模量，那么每一个晶体图对应一个弹性模量数值。模型训练的目标，就是学习“晶体图结构”与“目标性质”之间的关系。


该流程可概括为下列步骤，并可结合图 \ref{fig:crystal-graph-prediction} 理解：


\begin{lightquote}
\centering 晶体结构文件 → 原子节点 → 原子间边 → 节点特征和边特征 → 晶体图 → 性质预测
\end{lightquote}

\begin{figure}[H]
\centering
\BookFigureImage{images_11/2.png}
\caption{晶体图表示与性能预测示意图}
\label{fig:crystal-graph-prediction}
\end{figure}
\FloatBarrier

图 \ref{fig:crystal-graph-prediction} 所示流程中，图数据的组成可进一步概括为表 \ref{tab:crystal-graph-data-components}。


\begin{formalTable}[htbp]
\caption{晶体图数据组成示例}
\label{tab:crystal-graph-data-components}
\begin{tabularx}{\textwidth}{@{}YY@{}}
\toprule
\textbf{图数据组成} & \textbf{在 TiC 或高熵合金相关晶体结构中表示什么} \\
\midrule
节点 & Ti、C 或 Fe、Co、Cr、Ni、Ti 等原子 \\
节点特征 & 原子序数、电负性、原子半径、价电子数等 \\
边 & 截断半径内的原子邻近关系 \\
边特征 & 原子间距离或距离展开特征 \\
图标签 & 形成能、弹性模量、体模量或其他目标性质 \\
\bottomrule
\end{tabularx}
\end{formalTable}

图神经网络的核心是信息传递。每个原子节点在初始时只包含自己的元素属性，例如 Ti 节点知道自己是 Ti，C 节点知道自己是 C。经过一轮信息传递后，一个节点会接收周围邻居节点的信息。经过多轮信息传递后，节点就不仅包含自身信息，还包含局域原子环境信息。最后，模型把所有节点的信息汇总成整个晶体的表示，再输出目标性质。


这个过程可用材料语言理解：模型不只是看“这个材料含有哪些元素”，而是进一步看“这些元素在空间中如何排列、每个原子周围是什么邻居、原子间距离是多少、局域环境是否稳定”。这正是图神经网络相比传统描述符方法的重要优势。


为了让这个过程更具体，以下给出一个教学用的精简代码框架。它的目的不是提供完整可直接运行的项目，而是帮助研究者理解图神经网络建模的基本步骤。


\begin{lstlisting}[language=Python,escapeinside={(*@}{@*)}]
# (*@\textnormal{教学用简化框架：从晶体结构构建图并预测性质}@*)
# (*@\textnormal{实际运行时需要安装 pymatgen、torch、torch\_\allowbreak{}geometric 等库}@*)
from pymatgen.core import Structure
import torch
# (*@\textnormal{1. 读取晶体结构文件，例如 TiC.cif}@*)
structure = Structure.from_file("TiC.cif")
# (*@\textnormal{2. 设置截断半径}@*)
cutoff = 5.0
# (*@\textnormal{3. 构建节点特征：此处仅用原子序数作为示例}@*)
node_features = []
for site in structure:
    atomic_number = site.specie.Z
    node_features.append([atomic_number])
x = torch.tensor(node_features, dtype=torch.float)
# (*@\textnormal{4. 构建边：寻找截断半径内的邻居原子}@*)
edge_index = []
edge_attr = []
for i, site in enumerate(structure):
    neighbors = structure.get_neighbors(site, cutoff)
    for neighbor in neighbors:
        j = neighbor.index
        distance = neighbor.nn_distance
        edge_index.append([i, j])
        edge_attr.append([distance])
edge_index = torch.tensor(edge_index, dtype=torch.long).t().contiguous()
edge_attr = torch.tensor(edge_attr, dtype=torch.float)
# (*@\textnormal{5. 图标签：例如该结构的形成能或弹性模量}@*)
# (*@\textnormal{此处用一个占位数值表示，实际应来自数据库或计算结果}@*)
y = torch.tensor([target_property], dtype=torch.float)
# (*@\textnormal{至此，一个晶体结构已经被转化为图数据：}@*)
# (*@\textnormal{x 表示节点特征}@*)
# (*@\textnormal{edge\_\allowbreak{}index 表示边连接关系}@*)
# (*@\textnormal{edge\_\allowbreak{}attr 表示边特征}@*)
# (*@\textnormal{y 表示目标性质}@*)
\end{lstlisting}

这段代码展示了晶体图构建的基本思想。Structure.from\_\allowbreak{}file("TiC.cif") 用来读取晶体结构文件；node\_\allowbreak{}features 保存每个原子的节点特征；edge\_\allowbreak{}index 保存哪些原子之间存在边；edge\_\allowbreak{}attr 保存边特征，例如原子间距离；y 是模型要预测的目标性质。


实际使用时，节点特征不应只包含原子序数，还可加入电负性、原子半径、价电子数等元素属性。边特征也不一定只用一个距离数值，可把距离展开成一组函数，以便模型更好学习不同距离范围内的相互作用。对于不同材料任务，目标 y 也可不同。如果关注结构稳定性，可预测形成能；如果关注力学性能，可预测弹性模量或体模量；如果进一步结合实验数据，也可尝试预测与硬度相关的性质。


如果借助 AI 辅助构建 TiC 晶体图，提示词应突出“代表性晶体相”和“目标性质”两个信息，避免把图神经网络误写成直接预测整个涂层硬度的工具。示例任务可写成如下形式。

\begin{samepage}
\begin{userbox}
我想以高熵合金涂层中的 TiC 硬质相为例，构建一个用于性质预测的晶体图。请帮我设计从 TiC.cif 到图数据的处理流程，目标性质可以是形成能、弹性模量或体模量。

要求如下：
\begin{enumerate}
\item 说明如何读取 TiC.cif，并把 Ti 和 C 原子转化为节点；
\item 说明节点特征可选择哪些元素属性；
\item 根据 5 \AA{} 截断半径建立边，并说明为什么要考虑周期性边界条件；
\item 说明边特征和图标签分别是什么；
\item 提醒这个晶体图只能表示 TiC 这一代表性相，不宜直接代表整个高熵合金涂层硬度。
\end{enumerate}
\end{userbox}

\begin{aibox}
可按以下流程构建图数据：

\begin{enumerate}
\item \textbf{读取结构：}使用 pymatgen 读取 TiC.cif，检查晶格参数、Ti/C 原子坐标和结构完整性。
\item \textbf{构建节点：}把每个 Ti 原子和 C 原子作为节点。节点特征可包括原子序数、电负性、原子半径和价电子数。
\item \textbf{构建边：}使用 5 \AA{} 截断半径寻找邻近原子。由于 TiC 是周期晶体，邻居搜索必须考虑相邻晶胞，否则会漏掉晶胞边界附近的邻居。
\item \textbf{添加边特征：}可先使用原子间距离作为边特征，后续再根据模型需要使用距离展开特征。
\item \textbf{设置图标签：}如果训练目标是形成能，图标签就是形成能；如果目标是弹性模量或体模量，图标签应来自数据库、第一性原理计算或可靠实验/计算资料。
\end{enumerate}

这里需要专门说明，TiC 晶体图只能帮助理解涂层中某个硬质相的结构性质。真实高熵合金涂层硬度还受到晶粒尺寸、相分布、界面、孔隙和残余应力影响，不能由单个 TiC 晶体图直接决定。
\end{aibox}
\end{samepage}

与传统描述符模型相比，图神经网络的优势在于它能更直接地保留原子结构信息。传统模型需要人工预先设计“平均电负性”“原子半径差”等特征，而图神经网络可从原子种类、原子间距离和局域连接关系中自动学习结构信息。对于晶体材料来说，这种表示方式更接近材料本身。


但是，图神经网络也不是万能的。首先，它需要可靠的晶体结构输入。如果 CIF 文件错误、原子占位不清楚或结构不完整，构建出的图也会有问题。其次，它通常需要较多训练数据。如果只有少量结构样本，图神经网络也可能过拟合。再次，对于真实高熵合金涂层来说，整体硬度还受到晶粒尺寸、相分布、孔隙率、界面结合和残余应力等显微组织因素影响，单个晶体图不能完全代表整个涂层。


因此，在高熵合金涂层研究中，图神经网络更适合作为结构层面的补充工具。它可帮助研究者理解某个候选相、硬质相或代表性晶体结构的稳定性和力学相关性质；而整体涂层硬度仍然需要结合上一节中的表格特征模型、显微组织数据和实验测试结果共同分析。


\subsection{原子间势与熔点模拟}

前两节讨论的是材料性能预测。无论是传统机器学习模型，还是图神经网络模型，它们通常都是输入一个材料样品或晶体结构，然后输出一个性能数值，例如硬度、形成能、带隙或弹性模量。这类模型回答的是“这个材料可能具有什么性质”。


但是，材料科学中还有另一类重要问题：材料在原子尺度上如何运动、如何扩散、如何相变、如何熔化。例如，超高温陶瓷在高温下能否保持结构稳定？原子在升温过程中是否发生剧烈扩散？材料在什么温度附近开始熔化？这些问题不是简单预测一个静态数值就能完全回答的，而需要模拟原子随时间的运动过程。


分子动力学可用来研究这类原子尺度动态行为。在分子动力学中，研究者根据原子之间的相互作用计算每个原子的受力，再根据力更新原子位置，从而模拟材料随时间演化的过程。问题在于，原子间相互作用如何计算。第一性原理分子动力学精度较高，但计算成本高，难以处理大体系和长时间尺度；传统经验势速度快，但对复杂材料体系的适用性有限。机器学习势正是在这种背景下发展起来的。


机器学习势的核心思想，是用机器学习模型学习第一性原理计算得到的势能面。从概念上看，第一性原理计算可给出某一组原子结构的能量、原子受力和应力；机器学习势则利用这些数据学习“原子结构--能量/力”之间的关系。训练完成后，机器学习势可像传统势函数一样快速计算能量和力，从而支持更大规模、更长时间的分子动力学模拟。


\subsubsection*{\textbf{1. 案例示范：TiC 熔化行为模拟}}

本节只围绕一个任务展开：用机器学习势辅助模拟 TiC 超高温陶瓷的熔化行为。TiC 是典型的超高温陶瓷之一，具有高熔点、高硬度和良好的高温稳定性，也常作为高熵合金涂层或复合涂层中的硬质增强相。通过这个例子，可说明机器学习势如何从第一性原理数据出发，进一步服务于高温分子动力学模拟。


首先需要明确，机器学习势与前面性能预测模型的目标不同。前面随机森林模型预测的是硬度，图神经网络可预测形成能或弹性模量；而机器学习势要预测的是原子结构对应的总能量和每个原子的受力。它不是直接告诉研究者“TiC 的熔点是多少”，而是先学会在不同原子构型下计算能量和力，然后通过分子动力学模拟升温过程，再从模拟结果中判断熔化行为。


对于 TiC 熔点模拟，可按照以下流程理解。


第一步，准备 TiC 的初始晶体结构。研究者需要获得 TiC 的晶体结构文件，例如 CIF 文件。这个结构文件包含 Ti 和 C 原子的排列方式、晶格参数和空间结构。初始结构通常可来自数据库、文献或已有计算结果。


第二步，生成不同构型的训练数据。机器学习势不应只看一个理想晶体结构，因为真实高温模拟中原子会振动、偏移，甚至出现缺陷和局部无序。因此，需要构造一批不同的 TiC 原子构型，例如平衡晶体结构、轻微扰动结构、不同体积结构、不同温度下短时间分子动力学得到的结构等。这样，训练数据才能覆盖模型后续可能遇到的原子环境。


第三步，用第一性原理计算这些构型的能量和力。对于每一个 TiC 构型，需要通过第一性原理计算得到总能量、每个原子的力以及必要时的应力。这一步计算成本较高，但只需要对训练集中的有限构型进行。可理解为：第一性原理计算负责提供高质量“标准答案”，机器学习势负责学习这些标准答案中的规律。


第四步，训练机器学习势。训练时，模型输入原子种类和原子坐标，输出预测能量和原子力。训练目标是让模型预测的能量和力尽可能接近第一性原理计算结果。常见机器学习势包括 Behler--Parrinello 神经网络势、高斯近似势、DeepMD\footnote{Deep Potential Molecular Dynamics，深度势分子动力学。}、MACE\footnote{Message Passing Atomic Cluster Expansion，消息传递原子簇展开模型。}、NequIP\footnote{Neural Equivariant Interatomic Potentials，神经等变原子间势。}、M3GNet\footnote{Materials 3-body Graph Network，材料三体图网络。}、CHGNet\footnote{Crystal Hamiltonian Graph Neural Network，晶体哈密顿量图神经网络。} 等。教材入门阶段不需要详细展开每种模型的数学细节，重点是理解它们都在学习原子结构与能量/力之间的关系。


第五步，验证机器学习势是否可靠。训练完成后，不宜直接用于高温模拟，而要先在测试构型上验证预测误差。常见做法是比较模型预测能量与第一性原理能量之间的误差，以及模型预测力与第一性原理力之间的误差。如果误差较大，说明模型还没有学好 TiC 的势能面，需要增加训练数据、调整模型参数或补充高温构型。


第六步，用训练好的机器学习势进行分子动力学模拟。此时，机器学习势可快速计算每一步原子的能量和力，进而模拟 TiC 在不同温度下的原子运动。为了观察熔化行为，可逐步升高温度，例如从较低温度开始，逐步升温到高温区间，记录能量、体积、径向分布函数、均方位移等随温度变化的情况。


第七步，根据模拟结果判断熔化行为。材料熔化时，原子排列会从有序晶体逐渐转变为无序液态结构。模拟中可通过多个指标辅助判断：能量--温度曲线是否出现突变，体积是否发生明显变化，径向分布函数的长程有序峰是否消失，均方位移是否显著增加。通过这些现象，可估计 TiC 的熔化温度范围。


这个过程可概括为：


\begin{lightquote}
\centering TiC 晶体结构 → 多种原子构型 → 第一性原理能量和力 → 训练机器学习势 → 验证误差 → 分子动力学升温模拟 → 分析熔化行为
\end{lightquote}

为了让研究者更清楚地理解流程，以下给出一个教学用的简化代码框架。它不是完整可直接运行的工程代码，因为真实机器学习势训练通常需要专门软件和大量计算资源。此处的目的，是帮助理解每一步在做什么。


\begin{lstlisting}[language=Python,escapeinside={(*@}{@*)}]
# (*@\textnormal{教学用简化框架：机器学习势辅助 TiC 熔化模拟}@*)
# (*@\textnormal{实际研究中可使用 DeepMD、MACE、NequIP、CHGNet、M3GNet 等工具}@*)
# (*@\textnormal{1. 读取 TiC 初始晶体结构}@*)
structure = read_structure("TiC.cif")
# (*@\textnormal{2. 生成训练构型}@*)
# (*@\textnormal{包括平衡结构、扰动结构、不同体积结构和高温采样结构}@*)
train_structures = generate_configurations(
    structure,
    perturb=True,
    volume_changes=True,
    temperature_sampling=True
)
# (*@\textnormal{3. 对训练构型进行第一性原理计算}@*)
# (*@\textnormal{得到每个构型的能量、原子力和应力}@*)
dft_dataset = []
for config in train_structures:
    energy, forces, stress = run_dft_calculation(config)
    dft_dataset.append({
        "structure": config,
        "energy": energy,
        "forces": forces,
        "stress": stress
    })
# (*@\textnormal{4. 训练机器学习势}@*)
ml_potential = train_machine_learning_potential(
    dataset=dft_dataset,
    target_properties=["energy", "forces", "stress"]
)
# (*@\textnormal{5. 在测试构型上验证模型误差}@*)
test_results = evaluate_potential(
    model=ml_potential,
    test_dataset="test_configurations"
)
print(test_results)
# (*@\textnormal{6. 使用训练好的机器学习势进行分子动力学升温模拟}@*)
md_results = run_molecular_dynamics(
    structure=structure,
    potential=ml_potential,
    temperature_schedule=[1000, 1500, 2000, 2500, 3000, 3500, 4000],
    simulation_time="several ps or longer"
)
# (*@\textnormal{7. 分析熔化行为}@*)
plot_energy_temperature(md_results)
plot_radial_distribution_function(md_results)
plot_mean_square_displacement(md_results)
\end{lstlisting}

这段框架中，read\_\allowbreak{}structure() 表示读取 TiC 的晶体结构；generate\_\allowbreak{}configurations() 表示生成不同原子构型；run\_\allowbreak{}dft\_\allowbreak{}calculation() 表示用第一性原理计算能量和力；train\_\allowbreak{}machine\_\allowbreak{}learning\_\allowbreak{}potential() 表示训练机器学习势；run\_\allowbreak{}molecular\_\allowbreak{}dynamics() 表示用训练好的势进行分子动力学模拟；最后通过能量曲线、径向分布函数和均方位移分析熔化行为。


与前两节代码不同，此处给出的是流程框架，而不是几十行就能完成的入门脚本。原因是机器学习势涉及第一性原理计算、模型训练和分子动力学模拟，真实流程需要较多软件和计算资源。但这个框架可帮助研究者理解机器学习势在材料研究中的位置：它不是直接替代实验，也不是简单回归一个性能值，而是连接第一性原理计算和大尺度原子模拟的桥梁。


在机器学习势和分子动力学任务中，AI 辅助更适合用于解释文件字段、检查流程是否完整和整理输出指标，而不宜直接判断熔点结论。下面的写法更适合用于流程核查。

\begin{samepage}
\begin{userbox}
我想用机器学习势辅助模拟 TiC 的熔化行为。已有 TiC 初始结构，计划生成扰动构型和高温采样构型，用第一性原理计算能量与力，再训练机器学习势并进行升温分子动力学模拟。请帮我检查这个流程是否完整，并说明每类输出指标的物理意义。

要求如下：
\begin{enumerate}
\item 说明训练数据应包含哪些 TiC 原子构型；
\item 说明第一性原理计算需要记录哪些目标量；
\item 检查机器学习势训练后应如何验证能量和力误差；
\item 说明升温分子动力学中能量、体积、径向分布函数和均方位移分别反映什么；
\item 明确提醒不应只根据单一指标判断熔点。
\end{enumerate}
\end{userbox}

\begin{aibox}
这个流程可分为四个检查重点：

\begin{enumerate}
\item \textbf{训练构型覆盖范围：}除平衡 TiC 晶体结构外，还应包含轻微扰动结构、不同体积结构、不同温度采样结构，以及接近高温无序状态的构型。否则模型可能只适用于低温平衡结构。
\item \textbf{第一性原理标注目标：}每个构型至少应记录总能量和每个原子的力；如果后续模拟涉及体积变化，还应考虑应力信息。
\item \textbf{机器学习势验证：}训练后应在独立测试构型上比较预测能量、预测力与第一性原理结果的误差，不应只看训练集误差。
\item \textbf{熔化行为分析：}能量--温度曲线突变可能反映相变，体积突变可能反映结构膨胀，径向分布函数长程有序峰消失说明晶体有序性降低，均方位移显著增大说明原子扩散增强。
\end{enumerate}

熔点判断应综合多个指标，并结合体系大小、升温速率、模拟时间和训练数据覆盖范围共同分析。单独依靠 AI 对曲线的文字解释，不能作为可靠的材料结论。
\end{aibox}
\end{samepage}

需要特别强调的是，机器学习势的可靠性高度依赖训练数据。如果训练集中只包含低温平衡结构，而模拟时却用于高温熔化过程，模型很可能外推到训练分布之外，预测结果就不可靠。因此，在 TiC 熔点模拟中，训练数据应尽量包含高温扰动构型、不同体积构型和接近熔化状态的构型。否则，模型可能在低温结构上误差很小，却在高温模拟中产生错误行为。


此外，机器学习势不能完全替代第一性原理计算和实验验证。它的优势是加速模拟，但前提是训练数据可靠、验证误差合理、模拟条件设置正确。对于超高温陶瓷这类材料，熔点模拟还可能受到模拟体系大小、升温速率、时间尺度和缺陷状态影响。因此，模拟得到的熔化温度应被理解为计算预测结果，需要与实验数据或更高精度计算进行对照。

\subsubsection*{\textbf{小结}}

本节围绕硬度预测、图神经网络性质预测和机器学习势展开。
描述符模型用于表格化的宏观性能预测，图神经网络用于原子结构关联的性质预测，机器学习势用于原子尺度动力学模拟。
性能指标、数据泄漏、过拟合、结构文件质量和训练构型覆盖都会影响模型结论的可靠性。

\subsection*{习题}

\begin{enumerate}[leftmargin=*]
\item 比较描述符模型、图神经网络和机器学习势的适用任务。
\item 机器学习势适合哪些模拟场景？训练覆盖为何重要？
\item 为高熵合金涂层硬度预测设计输入特征，并说明依据。
\end{enumerate}

\clearpage
\section{模型解释与知识嵌入}


\begin{sectionkeypoints}
\item 可解释性分析关注模型训练之后的问题，即模型主要依赖哪些特征、这些特征如何影响预测结果，以及解释是否符合材料机理。
\item 特征重要性和 SHAP\footnote{SHapley Additive exPlanations，一种基于博弈论的可解释性分析方法。} 可帮助把机器学习输出转化为材料语言，但它们反映的是当前数据和当前模型中的关系，不等同于材料真理。
\item 物理知识嵌入强调在建模前和建模过程中主动引入材料规律，包括特征选择、数据筛选、目标泄漏检查和结果合理性判断。
\item 对高熵合金涂层硬度预测而言，原子半径差、晶粒尺寸、硬质相和孔隙率等特征应与固溶强化、细晶强化、第二相强化和组织致密化联系起来理解。
\item 一个可靠的材料人工智能模型，不仅要有较小预测误差，还要能被实验表征、物理机理和材料常识共同支持。
\end{sectionkeypoints}

\subsection{模型解释与特征归因}

前面几节已经介绍了材料性能预测的基本方法。通过描述符和机器学习模型，研究者可根据材料成分、结构和组织信息预测硬度、带隙、形成能、熔点等性能。对于工程应用来说，预测准确当然重要；但对于材料科学研究来说，仅仅得到一个预测数值还不够。研究者还需要进一步知道：模型为什么给出这样的预测？它主要依据了哪些材料特征？这些特征是否符合已有材料机理？


这正是模型可解释性要解决的问题。可解释性关注的是模型在预测时主要依赖哪些输入特征，以及这些特征如何影响预测结果。对于材料科学而言，可解释性尤其重要，因为材料研究的目标不仅是预测某个样品性能，更是提炼可理解、可验证、可迁移的材料规律。


本节继续围绕前面的高熵合金涂层硬度预测任务展开。假设研究者已经训练了一个随机森林模型，用来预测不同高熵合金涂层的硬度。模型的输入特征包括原子半径差、平均电负性、电负性差、价电子浓度、是否含硬质相和平均晶粒尺寸，输出目标是硬度/HV。现在的问题不再是“模型能不能预测硬度”，而是“模型主要根据什么判断某个涂层硬度较高”。


如果只看预测误差，例如 MAE、RMSE 和 R\textsuperscript{2}，研究者只能知道模型预测效果大致如何，却不知道模型内部学到了什么。例如，一个模型可能在测试集上误差较小，但它可能主要依赖了某个偶然相关的字段，而不是真正反映材料强化机制的特征。这样的模型即使暂时预测准确，也很难让研究者放心。因此，在材料人工智能中，模型解释往往和模型评价同样重要。


对于随机森林这类模型，最常见的解释方法是特征重要性分析。特征重要性可粗略告诉研究者：模型在预测硬度时，哪些输入特征贡献更大。假设训练后的模型输出表 \ref{tab:hardness-feature-importance} 所示结果。


\begin{formalTable}[htbp]
\caption{随机森林硬度预测特征重要性示例}
\label{tab:hardness-feature-importance}
\begin{tabularx}{\textwidth}{@{}YY@{}}
\toprule
\textbf{特征} & \textbf{重要性} \\
\midrule
是否含硬质相 & 0.34 \\
平均晶粒尺寸 & 0.26 \\
原子半径差 & 0.18 \\
价电子浓度 & 0.13 \\
电负性差 & 0.06 \\
平均电负性 & 0.03 \\
\bottomrule
\end{tabularx}
\end{formalTable}

这个结果说明，在当前数据集和模型中，“是否含硬质相”和“平均晶粒尺寸”对硬度预测影响最大。接下来，研究者需要把这个模型结果转化为材料语言。含硬质相的涂层往往可能通过第二相强化提高硬度；晶粒尺寸越小，晶界数量越多，可能带来细晶强化；原子半径差较大时，晶格畸变增强，也可能产生固溶强化。因此，如果模型认为这些特征重要，并且这种排序与材料强化机制基本一致，就说明模型结果具有一定可解释性。


但是，特征重要性只能告诉研究者“哪些特征重要”，不宜直接告诉研究者“这个特征让硬度升高还是降低”。例如，平均晶粒尺寸重要，并不意味着晶粒尺寸越大硬度越高，也不意味着晶粒尺寸越小一定硬度越高。它只能说明模型在预测时经常用到这个特征。为了进一步分析特征对预测结果的正负影响，可引入 SHAP 方法。


SHAP 可理解为一种更细致的解释工具。它不仅能分析某个特征整体上是否重要，还可分析在某一个样品中，某个特征是把预测硬度往高处推，还是往低处拉。例如，对于一个含 TiC 硬质相、晶粒较细的高熵合金涂层样品，SHAP 分析可能显示：“是否含硬质相”对预测硬度有正贡献，“平均晶粒尺寸较小”也对预测硬度有正贡献。这种结果就可和第二相强化、细晶强化联系起来解释。


为了更清楚地说明可解释性分析，以下给出一个精简代码框架。假设研究者已经在第 \ref{subsec:descriptor-hardness-prediction} 节中训练好了随机森林硬度预测模型 model，并且已经得到输入特征 X\_\allowbreak{}train、X\_\allowbreak{}test 和特征名称 feature\_\allowbreak{}cols。以下代码用于输出随机森林特征重要性，并进一步使用 SHAP 分析模型预测结果。


\begin{lstlisting}[language=Python,escapeinside={(*@}{@*)}]
# (*@\textnormal{1. 随机森林特征重要性}@*)
importance = pd.DataFrame({
    "feature": feature_cols,
    "importance": model.feature_importances_
}).sort_values(by="importance", ascending=False)
print(importance)
# (*@\textnormal{2. 绘制特征重要性柱状图}@*)
plt.figure(figsize=(6, 4))
plt.barh(importance["feature"], importance["importance"])
plt.xlabel("Feature importance")
plt.ylabel("Feature")
plt.title("Feature importance for hardness prediction")
plt.gca().invert_yaxis()
plt.tight_layout()
plt.show()
# (*@\textnormal{3. 使用 SHAP 解释模型}@*)
# (*@\textnormal{需要先安装 shap: pip install shap}@*)
import shap
explainer = shap.TreeExplainer(model)
shap_values = explainer.shap_values(X_test)
# (*@\textnormal{4. 绘制 SHAP 总结图}@*)
shap.summary_plot(
    shap_values,
    X_test,
    feature_names=feature_cols
)
\end{lstlisting}

这段代码分为两个部分。第一部分使用 model.feature\_\allowbreak{}importances\_\allowbreak{} 输出随机森林自身的特征重要性，并画出柱状图。这个图可帮助研究者快速判断模型最关注哪些特征。第二部分使用 SHAP 对模型进行解释，shap.summary\_\allowbreak{}plot() 可显示每个特征对预测结果的整体影响。


可解释性分析中的 AI 辅助，应服务于“把模型输出转化为材料语言”，而不是替代机理判断。提示词中需要同时给出模型结果和要求它说明证据限制。可按下面方式组织问题。

\begin{samepage}
\begin{userbox}
我用随机森林预测高熵合金涂层硬度，模型输出的特征重要性排序为：是否含硬质相 0.34，平均晶粒尺寸 0.26，原子半径差 0.18，价电子浓度 0.13，电负性差 0.06，平均电负性 0.03。SHAP 图显示，含硬质相通常对预测硬度有正贡献，晶粒尺寸较大时常出现负贡献。

请帮我把这些模型解释结果转化为材料语言。

要求如下：
\begin{enumerate}
\item 分别联系第二相强化、细晶强化、固溶强化或相稳定性解释这些特征；
\item 说明这些解释需要哪些实验表征支持，例如 XRD、SEM、EDS\footnote{Energy Dispersive Spectroscopy，能谱分析。} 或硬度测试；
\item 提醒特征重要性和 SHAP 只能说明当前模型中的关系，不宜直接证明因果关系；
\item 如果模型把样品编号、数据来源等无材料意义字段排在前面，也请说明可能存在什么问题。
\end{enumerate}
\end{userbox}

\begin{aibox}
可这样理解这组结果：

\begin{enumerate}
\item “是否含硬质相”重要性最高，说明模型在预测硬度时较多依赖硬质相信息。若涂层中确实存在 TiC、碳化物或硼化物等硬质相，可与第二相强化联系起来解释。
\item “平均晶粒尺寸”重要性较高，且较大晶粒尺寸对应负向 SHAP 贡献，可能说明模型学习到了晶粒细化有利于提高硬度的趋势，可与细晶强化机制联系起来。
\item “原子半径差”有一定重要性，可与多主元合金中的晶格畸变和固溶强化联系起来，但需要结合相结构和成分分布判断。
\item “价电子浓度”可能与 FCC/BCC 相稳定性有关，但它对硬度的影响通常需要通过相结构、组织形貌和其他描述符共同解释。
\end{enumerate}

这些解释仍然只是模型层面的线索。若要形成可靠材料结论，需要用 XRD 确认物相，用 SEM/EDS 观察硬质相和成分分布，用晶粒尺寸统计验证细晶强化趋势，并检查硬度测试条件是否一致。如果无材料意义的字段在模型中排在前面，可能说明数据整理、特征选择或数据划分存在问题。
\end{aibox}
\end{samepage}

模型解释并不等于材料真理。特征重要性高，只说明模型在当前数据集中较多依赖该特征；SHAP 值为正，也只说明在当前模型中该特征对预测值有正向贡献。这些结果可能受到数据量、数据分布、特征相关性和模型类型影响。例如，如果所有含硬质相的样品刚好也具有较小晶粒尺寸，模型可能难以区分到底是硬质相起主要作用，还是晶粒细化起主要作用。因此，可解释性分析必须与实验表征和材料机理结合，不能孤立理解。


对于高熵合金涂层硬度预测，可按照以下思路解读模型解释结果。首先看特征重要性，判断模型主要关注哪些因素。其次看 SHAP 分析，判断这些因素对预测硬度的正负影响。然后把模型结果与材料强化机制对应起来：硬质相对应第二相强化，晶粒尺寸对应细晶强化，原子半径差对应晶格畸变和固溶强化，价电子浓度可能与相结构稳定性有关。最后，再回到实验数据中检查这些解释是否被组织形貌、物相分析和硬度测试支持。


例如，如果模型认为“是否含硬质相”最重要，而 SEM 和 XRD 也显示高硬度样品中确实存在较多 TiC 等硬质相，那么模型解释与实验观察相互支持。相反，如果模型认为某个表面上无关的字段最重要，例如样品编号或数据来源，那么就需要警惕模型可能学到了错误关联。这种情况并不是模型真正理解了材料规律，而可能是数据整理或特征设计存在问题。


除了特征重要性和 SHAP，还有一些更偏研究生层次的方法可用于材料模型解释。例如，符号回归和 SISSO\footnote{Sure Independence Screening and Sparsifying Operator，一种用于筛选描述符并构建稀疏表达式的方法。} 试图从数据中寻找简洁的数学表达式或组合描述符，使模型结果更接近材料专业人员能够理解的公式关系。对于入门阶段的研究者来说，可先掌握特征重要性和 SHAP 的基本思想；对于研究生，可进一步了解这些能够自动发现简洁规律的可解释建模方法。


本节的重点不是让研究者记住某一个解释算法，而是理解材料人工智能中的一个基本原则：模型预测之后，必须回到材料问题本身去解释。一个好的材料机器学习模型，不仅应该预测误差较小，还应该能够帮助研究者理解哪些成分、结构或组织因素可能控制性能。只有当模型结果能够与材料机制相互印证时，它才更可能成为材料研究和材料设计的有效工具。


通过高熵合金涂层硬度预测这个例子可看到，可解释性分析把机器学习结果和材料科学知识连接起来。模型指出哪些特征重要，材料科学知识帮助研究者解释这些特征为什么重要。但最终是否形成可靠规律，仍然需要实验表征、物理机理和专业判断共同验证。


下一节将进一步讨论物理知识嵌入。可解释性更多是在模型训练之后分析“模型为什么这样预测”，而物理知识嵌入则是在建模之前或建模过程中，就主动把材料科学规律加入数据选择、特征设计和模型约束中，使模型尽量避免学习错误关系。


\subsection{物理知识与模型约束}

上一节讨论了模型可解释性，即模型训练完成之后，研究者如何分析模型主要依据哪些特征做出预测，并将这些特征与材料机理联系起来。可解释性强调的是“事后解释”：模型已经训练好了，研究者再分析它为什么这样预测。


但是，在材料科学中，仅仅依靠事后解释还不够。更理想的做法是，在数据整理、特征选择、模型训练和结果判断的过程中，就主动引入材料科学知识，使模型尽量在合理的物理和化学范围内学习规律。这就是物理知识嵌入的基本思想。


物理知识嵌入并不一定是复杂的公式或高级算法。对于入门阶段，可把它理解为：不要把所有能收集到的字段都机械地丢给模型，而是根据材料机理选择有意义的数据和特征；不要让模型学习明显不合理的关系，而是用物理、化学和实验常识约束建模过程。这样得到的模型，通常比完全黑箱的模型更可靠，也更容易被研究者接受。


以前面的高熵合金涂层硬度预测任务为例，研究者已经用原子半径差、平均电负性、电负性差、价电子浓度、是否含硬质相和平均晶粒尺寸等特征预测涂层硬度。接下来需要判断的是：这些特征为什么可作为输入？哪些字段不应该作为输入？数据应该如何筛选，才能让模型学习到更接近材料规律的关系？


首先，物理知识可用于指导特征选择。对于高熵合金涂层硬度来说，硬度提高通常可能来自固溶强化、细晶强化、第二相强化和组织致密化等机制。因此，输入特征应该尽量围绕这些机制设计。


例如，原子半径差可与晶格畸变和固溶强化联系起来。多种元素固溶在同一晶格中时，元素原子半径差异会导致晶格畸变，从而可能阻碍位错运动，提高材料硬度。平均晶粒尺寸可与细晶强化联系起来。晶粒越细，晶界数量越多，位错运动越容易受到阻碍，硬度可能提高。是否含硬质相可与第二相强化联系起来。如果涂层中生成 TiC、碳化物或硼化物等硬质相，这些硬质相可能作为强化相提高涂层硬度。


因此，这些特征不是随便选出来的数字，而是和材料强化机制存在对应关系。模型使用这些特征预测硬度，才有可能学习到具有材料意义的规律。


这种关系可整理为表 \ref{tab:mechanism-feature-mapping}。


\begin{formalTable}[htbp]
\caption{材料机制与候选特征对应关系}
\label{tab:mechanism-feature-mapping}
\begin{tabularx}{\textwidth}{@{}YYY@{}}
\toprule
\textbf{材料机制} & \textbf{可转化的模型特征} & \textbf{对硬度的可能影响} \\
\midrule
固溶强化 & 原子半径差、电负性差、价电子浓度 & 晶格畸变和相稳定性变化可能影响硬度 \\
细晶强化 & 平均晶粒尺寸 & 晶粒细化可能提高硬度 \\
第二相强化 & 是否含硬质相、硬质相含量 & 硬质相可能提高涂层硬度 \\
组织致密化 & 孔隙率、裂纹数量、涂层缺陷 & 孔隙和裂纹可能降低硬度和服役性能 \\
\bottomrule
\end{tabularx}
\end{formalTable}

这个表格体现了物理知识嵌入的第一层含义：把材料机理转化为模型特征。模型输入不是任意字段，而是由材料知识筛选出来的、有明确含义的特征。


在实际建模中，这一步应先由材料机理提出候选特征，再结合已有实验记录判断哪些特征可获得、哪些特征需要补充测试。特征设计不应只看字段是否存在，还要看字段是否能对应到明确的材料机制。


其次，物理知识可帮助排除不合理特征。并不是所有出现在数据表中的字段都适合作为模型输入。例如，样品编号、测试日期、数据来源、实验人员姓名等字段，通常不应该作为材料性能预测特征。它们可能与硬度在数据表中出现某种偶然相关，但这种相关并不具有材料意义。如果模型根据样品编号预测硬度，就说明它可能学到了数据排序或实验批次的偶然关系，而不是材料规律。


前文已经讨论过目标泄漏，这里只强调它与物理知识嵌入的关系：输入字段必须在预测之前可获得，且不能由目标结果推导而来。例如，“硬度等级”“是否高硬度样品”“硬度排序”这类字段不能作为硬度预测输入，因为它们已经包含了目标硬度的信息。研究者需要检查每一个输入字段：它是否具有材料意义？是否直接包含目标结果？如果无法回答这些问题，就不应该轻易放入模型。


第三，物理知识可指导数据筛选和标准化。材料实验数据往往受到测试条件影响。以硬度为例，不同载荷、不同保载时间、不同测试位置可能得到不同硬度值。对于高熵合金涂层，表面硬度、截面硬度和不同层位硬度也可能存在差异。如果把不同测试条件下的硬度值直接混合到一个数据集中，模型可能学到混乱关系。


在建立硬度预测模型时，应该尽量保证硬度数据具有可比性。例如，尽量使用相同测试方法、相近测试载荷、相同单位和相同样品位置的数据。如果不同数据来源不可避免，就应在数据表中记录测试条件，必要时把测试条件作为附加字段，或者只选择条件一致的数据进行建模。


这种数据筛选本身就是物理知识嵌入的一部分。它告诉模型：哪些数据可放在一起学习，哪些数据不宜简单混用。对于材料科学来说，数据标准化不是形式工作，而是保证模型学习到真实规律的前提。


检查数据可比性时，应重点核验测试载荷、测试单位、热处理条件、样品位置和数据来源等字段。最终是否保留某条数据，仍然需要研究者根据实验条件和研究目标判断。


第四，物理知识可作为模型结果的合理性检查。模型训练完成后，如果预测结果明显违背材料常识，就需要警惕。例如，如果模型预测孔隙率越高硬度越高，或者晶粒尺寸越大硬度越高，就需要进一步检查数据是否存在偏差、特征是否相关、样本数量是否太少，或者模型是否受异常样品影响。当然，材料规律并不总是简单单调关系，但当模型结果与常识明显冲突时，必须进行核查，而不是直接接受。


对于高熵合金涂层硬度预测，可设定一些基本的合理性检查。例如，含硬质相的样品不一定总是硬度最高，但如果模型完全认为硬质相无关，就需要检查硬质相字段是否准确；晶粒细化通常有利于硬度提高，如果模型得出相反趋势，也要检查晶粒尺寸统计是否一致；孔隙率增加通常不利于力学性能，如果模型认为孔隙率越高越好，就需要检查数据是否存在偶然相关或测量问题。


第五，物理知识还可用于更高级的模型约束。在入门阶段，研究者主要通过特征选择、数据筛选和结果检查嵌入材料知识。对于更深入的研究，还可把物理规律直接写进模型或损失函数中。例如，在原子尺度模型中，可要求模型满足平移、旋转和置换对称性；在热力学相关任务中，可引入稳定性约束；在力学模型中，可要求某些预测量满足基本物理关系。这类方法通常称为物理约束机器学习或物理知识嵌入机器学习。


对于本章而言，不需要展开复杂数学形式，但要让研究者理解一个基本思想：材料模型不应只追求拟合数据，还要尽量尊重材料科学规律。一个模型如果预测误差很小，但违背基本物理常识，或者依赖无法解释的虚假相关字段，那么它在材料研究中的价值是有限的。


可用如下流程概括物理知识嵌入在高熵合金涂层硬度预测中的作用：


\begin{lightquote}
\centering 研究目标 → 材料机理分析 → 特征选择 → 数据筛选 → 模型训练 → 结果解释 → 实验和机理验证
\end{lightquote}

其中，材料机理分析不是最后才出现，而是从一开始就参与建模。研究者首先根据固溶强化、细晶强化、第二相强化等机制选择特征；再根据测试条件筛选可比数据；然后训练模型；最后再用实验表征和材料机理检查模型解释是否合理。


为了帮助研究者理解这个流程，可设计一个简单的操作任务。假设已经有一张高熵合金涂层硬度数据表，其中包含以下字段：样品编号、化学成分、原子半径差、电负性差、价电子浓度、是否含硬质相、晶粒尺寸、孔隙率、硬度等级、硬度/HV、测试载荷、数据来源。研究者需要判断哪些字段可作为输入特征，哪些字段应该保留但不作为模型输入，哪些字段必须删除或谨慎使用。


按照物理知识嵌入的原则，可这样处理：原子半径差、电负性差、价电子浓度、是否含硬质相、晶粒尺寸和孔隙率可作为候选输入特征，因为它们具有材料意义；测试载荷和数据来源应保留，用于判断数据是否可比，但不一定直接作为性能特征；样品编号通常不应作为输入；硬度等级不能作为输入，因为它可能由硬度结果推导而来，存在目标泄漏风险；硬度/HV 是预测目标 y，不应同时作为输入特征 X。


可整理为表 \ref{tab:physics-informed-field-treatment}。


\begin{formalTable}[htbp]
\caption{物理知识嵌入建模过程中的字段处理}
\label{tab:physics-informed-field-treatment}
\begin{tabularx}{\textwidth}{@{}YYY@{}}
\toprule
\textbf{字段} & \textbf{是否作为模型输入} & \textbf{原因} \\
\midrule
原子半径差 & 可 & 与晶格畸变和固溶强化有关 \\
电负性差 & 可 & 与元素化学差异和相形成有关 \\
价电子浓度 & 可 & 可能影响相结构稳定性 \\
是否含硬质相 & 可 & 对应第二相强化 \\
晶粒尺寸 & 可 & 对应细晶强化 \\
孔隙率 & 可 & 反映组织致密性 \\
测试载荷 & 谨慎使用或用于筛选 & 影响硬度可比性 \\
数据来源 & 不作为材料特征 & 用于追溯和分组检查 \\
样品编号 & 不使用 & 无材料意义 \\
硬度等级 & 不使用 & 可能造成目标泄漏 \\
硬度/HV & 目标值 & 模型要预测的 y \\
\bottomrule
\end{tabularx}
\end{formalTable}

这类字段筛选任务适合用提示词形式进行辅助检查。提示词应要求同时判断材料意义、数据可比性和目标泄漏风险，而不是只问“哪些列可输入模型”。

\begin{samepage}
\begin{userbox}
我有一张高熵合金涂层硬度数据表，字段包括：样品编号、化学成分、原子半径差、电负性差、价电子浓度、是否含硬质相、晶粒尺寸、孔隙率、硬度等级、硬度/HV、测试载荷、数据来源。请帮我判断哪些字段适合作为模型输入，哪些字段应保留但不作为输入，哪些字段可能造成目标泄漏或虚假相关。

要求如下：
\begin{enumerate}
\item 按固溶强化、细晶强化、第二相强化和组织致密化解释可用特征；
\item 指出哪些字段用于追溯或筛选数据，但不宜作为材料性能特征；
\item 明确指出硬度/HV 和硬度等级为什么不能作为输入特征；
\item 提醒哪些测试条件会影响硬度数据可比性；
\item 不要只给机器学习角度判断，要结合材料机理说明理由。
\end{enumerate}
\end{userbox}

\begin{aibox}
可按三类处理这些字段：

\begin{enumerate}
\item 可作为候选输入特征的字段：原子半径差、电负性差、价电子浓度、是否含硬质相、晶粒尺寸、孔隙率。这些字段分别可与固溶强化、相结构稳定性、第二相强化、细晶强化和组织致密化联系起来。
\item 应保留但不宜直接作为材料性能特征的字段：测试载荷、数据来源、化学成分和样品编号。测试载荷和数据来源主要用于判断数据可比性和可追溯性；样品编号通常没有材料意义；化学成分可保留，但进入模型前通常需要转化为成分描述符。
\item 不能作为输入特征的字段：硬度/HV 是预测目标，硬度等级往往由硬度结果划分而来，二者如果进入输入特征，会造成目标泄漏。
\end{enumerate}

此外，硬度数据是否可比还取决于测试载荷、保载时间、测试位置、样品表面状态和测试方法。如果这些条件不一致，应先进行筛选、分组或记录为元数据，再决定是否合并建模。
\end{aibox}
\end{samepage}

需要强调的是，物理知识嵌入并不意味着模型一定更复杂。很多时候，它体现为更谨慎的数据选择、更合理的特征设计和更严格的结果检查。对于材料人工智能入门阶段的研究者来说，最重要的不是马上构建复杂的物理约束模型，而是养成一种意识：建模前先问材料问题，建模中保留物理意义，建模后回到实验和机理验证。


本节说明了物理知识嵌入的基本思想。可解释性帮助研究者在模型训练后理解模型，物理知识嵌入则帮助研究者在建模前和建模过程中约束模型。二者共同说明：人工智能用于材料科学，并不是把数据机械地交给模型，而是让数据驱动方法与材料科学知识相互配合。至此，第\ref{chap:materials-cognition-prediction}章完成了从材料数据、材料表示、性能预测、原子尺度模拟到模型解释的完整链条。下一章将进一步讨论逆向设计问题，即如何根据目标性能反向寻找和设计新材料。

\subsubsection*{\textbf{小结}}

本节说明模型解释和物理知识嵌入共同服务于材料判断。
特征重要性和 SHAP 可提示模型依据，但不能直接证明因果机理，仍需 XRD、SEM、EDS 等证据支持。
物理知识应贯穿数据筛选、特征选择、模型训练和结果检查，避免模型学习虚假相关。

\subsection*{习题}

\begin{enumerate}[leftmargin=*]
\item SHAP 为什么不能直接证明材料机理？还需哪些表征？
\item 可解释性分析如何辅助材料筛选？请举例说明。
\item 只有 FCC 数据时，能否预测 BCC 合金？请说明风险和验证方案。
\end{enumerate}

\section*{本章小结}

本章围绕材料认知与性能预测展开，形成了从数据到模型、再从模型回到材料问题的基本链条。材料数据整理强调字段完整、来源可追溯和条件可比较；材料表示强调把化学式、结构和局域环境转化为描述符或图结构；性能预测部分说明了描述符模型、图神经网络和机器学习势在不同尺度任务中的作用；模型解释与物理知识嵌入则进一步说明，模型预测结果必须经过材料机理、实验表征和物理约束的共同检验。由此可见，人工智能在材料科学中的价值不在于替代实验和理论，而在于提高数据利用效率、缩小候选范围并辅助形成可验证的材料判断。

\chapter{材料设计与智能发现}
\label{chap:materials-design-discovery}


第\ref{chap:materials-cognition-prediction}章讨论正向预测，即在材料成分、结构或实验条件已知的情况下，预测材料可能具有的性能。本章进一步讨论逆向设计与智能发现：当目标性能已经给定时，如何从大量候选材料中筛选出值得验证的方案。这个过程并不是简单地让模型“生成答案”，而是要把目标性能、稳定性、可合成性、成本、安全性和实验条件共同纳入考虑。

本章按照材料发现的一般流程展开，重点解决四个问题。第一节说明如何把应用需求转化为材料设计目标、候选变量和约束边界。第二节讨论候选筛选和主动学习如何帮助研究者逐层缩小搜索空间，并安排下一轮实验或计算。第三节介绍生成设计如何提出库外候选，以及生成候选为什么必须经过有效性、稳定性、新颖性和可合成性验证。第四节讨论可信发现、方法选择和闭环迭代，强调模型推荐必须与计算验证、实验反馈和应用约束共同使用。通过本章学习，研究者应理解人工智能如何参与材料设计决策，同时认识候选材料最终仍需要通过物理判断、计算验证和实验结果逐步确认。




\clearpage
\section{问题转换与设计边界}


\begin{sectionkeypoints}
\item 正向预测回答“已知材料可能具有什么性质”，逆向设计回答“为了获得目标性质应当寻找什么材料”。
\item 逆向设计本质上是受约束搜索，不是单纯追求最高预测分数，而是在稳定性、可合成性、成本、安全和设备条件内寻找可验证候选。
\item 材料设计对象可能位于成分、原子结构、微结构、工艺和服役环境等不同层级，层级选择必须与研究目标一致。
\item 候选筛选、主动学习和生成设计是本章三条核心路线，实际研究中常常组合使用。
\end{sectionkeypoints}

\subsection{正向预测与逆向设计}

第\ref{chap:materials-cognition-prediction}章讨论的是正向问题：给定材料的成分、结构、工艺或表征信息，预测形成能、带隙、硬度、折射率或熔点等性质。第\ref{chap:materials-design-discovery}章把问题方向反过来：当目标性质已经确定时，寻找哪些材料、配方、结构或工艺值得进一步尝试。前者回答“这种材料具有什么性质”，后者回答“具有目标性质的材料可能在哪里”。


材料科学长期依赖经验规则和实验试错，但多组元合金、复杂晶体和玻璃配方的候选数量远超实验能够逐一验证的范围。逆向设计的作用不是取消经验和实验，而是把大规模穷举转化为有方向的搜索：先明确目标和约束，再用模型提出或排序候选，最后通过计算和实验逐步确认。


从数学上看，逆向设计可理解为受约束优化。材料对象 x 可以是化学式、晶体结构、玻璃配方、微结构参数或制备条件；正向模型 f(x) 估计目标性能；约束 g(x) 描述稳定性、可合成性、成本、安全和设备边界。搜索的目的不是让某个预测分数无限增大，而是在满足约束的前提下找到一组可验证的候选。正向预测、逆向设计与闭环发现之间的关系见表 \ref{tab:forward-inverse-closed-loop}。


\begin{formalTable}[htbp]
\caption{正向预测、逆向设计与闭环发现}
\label{tab:forward-inverse-closed-loop}
\footnotesize
\begin{tabularx}{\textwidth}{@{}YYYYY@{}}
\toprule
\textbf{问题类型} & \textbf{输入} & \textbf{输出} & \textbf{材料例子} & \textbf{主要风险} \\
\midrule
正向预测 & 候选材料 x & 性质 f(x) & 由晶体结构预测带隙 & 训练集外误差被低估 \\
逆向设计 & 目标性质 y & 候选材料 x & 按目标折射率设计玻璃 & 候选不可合成或不稳定 \\
闭环发现 & 已有数据 + 目标 & 下一步实验 & 主动学习选择合金配比 & 实验噪声导致错误反馈 \\
\bottomrule
\end{tabularx}
\end{formalTable}

\subsubsection*{\textbf{1. 问题转换：从应用目标到材料任务}}

逆向设计的第一步，不是直接选择模型，而是把应用语言转化为材料科学语言。应用目标往往比较笼统，例如“希望玻璃折射率更高”“希望发光材料颜色更接近蓝光”“希望涂层在高温下更稳定”。这些表述还不宜直接进入模型，因为模型需要明确的输入变量、目标性能和约束条件。只有把应用目标拆解成可计算、可测量、可验证的材料任务，后续的筛选、优化和生成才有依据。

以 AR\footnote{Augmented Reality，增强现实。} 光波导玻璃为例，“高折射率”只是一个目标方向。真正进入设计流程时，还需要进一步明确折射率测量波长、透过率范围、色散要求、玻璃化转变温度（Tg\footnote{Glass transition temperature，玻璃化转变温度。}）、热膨胀系数、密度、原料安全性和熔制温度窗口。如果只追求折射率，模型可能推荐含高极化率重金属氧化物较多的配方，但这类配方可能带来密度升高、成本增加、透过率下降或环保风险。因此，逆向设计不是单目标求最大值，而是多目标、多约束条件下的候选选择。

其他材料任务也需要类似拆解。例如 microLED\footnote{Micro Light-Emitting Diode，微型发光二极管。} 发光材料可把目标颜色转化为发光波长或带隙范围，但仍需补充缺陷态、载流子复合效率、热稳定性和外延匹配等约束。这里不展开为新的主线案例，只用它说明：目标性能通常只是材料设计任务的一部分。

具体而言，问题转换通常包括四个问题：目标性能是什么？候选变量是什么？约束条件是什么？验证方式是什么？目标性能决定模型要预测什么；候选变量决定模型可搜索哪些材料空间；约束条件决定哪些候选必须排除；验证方式决定候选是否能从模型结果走向真实材料。表中的 DFT\footnote{Density Functional Theory，密度泛函理论。} 指密度泛函理论计算。典型应用目标到材料设计任务的转换见表 \ref{tab:application-to-design-task}。

\begin{formalTable}[htbp]
\caption{应用目标到材料设计任务的转换}
\label{tab:application-to-design-task}
\footnotesize
\begin{tabularx}{\textwidth}{@{}YYYYY@{}}
\toprule
\textbf{应用目标} & \textbf{材料任务} & \textbf{候选变量} & \textbf{关键约束} & \textbf{验证方式} \\
\midrule
高折射率玻璃 & 预测并优化折射率、Tg、热膨胀系数 & 氧化物组成比例 & 玻璃形成、透明性、熔制温度 & 熔制实验、折射率和透过率测试 \\
蓝光发光材料 & 筛选合适带隙和稳定结构 & 元素组成、晶体结构、掺杂 & 缺陷、外延兼容、热稳定 & DFT、发光测试、器件验证 \\
高硬度涂层 & 优化成分、硬质相和组织参数 & 合金成分、工艺、相组成 & 裂纹、孔隙、结合强度 & 显微硬度、SEM/XRD、服役测试 \\
超高温陶瓷 & 寻找高熔点和抗氧化候选 & 化学组成、晶体结构 & 高温稳定、氧化行为、可烧结性 & 高温实验、热分析、模拟验证 \\
\bottomrule
\end{tabularx}
\end{formalTable}

这一转换过程适合借助 AI 辅助梳理，但不能让 AI 直接给出最终材料结论。提示词应要求它把应用目标拆成材料变量、目标性能、约束条件和验证方式，而不是只列出几个候选材料名称。

\begin{samepage}
\begin{userbox}
我想设计一种用于 AR 光波导的高折射率玻璃。请帮我把这个应用目标转化为材料逆向设计任务。

要求如下：
\begin{enumerate}
\item 明确目标性能，例如折射率、透过率、色散、Tg、热膨胀系数等；
\item 说明候选变量可包括哪些氧化物组成或配方参数；
\item 区分必须满足的硬约束和可折中的软目标；
\item 给出候选筛选后需要进行哪些计算或实验验证；
\item 不要直接编造最终玻璃配方，只做任务拆解。
\end{enumerate}
\end{userbox}

\begin{aibox}
可将“高折射率 AR 光波导玻璃”拆解为以下设计任务：

\begin{enumerate}
\item \textbf{目标性能：}折射率是核心目标，但还应同时记录透过率、色散、Tg、热膨胀系数、密度和加工温度等指标。
\item \textbf{候选变量：}可把各氧化物组分比例作为设计变量，例如网络形成体、网络修饰体和提高折射率的氧化物比例。
\item \textbf{硬约束：}配方比例之和应满足组成约束；原料不能超出安全和环保边界；熔制温度应在设备可承受范围内；样品应能够形成玻璃而不是明显析晶。
\item \textbf{软目标：}折射率越高越好，但不能以明显牺牲透明性、热稳定性和加工性为代价。
\item 验证方式：模型筛选后，需要进行玻璃形成能力判断、熔制实验、折射率测试、透过率测试和热膨胀系数测试。
\end{enumerate}

这样处理后，逆向设计任务不再只是“寻找高折射率材料”，而是“在可熔制、可成玻、透明且热性质合适的约束下，寻找高折射率玻璃配方”。
\end{aibox}
\end{samepage}

问题转换应当在建模之前完成，而不是在模型给出候选之后再补充解释。如果目标一开始没有定义清楚，后面所有筛选和排序都会变得模糊。例如，若只写“设计高性能玻璃”，模型无法判断“高性能”指的是高折射率、高透过率、低热膨胀系数还是高热稳定性；若只写“设计蓝光材料”，模型也不知道重点是发光波长、发光效率、热稳定性还是外延兼容性。本章强调问题转换，是为了让研究者意识到：人工智能材料设计首先是材料问题定义，其次才是模型选择。

一个规范的逆向设计任务，宜用一句话概括为“在什么约束下，寻找什么类型的材料，使什么目标性能达到什么范围”。例如，“在不使用高风险元素、熔制温度不超过设备上限、能够形成透明玻璃的约束下，寻找折射率较高且热膨胀系数合适的氧化物玻璃配方”。这样的表述比“让模型设计高折射率玻璃”更清楚，也更便于后续建立数据表、选择描述符和设计验证实验。

\subsection{设计对象与约束边界}

\subsubsection*{\textbf{1. 设计边界：材料对象分层}}

逆向设计并不总是直接生成一个化学式。材料对象通常包含多个层级：成分决定由哪些元素或组分构成；原子结构决定空间群、占位、缺陷和局域配位；微结构包含晶粒、孔隙、第二相和界面；工艺包含熔融、退火、沉积、烧结和冷却条件。同一成分经过不同工艺可得到不同结构和性能，因此设计对象的层级必须与研究目标相匹配。材料逆向设计常见对象层级见表 \ref{tab:inverse-design-object-levels}。


\begin{formalTable}[htbp]
\caption{材料逆向设计的对象层级}
\label{tab:inverse-design-object-levels}
\begin{tabularx}{\textwidth}{@{}YYYY@{}}
\toprule
\textbf{设计层级} & \textbf{常见变量} & \textbf{适用任务} & \textbf{容易忽略的问题} \\
\midrule
成分 & 元素/氧化物比例、掺杂量 & 玻璃配方、合金筛选、催化剂组分 & 相形成与工艺窗口 \\
原子结构 & 晶格、坐标、缺陷、表面 & 带隙、离子迁移、稳定性 & 对称性和局域环境 \\
微结构 & 晶粒、孔隙、第二相、取向 & 力学、热学、疲劳与可靠性 & 制备历史强相关 \\
工艺 & 温度、气氛、压力、时间 & 可合成性与性能优化 & 元数据缺失会误导模型 \\
服役环境 & 光、电、热、湿度、载荷 & 器件可靠性和寿命 & 实验条件与应用条件不一致 \\
\bottomrule
\end{tabularx}
\end{formalTable}

除了设计层级，还需要区分约束的类型。材料设计中的约束可分为硬约束和软约束。硬约束是必须满足的条件，例如有毒元素不能使用、设备最高温度不能超过、配方比例必须归一化、候选结构不能明显违反价态和原子距离规则。软约束是希望尽量优化但可折中的条件，例如成本越低越好、密度越低越好、热膨胀系数越接近目标越好。硬约束用于排除候选，软约束用于排序候选。

如果不区分硬约束和软约束，逆向设计很容易变成一个不清楚的打分问题。例如，把折射率、成本、热稳定性和环保要求都压缩成一个总分，表面上方便排序，但不同目标之间的取舍会被隐藏。逆向设计应把硬约束、软目标和数据边界分开说明。

设计边界还包括数据边界。模型只能在训练数据覆盖的范围内比较可靠地工作。如果训练数据主要来自硅酸盐玻璃，那么把模型直接用于完全不同的硫系玻璃或氟化物玻璃，预测结果就可能不稳定。如果模型只学习了室温性质，就不宜直接推断高温服役行为。如果数据没有记录制备工艺，模型也很难判断同一成分在不同热处理条件下的性能差异。因此，逆向设计开始之前必须说明数据从哪里来、覆盖哪些材料体系、哪些变量没有记录。不同边界类型及其处理方式见表 \ref{tab:inverse-design-boundary-types}。

\begin{formalTable}[htbp]
\caption{材料逆向设计中的边界类型}
\label{tab:inverse-design-boundary-types}
\footnotesize
\begin{tabularx}{\textwidth}{@{}YYYY@{}}
\toprule
\textbf{边界类型} & \textbf{含义} & \textbf{材料例子} & \textbf{处理方式} \\
\midrule
硬约束 & 不满足就直接排除 & 有毒元素、不可成玻、设备温度超限 & 筛选前写入规则 \\
软目标 & 可折中但要排序 & 折射率、密度、成本、热膨胀系数 & 作为多目标评价项 \\
数据边界 & 模型训练数据覆盖范围 & 只含某类玻璃或某类晶体 & 明确适用范围和外推风险 \\
工艺边界 & 实际制备条件限制 & 熔制温度、气氛、烧结时间 & 作为实验可行性约束 \\
\bottomrule
\end{tabularx}
\end{formalTable}

在教学中，可让研究者首先写出设计边界，再选择模型。这样的顺序比“先选一个算法再想怎么用”更合理。算法只能在已经定义清楚的问题上发挥作用；如果目标、变量和边界本身模糊，模型输出再复杂也难以转化为真实材料方案。

这一类边界梳理同样可使用 AI 辅助，但提示词必须要求区分硬约束、软目标和数据边界。

\begin{samepage}
\begin{userbox}
我准备做一个材料逆向设计任务，目标是寻找候选高折射率玻璃。请帮我列出这个任务的设计边界。

要求如下：
\begin{enumerate}
\item 区分硬约束、软目标、数据边界和工艺边界；
\item 每一类边界都要说明它为什么重要；
\item 不要把所有条件合并成一个总分；
\item 说明哪些信息需要在实验记录或数据库中保留；
\item 不要直接给出候选配方。
\end{enumerate}
\end{userbox}

\begin{aibox}
可从四类边界整理：

\begin{enumerate}
\item \textbf{硬约束：}原料安全、环保限制、组成比例合法、可形成玻璃、设备温度可达。这类条件不满足时，候选应直接排除。
\item \textbf{软目标：}折射率、透过率、密度、热膨胀系数、成本和加工性能。这些目标之间可能存在折中，应分别记录和排序。
\item \textbf{数据边界：}需要说明训练数据主要来自哪些玻璃体系、性质是在什么条件下测得、是否包含失败样品和析晶样品。超出数据覆盖范围时，模型预测应谨慎使用。
\item \textbf{工艺边界：}熔制温度、保温时间、冷却方式、气氛、坩埚材料等都会影响玻璃形成和性能，应作为实验元数据保存。
\end{enumerate}

这些边界不是附加说明，而是逆向设计任务的一部分。只有先限定边界，模型推荐的候选才有可能进入真实实验。
\end{aibox}
\end{samepage}

边界条件还应随研究阶段动态更新。早期阶段可设置较宽的组成范围，用来观察模型是否能发现潜在趋势；进入实验阶段后，边界应逐渐收紧，例如限定可采购原料、设备温度、样品尺寸和测试条件；进入应用阶段后，还要加入成本、安全、长期稳定性和批量制备要求。换言之，边界不是一次写完就不变，而是随着数据积累和实验反馈不断修正。

对于研究者来说，最容易忽略的是“不可见边界”。有些限制并不直接写在数据表中，例如某些配方虽然数据库中有记录，但实验室没有相应原料；某些结构计算上稳定，但现有设备无法提供所需压力或气氛；某些材料性能很好，但含有不适合应用场景的元素。这些不可见边界如果不提前说明，模型推荐就会表面上合理、实际不可执行。因此，材料逆向设计的边界既包括数据中的字段，也包括实验室条件和应用场景中的真实限制。

本章讨论三条核心路线。候选筛选适合从已有数据库或枚举空间中逐层过滤；主动学习适合实验昂贵、每轮只能验证少量样品的任务；生成设计用于突破已有候选库，在目标条件下提出新结构或新配方。三者并非互相替代，实际研究常按“筛选--优化--生成--验证”的方式组合使用。

\subsubsection*{\textbf{小结}}

本节说明逆向设计不是直接生成材料，而是把应用需求转化为目标、变量和约束。
正向预测回答已知材料的性能，逆向设计回答为了目标性能应寻找哪些候选。
材料对象层级、硬约束、软目标、数据边界和工艺边界需要在建模前明确。

\subsection*{习题}

\begin{enumerate}[leftmargin=*]
\item 正向预测和逆向设计在高折射率玻璃任务中如何配合？
\item 逆向设计为什么必须先明确目标和约束？
\item 如何把高折射率玻璃需求转化为计算目标？
\end{enumerate}

\clearpage
\section{候选筛选与主动学习}


\begin{sectionkeypoints}
\item 候选筛选的基本逻辑是先建立候选库，再按照成本由低到高、可信度由粗到细逐层过滤。
\item 筛选漏斗通常包含规则过滤、机器学习代理模型、高精度计算和实验验证等层级，每一级都应记录通过和排除理由。
\item 主动学习适合实验昂贵、小样本和需要多轮优化的材料任务，核心是同时考虑预测性能和模型不确定性。
\item 实验噪声、批次效应和失败样品都应进入数据记录，否则模型可能把制备差异误认为材料本征规律。
\item 候选排序和主动学习评分只负责辅助决策，最终实验批次仍需要材料约束、工艺可行性和人工复核。
\end{sectionkeypoints}

\subsection{筛选漏斗与分层验证}

高通量虚拟筛选的基本思路，是先建立候选库，再按照成本由低到高、可信度由粗到细逐层过滤。最外层使用元素排除、价态平衡、成分范围、成本或毒性等规则；中间层使用第\ref{chap:materials-cognition-prediction}章建立的正向模型快速估计形成能、带隙、弹性模量或折射率；内层使用 DFT、分子动力学或热力学计算精算；最后才进入实验合成与表征。


\begin{figure}[H]
\centering
\BookFigureImage{images_11/3.png}
\caption{候选筛选与验证层级示意图}
\label{fig:candidate-screening-validation}
\end{figure}
\FloatBarrier

如图 \ref{fig:candidate-screening-validation} 所示，筛选漏斗利用了不同方法之间的成本差异。规则过滤几乎不需要计算，机器学习预测可快速处理大量候选，而高精度计算和真实实验只能用于较小范围。把便宜但粗略的方法放在前面，把昂贵但可靠的方法放在后面，能够避免把有限资源消耗在明显不合理的候选上。


筛选的第一个常见门槛是稳定性。无机晶体可参考形成能和凸包能量，合金和陶瓷还要考虑相分离、氧化和高温相变，玻璃则需要考虑玻璃形成能力、析晶倾向和熔制窗口。计算稳定不等于实验一定能够合成，热力学判断还要与动力学路径、前驱体和设备条件结合。


第二个门槛是多目标要求。AR 光波导玻璃不仅需要较高折射率，还要兼顾透过率、色散、玻璃化转变温度、热膨胀系数、密度和加工性能；microLED 发光材料除带隙外，还要考虑缺陷、载流子复合、热稳定和外延兼容性。相应地，候选排序不应只依据单一预测值，而应保留不同目标之间的折中关系。筛选漏斗中不同层级的作用见表 \ref{tab:screening-funnel-levels}。


\begin{formalTable}[htbp]
\caption{筛选漏斗中不同层级的作用}
\label{tab:screening-funnel-levels}
\begin{tabularx}{\textwidth}{@{}YYYY@{}}
\toprule
\textbf{筛选层级} & \textbf{常见工具} & \textbf{适合回答的问题} & \textbf{学习要点} \\
\midrule
规则层 & 元素价态、毒性、成本、工艺窗口 & 哪些候选明显不可取？ & 规则便宜但粗糙，适合第一轮排除 \\
代理模型层 & 回归/分类模型、图神经网络、集成模型 & 哪些候选可能性能较好？ & 模型用于排序，不等于最终证明 \\
物理计算层 & DFT、相图、分子动力学、热力学模型 & 候选是否稳定、机制是否可信？ & 精度更高但成本也更高 \\
实验层 & 合成、表征、服役测试 & 候选能否成为真实材料？ & 实验验证是材料发现的终点而非附属步骤 \\
\bottomrule
\end{tabularx}
\end{formalTable}

\subsubsection*{\textbf{1. 案例示范：高折射率玻璃候选筛选}}

以高折射率玻璃为例，首先把应用要求转化为可计算指标；然后限定各氧化物比例之和、原料安全和设备温度等边界；代理模型预测折射率、Tg、热膨胀系数和析晶倾向；物理或经验模型检查玻璃形成与黏度窗口；最后制备少量代表性配方并测量真实光学损耗。每一级都应记录排除理由，使候选清单可复核。

具体而言，第一层规则可先排除明显不合理的配方。例如，各氧化物比例之和必须为 100\%，单一组分含量不能超过玻璃形成的经验范围，含有明显不适合教学或工程应用的高风险元素时应直接排除。第二层可用代理模型预测折射率和热性质，把候选从几千个缩小到几十个。第三层再用更可靠的热力学判断或经验模型检查玻璃形成能力、析晶风险和黏度窗口。最后才选择少量配方进入熔制实验。

筛选漏斗的一个重要原则是“保留失败理由”。如果一个候选被排除，应记录它是因为折射率不够、Tg 太低、热膨胀系数过高、含有风险元素，还是因为预测不确定性太大。这样做有两个好处。第一，后续可复核筛选规则是否过于严格；第二，失败信息可反馈给下一轮模型或实验设计。如果只保留最后通过的候选，而删除中间淘汰信息，材料发现过程就会失去可追溯性。

在初学阶段，研究者可将筛选漏斗理解为一个逐层提问的过程：第一问，这个候选是否明显违反基本规则？第二问，它的目标性能是否接近要求？第三问，它是否可能稳定存在或可制备？第四问，它是否值得消耗实验资源？每一层都不是绝对正确的判断，而是在不同成本下尽量减少无效尝试。筛选过程中需要保留的判断信息见表 \ref{tab:screening-funnel-records}。

\begin{formalTable}[htbp]
\caption{筛选漏斗中需要保留的判断信息}
\label{tab:screening-funnel-records}
\footnotesize
\begin{tabularx}{\textwidth}{@{}YYYY@{}}
\toprule
\textbf{筛选问题} & \textbf{排除依据} & \textbf{保留信息} & \textbf{后续用途} \\
\midrule
是否违反基本规则 & 比例错误、风险元素、明显价态异常 & 排除原因和原始配方 & 修正规则或候选生成空间 \\
性能是否接近目标 & 折射率、带隙、硬度等预测值偏离 & 预测值和模型版本 & 比较不同代理模型 \\
是否可能稳定 & 凸包能、析晶倾向、相分离风险 & 稳定性指标和计算方法 & 决定是否进入高精度计算 \\
是否值得实验 & 成本、设备温度、样品数量限制 & 排序分数和实验批次 & 形成可执行实验计划 \\
\bottomrule
\end{tabularx}
\end{formalTable}

这一过程可借助 AI 辅助把筛选规则写成清单，但规则本身必须由材料任务决定。可用如下提示词把流程规范化。

\begin{samepage}
\begin{userbox}
我正在设计一个高折射率玻璃候选筛选漏斗。候选数据表中有氧化物比例、预测折射率、预测 Tg、热膨胀系数、是否含风险元素、预测析晶风险等字段。请帮我把筛选流程整理成逐层规则。

要求如下：
\begin{enumerate}
\item 先列出规则过滤层、代理模型层、物理/经验验证层和实验层；
\item 每一层说明主要排除什么候选；
\item 明确哪些排除理由需要记录；
\item 提醒候选排序不应只看折射率；
\item 不要编造真实配方数据。
\end{enumerate}
\end{userbox}

\begin{aibox}
可将筛选漏斗设计为四层：

\begin{enumerate}
\item \textbf{规则过滤层：}检查氧化物比例是否合理、是否含风险元素、熔制温度是否超出设备范围。被排除的候选应记录具体违反哪条规则。
\item \textbf{代理模型层：}预测折射率、Tg 和热膨胀系数。折射率高但 Tg 过低或热膨胀系数过高的候选不宜直接排在最前。
\item \textbf{物理/经验验证层：}进一步检查玻璃形成能力、析晶倾向、黏度窗口和配方稳定性。此处的目标是排除“预测性能好但难以形成玻璃”的候选。
\item \textbf{实验层：}选择少量代表性配方进行熔制、退火、折射率和透过率测试，并记录失败样品原因。
\end{enumerate}

筛选报告中应保留每一级候选数量、通过率和排除理由。这样可判断筛选规则是否过严，也便于后续主动学习利用失败信息。
\end{aibox}
\end{samepage}

筛选漏斗还有一个容易被忽视的问题：前一级筛选规则会影响后一级能看到什么候选。如果第一层规则过于严格，可能会把一些有潜力但不符合传统经验的候选提前排除；如果第一层规则过于宽松，又会把大量明显不可行的候选留给后续高成本计算。为此，筛选规则需要根据任务目标反复校准，而不是一次确定后不再修改。

在课堂练习中，可让研究者比较两种筛选策略。一种策略偏保守，优先排除风险候选，适合实验资源很有限的任务；另一种策略偏探索，保留更多边界候选，适合希望发现新体系的任务。两种策略没有绝对优劣，关键在于是否说明了筛选目的、资源条件和验证能力。这样可避免把筛选漏斗理解为机械流程，而是理解为材料决策工具。


以下给出一个最小候选排序代码示例。代码把目标带隙、凸包能量和元素风险写成显式评分规则，用于说明筛选假设如何转化为可复核的计算步骤。评分只负责初步排序，不宜代替稳定性复算和实验验证。


\begin{lstlisting}[language=Python,escapeinside={(*@}{@*)}]
import pandas as pd
# (*@\textnormal{gap 的单位为 eV，hull 的单位为 eV/atom}@*)
cands = pd.DataFrame({
    "id": ["mp-1", "mp-2", "mp-3", "mp-4"],
    "formula": ["M1O2", "M2O3", "M3S", "PbM4O"],
    "gap": [2.10, 1.55, 2.35, 2.00],
    "hull": [0.018, 0.006, 0.065, 0.012],
    "elements": ["M1,O", "M2,O", "M3,S", "Pb,M4,O"],
})
target_gap = 2.0
cands["gap_score"] = 1 - (cands["gap"] - target_gap).abs() / target_gap
cands["stability_score"] = (0.08 - cands["hull"]).clip(lower=0) / 0.08
cands["element_penalty"] = cands["elements"].str.contains("Pb|Cd|Hg").astype(int)
cands["priority"] = (
    0.45 * cands["gap_score"]
    + 0.45 * cands["stability_score"]
    - 0.50 * cands["element_penalty"]
)
cols = ["id", "formula", "gap", "hull", "priority"]
print(cands.sort_values("priority", ascending=False)[cols])
\end{lstlisting}

\manualfigcaption{\textbf{代码示例：}候选材料初步排序示例}

\subsection{主动学习与实验决策}

高通量筛选是在固定候选库中做过滤，主动学习则是在搜索过程中不断更新模型。研究者首先用少量实验或计算训练代理模型，模型同时给出候选性能和不确定性，再选择下一批最值得验证的候选；新结果返回后重新训练模型，形成序贯迭代。


\begin{figure}[H]
\centering
\BookFigureImage{images_11/4.png}
\caption{主动学习与生成设计闭环示意图}
\label{fig:active-learning-generative-loop}
\end{figure}
\FloatBarrier

图 \ref{fig:active-learning-generative-loop} 概括了主动学习与生成设计之间的闭环关系。贝叶斯优化是常见的主动学习方法，适合小样本和高成本实验。它通常包含代理模型和采集函数两个部分：代理模型近似真实实验或高精度计算，采集函数综合预测均值与不确定性，决定下一步优先验证哪些候选。其核心不是总选择当前分数最高的样品，而是在“利用”和“探索”之间取得平衡。


利用型候选靠近模型认为的高性能区域，有利于尽快得到改进材料；探索型候选位于模型不确定但物理上仍合理的区域，用于扩展知识边界；重复型或近重复型候选用于估计制备和测量噪声。真实实验批次通常同时包含这几类样品，而不是只给出一个单一的“最佳配方”。主动学习批次中的候选类型见表 \ref{tab:active-learning-batch-types}。


\begin{formalTable}[htbp]
\caption{主动学习批次中的候选类型}
\label{tab:active-learning-batch-types}
\begin{tabularx}{\textwidth}{@{}YYY@{}}
\toprule
\textbf{候选类型} & \textbf{选择依据} & \textbf{在实验轮次中的作用} \\
\midrule
利用型 & 预测性能高且约束满足 & 尽快逼近可用材料 \\
探索型 & 模型不确定但物理上可行 & 扩展模型对未知区域的认识 \\
边界型 & 靠近工艺或相稳定边界 & 识别可制备窗口 \\
重复型 & 已有样品或近似重复配方 & 估计实验噪声和批次误差 \\
\bottomrule
\end{tabularx}
\end{formalTable}

实验噪声和批次效应是主动学习能否落地的关键。同一成分可能因混料、升温速率、保温时间、气氛、坩埚污染或仪器校准而表现不同。如果模型把这些差异全部当成材料本征规律，就会错误更新搜索方向。这意味着，样品批次、设备、制备程序、测量条件和失败原因都应进入数据记录。


失败样品同样提供信息。某类成分若反复析晶、开裂或产生杂相，说明该区域可能超出可制备边界。例如，某轮玻璃熔制实验中如果发现 La\textsubscript{2}O\textsubscript{3} 含量高于 12\% 的配方全部析晶，下一轮模型排序就应把这一区域标记为高风险，并降低相邻候选的推荐优先级。主动学习不应只学习“哪些候选性能高”，还应学习“哪些区域无法稳定制备”。这使下一轮实验决策同时考虑性能、不确定性、工艺风险和已有负结果。


\subsubsection*{\textbf{1. 案例示范：AR 光波导玻璃主动学习}}

AR 光波导玻璃可作为连续案例。初始数据由已有氧化物配方及其折射率、Tg、热膨胀系数和析晶情况构成；代理模型预测候选性质和不确定性；每轮选择若干高性能候选、若干未知区域候选和少量重复样品；实验结果再用于修正模型和玻璃形成边界。这样，模型输出被转化为可执行的实验批次。

主动学习与普通筛选最大的不同，是它把“下一步做什么实验”作为核心问题。普通筛选常常一次性给出候选排序，而主动学习会根据已经完成的实验不断调整模型。第一轮实验后，模型可能发现某个组成区域折射率较高但容易析晶；第二轮就应减少该区域中高风险候选的比例，同时在相邻区域选择少量探索样品。这样，实验不是简单验证模型，而是在持续修正模型。

在真实实验中，一轮实验通常不应只选择预测性能最高的候选。如果全部选择利用型样品，模型可能很快陷入局部区域，对未知空间了解不足；如果全部选择探索型样品，短期内又可能得不到性能改进。为此，一个合理批次通常同时包含高性能候选、模型不确定候选、边界候选和重复样品。重复样品表面上浪费名额，但它可估计实验噪声和批次误差，对判断模型是否真的改进很重要。

主动学习还需要记录失败类型。对于玻璃配方来说，失败可能表现为不能熔清、冷却后析晶、样品开裂、气泡过多、折射率测试失败或热膨胀系数不满足要求。不同失败类型对应不同边界：不能熔清可能是工艺温度问题，析晶可能是玻璃形成能力不足，开裂可能与热膨胀系数或退火制度有关。把这些失败统一写成“失败”会损失大量材料信息。主动学习实验结果的记录方式见表 \ref{tab:active-learning-result-records}。

\begin{formalTable}[htbp]
\caption{主动学习实验结果的记录方式}
\label{tab:active-learning-result-records}
\footnotesize
\begin{tabularx}{\textwidth}{@{}YYYY@{}}
\toprule
\textbf{实验结果类型} & \textbf{记录内容} & \textbf{对模型的作用} & \textbf{下一轮决策} \\
\midrule
性能达标 & 折射率、透过率、Tg、热膨胀系数 & 提供高价值正样本 & 在附近细化搜索 \\
性能一般但可成玻 & 完整性质数据和配方 & 扩展模型覆盖范围 & 可作为训练样本保留 \\
析晶或开裂 & 失败现象、热处理条件 & 标记不可制备边界 & 减少相邻高风险候选 \\
重复样品 & 批次、设备、测试误差 & 估计噪声水平 & 修正不确定性判断 \\
\bottomrule
\end{tabularx}
\end{formalTable}

AI 可辅助把主动学习输出整理成实验批次说明，但提示词中必须要求它保留利用、探索、边界和重复样品的比例，并提醒不能自动执行模型给出的批次。

\begin{samepage}
\begin{userbox}
我正在做 AR 光波导玻璃配方的主动学习优化。模型已经给出候选配方的预测折射率、不确定性、Tg、热膨胀系数、距离已有样本的远近以及是否满足基本约束。请帮我整理下一轮实验批次选择原则。

要求如下：
\begin{enumerate}
\item 区分利用型、探索型、边界型和重复型样品；
\item 说明每类样品为什么需要进入实验批次；
\item 提醒哪些实验元数据必须记录；
\item 说明失败样品如何反馈到下一轮模型；
\item 不要直接把模型最高分样品全部选为下一轮实验。
\end{enumerate}
\end{userbox}

\begin{aibox}
下一轮实验批次可按四类候选组合：

\begin{enumerate}
\item \textbf{利用型样品：}预测折射率较高、Tg 和热膨胀系数满足约束，用于尽快获得性能改进。
\item \textbf{探索型样品：}预测不确定性较高但物理上仍可行，用于扩展模型对未知组成区域的认识。
\item \textbf{边界型样品：}靠近玻璃形成或工艺边界，用于判断哪些区域容易析晶、开裂或无法熔制。
\item \textbf{重复型样品：}选择已有样品或近似重复样品，用于估计实验噪声、批次差异和测试误差。
\end{enumerate}

每个样品都应记录配方、原料批次、熔制温度、保温时间、冷却制度、样品状态、测试设备和失败原因。失败样品不应删除，而应作为不可制备区域或高风险区域的信息反馈给下一轮模型。
\end{aibox}
\end{samepage}

以下代码示例保留了主动学习最核心的计算关系：预测值体现“利用”，不确定性体现“探索”，Tg 和热膨胀系数等惩罚项体现工艺约束。代码输出的是待人工复核的实验批次，而不是可自动执行的最终方案。


\begin{lstlisting}[language=Python,escapeinside={(*@}{@*)}]
import pandas as pd
cand = pd.DataFrame({
    "id": ["G01", "G02", "G03", "G04", "G05", "G06"],
    "n_pred": [1.86, 1.91, 1.88, 1.83, 1.94, 1.89],
    "sigma": [0.012, 0.010, 0.006, 0.030, 0.018, 0.026],
    "Tg": [565, 548, 570, 535, 505, 552],
    "cte": [8.1, 8.7, 8.4, 8.9, 9.8, 8.0],
    "dist": [0.25, 0.18, 0.35, 0.12, 0.08, 0.10],
    "ok": [True, True, True, True, False, True],
})
target_cte = 8.3
cand["penalty"] = (
    (cand["Tg"] < 520).astype(int) * 0.20
    + (cand["cte"] - target_cte).abs() * 0.03
    + (~cand["ok"]).astype(int) * 10
)
cand["score"] = cand["n_pred"] + 0.4 * cand["sigma"] - cand["penalty"]
feasible = cand[cand["ok"]].copy()
use = feasible.sort_values("n_pred", ascending=False).head(2)
unused = feasible.drop(use.index)
explore = unused.sort_values("sigma", ascending=False).head(2)
unused = unused.drop(explore.index)
boundary = unused.sort_values("dist").head(1)
# (*@\textnormal{保留 G03 作为重复样品，用于评估实验重复性}@*)
repeat = feasible[feasible["id"].eq("G03")].copy()
batch = pd.concat([
    use.assign(kind="use"),
    explore.assign(kind="explore"),
    boundary.assign(kind="boundary"),
    repeat.assign(kind="repeat"),
])
cols = ["id", "kind", "n_pred", "sigma", "Tg", "cte", "score"]
print(batch[cols])
\end{lstlisting}

\manualfigcaption{\textbf{代码示例：}主动学习实验批次评分示例}

\subsubsection*{\textbf{小结}}

本节说明候选筛选应按照由低成本到高成本、由粗判断到精验证的漏斗推进。
规则过滤、代理模型、物理计算和实验验证各有作用，每一级都应保留通过和排除理由。
主动学习把预测值、不确定性和实验反馈结合起来，失败样品和噪声记录同样是有价值的数据。

\subsection*{习题}

\begin{enumerate}[leftmargin=*]
\item 高通量筛选为什么常采用分层漏斗？
\item 比较规则过滤、代理模型、DFT 和实验验证的作用。
\item 主动学习中的“利用”和“探索”分别指什么？
\end{enumerate}

\clearpage
\section{生成设计与候选验证}


\begin{sectionkeypoints}
\item 生成设计改变的是候选产生方式，可在目标条件下提出新的组成、结构或配方，但不能取消材料约束和实验验证。
\item 材料表示决定生成模型能够处理的对象。晶体、玻璃、分子和聚合物需要不同表示方式，不宜简单套用同一种输入格式。
\item 条件生成需要把目标性能转化为可计算条件，例如带隙范围、形成能、元素体系、成本边界或供应链限制。
\item 生成候选必须经过有效性、稳定性、新颖性、可合成性和应用性等多层验证，不应只看生成数量。
\item 生成设计通常需要与筛选、主动学习、DFT 计算和实验验证串联使用，才能形成可信的材料发现流程。
\end{sectionkeypoints}

\subsection{条件生成与验证流程}

候选筛选只能在已经列出的空间中寻找材料。如果候选库没有包含某类新结构或新配方，筛选再快也无法发现它。生成设计尝试学习已有材料的分布，并在目标条件下提出新的组成、结构或配方。它改变的是候选产生方式，而不是取消物理约束和实验验证。


材料表示决定生成模型能够处理什么。分子可使用字符串或分子图，晶体需要同时表示元素种类、周期性晶格和原子坐标，玻璃与非晶材料常以成分和短程有序统计表示，聚合物还要考虑单体序列和拓扑。生成表示若不能表达周期性、价态和局域配位，模型容易产生形式正确但没有材料意义的候选。表中的 VAE\footnote{Variational Autoencoder，变分自编码器。} 指变分自编码器，GAN\footnote{Generative Adversarial Network，生成对抗网络。} 指生成对抗网络，自回归模型\footnote{Autoregressive Model，自回归模型，按序列顺序逐步生成数据。} 按步骤生成序列或结构片段，扩散模型\footnote{Diffusion Model，扩散模型，从噪声出发逐步去噪生成数据。} 从噪声中逐步生成结构，SMILES\footnote{Simplified Molecular Input Line Entry System，简化分子线性输入系统。} 是分子线性表示方法。不同生成模型路线与材料对象的对应关系见表 \ref{tab:generative-model-routes}。


\begin{formalTable}[htbp]
\caption{生成模型路线与材料对象}
\label{tab:generative-model-routes}
\begin{tabularx}{\textwidth}{@{}YYYY@{}}
\toprule
\textbf{模型路线} & \textbf{材料生成中的直觉} & \textbf{典型对象} & \textbf{需要特别注意} \\
\midrule
VAE & 把材料压缩到连续潜空间并在其中优化 & 分子、聚合物、部分晶体表示 & 潜空间是否有物理意义 \\
GAN & 生成器与判别器对抗，逼近真实样本分布 & 显微组织、分子、结构图像 & 训练稳定性与约束不足 \\
自回归模型 & 按序列或步骤逐步生成 & SMILES、合成路线、结构 token & 语法有效性和错误传播 \\
扩散模型 & 从噪声逐步去噪生成结构 & 晶体、分子三维结构 & 周期边界、对称性和坐标处理 \\
混合工作流 & 生成、筛选、DFT、实验闭环结合 & 无机晶体、玻璃、催化剂 & 验证流程比模型名称更重要 \\
\bottomrule
\end{tabularx}
\end{formalTable}

VAE、GAN、自回归模型和扩散模型可采用不同方式学习材料分布。VAE 更适合把分子或配方压缩到连续潜空间中再优化；GAN 常用于生成图像式显微组织或分子结构，但训练稳定性和约束控制需要特别检查；自回归模型适合 SMILES、合成路线或结构 token 这类有顺序的表示；扩散模型适合逐步生成分子或晶体三维结构，但必须处理周期性边界、坐标对称性和原子距离合理性。本章不展开各模型的数学细节，只保留共同逻辑：模型从已有材料中学习可行区域，目标条件用于引导生成方向，化学和结构约束用于排除无效样本，正向模型和高精度计算用于评价生成结果。


\subsubsection*{\textbf{1. 案例示范：发光材料条件生成}}

条件生成是逆向设计的关键。目标条件可以是带隙范围、形成能、体模量、磁性、元素体系或供应链限制。以发光材料为例，目标波长可转化为带隙范围，生成模型提出候选组成和结构，再由第\ref{chap:materials-cognition-prediction}章的带隙模型、DFT 和实验逐级验证。目标带隙只是入口，缺陷、掺杂、热稳定和器件兼容性仍需单独检查。

对于发光材料设计，条件生成可从两个层次理解。第一层是性能条件，例如目标发光波长、带隙范围、热稳定性和缺陷容忍度。第二层是材料边界，例如元素体系、晶体结构类型、掺杂元素、成本和安全要求。如果只给生成模型一个带隙条件，模型可能提出带隙合适但不稳定、含风险元素或难以合成的候选。因此，条件生成需要同时输入目标和边界。

发光材料还具有明显的“性质链条”。目标颜色对应发光波长，发光波长与带隙有关，但材料是否真正发光还取决于缺陷、能级、载流子复合和非辐射损失。换言之，带隙只是候选进入下一步验证的第一道门槛。一个候选材料即使带隙接近目标，如果存在深能级缺陷、热稳定性差或与器件结构不兼容，也不宜直接认为适合应用。

由此可见，生成设计更适合作为“候选扩展器”。它可帮助研究者跳出已有数据库或人工枚举空间，提出一些值得进一步检查的新组成或新结构。但候选是否有意义，要经过规则检查、正向预测、结构弛豫、高精度计算和实验验证。生成数量本身没有价值，真正有价值的是通过层层验证后仍然保留下来的候选。发光材料条件生成中的约束类型见表 \ref{tab:luminescent-material-constraints}。

\begin{formalTable}[htbp]
\caption{发光材料条件生成中的约束类型}
\label{tab:luminescent-material-constraints}
\footnotesize
\begin{tabularx}{\textwidth}{@{}YYYY@{}}
\toprule
\textbf{条件类型} & \textbf{发光材料中的含义} & \textbf{可能数据字段} & \textbf{验证方式} \\
\midrule
目标条件 & 希望获得的发光颜色或带隙范围 & 目标波长、带隙范围 & 光谱测试、带隙预测、DFT \\
结构条件 & 候选结构是否合理 & 空间群、原子坐标、局域配位 & 结构弛豫、距离检查、声子计算 \\
稳定条件 & 候选是否可能存在 & 形成能、凸包能、热稳定性 & 高精度计算、热分析、实验合成 \\
应用条件 & 是否适合器件环境 & 缺陷、热稳定、外延兼容 & 发光测试、器件测试、可靠性测试 \\
\bottomrule
\end{tabularx}
\end{formalTable}

在使用 AI 辅助生成设计时，提示词应强调“提出候选后的验证链条”，而不是让 AI 直接给出材料名称。示例提示词可这样组织。

\begin{samepage}
\begin{userbox}
我想设计一种目标发光波长在蓝光范围附近的无机发光材料。请帮我把这个目标转化为条件生成设计流程。

要求如下：
\begin{enumerate}
\item 说明目标波长如何转化为带隙或能级相关条件；
\item 说明生成模型可提出哪些类型的候选，例如组成、掺杂或晶体结构；
\item 列出生成后必须检查的化学、结构、稳定性和应用约束；
\item 说明为什么不应只根据带隙判断材料是否适合发光器件；
\item 不要直接虚构某个“最佳材料”。
\end{enumerate}
\end{userbox}

\begin{aibox}
可把任务拆成以下流程：

\begin{enumerate}
\item 目标条件设定：将目标发光波长转化为近似带隙范围，同时记录发光效率、热稳定性和器件兼容性等要求。
\item \textbf{候选生成：}生成模型可提出候选组成、掺杂元素、晶体结构或局域配位环境，但这些结果只是候选，不是最终材料。
\item \textbf{有效性检查：}检查化学式、价态、原子距离、晶格体积和结构格式是否合理，排除明显无意义候选。
\item 稳定性验证：用形成能、凸包能、结构弛豫和必要的声子或热稳定性分析判断候选是否可能存在。
\item \textbf{应用验证：}进一步检查缺陷、非辐射复合、热稳定、外延匹配和器件环境适配性。
\end{enumerate}

带隙合适只能说明候选可能满足发光颜色要求，不能说明它一定具有高发光效率或可制备性。生成模型的作用是扩展候选空间，最终判断仍然需要计算和实验验证。
\end{aibox}
\end{samepage}

条件生成还需要保留“负面条件”。很多时候，研究者不仅知道想要什么，也知道不想要什么。例如，不希望候选含有某些高风险元素，不希望晶体结构出现过短原子距离，不希望配方超出已知可制备范围，不希望候选与训练集中已有材料完全重复。这些负面条件如果不写入生成和筛选流程，模型可能不断提出形式上满足目标性能、但材料上明显不可接受的候选。

具体而言，条件生成的输入可分为三类：目标条件、边界条件和排除条件。目标条件推动模型朝目标性能靠近；边界条件保证候选处于可研究范围；排除条件提前去掉明显不可接受的结果。对于材料科学任务，三类条件同样重要。只写目标条件，容易得到“性能诱人但不可用”的候选；只写排除条件，又可能使搜索空间过窄，失去发现新材料的机会。


\subsection{候选验证与路线组合}

\subsubsection*{\textbf{1. 验证流程：从有效性到应用性}}

生成候选首先要检查有效性：化学式、价态、原子距离、晶格体积和周期性是否合理；其次检查稳定性：结构弛豫后是否保持、形成能和声子是否可接受；再次检查新颖性：候选是否只是训练集或数据库中的重复结构；最后检查可合成性和应用条件。任何一级不通过，都不应把候选直接称为材料发现。生成设计的主要验证维度见表 \ref{tab:generative-design-validation}。


\begin{formalTable}[htbp]
\caption{生成设计的验证维度}
\label{tab:generative-design-validation}
\begin{tabularx}{\textwidth}{@{}YYYY@{}}
\toprule
\textbf{验证指标} & \textbf{问题} & \textbf{常用判断} & \textbf{为什么重要} \\
\midrule
有效性 & 结构是否符合基本规则？ & 价态、距离、格式、晶格范围 & 避免无意义候选 \\
稳定性 & 候选是否可能存在？ & 形成能、凸包能、声子、动力学 & 减少纸上材料 \\
新颖性 & 是否只是训练集复述？ & 结构匹配、组成相似度、数据库检索 & 证明发现价值 \\
可合成性 & 实验能否做出来？ & 合成路线、温度窗口、前驱体、杂相 & 连接真实材料 \\
应用性 & 能否满足场景要求？ & 服役环境、成本、可靠性、规模化 & 决定工程价值 \\
\bottomrule
\end{tabularx}
\end{formalTable}

可靠的验证链条通常分为四层。第一层是格式、几何和重复性检查；第二层是快速正向模型打分；第三层是 DFT、分子动力学或热力学精算；第四层是实验合成、结构表征和性能测试。研究报告还应说明每一级候选的通过率，区分生成数量、有效数量、稳定数量和实验验证数量。


筛选、主动学习和生成设计可串联使用：生成模型扩展候选空间，规则和代理模型快速排除明显不合理候选，主动学习决定哪些候选最值得进入下一轮高精度计算或实验。最终判断仍然来自材料约束和验证证据，而不是模型名称。

可合成性是生成设计最容易被低估的环节。一个结构在计算上稳定，不代表实验上容易合成。实验可达性受到前驱体、温度、压力、气氛、反应路径、动力学障碍和杂相竞争影响。对于玻璃材料，候选配方还必须能熔制、澄清、冷却成玻璃并避免严重析晶；对于晶体材料，候选结构还要考虑相竞争和合成窗口。可合成性判断不足，是“纸上材料”大量出现的重要原因。生成候选在不同证据阶段的表述边界见表 \ref{tab:generated-candidate-evidence}。

\begin{formalTable}[htbp]
\caption{生成候选在不同证据阶段的表述边界}
\label{tab:generated-candidate-evidence}
\footnotesize
\begin{tabularx}{\textwidth}{@{}YYYY@{}}
\toprule
\textbf{候选阶段} & \textbf{可采用的表述} & \textbf{不能怎样表述} & \textbf{需要补充证据} \\
\midrule
模型生成 & 生成了一批候选结构或配方 & 发现了一批新材料 & 格式、价态和几何检查 \\
模型打分 & 获得高优先级候选 & 证明候选性能优异 & 正向模型误差和适用范围 \\
高精度计算 & 候选在计算上接近稳定 & 候选一定可合成 & 结构弛豫、凸包能、相竞争 \\
实验合成 & 候选已被制备并测试 & 可直接工程应用 & 服役、可靠性和规模化验证 \\
\bottomrule
\end{tabularx}
\end{formalTable}

这一部分也适合用 AI 辅助整理验证清单。关键是要求 AI 区分证据层级，避免把模型输出写成最终结论。

\begin{samepage}
\begin{userbox}
我用生成模型得到了一批候选无机晶体结构。请帮我设计一个候选验证清单，判断这些候选能否进入后续计算或实验。

要求如下：
\begin{enumerate}
\item 按有效性、稳定性、新颖性、可合成性和应用性分层；
\item 每一层说明需要检查哪些信息；
\item 区分“模型推荐候选”和“已验证材料”的表述；
\item 提醒如何记录每一级候选数量和通过率；
\item 不要把生成模型输出直接称为材料发现。
\end{enumerate}
\end{userbox}

\begin{aibox}
可建立五层验证清单：

\begin{enumerate}
\item \textbf{有效性：}检查化学式、价态、原子距离、晶格体积、周期性和文件格式是否合理。
\item 稳定性：进行结构弛豫，检查形成能、凸包能、声子稳定性或分子动力学稳定性。
\item \textbf{新颖性：}与训练集、材料数据库和文献记录比对，判断是否为已有材料或高度相似结构。
\item \textbf{可合成性：}分析前驱体、温度窗口、相竞争、动力学路径和实验设备条件。
\item \textbf{应用性：}结合目标场景检查性能、成本、安全、服役环境和长期可靠性。
\end{enumerate}

报告中应分别写出生成数量、有效数量、稳定数量、新颖数量、进入实验数量和实验验证数量。模型生成和模型打分只能说明候选值得进一步检查，不宜直接等同于材料发现。
\end{aibox}
\end{samepage}

在实际写作中，还应避免把验证流程写成一条单向直线。候选在任何一级验证失败，都可能需要返回前面的步骤。例如，结构弛豫后候选发生明显畸变，说明生成结构可能不合理；DFT 显示候选远离稳定区域，说明筛选规则需要加强稳定性约束；实验中反复出现杂相，说明可合成性判断不足。这些失败不是无效信息，而是帮助修正候选生成和筛选规则的重要反馈。

因此，候选验证更像是一个带反馈的流程。生成模型提出候选，验证流程排除或确认候选，失败原因再返回去修改约束、训练数据或生成条件。只有把反馈写清楚，生成设计才不是“模型输出一批材料”，而是“模型、计算和实验共同收敛到更可靠候选”的过程。

\subsubsection*{\textbf{小结}}

本节说明生成设计的作用是扩展候选空间，而不是直接给出最终材料。
条件生成需要同时设置目标条件、边界条件和排除条件，并选择适合材料对象的表示方式。
生成候选必须经过有效性、稳定性、新颖性、可合成性和应用性验证，失败结果也应回流流程。

\subsection*{习题}

\begin{enumerate}[leftmargin=*]
\item 生成候选为什么不能直接称为材料发现？
\item 为发光材料生成任务设计三类验证指标。
\end{enumerate}

\clearpage
\section{可信发现与路线选择}


\begin{sectionkeypoints}
\item 可信材料发现需要区分数据、模型、材料、验证和应用等不同证据层级，不应把模型推荐直接等同于真实材料发现。
\item 方法选择应由材料问题决定。候选库明确时优先筛选，实验昂贵时优先主动学习，候选库不足时再考虑生成模型。
\item 安全、环境、供应链、成本、可加工性和长期可靠性应作为并列约束保留，不能被压缩成无法解释的单一分数。
\item 第\ref{chap:materials-cognition-prediction}章的材料表示、性质预测、不确定性、机器学习势和可解释性，是第\ref{chap:materials-design-discovery}章候选筛选、生成验证和闭环优化的基础。
\item 全模块的完整逻辑是：定义目标、限定边界、提出候选、逐层验证、实验反馈，并在新数据基础上继续更新模型。
\end{sectionkeypoints}

\subsection{可靠评价与方法选择}

人工智能能够扩大材料搜索范围，也会放大数据偏差和模型外推风险。材料数据库中的计算参数、实验条件和样品质量并不完全一致，成功结果通常多于失败结果，热门体系的数据也远多于冷门体系。模型会继承这些差异，因此候选排名并不是客观真理，而是特定数据和模型边界下的判断。


材料性能还高度依赖结构、工艺和服役环境。同一化学式在不同缺陷、微结构和热处理条件下可能表现迥异，如果数据只记录成分和最终性能，模型就会把未记录的工艺差异当成噪声。逆向设计因此必须明确材料对象和数据适用范围，不应把材料简化为一个化学式或单一分数。


“纸上材料”指计算或模型预测较好，但缺少可合成性、规模化或应用验证的候选。真正的材料发现需要区分不同证据层级：模型预测提供线索，高精度计算提高物理可信度，实验合成证明材料可达，器件测试和长期稳定性才说明应用潜力。材料发现结果的分层评价见表 \ref{tab:materials-discovery-evaluation}。


\begin{formalTable}[htbp]
\caption{材料发现结果的分层评价}
\label{tab:materials-discovery-evaluation}
\begin{tabularx}{\textwidth}{@{}YYY@{}}
\toprule
\textbf{评价层面} & \textbf{核心判断} & \textbf{可靠表述示例} \\
\midrule
数据 & 来源、去重、划分是否清楚 & 在某数据库和划分方式下模型表现良好 \\
模型 & 是否有基线、不确定性和边界说明 & 模型推荐了若干高优先级候选 \\
材料 & 是否满足化学、结构和稳定性约束 & 候选在计算上接近稳定且结构合理 \\
验证 & 是否经过 DFT、实验或器件测试 & 某候选已合成并测得目标性质 \\
应用 & 是否考虑成本、安全和可靠性 & 候选具备进一步器件验证价值 \\
\bottomrule
\end{tabularx}
\end{formalTable}

可信评价的关键，是把“模型表现好”和“材料真的好”分开。模型表现好，通常指在某个数据集、某种划分方式和某个评价指标下，预测误差较小或排序效果较好。材料真的好，则意味着候选在真实制备、结构稳定、性能测试和服役环境中都能经受验证。二者之间存在距离。材料人工智能研究中常见的问题，是把模型层面的高分候选直接写成材料发现结果，忽略了验证证据不足。

评价结果时，还要说明模型适用范围。例如，模型是在无机晶体数据库上训练的，还是在某类玻璃配方上训练的；数据是计算数据、实验数据，还是二者混合；测试集划分是随机划分，还是按材料体系、时间或组成空间划分。如果训练集和测试集高度相似，模型指标可能看起来很好，但对真正未知材料的预测能力并不一定强。对于逆向设计来说，外推风险比普通插值预测更突出，因为候选往往位于已有数据边界附近。

不确定性也是可信评价的重要组成部分。一个候选预测性能高，但模型不确定性也高，说明它可能值得探索，但不宜直接排在确定性较高的候选之前。相反，一个候选预测性能略低但不确定性小、约束满足、可合成性强，也可能更适合作为近期实验对象。候选报告中应同时给出预测值、不确定性、约束是否满足和推荐理由。材料发现结果的可靠表述方式见表 \ref{tab:reliable-discovery-expression}。

\begin{formalTable}[htbp]
\caption{材料发现结果的可靠表述方式}
\label{tab:reliable-discovery-expression}
\footnotesize
\begin{tabularx}{\textwidth}{@{}YYYY@{}}
\toprule
\textbf{常见表述} & \textbf{风险} & \textbf{更稳妥的写法} & \textbf{还需补充} \\
\midrule
模型发现了某材料 & 夸大模型作用 & 模型推荐该候选进入验证 & 计算或实验验证 \\
该候选性能优异 & 忽略不确定性和边界 & 在当前模型和数据范围内预测值较高 & 不确定性、误差范围 \\
生成了大量新材料 & 混淆生成与发现 & 生成了若干待验证候选 & 去重、新颖性和稳定性检查 \\
候选可用于器件 & 跳过应用验证 & 候选具备进一步器件验证价值 & 器件测试和长期可靠性 \\
\bottomrule
\end{tabularx}
\end{formalTable}

候选报告还需要特别注意“数量表述”。生成模型可能一次输出几千个候选，但这并不意味着发现了几千种新材料。更准确的写法应区分：生成候选数量、格式有效数量、结构合理数量、计算稳定数量、去重后新颖数量、进入实验数量和实验验证成功数量。只有经过实验合成和性能测试的候选，才能接近“发现”的表述；只经过模型打分的候选，只能称为“推荐候选”或“高优先级候选”。

新颖性检查也属于可信评价的一部分。生成模型可能复述训练集中已经存在的材料，或者只对已有结构做很小改动。如果不做数据库检索和结构匹配，就可能把已有材料误写成新发现。对于晶体材料，可比较化学组成、空间群、晶格参数和结构相似度；对于玻璃或配方材料，可比较组分比例和已有专利、文献或实验记录。新颖性不是单靠模型输出判断的，而需要与数据库和文献记录对照。

\subsubsection*{\textbf{1. 路线选择：问题决定方法}}

方法选择应由问题决定，而不是由模型新旧决定。候选空间已经明确时，高通量筛选通常最直接；实验成本高且每轮样品有限时，主动学习更适合；需要突破已有候选库时，生成模型才体现优势。玻璃形成边界、相稳定、毒性和设备温度等硬约束，应尽可能在搜索前加入，而不是在模型给出结果后再补救。不同问题情形下的方法选择见表 \ref{tab:inverse-design-method-selection}。


\begin{formalTable}[htbp]
\caption{材料逆向设计的方法选择}
\label{tab:inverse-design-method-selection}
\begin{tabularx}{\textwidth}{@{}YYYY@{}}
\toprule
\textbf{问题情形} & \textbf{优先方法} & \textbf{理由} & \textbf{典型材料场景} \\
\midrule
候选库明确 & 高通量筛选 & 流程清楚、可解释、验证成本低 & 从数据库筛选稳定光电材料 \\
实验昂贵且小样本 & 贝叶斯优化/主动学习 & 每一轮实验都要最大化信息量 & 玻璃配方、合金热处理、催化条件 \\
候选库不足 & 生成模型 + 后筛选 & 需要提出库外候选 & 新晶体结构、功能分子、聚合物序列 \\
安全或法规约束强 & 规则约束 + 人工审核 & 先排除不可接受候选 & 含毒元素、稀缺元素或高风险反应 \\
\bottomrule
\end{tabularx}
\end{formalTable}

可靠评价还应考虑安全、环境和供应链。候选若含有有毒、稀缺或高风险元素，即使单项性能突出，也可能不适合真实应用。相应地，性能、成本、环境影响、可加工性和生命周期不应被压缩成一个无法解释的总分，而应作为并列约束保留在候选报告中。

方法选择还应考虑课程或项目阶段。对于入门学习，最合适的方法往往不是最新的生成模型，而是可解释、可复核的筛选流程。因为研究者需要先理解数据字段、模型输入、约束条件和验证证据之间的关系。对于研究生或真实科研项目，如果已有稳定的数据积累和实验反馈机制，再引入主动学习或生成模型才更合理。否则，复杂模型可能只是增加技术难度，并不能提高材料发现质量。

在写作时，可把方法选择写成一组判断问题：候选空间是否已经明确？每次实验成本是否很高？是否需要提出已有数据库之外的新候选？是否存在严格安全或工艺边界？是否有能力进行后续高精度计算和实验验证？这些问题比“选择哪个先进算法”更重要。

这一节中的方法选择和可信评价也可借助 AI 辅助整理，但提示词应要求它给出理由和证据边界，不应只推荐一个模型名称。

\begin{samepage}
\begin{userbox}
我正在写材料逆向设计方案，候选空间包括已有数据库中的无机晶体和少量可枚举组成。实验成本较高，每轮只能验证 5～10 个样品。请帮我判断优先使用高通量筛选、主动学习还是生成模型。

要求如下：
\begin{enumerate}
\item 根据候选空间、实验成本、数据量和验证能力进行判断；
\item 说明每种方法适合和不适合的情况；
\item 给出推荐路线，但必须说明理由；
\item 提醒结果报告中应如何避免夸大模型作用；
\item 不要只因为生成模型较新就优先推荐它。
\end{enumerate}
\end{userbox}

\begin{aibox}
在这个任务中，候选空间已有数据库和可枚举组成，说明不一定需要首先使用生成模型。更合理的路线可以是：

\begin{enumerate}
\item 第一阶段使用高通量筛选：先用规则、正向模型和稳定性判断从已有候选库中排除明显不合理候选。
\item \textbf{第二阶段使用主动学习：}由于每轮实验只能验证 5～10 个样品，应让模型同时考虑预测性能和不确定性，选择高价值实验批次。
\item 生成模型暂时作为补充：只有当已有候选库不能覆盖目标性能区域，或者需要提出库外结构时，再考虑生成模型。
\end{enumerate}

报告中应写成“模型推荐若干高优先级候选，并通过后续计算或实验验证”，而不是“模型发现了最终材料”。如果没有实验或高精度计算支持，应明确说明候选仍处于待验证阶段。
\end{aibox}
\end{samepage}

\subsection{能力复用与闭环迭代}

\subsubsection*{\textbf{1. 两章闭环：正向模型服务逆向设计}}

两章之间形成正反闭环。第\ref{chap:materials-cognition-prediction}章把材料数据转化为描述符或结构图，训练正向模型并解释模型边界；第\ref{chap:materials-design-discovery}章把这些模型作为候选筛选和生成验证的评价器，再用新计算和实验结果更新数据。没有正向预测，逆向搜索缺少可靠评价；没有逆向设计，正向模型容易停留在被动打分。第\ref{chap:materials-cognition-prediction}章能力在第\ref{chap:materials-design-discovery}章中的复用关系见表 \ref{tab:chapter-one-capability-reuse}。


\begin{formalTable}[htbp]
\caption{第\ref{chap:materials-cognition-prediction}章能力在第\ref{chap:materials-design-discovery}章中的复用}
\label{tab:chapter-one-capability-reuse}
\begin{tabularx}{\textwidth}{@{}YYY@{}}
\toprule
\textbf{第\ref{chap:materials-cognition-prediction}章能力} & \textbf{第\ref{chap:materials-design-discovery}章中的用途} & \textbf{材料例子} \\
\midrule
性质预测 & 作为筛选和生成后的快速打分器 & 带隙筛选 microLED 候选 \\
材料表示 & 决定候选能否被模型有效评价 & 晶体图、玻璃成分、分子图 \\
不确定性量化 & 指导主动学习下一步实验 & 高折射率玻璃小样本优化 \\
机器学习势 & 验证候选结构和高温行为 & 超高温陶瓷和高熵合金 \\
可解释性 & 把模型推荐转化为材料规律 & 形成能力、熔点或缺陷趋势分析 \\
\bottomrule
\end{tabularx}
\end{formalTable}

完整逻辑可概括为：目标性质定义问题，材料知识限定边界，筛选、主动学习或生成模型提出候选，正向模型和物理计算逐层评价，实验结果决定最终结论并反馈到数据集。人工智能的价值在于提高搜索与决策效率，材料科学的责任则是保证问题、约束和证据真实可靠。

闭环迭代要落实到可操作的数据管理上。每一轮推荐都应记录模型版本、训练数据版本、候选来源、预测值、不确定性、约束检查结果和推荐理由；每一轮验证后，还要把成功样品、失败样品、测试条件和失败原因按同一格式回写数据表。这样，下一轮模型更新才知道哪些变化来自材料本身，哪些变化来自数据版本、实验批次或测试条件。

例如，在高折射率玻璃任务中，第\ref{chap:materials-cognition-prediction}章中讨论的数据整理和描述符构建可帮助建立折射率、Tg 和热膨胀系数预测模型；第\ref{chap:materials-design-discovery}章则利用这些模型筛选候选配方，并通过主动学习选择下一轮熔制实验。在高熵合金涂层任务中，第\ref{chap:materials-cognition-prediction}章的硬度预测、图神经网络和机器学习势可提供性能估计与结构理解；第\ref{chap:materials-design-discovery}章则进一步用这些模型决定哪些成分或硬质相值得验证。这样的闭环，才是人工智能真正进入材料发现流程的方式。

闭环迭代还要求新数据及时回流。实验失败、性能一般、模型预测错误的样品，都应被记录下来。只把成功样品加入数据库，会使模型越来越偏向成功区域，无法学习失败边界；只记录最终性能而不记录工艺条件，也会使模型无法区分成分效应和工艺效应。闭环发现不仅是“模型推荐--实验验证”，还包括“实验记录--数据更新--模型修正”。材料智能发现闭环中的数据回流内容见表 \ref{tab:closed-loop-data-feedback}。

\begin{formalTable}[htbp]
\caption{材料智能发现闭环中的数据回流}
\label{tab:closed-loop-data-feedback}
\footnotesize
\begin{tabularx}{\textwidth}{@{}YYYY@{}}
\toprule
\textbf{闭环环节} & \textbf{主要任务} & \textbf{需要记录的信息} & \textbf{常见问题} \\
\midrule
模型推荐 & 根据目标和约束提出候选 & 预测值、不确定性、推荐理由 & 只给排名，不说明边界 \\
计算验证 & 用更高精度方法复核候选 & 结构、能量、稳定性、失败原因 & 只保留通过候选 \\
实验验证 & 合成、表征和性能测试 & 工艺、样品状态、测试条件 & 缺少失败样品记录 \\
数据回流 & 更新数据集并重训模型 & 新数据、元数据、版本信息 & 数据来源不可追溯 \\
\bottomrule
\end{tabularx}
\end{formalTable}

AI 可辅助整理闭环报告结构，尤其是检查是否遗漏数据来源、验证证据和失败记录。但同样需要明确：AI 只能帮助组织报告，不宜替代实验和模型审查。

\begin{samepage}
\begin{userbox}
我想把一个材料智能发现项目写成闭环流程报告。项目包括正向模型预测、候选筛选、实验验证和数据回流。请帮我设计报告结构。

要求如下：
\begin{enumerate}
\item 说明第\ref{chap:materials-cognition-prediction}章中的材料表示、性能预测和模型解释如何服务第\ref{chap:materials-design-discovery}章的逆向设计；
\item 按模型推荐、计算验证、实验验证和数据回流四个环节组织；
\item 提醒每个环节需要记录哪些信息；
\item 特别说明失败样品和预测错误样品为什么也要记录；
\item 不要把闭环写成模型自动发现材料。
\end{enumerate}
\end{userbox}

\begin{aibox}
闭环流程报告可包括四部分：

\begin{enumerate}
\item \textbf{模型推荐：}说明使用了哪些材料表示和正向预测模型，给出候选的预测值、不确定性和推荐理由。
\item \textbf{计算验证：}记录结构弛豫、稳定性判断、形成能或其他高精度计算结果，并保留未通过候选的原因。
\item \textbf{实验验证：}记录制备工艺、样品状态、表征结果、性能测试条件和失败现象，不应只记录成功样品。
\item \textbf{数据回流：}说明新数据如何加入数据库，如何更新训练集，模型是否重新训练，以及新模型与旧模型相比是否改进。
\end{enumerate}

失败样品和预测错误样品同样重要，因为它们可帮助模型学习不可制备区域、实验噪声和模型外推边界。材料智能发现不是模型自动给出答案，而是模型、计算、实验和数据更新共同形成的迭代过程。
\end{aibox}
\end{samepage}

如果把两章放在一个完整模块中理解，第\ref{chap:materials-cognition-prediction}章更像是“建立评价能力”，第\ref{chap:materials-design-discovery}章更像是“利用评价能力做决策”。没有第\ref{chap:materials-cognition-prediction}章的数据整理、表示学习和模型评估，第\ref{chap:materials-design-discovery}章的筛选和生成就缺少可靠判断依据；没有第\ref{chap:materials-design-discovery}章的候选选择和实验反馈，第\ref{chap:materials-cognition-prediction}章的模型又容易停留在已有数据上的被动预测。二者结合起来，才能形成从材料认知到材料发现的完整教学闭环。

最后一节不只是总结第\ref{chap:materials-design-discovery}章，也是在说明整个模块的逻辑边界：人工智能可帮助研究者更快地整理材料数据、构建候选、预测性能、安排验证和更新模型，但不应跳过材料科学中的基本问题。材料是否真实存在，能否制备出来，是否满足应用场景，仍然要靠实验、计算和工程约束共同判断。

\subsubsection*{\textbf{小结}}

本节强调可信材料发现需要区分模型推荐、计算证据、实验验证和应用潜力。
方法选择应由问题决定，候选库明确时优先筛选，实验昂贵时优先主动学习，候选不足时再考虑生成。
第\ref{chap:materials-cognition-prediction}章建立评价能力，第\ref{chap:materials-design-discovery}章利用评价能力做决策，数据回流使整个模块形成闭环。

\subsection*{习题}

\begin{enumerate}[leftmargin=*]
\item 数据回流为什么是材料智能发现闭环的关键？
\item 如何根据候选库、实验成本和数据量选择材料发现路线？
\end{enumerate}

\section*{本章小结}

本章讨论了材料设计与智能发现的基本思路。逆向设计首先要求将应用需求转化为可计算目标、变量和约束；候选筛选通过规则、代理模型、高精度计算和实验验证逐层缩小范围；主动学习利用预测值和不确定性安排下一批实验或计算任务；生成设计能够提出新的结构或配方，但必须接受有效性、稳定性、新颖性、可合成性和应用条件等多层验证。全章强调，智能发现不是单次模型输出，而是目标定义、候选构建、分层验证和数据回流共同组成的闭环过程。

% 文件中不加载任何宏包，所有样式均调用项目现有的 package.tex。
\part*{微模块5\quad 人工智能在材料科学中的应用}
\addcontentsline{toc}{part}{微模块5\texorpdfstring{\quad}{ }人工智能在材料科学中的应用}

\chapter{材料认知与性能预测}
\label{chap:materials-cognition-prediction}


理解材料性能的来源，是材料科学研究的核心任务。材料的力学、热学、电学、磁学和光学性能，往往不是由单一因素决定的，而是受到成分组成、微观结构和制备工艺的共同影响。随着研究对象从简单合金、陶瓷和半导体扩展到高熵合金、复合材料、低维材料和复杂功能材料，材料候选空间迅速增大，单纯依靠实验试错已经难以系统揭示其中的规律。在这一背景下，人工智能方法的价值，主要体现在对分散实验数据、计算数据和文献知识的组织与学习上，使模型能够从中提取“成分--结构--工艺--性能”之间的关联，从而辅助材料性能预测和后续实验设计。

本章围绕“材料认知与性能预测”展开，重点解决四个问题。第一节说明材料人工智能为什么需要可靠数据基础，以及实验、计算、文献和数据库数据应如何整理。第二节讨论材料成分和晶体结构如何转化为描述符或图表示，使模型能够读取材料信息。第三节介绍描述符模型、图神经网络和机器学习原子间势（下文简称“机器学习势”）如何服务硬度预测、性质预测和熔点模拟。第四节进一步说明模型解释与物理知识嵌入的作用，帮助研究者判断模型结果是否符合材料机理和实验事实。通过本章学习，研究者需要建立一个基本认识：人工智能并不是替代材料理论和实验验证，而是与数据、计算和材料机理共同构成辅助材料认知的新工具。




\clearpage
\section{范式变革与数据基础}


\begin{sectionkeypoints}
\item 材料人工智能中的数据并不只是单个性能数值，而是包含材料组成、结构信息、制备工艺、测试条件、数据来源和目标性能的完整记录。
\item 材料数据主要来自实验、计算、文献和数据库四类渠道。不同来源的数据各有优势，也各有局限，使用时需要结合具体研究任务进行判断。
\item 材料数据库为模型训练、性能预测和候选筛选提供了重要的数据基础，但数据库中的数据不宜直接等同于最终模型输入，还需要经过字段筛选、来源核验和格式整理。
\item 在正式建模之前，应先整理原始数据清单，明确记录材料是什么、结构是什么、性能数值是多少、数据来自哪里，以及采用了什么计算方法或实验条件。
\item AI 可辅助解释数据库字段、设计表格框架和检查数据规范性，但不宜替代数据库原始记录、实验记录和论文原文核验。材料数据整理的核心要求是清楚、可追溯、可检查。
\end{sectionkeypoints}

\subsection{实验试错与数据驱动}

材料科学的核心任务之一，是理解和建立材料成分、结构、工艺与性能之间的关联。研究者希望理解：为什么某种成分的合金具有更高强度？为什么某种陶瓷能够承受更高温度？为什么某种半导体材料具有合适的带隙并能用于发光器件？这些问题表面上属于不同材料体系，但本质上都指向同一个核心：材料的内部组成与结构如何决定其外部性能。


在传统材料研发中，研究者通常依靠已有理论、文献经验和实验试错来提出材料方案，再通过制备、表征和性能测试不断调整成分与工艺。例如，为了提高合金强度，可尝试加入不同合金元素；为了改善陶瓷耐高温性能，可调节烧结温度、晶粒尺寸或相组成；为了优化光电材料性能，可改变材料的元素组成和晶体结构。这种方式直接可靠，也符合材料实验研究的基本规律。它的优势在于结论来自真实样品和真实测试条件，缺点是往往需要大量重复实验，研发周期较长，成本也较高。


当材料体系从简单单组元或二元材料扩展到多组元合金、复合材料、异质结构、低维材料和复杂光电材料时，实验试错面对的组合空间会迅速扩大。以高熵合金为例，四种、五种甚至更多主要元素的组合、比例和热处理工艺都可能改变组织结构和性能；同时，强度、硬度、韧性、耐腐蚀性、导电性或发光性能又受到晶体结构、缺陷、晶粒尺寸、相组成、界面状态和服役环境共同影响。传统经验仍然重要，但单纯依靠人工逐一尝试，已经难以在有限时间内完成系统筛选。


在这一背景下，材料研发逐渐从单纯依赖实验试错，转向实验、计算、数据和人工智能相结合的新方式。实验能够提供真实可靠的性能结果，计算能够在实验之前预测部分结构和能量信息，数据库能够积累已有材料的成分、结构和性能数据，而人工智能则能够从这些数据中学习复杂的材料规律。换言之，人工智能并不是脱离实验和理论单独发挥作用，而是建立在已有材料知识、实验数据和计算数据基础上的辅助工具。


数据驱动并不是否定实验和理论，而是把已有实验数据、计算数据、文献数据和数据库资源转化为可分析、可学习和可预测的信息。人工智能模型可从大量材料数据中学习成分、结构与性能之间的关系，并用于预测未知材料的性质。这样，研究者不再需要在庞大的候选空间中完全盲目试错，而可借助数据和模型缩小研究范围，把实验资源集中到更有希望的材料体系上。


\subsubsection*{\textbf{1. 人工智能用于材料科学的基本流程}}

在人工智能用于材料科学的过程中，研究者通常需要依次解决几个关键环节。第一个环节是材料数据整理。需要收集和规范实验数据、计算数据、文献数据和数据库资源，并记录材料的成分、结构、工艺、性能以及数据来源等信息。只有数据来源清楚、含义明确，后续模型学习才具有可靠基础。

第二个环节是材料表示。机器学习模型不宜直接理解“合金”“晶体结构”或“热处理工艺”等材料概念，因此需要把材料转化为机器可处理的形式。例如，可用描述符表示化学成分和结构信息，也可用图结构表示原子之间的连接关系。

第三个环节是模型预测。在材料被合理表示之后，可建立机器学习模型，预测材料的形成能、带隙、弹性模量、熔点、硬度等性能。对于涉及原子运动和微观演化的问题，还可进一步构建机器学习势（利用机器学习模型学习原子结构与能量、力之间关系的势函数），用于模拟扩散、相变和熔化等原子尺度行为。

最后，模型预测结果还需要回到材料科学问题本身。通过可解释性方法，可分析模型为什么做出某种预测，识别哪些成分、结构或工艺因素对性能影响更大，从而把模型结果进一步转化为可理解的材料规律。


在材料研究中还必须明确，人工智能不宜替代材料专业知识、实验验证和物理判断。人工智能的作用更像是一种增强手段：它可帮助研究者整理数据、发现规律、提出候选、预测性能和辅助决策，但最终结论仍然需要通过材料理论分析和实验结果进行验证。尤其是在材料科学中，模型给出的预测结果并不等同于真实材料性能。材料是否能够被实际制备出来，结构是否稳定，以及能否在实际服役环境中长期可靠工作，仍然需要进一步验证。


围绕材料性能预测和材料设计两类问题，人工智能用于材料科学可概括为两条基本主线。第一条是正向预测，即在材料成分、结构或实验条件已经给定的情况下，预测材料可能具有的性能和行为。例如，已知一种晶体材料的结构，可预测它的形成能、带隙、弹性模量或熔点。第二条是逆向设计，即先给出目标性能，再反过来寻找可能满足要求的材料。从概念上看，正向预测关注已知材料可能表现出什么性能，逆向设计关注为了获得目标性能应当寻找什么材料。


本章主要讨论第一条主线，即材料科学中的正向预测。本章将从材料数据出发，进一步讨论材料如何被表示成机器可理解的形式，机器学习如何预测材料性能，机器学习势如何模拟原子尺度行为，以及可解释人工智能如何帮助研究者从模型结果中提炼材料规律。


\subsection{数据来源与材料数据库}

人工智能要预测材料性能，首先需要数据。对于材料科学而言，数据并不只是一个性能数值，而是对材料组成、制备过程、测试条件和性能结果的完整记录。一个可靠的材料数据条目，至少应该说明材料的化学组成、结构信息、制备工艺、测试条件、数据来源和目标性能。缺少这些基本信息，模型就难以建立清晰的输入与输出对应关系，所得预测结果也难以作为可靠依据。


材料数据主要来自实验、计算、文献和数据库四类渠道。


第一类是实验数据。实验数据来自材料制备、显微表征和性能测试过程，例如硬度、拉伸强度、腐蚀电流密度、熔点、热导率、带隙、发光波长等。实验数据最接近真实材料和真实服役环境，但获取成本较高，也容易受到制备工艺、测试条件和仪器参数的影响。例如，同一种合金在不同热处理温度下可能得到不同硬度，同一种涂层在不同腐蚀介质中可能表现出不同耐蚀性能。因此，实验数据不应只记录结果，还要尽量记录测试条件。


第二类是计算数据。计算数据主要来自第一性原理计算、分子动力学模拟、相场模拟等方法。例如，第一性原理计算可得到形成能、能带结构、弹性常数、态密度、晶体结构稳定性等信息；分子动力学可模拟扩散、相变和高温行为。计算数据的优点是容易标准化，也适合大规模生成；但计算数据依赖具体计算方法和参数，不宜简单等同于实验结果。


第三类是文献数据。大量材料知识存在于论文、专利、实验报告和技术资料中。例如，一篇论文可能同时记录材料成分、制备工艺、热处理条件、显微组织和性能结果。文献数据的优点是信息丰富，缺点是格式不统一，很多信息隐藏在正文、表格、图片或补充材料中，需要人工整理或借助文本处理工具提取。


第四类是材料数据库。材料数据库可理解为经过整理的材料信息资源库，它把大量材料的化学组成、晶体结构、计算性质、实验性质和文献信息集中到一个可检索的平台中。对于材料人工智能而言，数据库是重要数据来源，因为它可为模型训练、性能预测和材料筛选提供相对规范的数据基础。


常见材料数据库包括 Materials Project、OQMD\footnote{Open Quantum Materials Database，开放量子材料数据库。}、AFLOW\footnote{Automatic Flow Framework for Materials Discovery，自动化材料发现计算框架。}、NOMAD\footnote{Novel Materials Discovery，材料数据管理与共享平台。}、ICSD\footnote{Inorganic Crystal Structure Database，无机晶体结构数据库。}、JARVIS\footnote{Joint Automated Repository for Various Integrated Simulations，材料计算数据库与工具平台。} 等。不同数据库的侧重点并不完全相同：有的偏向无机晶体结构和稳定性数据，有的偏向大规模热力学计算结果，有的强调计算数据管理和共享，还有的主要提供实验晶体结构信息。教材中列举数据库名称，目的不是要求逐一掌握平台功能，而是说明不同数据来源对应不同研究任务。

对于入门阶段的研究者而言，不必一开始深入掌握每个数据库的所有功能，但需要知道不同数据库主要提供哪类材料数据，以及这些数据适合用于哪些研究任务。只有明确数据来源和适用范围，才能更合理地将数据库用于材料性能预测、候选筛选和模型训练。


使用材料数据库时，首先要明确自己的研究任务。不同任务需要的数据字段不同。例如，如果任务是预测无机光电材料的带隙，就需要关注化学式、晶体结构、空间群、带隙、形成能、数据来源和计算方法等字段；如果任务是预测合金硬度，就需要关注成分、相结构、制备工艺、热处理条件、硬度值和测试方法等字段；如果任务是筛选超高温陶瓷，就需要关注化学组成、晶体结构、熔点、密度、热稳定性和抗氧化相关数据。


因此，数据库只是提供了材料信息的来源，真正进入模型训练之前，还需要把这些信息整理成结构清楚、字段明确、来源可追溯的原始数据清单。下面以无机光电材料带隙数据整理为例，说明从数据库原始记录到可建模数据清单的具体做法。

\subsubsection*{\textbf{1. 案例示范：无机光电材料带隙数据整理}}

以“无机光电材料带隙预测”为例，数据准备的第一步不是训练模型，而是从数据库或文献中整理出一张原始数据清单。这个清单应该回答几个基本问题：材料是什么？结构是什么？带隙是多少？这个带隙来自哪里？是实验值还是计算值？如果是计算值，采用了什么计算方法？如果是实验值，是否记录了测试条件？


一张初步整理后的原始数据清单可按表 \ref{tab:raw-bandgap-data} 设计。表中的 eV\footnote{Electron volt，电子伏特，常用于表示电子能量和带隙大小。} 是带隙常用单位。


\begin{formalTable}[htbp]
\caption{原始带隙数据清单示例}
\label{tab:raw-bandgap-data}
\scriptsize
\setlength{\tabcolsep}{2.8pt}
\begin{tabularx}{\textwidth}{@{}YYYYYYYYY@{}}
\toprule
\textbf{材料编号} & \textbf{化学式} & \textbf{晶体结构} & \textbf{空间群} & \textbf{带隙/eV} & \textbf{数据类型} & \textbf{数据来源} & \textbf{方法或条件} & \textbf{备注} \\
\midrule
001 & GaN & 纤锌矿 & P6\textsubscript{3}mc & 3.40 & 实验 & 文献 & 室温测试 & 直接带隙 \\
002 & ZnO & 纤锌矿 & P6\textsubscript{3}mc & 3.37 & 实验 & 文献 & 室温测试 & 直接带隙 \\
003 & SrTiO\textsubscript{3} & 钙钛矿 & Pm-3m & 3.20 & 数据库 & 数据库记录 & 需注明方法 & 需核对来源 \\
\bottomrule
\end{tabularx}
\end{formalTable}

表 \ref{tab:raw-bandgap-data} 的作用是保存原始材料信息，而不是直接代表最终模型输入。它的重点在于把分散的信息整理成清楚、可追溯、可检查的记录。后续是否把化学式转化为成分特征，是否把晶体结构转化为结构描述符，属于下一节要讨论的材料表示问题。


在整理这类原始数据清单时，可借助 AI 辅助完成部分工作。此处的重点不是让 AI 代替研究者判断数据是否可靠，而是让 AI 帮助理解字段含义、设计表格框架和检查数据规范性。以下以无机光电材料带隙数据整理为例给出提示词示范。

\begin{samepage}
\begin{userbox}
我正在整理无机光电材料的带隙数据，准备用于后续机器学习建模。数据可能来自 Materials Project、JARVIS、文献或实验记录。请帮我设计一张原始数据清单的表头，并说明每一列的作用。

要求如下：
\begin{enumerate}
\item 表格要保留材料编号、化学式、晶体结构、空间群、带隙数值、带隙单位、数据类型、数据来源、计算方法或实验条件等信息；
\item 区分实验值和计算值；
\item 提醒哪些字段是后续建模前必须核验的；
\item 不要虚构具体材料数据，只需要给出表头设计和填写说明。
\end{enumerate}
\end{userbox}

\begin{aibox}
可将原始数据清单设计为以下字段：

\begin{itemize}
\item 材料编号：用于唯一标识每一条材料记录，便于后续追踪和检查。
\item 化学式：记录材料的元素组成，例如 GaN、ZnO、SrTiO$_3$ 等。
\item \textbf{晶体结构：}记录材料所属结构类型，例如纤锌矿、钙钛矿等。
\item \textbf{空间群：}记录晶体结构的对称性信息。
\item 带隙数值：记录材料的带隙大小。
\item 带隙单位：通常为 eV，避免不同来源数据混用时产生单位错误。
\item \textbf{数据类型：}说明该带隙是实验值还是计算值。
\item \textbf{数据来源：}记录数据来自数据库、论文、实验报告还是其他资料。
\item \textbf{方法或条件：}如果是计算值，应注明计算方法；如果是实验值，应注明测试温度、测试方法等条件。
\item \textbf{备注：}记录是否为直接带隙、是否需要进一步核验等补充信息。
\end{itemize}

在正式建模前，尤其需要核验带隙数值、数据类型、数据来源和方法或条件等字段。同一种材料在不同晶体结构、不同计算方法或不同测试条件下可能具有不同带隙，因此这些信息不应省略。
\end{aibox}
\end{samepage}

上面的提示词示范说明，AI 可帮助研究者快速形成数据表框架，但 AI 输出并不宜直接作为最终数据结果。研究者还需要回到数据库原始页面或论文原文，核对具体数值、计算方法、实验条件和材料结构是否准确。

因此，在材料数据整理阶段，AI 更适合充当数据整理助手，而不是数据来源本身。它可帮助设计表格、解释字段和生成检查规则，但具体数据是否准确，仍然需要根据数据库原始记录、实验记录或论文原文进行确认。尤其在材料科学中，同一种化学式可能对应不同晶体结构，不同计算方法也可能得到不同带隙值，因此数据来源、计算方法和实验条件必须保留。


为了提高数据质量，原始数据整理时应注意几个问题。第一，要统一单位。例如带隙通常用 eV，硬度可能用 HV\footnote{Vickers Hardness，维氏硬度单位。} 或 GPa\footnote{Gigapascal，吉帕，常用于表示压力、应力或弹性模量。}，能量可能用 eV/atom 或 kJ/mol。第二，要区分实验数据和计算数据。二者可同时记录，但不能在没有说明的情况下混用。第三，要保留数据来源。每一条数据宜能追溯到数据库编号、论文或实验记录。第四，要检查缺失值和重复值。同一个材料如果重复出现，需要判断它们是否对应不同结构、不同测试条件或不同计算方法。第五，要记录必要条件。对于实验数据，应尽量记录温度、压力、测试环境和样品状态；对于计算数据，应记录计算方法和参数来源。

\subsubsection*{\textbf{小结}}

本节说明材料研发从实验试错走向数据、计算和实验协同的基本背景。
材料数据应包含成分、结构、工艺、性能、来源和测试条件，而不只是单一性能数值。
数据库和 AI 可提高整理效率，但数据核验、条件记录和材料判断仍需由研究者完成。

\subsection*{习题}

\begin{enumerate}[leftmargin=*]
\item 材料数据为什么不能只记录性能值？还应补充哪些字段？
\item 比较四类材料数据来源，并说明选择依据。
\end{enumerate}

\clearpage
\section{特征构建与结构表示}


\begin{sectionkeypoints}
\item 材料特征构建的目的，是把化学式、晶体结构、空间群和性能数值等原始记录转化为机器学习模型能够读取的数值特征。
\item 成分描述符可由元素组成和元素属性计算得到，例如平均原子序数、平均电负性、电负性差和平均原子半径等。
\item 结构描述符用于补充原子排列信息，常见字段包括空间群编号、晶胞体积、密度和晶格常数等。
\item 图表示把晶体看成原子节点和原子间边组成的结构，可更直接地保留局域配位、原子间距离和周期性信息。
\item 描述符和图表示都需要材料判断。特征不是越多越好，应避免目标泄漏，并检查所选特征是否具有明确的材料意义。
\end{sectionkeypoints}

\subsection{成分与结构描述符}

上一节介绍了材料数据的来源和材料数据库的作用。通过实验、计算、文献或数据库，研究者可获得材料的化学式、晶体结构、空间群、性能数值和数据来源等信息。但是，这些信息只是原始材料记录，还不宜直接被大多数机器学习模型使用。机器学习模型真正能够处理的，通常是一组数值特征。因此，在进行性能预测之前，需要把材料信息进一步转化为模型可读取的数值表达。


这一过程称为材料的特征化，也常称为描述符构建。描述符指用数字表达材料中可能影响性能的关键信息。描述符不是随意选择的一组数字，而应该尽量反映材料的成分、结构和物理化学特征。例如，对于光电材料带隙预测，元素电负性、原子半径、晶体结构、空间群和晶胞参数等信息都可能与电子结构有关，因此可作为候选描述符。


\subsubsection*{\textbf{1. 案例示范：无机光电材料特征构建}}

本节围绕一个任务展开：无机光电材料带隙预测。假设研究者已经从材料数据库或文献中整理出表 \ref{tab:bandgap-raw-records} 所示的原始数据清单。


\begin{formalTable}[htbp]
\caption{无机光电材料原始记录示例}
\label{tab:bandgap-raw-records}
\footnotesize
\begin{tabularx}{\textwidth}{@{}YYYYY@{}}
\toprule
\textbf{材料} & \textbf{化学式} & \textbf{晶体结构} & \textbf{空间群} & \textbf{带隙/eV} \\
\midrule
氮化镓 & GaN & 纤锌矿 & P6\textsubscript{3}mc & 3.40 \\
氧化锌 & ZnO & 纤锌矿 & P6\textsubscript{3}mc & 3.37 \\
钛酸锶 & SrTiO\textsubscript{3} & 钙钛矿 & Pm-3m & 3.20 \\
\bottomrule
\end{tabularx}
\end{formalTable}

这个表格适合人阅读，但模型不宜直接理解“GaN”“纤锌矿”或“P6\textsubscript{3}mc”这些符号和文字。因此，第一步是把化学式转化为元素组成。


以 GaN 为例，它由 Ga 和 N 两种元素组成，原子比例为 1:1。为了便于模型处理，可把这个比例归一化为原子分数，即 Ga 为 0.50，N 为 0.50。ZnO 同理，Zn 为 0.50，O 为 0.50。SrTiO\textsubscript{3} 中 Sr、Ti 和 O 的比例为 1:1:3，归一化后分别为 0.20、0.20 和 0.60。


经过这一步，原始化学式可转化为表 \ref{tab:formula-composition-example} 所示形式。


\begin{formalTable}[htbp]
\caption{化学式到元素组成的转换示例}
\label{tab:formula-composition-example}
\begin{tabularx}{\textwidth}{@{}YYY@{}}
\toprule
\textbf{化学式} & \textbf{元素组成} & \textbf{原子分数} \\
\midrule
GaN & Ga, N & Ga: 0.50；N: 0.50 \\
ZnO & Zn, O & Zn: 0.50；O: 0.50 \\
SrTiO\textsubscript{3} & Sr, Ti, O & Sr: 0.20；Ti: 0.20；O: 0.60 \\
\bottomrule
\end{tabularx}
\end{formalTable}

这一步表面上简单，但当数据表中有几百或几千个化学式时，人工逐个拆解就很低效。实际建模时，通常需要用程序批量解析化学式，并检查程序是否能够正确处理下标、括号、掺杂和复杂化学式。只有化学式解析正确，后续成分描述符才有意义。


得到元素组成后，第二步是引入元素属性。机器学习模型并不知道 Ga、N、Zn、O、Sr、Ti 这些元素各自有什么物理化学特征，因此需要查找或建立元素属性表。这个属性表可包括原子序数、原子半径、电负性、价电子数、元素熔点等信息。


例如，GaN 的两个组成元素是 Ga 和 N。如果需要计算平均电负性，就需要知道 Ga 和 N 各自的电负性，再按照原子分数加权平均。一般来说，若第 i 种元素的原子分数为 xi，对应元素属性为 Pi，则该材料的平均属性可表示为：


\begin{equation*}\text{平均属性}=\sum_i x_i P_i\end{equation*}

这个公式可用于计算平均原子序数、平均电负性、平均原子半径、平均价电子数等特征。除了平均值，还可计算最大值、最小值、极差和方差。例如，电负性差可用组成元素中最大电负性与最小电负性的差值表示，用来描述元素之间化学性质差异。


对于带隙预测来说，电负性和电负性差具有直观意义。带隙与材料的电子结构和成键特征有关，而元素之间电负性差异会影响电子分布和键合方式。因此，平均电负性、电负性差、平均原子半径等都可作为初步描述符。当然，这些描述符并不能完全决定带隙，但它们为模型提供了可学习的材料信息。


经过这一步，数据表可进一步转化为表 \ref{tab:composition-descriptor-example} 所示的成分描述符表。


\begin{formalTable}[htbp]
\caption{成分描述符表构建示例}
\label{tab:composition-descriptor-example}
\footnotesize
\setlength{\tabcolsep}{2.8pt}
\begin{tabularx}{\textwidth}{@{}YYYYYY@{}}
\toprule
\textbf{化学式} & \textbf{平均原子序数} & \textbf{平均电负性} & \textbf{电负性差} & \textbf{平均原子半径} & \textbf{带隙/eV} \\
\midrule
GaN & 数值 & 数值 & 数值 & 数值 & 3.40 \\
ZnO & 数值 & 数值 & 数值 & 数值 & 3.37 \\
SrTiO\textsubscript{3} & 数值 & 数值 & 数值 & 数值 & 3.20 \\
\bottomrule
\end{tabularx}
\end{formalTable}

此处的“数值”需要根据元素属性表计算得到。实际操作时，可先准备一个元素属性 CSV\footnote{Comma-Separated Values，逗号分隔值文件。} 文件，其中包含元素符号、原子序数、电负性、原子半径等字段。然后将化学式解析结果与元素属性表关联起来，自动计算每个材料的平均值和差值。这样可减少重复计算，但计算结果仍然需要通过抽样检查来确认。


第三步，是加入结构描述符。仅有成分描述符还不够，因为材料性能不仅取决于由哪些元素组成，还取决于这些原子如何排列。GaN 和 ZnO 都是纤锌矿结构，而 SrTiO\textsubscript{3} 是钙钛矿结构。不同晶体结构对应不同的原子排列、局域环境和电子能带特征，因此结构信息对于带隙预测也很重要。


对于入门阶段的研究者，可先使用比较容易获得的结构描述符，例如空间群编号、晶胞体积、密度、晶格常数等。以空间群为例，P6\textsubscript{3}mc 可对应空间群编号 186，Pm-3m 可对应空间群编号 221。这样，文字形式的空间群就可转化为数值字段。晶胞体积和密度则可从晶体结构文件或数据库记录中获得。


加入结构描述符后，表格可进一步扩展为表 \ref{tab:structure-descriptor-example}。


\begin{formalTable}[htbp]
\caption{成分与结构描述符表构建示例}
\label{tab:structure-descriptor-example}
\scriptsize
\setlength{\tabcolsep}{2.8pt}
\begin{tabularx}{\textwidth}{@{}YYYYYYYY@{}}
\toprule
\textbf{化学式} & \textbf{平均原子序数} & \textbf{平均电负性} & \textbf{电负性差} & \textbf{平均原子半径} & \textbf{空间群编号} & \textbf{晶胞体积} & \textbf{带隙/eV} \\
\midrule
GaN & 数值 & 数值 & 数值 & 数值 & 186 & 数值 & 3.40 \\
ZnO & 数值 & 数值 & 数值 & 数值 & 186 & 数值 & 3.37 \\
SrTiO\textsubscript{3} & 数值 & 数值 & 数值 & 数值 & 221 & 数值 & 3.20 \\
\bottomrule
\end{tabularx}
\end{formalTable}

在这个表中，平均原子序数、平均电负性、电负性差、平均原子半径、空间群编号和晶胞体积等列就是模型输入特征，带隙/eV 是模型要预测的目标值。这样，原始的材料记录就被转化为了机器学习可处理的特征表。


该流程可概括为下列步骤，并可结合图 \ref{fig:materials-data-feature} 理解：


\begin{lightquote}
\centering 原始材料记录 → 化学式解析 → 元素属性计算 → 结构字段数值化 → 特征表
\end{lightquote}

\begin{figure}[H]
\centering
\BookFigureImage{images_11/1.png}
\caption{材料数据整理与特征构建示意图}
\label{fig:materials-data-feature}
\end{figure}
\FloatBarrier

在这一环节，如果借助 AI 辅助处理，应像数据整理阶段一样把任务、已有字段、期望输出和限制条件说明清楚。可按下面方式提出任务。

\begin{samepage}
\begin{userbox}
我已经整理了一张无机光电材料带隙数据表，字段包括化学式、晶体结构、空间群、带隙/eV。现在准备用于机器学习建模，请帮我设计从原始表格到特征表的处理流程。

要求如下：
\begin{enumerate}
\item 说明如何把化学式解析为元素组成和原子分数；
\item 说明如何利用元素属性表计算平均原子序数、平均电负性、电负性差和平均原子半径；
\item 指出晶体结构、空间群等文字字段哪些需要进一步数值化；
\item 提醒哪些字段不宜直接作为模型输入，以及如何避免目标泄漏；
\item 不要虚构具体数值，只给出处理步骤和检查要点。
\end{enumerate}
\end{userbox}

\begin{aibox}
可把处理流程分为五步：

\begin{enumerate}
\item 化学式解析：将 GaN、ZnO、SrTiO\textsubscript{3} 等化学式拆分为元素种类和原子比例，再归一化为原子分数。
\item \textbf{元素属性匹配：}准备元素属性表，至少包含元素符号、原子序数、电负性、原子半径等字段，并与解析后的元素组成对应。
\item \textbf{成分描述符计算：}按原子分数加权计算平均原子序数、平均电负性、平均原子半径，并计算电负性差等差异类特征。
\item \textbf{结构字段转换：}晶体结构名称和空间群符号不宜直接输入多数模型，应转换为空间群编号、晶胞体积、密度或其他可计算结构描述符。
\item \textbf{建模前检查：}带隙/eV 是预测目标，不能作为输入特征；发光颜色等由带隙直接决定的字段也要谨慎使用，避免目标泄漏。
\end{enumerate}

最终得到的特征表应做到字段清楚、单位一致、输入特征和预测目标分离，并保留原始数据来源，便于后续核验。
\end{aibox}
\end{samepage}

在实践中，描述符不是越多越好。如果数据样本较少，而输入特征过多，模型可能会过拟合。对于入门任务，可先选择少量具有明确材料意义、且容易从数据库获得的描述符，例如平均原子序数、平均电负性、电负性差、平均原子半径、空间群编号和晶胞体积。等研究者理解基本流程后，再逐步加入更复杂的局域环境描述符。


还要避免目标泄漏。目标泄漏指的是把与预测目标直接相关、甚至由目标结果推导出来的信息放入输入特征中。例如，如果目标是预测带隙，就不应把“发光颜色”作为输入特征，因为发光颜色本身往往已经反映了带隙大小。这样的模型即使看起来预测准确，也没有真正学习材料成分和结构对带隙的影响。


\subsection{图表示与物理对称性}

上一节介绍了成分描述符和结构描述符。通过这些描述符，研究者可将化学式、元素属性、空间群和晶胞体积等信息转化为数值特征，用于机器学习模型预测材料性能。这种方法直观、容易理解，也适合入门学习。但它仍然有一个明显问题：很多描述符是人工设计的，模型看到的是人工预先整理好的特征，而不是材料结构本身。


成分描述符依赖人工设计，适合快速形成特征表，但难以完整保留原子排列、局域配位和相互作用距离等结构信息。带隙、形成能、弹性性质、离子扩散能力等，都与原子之间的连接关系、距离和局域环境密切相关。如果模型只使用平均电负性、平均原子半径、空间群编号等手工描述符，可能会丢失一部分精细结构信息。


图表示为这一问题提供了另一种思路：把材料看作原子节点与原子间边组成的图结构。在这张图中，原子是节点，原子之间的邻近关系或相互作用是边。这样，模型就不再只是读取一组人工设计的表格特征，而是可直接从材料的原子结构中学习信息。


继续以无机光电材料带隙预测为例，假设研究者要预测 GaN 的带隙。上一节中，研究者可使用平均电负性、电负性差、平均原子半径、空间群编号等描述符表示 GaN。但如果采用图表示，问题就变为：Ga 原子和 N 原子在晶体中如何排列，它们之间的距离是多少，每个原子周围有多少邻居。


在晶体图中，节点通常代表原子。对于 GaN 来说，图中的节点可以是 Ga 原子和 N 原子。每个节点需要带有节点特征，用来描述这个原子的基本信息。例如，一个 Ga 节点可包含 Ga 的原子序数、元素类型、电负性、原子半径、价电子数等信息；一个 N 节点也可包含 N 的对应元素属性。这样，模型就知道每个节点代表什么元素。


边通常代表原子之间的邻近关系。晶体中原子不是孤立存在的，材料性能往往取决于原子周围的局域环境。因此，需要根据一定规则判断哪些原子之间建立边。最常见的方法是设置一个截断半径：如果两个原子之间的距离小于某个阈值，就认为它们在图中相连。例如，可规定距离小于 4 \AA{} 的原子之间建立边。这样，Ga 原子周围一定距离内的 N 原子会与它相连，形成局域原子环境。


每条边也可带有边特征。边特征可包括原子间距离、相对位置、键类型或距离展开后的数值表示。对于带隙预测来说，原子间距离和局域配位环境会影响轨道重叠、成键方式和电子能带结构。因此，图表示比简单成分描述符更接近材料真实结构。


\subsubsection*{\textbf{1. 案例示范：GaN 晶体图构建}}

以 GaN 为例，构建晶体图可按照以下步骤进行。


第一步，获得 GaN 的晶体结构文件。这个文件可来自材料数据库，常见格式包括 CIF\footnote{Crystallographic Information File，晶体学信息文件。} 文件或其他晶体结构格式。结构文件中记录了晶格参数、原子种类和原子坐标。对于图表示来说，晶体结构文件是关键输入，因为它提供了原子在空间中的排列信息。


第二步，读取原子信息。程序需要从结构文件中读取每个原子的元素类型和坐标。例如，一个晶胞中可能包含 Ga 原子和 N 原子，每个原子都有对应的空间位置。此时，每个原子就可被看作图中的一个节点。


第三步，为节点添加特征。对于每个原子节点，可用元素属性作为节点特征。例如，Ga 节点可用原子序数、电负性、原子半径、价电子数等表示；N 节点也用同样类型的属性表示。这样，不同元素虽然都是“节点”，但节点特征不同，模型可区分它们的化学性质。


第四步，根据原子间距离建立边。设定一个截断半径，例如 4 \AA{}。程序计算所有原子之间的距离，如果两个原子距离小于截断半径，就在它们之间建立一条边。由于晶体具有周期性，计算距离时还要考虑相邻晶胞中的原子，而不是只看一个孤立晶胞。否则，边界附近原子的邻居关系会被错误截断。


第五步，为边添加特征。最简单的边特征是两个原子之间的距离。例如，Ga--N 距离可作为边特征输入模型。更复杂的方法会把距离展开成一组函数值，使模型更容易学习不同距离对材料性质的影响。


第六步，将整张图与目标性能对应起来。对于带隙预测任务，GaN 的晶体图作为模型输入，GaN 的带隙数值作为模型输出。模型学习的目标，就是从晶体图中预测带隙。换言之，模型不再只学习“平均电负性--带隙”的关系，而是学习“原子种类、局域连接、原子距离和晶体结构--带隙”的关系。


用表格形式表示，图数据可概括为表 \ref{tab:gan-crystal-graph-data-components} 所示内容。


\begin{formalTable}[htbp]
\caption{GaN 晶体图数据组成示例}
\label{tab:gan-crystal-graph-data-components}
\begin{tabularx}{\textwidth}{@{}YYY@{}}
\toprule
\textbf{图的组成} & \textbf{在 GaN 中对应什么} & \textbf{作用} \\
\midrule
节点 & Ga 原子、N 原子 & 表示晶体中的原子 \\
节点特征 & 原子序数、电负性、原子半径等 & 描述每个原子的元素属性 \\
边 & 截断半径内相邻原子之间的连接 & 表示局域原子关系 \\
边特征 & 原子间距离、相对位置等 & 描述原子之间的空间关系 \\
图标签 & GaN 的带隙值 & 模型需要预测的目标性能 \\
\bottomrule
\end{tabularx}
\end{formalTable}

晶体图构建涉及结构文件读取、周期性边界条件和邻接关系检查，适合借助 AI 辅助梳理代码流程。但提示词不应只写“构建晶体图”，而要明确输入文件、截断半径、输出内容和检查要求。

\begin{samepage}
\begin{userbox}
我有一个 GaN 的 CIF 晶体结构文件，准备用它构建用于带隙预测的晶体图。请帮我设计一段 Python 处理流程，用 pymatgen 读取结构，并根据 4 \AA{} 截断半径建立原子邻接关系。

要求如下：
\begin{enumerate}
\item 输出每个原子的元素种类和坐标；
\item 将原子作为节点，并说明可使用哪些节点特征；
\item 根据 4 \AA{} 截断半径生成边列表，并说明是否考虑周期性边界条件；
\item 输出节点数量、边数量和部分邻居关系，便于检查图是否合理；
\item 不要假设 CIF 文件一定正确，提醒需要核对哪些结构信息。
\end{enumerate}
\end{userbox}

\begin{aibox}
可按以下流程处理：

\begin{enumerate}
\item 用 pymatgen 读取 GaN.cif，先检查晶格参数、原子种类、原子坐标和空间群信息是否完整。
\item 将每个原子作为图中的一个节点。节点特征可从元素属性中选取，例如原子序数、电负性、原子半径和价电子数。
\item 使用结构对象的邻居搜索功能，在 4 \AA{} 截断半径内寻找相邻原子。对于晶体材料，应使用考虑周期性边界条件的方法，而不是只计算单个晶胞内部距离。
\item 将邻居关系保存为边列表，边特征可先使用原子间距离。后续如果用于图神经网络，还可把距离展开为多个函数值。
\item 检查节点数、边数和每类原子的邻居数量。如果某些原子没有邻居，或者边数量明显异常，应重新检查截断半径、CIF 文件和周期性设置。
\end{enumerate}

这段流程只能帮助完成结构到图的转换。截断半径是否合理、节点特征是否足够、图结构是否符合 GaN 的晶体环境，仍然需要结合晶体结构知识判断。
\end{aibox}
\end{samepage}

图表示之所以适合材料结构，还因为它能够更自然地满足材料中的一些物理要求。材料性质不应该因为坐标系移动、整体旋转或原子编号改变而发生变化。例如，同一个 GaN 晶体，如果研究者把整个晶体平移一段距离，它的带隙不应该改变；如果研究者重新编号原子，材料性质也不应该改变；如果只是改变观察角度，材料本身也没有发生变化。


这些要求可概括为物理对称性。对于材料机器学习，常见的物理对称性包括平移不变性、旋转不变性和置换不变性。


平移不变性是指材料整体移动后，性质不应改变。换言之，原子坐标整体加上同一个位移，材料的带隙、形成能等性质不应变化。因此，模型不应只依赖原子的绝对坐标，而应该更多使用原子之间的相对位置、距离或局域环境。


旋转不变性是指材料整体旋转后，性质不应改变。对于带隙、形成能、总能量等标量性质，材料旋转一个角度后，预测值应该保持不变。因此，很多模型会使用原子间距离这类旋转不变的量作为边特征。


置换不变性是指相同类型原子的编号顺序改变后，性质不应改变。在晶体结构文件中，原子编号只是人为排列顺序，不代表真实物理差异。例如，同一个晶体中两个等价 Ga 原子交换编号后，化学组成、局域环境和带隙都没有改变，模型输出也不应改变。图神经网络通过对邻居信息进行聚合，可在一定程度上满足这种要求。


除了不变性，还有一个更进一步的概念叫等变性。从概念上看，不变性要求某些输出在变换后保持不变，而等变性要求输出随着输入变换而按照物理规律相应变化。例如，能量是标量，材料整体旋转后能量不变；力是矢量，如果结构旋转，力的方向也应该随之旋转。可以把它理解为：同一个原子环境换了观察方向，能量数值不变，但每个原子受力的方向要跟着坐标系一起转动。对于带隙预测这类标量任务，理解不变性已经足够；对于后续机器学习势、力预测和分子动力学模拟，等变性会更加重要。


因此，图表示不仅是一种数据格式，更是一种更符合材料物理特征的表示方式。它把材料看作原子节点和相互作用边组成的结构，同时尽量尊重材料中的平移、旋转和置换对称性。相比只使用手工描述符，图表示能够保留更多局域结构信息，也为后续图神经网络模型奠定基础。


图表示并不意味着完全不需要人工判断。构建晶体图时，截断半径、节点特征、边特征表示和周期性边界条件处理都会影响模型结果。即使借助工具生成代码并检查图结构，上述关键设置是否合理，仍然需要结合材料结构和预测任务来判断。

\subsubsection*{\textbf{小结}}

本节说明材料信息需要转化为描述符、结构字段或晶体图，才能进入机器学习模型。
成分和结构描述符适合表格建模，图表示则更适合保留原子连接、距离和局域环境。
特征构建必须避免目标泄漏，并结合物理对称性、截断半径和边特征检查输入是否合理。

\subsection*{习题}

\begin{enumerate}[leftmargin=*]
\item 比较成分描述符和结构描述符的作用及适用场景。
\item 图表示如何保留比成分描述符更多的结构信息？
\item 成分描述符和图表示各有什么优势与局限？请结合带隙预测举例说明。
\end{enumerate}

\clearpage
\section{性能预测与尺度模拟}


\begin{sectionkeypoints}
\item 材料性能预测的核心，是让模型学习材料特征与目标性能之间的对应关系。硬度、带隙、形成能、弹性模量和熔点都可成为预测目标。
\item 描述符模型适合表格数据和入门任务，关键步骤包括整理特征表、划分训练测试集、训练模型、评价误差和检查特征重要性。
\item 图神经网络适合处理晶体结构数据，能够从原子节点、原子间边和局域环境中学习材料性质。
\item 机器学习势学习的是原子结构与能量、力之间的关系，可连接第一性原理计算和分子动力学模拟。
\item 模型预测和尺度模拟都必须进行可靠性检查。训练集外推、过拟合、数据泄漏、结构文件错误和模拟参数不合理都会影响结论。
\end{sectionkeypoints}

\subsection{描述符模型与硬度预测}
\label{subsec:descriptor-hardness-prediction}

有了材料描述符特征表之后，下一步是建立特征与性能之间的映射关系。也就是说，模型需要根据成分、结构、工艺或显微组织等输入特征，预测硬度、带隙、形成能、熔点等目标性能。


材料性能预测的基本思想，是让模型学习“材料特征”和“目标性能”之间的对应关系。材料特征可来自成分、结构、工艺和显微组织，目标性能可以是硬度、强度、带隙、形成能、熔点或耐腐蚀性能等。对于传统机器学习方法而言，材料首先要被转化成一张特征表，每一行代表一个材料样品，每一列代表一个可计算的特征，最后一列通常是需要预测的性能。


\subsubsection*{\textbf{1. 案例示范：高熵合金涂层硬度预测}}

本节只围绕一个具体任务展开：高熵合金涂层硬度预测。假设研究者希望根据高熵合金涂层的成分和组织信息，预测涂层硬度。这个例子适合说明传统机器学习预测流程，因为硬度是材料实验中常见的数值型性能，既受成分影响，也受相结构、晶粒尺寸和硬质相等组织因素影响。


对于一个高熵合金涂层样品，原始实验记录可能包括名义成分、制备工艺、XRD\footnote{X-ray Diffraction，X 射线衍射。} 相分析、显微组织观察和显微硬度测试结果。经过前面的描述符构建过程后，这些信息可整理为表 \ref{tab:hea-hardness-features} 所示形式。


\begin{formalTable}[htbp]
\caption{高熵合金涂层硬度预测特征表示示例}
\label{tab:hea-hardness-features}
\scriptsize
\setlength{\tabcolsep}{2.8pt}
\begin{tabularx}{\textwidth}{@{}YYYYYYYY@{}}
\toprule
\textbf{样品编号} & \textbf{原子半径差} & \textbf{平均电负性} & \textbf{电负性差} & \textbf{价电子浓度} & \textbf{是否含硬质相} & \textbf{平均晶粒尺寸/\(\mu\)m} & \textbf{硬度/HV} \\
\midrule
S1 & 数值 & 数值 & 数值 & 数值 & 0 & 数值 & 数值 \\
S2 & 数值 & 数值 & 数值 & 数值 & 1 & 数值 & 数值 \\
S3 & 数值 & 数值 & 数值 & 数值 & 1 & 数值 & 数值 \\
S4 & 数值 & 数值 & 数值 & 数值 & 1 & 数值 & 数值 \\
\bottomrule
\end{tabularx}
\end{formalTable}

表 \ref{tab:hea-hardness-features} 中，“原子半径差、平均电负性、电负性差、价电子浓度、是否含硬质相、平均晶粒尺寸”等列是模型输入特征，“硬度/HV”是模型需要预测的目标值。换成机器学习语言，就是把输入特征组成矩阵 X，把硬度组成目标向量 y。模型训练的目的，就是从 X 中学习到 y 的变化规律。


在这个任务中，硬度是一个连续数值，因此它属于回归任务。回归任务的目标不是判断材料属于哪一类，而是预测一个具体数值。例如，模型输入某个涂层样品的成分和组织特征后，输出预测硬度是多少 HV。与之相对，如果任务是判断材料是否耐腐蚀、是否稳定、是否属于某种相结构，那就更接近分类任务。


为了让这个过程更加具体，以下给出一个精简的机器学习代码框架。假设研究者已经整理好了一个高熵合金涂层硬度数据表，文件名为 hea\_\allowbreak{}coating\_\allowbreak{}hardness.csv。表格中包含若干材料描述符和目标性能“硬度/HV”。代码的目标是读取数据、划分训练集和测试集、训练随机森林回归模型，并评价预测误差。代码采用平均绝对误差（Mean Absolute Error，MAE）、均方根误差（Root Mean Squared Error，RMSE）和决定系数（Coefficient of Determination，\(R^2\)）评价预测误差与拟合程度。


\begin{lstlisting}[language=Python,escapeinside={(*@}{@*)}]
# (*@\textnormal{1. 导入常用库}@*)
import pandas as pd
from sklearn.model_selection import train_test_split
from sklearn.ensemble import RandomForestRegressor
from sklearn.metrics import mean_absolute_error, mean_squared_error, r2_score
# (*@\textnormal{2. 读取高熵合金涂层硬度数据表}@*)
data = pd.read_csv("hea_coating_hardness.csv")
# (*@\textnormal{3. 指定输入特征和预测目标}@*)
# (*@\textnormal{这些特征来自前面已经构建好的材料描述符}@*)
feature_cols = [
    "atomic_radius_diff",             # (*@\textnormal{原子半径差}@*)
    "mean_electronegativity",         # (*@\textnormal{平均电负性}@*)
    "electronegativity_diff",         # (*@\textnormal{电负性差}@*)
    "valence_electron_concentration", # (*@\textnormal{价电子浓度}@*)
    "has_hard_phase",                 # (*@\textnormal{是否含硬质相：0 表示否，1 表示是}@*)
    "grain_size"                      # (*@\textnormal{平均晶粒尺寸}@*)
]
X = data[feature_cols]
y = data["hardness_HV"]
# (*@\textnormal{4. 划分训练集和测试集}@*)
X_train, X_test, y_train, y_test = train_test_split(
    X,
    y,
    test_size=0.2,
    random_state=42
)
# (*@\textnormal{5. 建立并训练随机森林回归模型}@*)
model = RandomForestRegressor(
    n_estimators=200,
    random_state=42
)
model.fit(X_train, y_train)
# (*@\textnormal{6. 在测试集上进行硬度预测}@*)
y_pred = model.predict(X_test)
# (*@\textnormal{7. 评价模型预测效果}@*)
mae = mean_absolute_error(y_test, y_pred)
rmse = mean_squared_error(y_test, y_pred) ** 0.5
r2 = r2_score(y_test, y_pred)
print("MAE:", mae)
print("RMSE:", rmse)
print("R2:", r2)
# (*@\textnormal{8. 输出特征重要性}@*)
importance = pd.DataFrame({
    "feature": feature_cols,
    "importance": model.feature_importances_
}).sort_values(by="importance", ascending=False)
print(importance)
\end{lstlisting}

这段代码对应了材料性能预测的基本流程。首先，pd.read\_\allowbreak{}csv() 用于读取已经整理好的材料数据表。此处的数据表不是原始实验记录，而是已经经过描述符构建后的特征表。换言之，代码中的输入不是“FeCoCrNiTi”这样的化学式，而是“原子半径差、平均电负性、电负性差、价电子浓度、是否含硬质相、平均晶粒尺寸”等数值特征。


其中，X = data[feature\_\allowbreak{}cols] 表示模型输入，y = data["hardness\_\allowbreak{}HV"] 表示模型输出。对于本例来说，X 是涂层样品的成分和组织特征，y 是实验测得的硬度值。模型训练时要学习的，就是这些特征与硬度之间的关系。


train\_\allowbreak{}test\_\allowbreak{}split() 用于划分训练集和测试集。训练集用于让模型学习规律，测试集用于检验模型对新样品的预测能力。此处设置 test\_\allowbreak{}size=0.2，表示 80\% 的样品用于训练，20\% 的样品用于测试。对于材料数据来说，样本数量往往不大，因此测试集必须严格独立，不能参与模型训练。否则，模型评估结果会偏高，不能代表真实预测能力。


本例使用的模型是随机森林回归模型。随机森林由多棵决策树组成，能够处理非线性关系，对小规模表格数据比较友好，也可输出特征重要性。对于高熵合金涂层硬度预测来说，成分描述符、硬质相和晶粒尺寸与硬度之间通常不是简单线性关系，因此随机森林适合作为入门模型。此处选择随机森林并不意味着它一定最优，而是因为它便于研究者理解完整建模流程。


模型训练完成后，model.predict(X\_\allowbreak{}test) 会对测试集样品进行硬度预测。此时模型只看到测试样品的输入特征，不知道真实硬度。随后，研究者用真实硬度和预测硬度进行比较，计算 MAE、RMSE 和 R\textsuperscript{2}。


MAE 是平均绝对误差，表示模型预测值与真实值平均相差多少。例如，如果 MAE 为 30 HV，可理解为模型对测试样品的硬度平均预测误差约为 30 HV。RMSE 是均方根误差，对较大的预测误差更加敏感。如果少数样品预测偏差很大，RMSE 会明显增大。R\textsuperscript{2} 是决定系数，用来衡量预测值与真实值变化趋势的一致程度，越接近 1，说明模型预测效果越好。


对于材料研究来说，仅有误差指标还不够。宜进一步绘制“真实硬度--预测硬度”散点图。如果模型预测较好，散点应大致分布在 y = x 参考线附近。图中的坐标轴、单位、图例和参考线都要清楚，否则模型结果不便于与实验报告或论文表达衔接。


除了预测误差，随机森林模型还可输出特征重要性。代码中的 model.feature\_\allowbreak{}importances\_\allowbreak{} 可显示不同输入特征对硬度预测的贡献大小。如果结果显示“是否含硬质相”“平均晶粒尺寸”“原子半径差”等特征重要性较高，就可结合材料强化机制进行分析。例如，硬质相可能带来第二相强化，晶粒细化可能带来细晶强化，原子半径差可能与固溶强化和晶格畸变有关。


在代码实现阶段，AI 更适合承担“字段核对、代码改写、报错解释和绘图补充”等辅助工作，而不是替研究者判断哪些特征一定正确。可用下面的提示词约束任务边界。

\begin{samepage}
\begin{userbox}
我正在用随机森林预测高熵合金涂层硬度，CSV 表头包括原子半径差、平均电负性、电负性差、价电子浓度、是否含硬质相、晶粒尺寸和硬度/HV。请帮我检查这段建模流程是否清楚，并补充测试集真实硬度与预测硬度散点图的绘图步骤。

要求如下：
\begin{enumerate}
\item 判断哪些列可作为输入特征，哪一列是预测目标；
\item 提醒哪些字段需要确认数据来源或实验条件；
\item 说明如果出现列名不存在、缺失值或字符串字段无法输入模型，应如何检查；
\item 补充绘制真实值--预测值散点图和 y=x 参考线的思路；
\item 不要替我判断模型一定可靠，只给出代码检查和结果呈现建议。
\end{enumerate}
\end{userbox}

\begin{aibox}
表 \ref{tab:hea-hardness-features} 中，原子半径差、平均电负性、电负性差、价电子浓度、是否含硬质相和晶粒尺寸可作为候选输入特征，硬度/HV 是预测目标，不应同时放入输入特征。

建模前需要重点检查三类问题：
\begin{enumerate}
\item \textbf{字段是否存在：}代码中指定的输入特征名称必须与 CSV 表头完全一致，否则会出现列名不存在的错误。
\item \textbf{数据是否可输入模型：}输入特征应为数值型；如果存在空值、文字标签或单位混用，需要先清洗或编码。
\item \textbf{材料含义是否可靠：}是否含硬质相应来自 XRD、SEM\footnote{Scanning Electron Microscopy，扫描电子显微镜。} 或其他表征结果，晶粒尺寸应说明统计方法，硬度值应记录测试载荷和保载时间。
\end{enumerate}

模型评价时，除 MAE、RMSE 和 R\textsuperscript{2} 外，可绘制测试集真实硬度与预测硬度散点图，并添加 y=x 参考线。如果散点明显偏离参考线，说明模型在部分样品上的预测误差较大，需要回到数据质量、样本数量和特征设计中继续检查。
\end{aibox}
\end{samepage}

在实际建模时，还应特别注意过拟合。材料数据往往样本少、特征多，如果只有十几个样品却加入几十个描述符，模型可能在训练集上表现很好，但在测试集上预测很差。这说明模型可能只是记住了已有样品，而没有真正学到可推广的规律。为了降低过拟合风险，可减少无关特征，采用交叉验证，增加样本数量，或者选择更简单、更稳定的模型。是否属于过拟合，仍然要结合训练误差、测试误差、数据规模和材料任务共同判断。


数据泄漏同样需要警惕。预测硬度时，不应把由硬度结果直接推导出的字段作为输入特征。例如，“是否高硬度样品”“硬度等级”这类字段不能作为输入，因为它们本身已经包含了目标信息。否则模型看起来预测很准，实际上并没有从成分和组织中学习硬度规律。


\subsection{图神经网络与性质预测}

上一节以高熵合金涂层硬度预测为例，介绍了传统机器学习的基本流程。那种方法的特点是：研究者首先根据材料知识人工构造描述符，例如原子半径差、电负性差、价电子浓度、是否含硬质相和晶粒尺寸，然后再把这些描述符输入随机森林等模型中预测硬度。这种方法直观、容易操作，也适合材料数据量较少的入门任务。


但是，传统描述符方法也有局限。材料性能不仅取决于几个平均特征，还与原子在晶体中的排列方式、局域配位环境、原子间距离和结构稳定性有关。尤其对于高熵合金涂层这类复杂体系，材料中可能同时存在 FCC\footnote{Face-Centered Cubic，面心立方结构。} 固溶体、BCC\footnote{Body-Centered Cubic，体心立方结构。} 相、碳化物或硼化物等不同结构。仅用“是否含硬质相”“平均晶粒尺寸”这样的表格特征，很难完整表达原子尺度结构对性能的影响。


图神经网络提供了另一种建模思路。它不再把材料完全压缩成一行人工描述符，而是把材料看作由原子节点和原子间关系组成的图结构。在晶体图中，每个原子是一个节点，原子之间的邻近关系是边，节点和边都可带有特征。模型通过在图上进行信息传递，学习局域原子环境与材料性质之间的关系。


这里要限定问题边界：高熵合金涂层是一个多尺度、多相组织体系，真实涂层中包含晶粒、晶界、第二相、孔隙、界面和残余应力等复杂因素。图神经网络不宜简单地把一张 SEM 图或一个涂层样品直接转化成完整晶体图，然后准确预测整体硬度。更合理的入门方式是：选取涂层中的代表性晶体相或候选晶体结构，例如 FCC 固溶体、BCC 固溶体或 TiC 等硬质相，把它们转化为晶体图，学习结构与形成能、弹性模量或稳定性等性质之间的关系。然后再把这些原子尺度信息作为理解涂层硬度和强化机制的补充。


这里的重点不是用图神经网络直接预测整个涂层硬度，而是说明如何把涂层中的代表性晶体结构转化为图，并用图神经网络预测与性能相关的材料性质。


假设研究者研究的是一种 FeCoCrNiTi 高熵合金涂层。实验表明，涂层中可能存在 FCC 固溶体、BCC 相和 TiC 等硬质相。上一节中，研究者用人工描述符预测涂层硬度；现在研究者从更微观的角度思考：如果想理解某个相为什么可能更稳定、为什么某种硬质相会提高硬度，就需要进一步关注原子尺度结构。图神经网络正适合处理这类晶体结构信息。


对于一个代表性晶体结构，原始数据通常来自 CIF 文件或数据库结构文件。一个结构文件中记录了晶格参数、原子种类和原子坐标。图神经网络的第一步，就是把这个晶体结构转化为晶体图。


\subsubsection*{\textbf{1. 案例示范：TiC 晶体图性质预测}}

以 TiC 硬质相为例，它在高熵合金涂层中常被视为提高硬度的重要相之一。构建 TiC 晶体图时，可按照以下步骤进行。


第一步，读取晶体结构文件。结构文件中包含 Ti 和 C 原子的种类、坐标以及晶格参数。对于模型来说，这些信息是图结构的基础。


第二步，把原子转化为节点。TiC 中的 Ti 原子和 C 原子都可看作图中的节点。每个节点需要包含节点特征，例如原子序数、电负性、原子半径、价电子数等。这样，模型能够区分 Ti 节点和 C 节点，并理解它们具有不同化学属性。


第三步，根据原子间距离建立边。在晶体中，一个原子的性质受到周围邻居原子的影响。因此，可设置一个截断半径，例如 5 \AA{}。如果两个原子之间的距离小于这个截断半径，就在它们之间建立一条边。这样，模型就可知道哪些原子彼此相邻。


第四步，为边添加特征。最常见的边特征是原子间距离。距离越近，原子之间相互作用通常越强；距离不同，键合环境也不同。因此，边特征可帮助模型学习局域结构对性质的影响。


第五步，把整个晶体图与目标性质对应起来。如果目标是预测形成能，那么每一个晶体图对应一个形成能数值；如果目标是预测弹性模量，那么每一个晶体图对应一个弹性模量数值。模型训练的目标，就是学习“晶体图结构”与“目标性质”之间的关系。


该流程可概括为下列步骤，并可结合图 \ref{fig:crystal-graph-prediction} 理解：


\begin{lightquote}
\centering 晶体结构文件 → 原子节点 → 原子间边 → 节点特征和边特征 → 晶体图 → 性质预测
\end{lightquote}

\begin{figure}[H]
\centering
\BookFigureImage{images_11/2.png}
\caption{晶体图表示与性能预测示意图}
\label{fig:crystal-graph-prediction}
\end{figure}
\FloatBarrier

图 \ref{fig:crystal-graph-prediction} 所示流程中，图数据的组成可进一步概括为表 \ref{tab:crystal-graph-data-components}。


\begin{formalTable}[htbp]
\caption{晶体图数据组成示例}
\label{tab:crystal-graph-data-components}
\begin{tabularx}{\textwidth}{@{}YY@{}}
\toprule
\textbf{图数据组成} & \textbf{在 TiC 或高熵合金相关晶体结构中表示什么} \\
\midrule
节点 & Ti、C 或 Fe、Co、Cr、Ni、Ti 等原子 \\
节点特征 & 原子序数、电负性、原子半径、价电子数等 \\
边 & 截断半径内的原子邻近关系 \\
边特征 & 原子间距离或距离展开特征 \\
图标签 & 形成能、弹性模量、体模量或其他目标性质 \\
\bottomrule
\end{tabularx}
\end{formalTable}

图神经网络的核心是信息传递。每个原子节点在初始时只包含自己的元素属性，例如 Ti 节点知道自己是 Ti，C 节点知道自己是 C。经过一轮信息传递后，一个节点会接收周围邻居节点的信息。经过多轮信息传递后，节点就不仅包含自身信息，还包含局域原子环境信息。最后，模型把所有节点的信息汇总成整个晶体的表示，再输出目标性质。


这个过程可用材料语言理解：模型不只是看“这个材料含有哪些元素”，而是进一步看“这些元素在空间中如何排列、每个原子周围是什么邻居、原子间距离是多少、局域环境是否稳定”。这正是图神经网络相比传统描述符方法的重要优势。


为了让这个过程更具体，以下给出一个教学用的精简代码框架。它的目的不是提供完整可直接运行的项目，而是帮助研究者理解图神经网络建模的基本步骤。


\begin{lstlisting}[language=Python,escapeinside={(*@}{@*)}]
# (*@\textnormal{教学用简化框架：从晶体结构构建图并预测性质}@*)
# (*@\textnormal{实际运行时需要安装 pymatgen、torch、torch\_\allowbreak{}geometric 等库}@*)
from pymatgen.core import Structure
import torch
# (*@\textnormal{1. 读取晶体结构文件，例如 TiC.cif}@*)
structure = Structure.from_file("TiC.cif")
# (*@\textnormal{2. 设置截断半径}@*)
cutoff = 5.0
# (*@\textnormal{3. 构建节点特征：此处仅用原子序数作为示例}@*)
node_features = []
for site in structure:
    atomic_number = site.specie.Z
    node_features.append([atomic_number])
x = torch.tensor(node_features, dtype=torch.float)
# (*@\textnormal{4. 构建边：寻找截断半径内的邻居原子}@*)
edge_index = []
edge_attr = []
for i, site in enumerate(structure):
    neighbors = structure.get_neighbors(site, cutoff)
    for neighbor in neighbors:
        j = neighbor.index
        distance = neighbor.nn_distance
        edge_index.append([i, j])
        edge_attr.append([distance])
edge_index = torch.tensor(edge_index, dtype=torch.long).t().contiguous()
edge_attr = torch.tensor(edge_attr, dtype=torch.float)
# (*@\textnormal{5. 图标签：例如该结构的形成能或弹性模量}@*)
# (*@\textnormal{此处用一个占位数值表示，实际应来自数据库或计算结果}@*)
y = torch.tensor([target_property], dtype=torch.float)
# (*@\textnormal{至此，一个晶体结构已经被转化为图数据：}@*)
# (*@\textnormal{x 表示节点特征}@*)
# (*@\textnormal{edge\_\allowbreak{}index 表示边连接关系}@*)
# (*@\textnormal{edge\_\allowbreak{}attr 表示边特征}@*)
# (*@\textnormal{y 表示目标性质}@*)
\end{lstlisting}

这段代码展示了晶体图构建的基本思想。Structure.from\_\allowbreak{}file("TiC.cif") 用来读取晶体结构文件；node\_\allowbreak{}features 保存每个原子的节点特征；edge\_\allowbreak{}index 保存哪些原子之间存在边；edge\_\allowbreak{}attr 保存边特征，例如原子间距离；y 是模型要预测的目标性质。


实际使用时，节点特征不应只包含原子序数，还可加入电负性、原子半径、价电子数等元素属性。边特征也不一定只用一个距离数值，可把距离展开成一组函数，以便模型更好学习不同距离范围内的相互作用。对于不同材料任务，目标 y 也可不同。如果关注结构稳定性，可预测形成能；如果关注力学性能，可预测弹性模量或体模量；如果进一步结合实验数据，也可尝试预测与硬度相关的性质。


如果借助 AI 辅助构建 TiC 晶体图，提示词应突出“代表性晶体相”和“目标性质”两个信息，避免把图神经网络误写成直接预测整个涂层硬度的工具。示例任务可写成如下形式。

\begin{samepage}
\begin{userbox}
我想以高熵合金涂层中的 TiC 硬质相为例，构建一个用于性质预测的晶体图。请帮我设计从 TiC.cif 到图数据的处理流程，目标性质可以是形成能、弹性模量或体模量。

要求如下：
\begin{enumerate}
\item 说明如何读取 TiC.cif，并把 Ti 和 C 原子转化为节点；
\item 说明节点特征可选择哪些元素属性；
\item 根据 5 \AA{} 截断半径建立边，并说明为什么要考虑周期性边界条件；
\item 说明边特征和图标签分别是什么；
\item 提醒这个晶体图只能表示 TiC 这一代表性相，不宜直接代表整个高熵合金涂层硬度。
\end{enumerate}
\end{userbox}

\begin{aibox}
可按以下流程构建图数据：

\begin{enumerate}
\item \textbf{读取结构：}使用 pymatgen 读取 TiC.cif，检查晶格参数、Ti/C 原子坐标和结构完整性。
\item \textbf{构建节点：}把每个 Ti 原子和 C 原子作为节点。节点特征可包括原子序数、电负性、原子半径和价电子数。
\item \textbf{构建边：}使用 5 \AA{} 截断半径寻找邻近原子。由于 TiC 是周期晶体，邻居搜索必须考虑相邻晶胞，否则会漏掉晶胞边界附近的邻居。
\item \textbf{添加边特征：}可先使用原子间距离作为边特征，后续再根据模型需要使用距离展开特征。
\item \textbf{设置图标签：}如果训练目标是形成能，图标签就是形成能；如果目标是弹性模量或体模量，图标签应来自数据库、第一性原理计算或可靠实验/计算资料。
\end{enumerate}

这里需要专门说明，TiC 晶体图只能帮助理解涂层中某个硬质相的结构性质。真实高熵合金涂层硬度还受到晶粒尺寸、相分布、界面、孔隙和残余应力影响，不能由单个 TiC 晶体图直接决定。
\end{aibox}
\end{samepage}

与传统描述符模型相比，图神经网络的优势在于它能更直接地保留原子结构信息。传统模型需要人工预先设计“平均电负性”“原子半径差”等特征，而图神经网络可从原子种类、原子间距离和局域连接关系中自动学习结构信息。对于晶体材料来说，这种表示方式更接近材料本身。


但是，图神经网络也不是万能的。首先，它需要可靠的晶体结构输入。如果 CIF 文件错误、原子占位不清楚或结构不完整，构建出的图也会有问题。其次，它通常需要较多训练数据。如果只有少量结构样本，图神经网络也可能过拟合。再次，对于真实高熵合金涂层来说，整体硬度还受到晶粒尺寸、相分布、孔隙率、界面结合和残余应力等显微组织因素影响，单个晶体图不能完全代表整个涂层。


因此，在高熵合金涂层研究中，图神经网络更适合作为结构层面的补充工具。它可帮助研究者理解某个候选相、硬质相或代表性晶体结构的稳定性和力学相关性质；而整体涂层硬度仍然需要结合上一节中的表格特征模型、显微组织数据和实验测试结果共同分析。


\subsection{原子间势与熔点模拟}

前两节讨论的是材料性能预测。无论是传统机器学习模型，还是图神经网络模型，它们通常都是输入一个材料样品或晶体结构，然后输出一个性能数值，例如硬度、形成能、带隙或弹性模量。这类模型回答的是“这个材料可能具有什么性质”。


但是，材料科学中还有另一类重要问题：材料在原子尺度上如何运动、如何扩散、如何相变、如何熔化。例如，超高温陶瓷在高温下能否保持结构稳定？原子在升温过程中是否发生剧烈扩散？材料在什么温度附近开始熔化？这些问题不是简单预测一个静态数值就能完全回答的，而需要模拟原子随时间的运动过程。


分子动力学可用来研究这类原子尺度动态行为。在分子动力学中，研究者根据原子之间的相互作用计算每个原子的受力，再根据力更新原子位置，从而模拟材料随时间演化的过程。问题在于，原子间相互作用如何计算。第一性原理分子动力学精度较高，但计算成本高，难以处理大体系和长时间尺度；传统经验势速度快，但对复杂材料体系的适用性有限。机器学习势正是在这种背景下发展起来的。


机器学习势的核心思想，是用机器学习模型学习第一性原理计算得到的势能面。从概念上看，第一性原理计算可给出某一组原子结构的能量、原子受力和应力；机器学习势则利用这些数据学习“原子结构--能量/力”之间的关系。训练完成后，机器学习势可像传统势函数一样快速计算能量和力，从而支持更大规模、更长时间的分子动力学模拟。


\subsubsection*{\textbf{1. 案例示范：TiC 熔化行为模拟}}

本节只围绕一个任务展开：用机器学习势辅助模拟 TiC 超高温陶瓷的熔化行为。TiC 是典型的超高温陶瓷之一，具有高熔点、高硬度和良好的高温稳定性，也常作为高熵合金涂层或复合涂层中的硬质增强相。通过这个例子，可说明机器学习势如何从第一性原理数据出发，进一步服务于高温分子动力学模拟。


首先需要明确，机器学习势与前面性能预测模型的目标不同。前面随机森林模型预测的是硬度，图神经网络可预测形成能或弹性模量；而机器学习势要预测的是原子结构对应的总能量和每个原子的受力。它不是直接告诉研究者“TiC 的熔点是多少”，而是先学会在不同原子构型下计算能量和力，然后通过分子动力学模拟升温过程，再从模拟结果中判断熔化行为。


对于 TiC 熔点模拟，可按照以下流程理解。


第一步，准备 TiC 的初始晶体结构。研究者需要获得 TiC 的晶体结构文件，例如 CIF 文件。这个结构文件包含 Ti 和 C 原子的排列方式、晶格参数和空间结构。初始结构通常可来自数据库、文献或已有计算结果。


第二步，生成不同构型的训练数据。机器学习势不应只看一个理想晶体结构，因为真实高温模拟中原子会振动、偏移，甚至出现缺陷和局部无序。因此，需要构造一批不同的 TiC 原子构型，例如平衡晶体结构、轻微扰动结构、不同体积结构、不同温度下短时间分子动力学得到的结构等。这样，训练数据才能覆盖模型后续可能遇到的原子环境。


第三步，用第一性原理计算这些构型的能量和力。对于每一个 TiC 构型，需要通过第一性原理计算得到总能量、每个原子的力以及必要时的应力。这一步计算成本较高，但只需要对训练集中的有限构型进行。可理解为：第一性原理计算负责提供高质量“标准答案”，机器学习势负责学习这些标准答案中的规律。


第四步，训练机器学习势。训练时，模型输入原子种类和原子坐标，输出预测能量和原子力。训练目标是让模型预测的能量和力尽可能接近第一性原理计算结果。常见机器学习势包括 Behler--Parrinello 神经网络势、高斯近似势、DeepMD\footnote{Deep Potential Molecular Dynamics，深度势分子动力学。}、MACE\footnote{Message Passing Atomic Cluster Expansion，消息传递原子簇展开模型。}、NequIP\footnote{Neural Equivariant Interatomic Potentials，神经等变原子间势。}、M3GNet\footnote{Materials 3-body Graph Network，材料三体图网络。}、CHGNet\footnote{Crystal Hamiltonian Graph Neural Network，晶体哈密顿量图神经网络。} 等。教材入门阶段不需要详细展开每种模型的数学细节，重点是理解它们都在学习原子结构与能量/力之间的关系。


第五步，验证机器学习势是否可靠。训练完成后，不宜直接用于高温模拟，而要先在测试构型上验证预测误差。常见做法是比较模型预测能量与第一性原理能量之间的误差，以及模型预测力与第一性原理力之间的误差。如果误差较大，说明模型还没有学好 TiC 的势能面，需要增加训练数据、调整模型参数或补充高温构型。


第六步，用训练好的机器学习势进行分子动力学模拟。此时，机器学习势可快速计算每一步原子的能量和力，进而模拟 TiC 在不同温度下的原子运动。为了观察熔化行为，可逐步升高温度，例如从较低温度开始，逐步升温到高温区间，记录能量、体积、径向分布函数、均方位移等随温度变化的情况。


第七步，根据模拟结果判断熔化行为。材料熔化时，原子排列会从有序晶体逐渐转变为无序液态结构。模拟中可通过多个指标辅助判断：能量--温度曲线是否出现突变，体积是否发生明显变化，径向分布函数的长程有序峰是否消失，均方位移是否显著增加。通过这些现象，可估计 TiC 的熔化温度范围。


这个过程可概括为：


\begin{lightquote}
\centering TiC 晶体结构 → 多种原子构型 → 第一性原理能量和力 → 训练机器学习势 → 验证误差 → 分子动力学升温模拟 → 分析熔化行为
\end{lightquote}

为了让研究者更清楚地理解流程，以下给出一个教学用的简化代码框架。它不是完整可直接运行的工程代码，因为真实机器学习势训练通常需要专门软件和大量计算资源。此处的目的，是帮助理解每一步在做什么。


\begin{lstlisting}[language=Python,escapeinside={(*@}{@*)}]
# (*@\textnormal{教学用简化框架：机器学习势辅助 TiC 熔化模拟}@*)
# (*@\textnormal{实际研究中可使用 DeepMD、MACE、NequIP、CHGNet、M3GNet 等工具}@*)
# (*@\textnormal{1. 读取 TiC 初始晶体结构}@*)
structure = read_structure("TiC.cif")
# (*@\textnormal{2. 生成训练构型}@*)
# (*@\textnormal{包括平衡结构、扰动结构、不同体积结构和高温采样结构}@*)
train_structures = generate_configurations(
    structure,
    perturb=True,
    volume_changes=True,
    temperature_sampling=True
)
# (*@\textnormal{3. 对训练构型进行第一性原理计算}@*)
# (*@\textnormal{得到每个构型的能量、原子力和应力}@*)
dft_dataset = []
for config in train_structures:
    energy, forces, stress = run_dft_calculation(config)
    dft_dataset.append({
        "structure": config,
        "energy": energy,
        "forces": forces,
        "stress": stress
    })
# (*@\textnormal{4. 训练机器学习势}@*)
ml_potential = train_machine_learning_potential(
    dataset=dft_dataset,
    target_properties=["energy", "forces", "stress"]
)
# (*@\textnormal{5. 在测试构型上验证模型误差}@*)
test_results = evaluate_potential(
    model=ml_potential,
    test_dataset="test_configurations"
)
print(test_results)
# (*@\textnormal{6. 使用训练好的机器学习势进行分子动力学升温模拟}@*)
md_results = run_molecular_dynamics(
    structure=structure,
    potential=ml_potential,
    temperature_schedule=[1000, 1500, 2000, 2500, 3000, 3500, 4000],
    simulation_time="several ps or longer"
)
# (*@\textnormal{7. 分析熔化行为}@*)
plot_energy_temperature(md_results)
plot_radial_distribution_function(md_results)
plot_mean_square_displacement(md_results)
\end{lstlisting}

这段框架中，read\_\allowbreak{}structure() 表示读取 TiC 的晶体结构；generate\_\allowbreak{}configurations() 表示生成不同原子构型；run\_\allowbreak{}dft\_\allowbreak{}calculation() 表示用第一性原理计算能量和力；train\_\allowbreak{}machine\_\allowbreak{}learning\_\allowbreak{}potential() 表示训练机器学习势；run\_\allowbreak{}molecular\_\allowbreak{}dynamics() 表示用训练好的势进行分子动力学模拟；最后通过能量曲线、径向分布函数和均方位移分析熔化行为。


与前两节代码不同，此处给出的是流程框架，而不是几十行就能完成的入门脚本。原因是机器学习势涉及第一性原理计算、模型训练和分子动力学模拟，真实流程需要较多软件和计算资源。但这个框架可帮助研究者理解机器学习势在材料研究中的位置：它不是直接替代实验，也不是简单回归一个性能值，而是连接第一性原理计算和大尺度原子模拟的桥梁。


在机器学习势和分子动力学任务中，AI 辅助更适合用于解释文件字段、检查流程是否完整和整理输出指标，而不宜直接判断熔点结论。下面的写法更适合用于流程核查。

\begin{samepage}
\begin{userbox}
我想用机器学习势辅助模拟 TiC 的熔化行为。已有 TiC 初始结构，计划生成扰动构型和高温采样构型，用第一性原理计算能量与力，再训练机器学习势并进行升温分子动力学模拟。请帮我检查这个流程是否完整，并说明每类输出指标的物理意义。

要求如下：
\begin{enumerate}
\item 说明训练数据应包含哪些 TiC 原子构型；
\item 说明第一性原理计算需要记录哪些目标量；
\item 检查机器学习势训练后应如何验证能量和力误差；
\item 说明升温分子动力学中能量、体积、径向分布函数和均方位移分别反映什么；
\item 明确提醒不应只根据单一指标判断熔点。
\end{enumerate}
\end{userbox}

\begin{aibox}
这个流程可分为四个检查重点：

\begin{enumerate}
\item \textbf{训练构型覆盖范围：}除平衡 TiC 晶体结构外，还应包含轻微扰动结构、不同体积结构、不同温度采样结构，以及接近高温无序状态的构型。否则模型可能只适用于低温平衡结构。
\item \textbf{第一性原理标注目标：}每个构型至少应记录总能量和每个原子的力；如果后续模拟涉及体积变化，还应考虑应力信息。
\item \textbf{机器学习势验证：}训练后应在独立测试构型上比较预测能量、预测力与第一性原理结果的误差，不应只看训练集误差。
\item \textbf{熔化行为分析：}能量--温度曲线突变可能反映相变，体积突变可能反映结构膨胀，径向分布函数长程有序峰消失说明晶体有序性降低，均方位移显著增大说明原子扩散增强。
\end{enumerate}

熔点判断应综合多个指标，并结合体系大小、升温速率、模拟时间和训练数据覆盖范围共同分析。单独依靠 AI 对曲线的文字解释，不能作为可靠的材料结论。
\end{aibox}
\end{samepage}

需要特别强调的是，机器学习势的可靠性高度依赖训练数据。如果训练集中只包含低温平衡结构，而模拟时却用于高温熔化过程，模型很可能外推到训练分布之外，预测结果就不可靠。因此，在 TiC 熔点模拟中，训练数据应尽量包含高温扰动构型、不同体积构型和接近熔化状态的构型。否则，模型可能在低温结构上误差很小，却在高温模拟中产生错误行为。


此外，机器学习势不能完全替代第一性原理计算和实验验证。它的优势是加速模拟，但前提是训练数据可靠、验证误差合理、模拟条件设置正确。对于超高温陶瓷这类材料，熔点模拟还可能受到模拟体系大小、升温速率、时间尺度和缺陷状态影响。因此，模拟得到的熔化温度应被理解为计算预测结果，需要与实验数据或更高精度计算进行对照。

\subsubsection*{\textbf{小结}}

本节围绕硬度预测、图神经网络性质预测和机器学习势展开。
描述符模型用于表格化的宏观性能预测，图神经网络用于原子结构关联的性质预测，机器学习势用于原子尺度动力学模拟。
性能指标、数据泄漏、过拟合、结构文件质量和训练构型覆盖都会影响模型结论的可靠性。

\subsection*{习题}

\begin{enumerate}[leftmargin=*]
\item 比较描述符模型、图神经网络和机器学习势的适用任务。
\item 机器学习势适合哪些模拟场景？训练覆盖为何重要？
\item 为高熵合金涂层硬度预测设计输入特征，并说明依据。
\end{enumerate}

\clearpage
\section{模型解释与知识嵌入}


\begin{sectionkeypoints}
\item 可解释性分析关注模型训练之后的问题，即模型主要依赖哪些特征、这些特征如何影响预测结果，以及解释是否符合材料机理。
\item 特征重要性和 SHAP\footnote{SHapley Additive exPlanations，一种基于博弈论的可解释性分析方法。} 可帮助把机器学习输出转化为材料语言，但它们反映的是当前数据和当前模型中的关系，不等同于材料真理。
\item 物理知识嵌入强调在建模前和建模过程中主动引入材料规律，包括特征选择、数据筛选、目标泄漏检查和结果合理性判断。
\item 对高熵合金涂层硬度预测而言，原子半径差、晶粒尺寸、硬质相和孔隙率等特征应与固溶强化、细晶强化、第二相强化和组织致密化联系起来理解。
\item 一个可靠的材料人工智能模型，不仅要有较小预测误差，还要能被实验表征、物理机理和材料常识共同支持。
\end{sectionkeypoints}

\subsection{模型解释与特征归因}

前面几节已经介绍了材料性能预测的基本方法。通过描述符和机器学习模型，研究者可根据材料成分、结构和组织信息预测硬度、带隙、形成能、熔点等性能。对于工程应用来说，预测准确当然重要；但对于材料科学研究来说，仅仅得到一个预测数值还不够。研究者还需要进一步知道：模型为什么给出这样的预测？它主要依据了哪些材料特征？这些特征是否符合已有材料机理？


这正是模型可解释性要解决的问题。可解释性关注的是模型在预测时主要依赖哪些输入特征，以及这些特征如何影响预测结果。对于材料科学而言，可解释性尤其重要，因为材料研究的目标不仅是预测某个样品性能，更是提炼可理解、可验证、可迁移的材料规律。


本节继续围绕前面的高熵合金涂层硬度预测任务展开。假设研究者已经训练了一个随机森林模型，用来预测不同高熵合金涂层的硬度。模型的输入特征包括原子半径差、平均电负性、电负性差、价电子浓度、是否含硬质相和平均晶粒尺寸，输出目标是硬度/HV。现在的问题不再是“模型能不能预测硬度”，而是“模型主要根据什么判断某个涂层硬度较高”。


如果只看预测误差，例如 MAE、RMSE 和 R\textsuperscript{2}，研究者只能知道模型预测效果大致如何，却不知道模型内部学到了什么。例如，一个模型可能在测试集上误差较小，但它可能主要依赖了某个偶然相关的字段，而不是真正反映材料强化机制的特征。这样的模型即使暂时预测准确，也很难让研究者放心。因此，在材料人工智能中，模型解释往往和模型评价同样重要。


对于随机森林这类模型，最常见的解释方法是特征重要性分析。特征重要性可粗略告诉研究者：模型在预测硬度时，哪些输入特征贡献更大。假设训练后的模型输出表 \ref{tab:hardness-feature-importance} 所示结果。


\begin{formalTable}[htbp]
\caption{随机森林硬度预测特征重要性示例}
\label{tab:hardness-feature-importance}
\begin{tabularx}{\textwidth}{@{}YY@{}}
\toprule
\textbf{特征} & \textbf{重要性} \\
\midrule
是否含硬质相 & 0.34 \\
平均晶粒尺寸 & 0.26 \\
原子半径差 & 0.18 \\
价电子浓度 & 0.13 \\
电负性差 & 0.06 \\
平均电负性 & 0.03 \\
\bottomrule
\end{tabularx}
\end{formalTable}

这个结果说明，在当前数据集和模型中，“是否含硬质相”和“平均晶粒尺寸”对硬度预测影响最大。接下来，研究者需要把这个模型结果转化为材料语言。含硬质相的涂层往往可能通过第二相强化提高硬度；晶粒尺寸越小，晶界数量越多，可能带来细晶强化；原子半径差较大时，晶格畸变增强，也可能产生固溶强化。因此，如果模型认为这些特征重要，并且这种排序与材料强化机制基本一致，就说明模型结果具有一定可解释性。


但是，特征重要性只能告诉研究者“哪些特征重要”，不宜直接告诉研究者“这个特征让硬度升高还是降低”。例如，平均晶粒尺寸重要，并不意味着晶粒尺寸越大硬度越高，也不意味着晶粒尺寸越小一定硬度越高。它只能说明模型在预测时经常用到这个特征。为了进一步分析特征对预测结果的正负影响，可引入 SHAP 方法。


SHAP 可理解为一种更细致的解释工具。它不仅能分析某个特征整体上是否重要，还可分析在某一个样品中，某个特征是把预测硬度往高处推，还是往低处拉。例如，对于一个含 TiC 硬质相、晶粒较细的高熵合金涂层样品，SHAP 分析可能显示：“是否含硬质相”对预测硬度有正贡献，“平均晶粒尺寸较小”也对预测硬度有正贡献。这种结果就可和第二相强化、细晶强化联系起来解释。


为了更清楚地说明可解释性分析，以下给出一个精简代码框架。假设研究者已经在第 \ref{subsec:descriptor-hardness-prediction} 节中训练好了随机森林硬度预测模型 model，并且已经得到输入特征 X\_\allowbreak{}train、X\_\allowbreak{}test 和特征名称 feature\_\allowbreak{}cols。以下代码用于输出随机森林特征重要性，并进一步使用 SHAP 分析模型预测结果。


\begin{lstlisting}[language=Python,escapeinside={(*@}{@*)}]
# (*@\textnormal{1. 随机森林特征重要性}@*)
importance = pd.DataFrame({
    "feature": feature_cols,
    "importance": model.feature_importances_
}).sort_values(by="importance", ascending=False)
print(importance)
# (*@\textnormal{2. 绘制特征重要性柱状图}@*)
plt.figure(figsize=(6, 4))
plt.barh(importance["feature"], importance["importance"])
plt.xlabel("Feature importance")
plt.ylabel("Feature")
plt.title("Feature importance for hardness prediction")
plt.gca().invert_yaxis()
plt.tight_layout()
plt.show()
# (*@\textnormal{3. 使用 SHAP 解释模型}@*)
# (*@\textnormal{需要先安装 shap: pip install shap}@*)
import shap
explainer = shap.TreeExplainer(model)
shap_values = explainer.shap_values(X_test)
# (*@\textnormal{4. 绘制 SHAP 总结图}@*)
shap.summary_plot(
    shap_values,
    X_test,
    feature_names=feature_cols
)
\end{lstlisting}

这段代码分为两个部分。第一部分使用 model.feature\_\allowbreak{}importances\_\allowbreak{} 输出随机森林自身的特征重要性，并画出柱状图。这个图可帮助研究者快速判断模型最关注哪些特征。第二部分使用 SHAP 对模型进行解释，shap.summary\_\allowbreak{}plot() 可显示每个特征对预测结果的整体影响。


可解释性分析中的 AI 辅助，应服务于“把模型输出转化为材料语言”，而不是替代机理判断。提示词中需要同时给出模型结果和要求它说明证据限制。可按下面方式组织问题。

\begin{samepage}
\begin{userbox}
我用随机森林预测高熵合金涂层硬度，模型输出的特征重要性排序为：是否含硬质相 0.34，平均晶粒尺寸 0.26，原子半径差 0.18，价电子浓度 0.13，电负性差 0.06，平均电负性 0.03。SHAP 图显示，含硬质相通常对预测硬度有正贡献，晶粒尺寸较大时常出现负贡献。

请帮我把这些模型解释结果转化为材料语言。

要求如下：
\begin{enumerate}
\item 分别联系第二相强化、细晶强化、固溶强化或相稳定性解释这些特征；
\item 说明这些解释需要哪些实验表征支持，例如 XRD、SEM、EDS\footnote{Energy Dispersive Spectroscopy，能谱分析。} 或硬度测试；
\item 提醒特征重要性和 SHAP 只能说明当前模型中的关系，不宜直接证明因果关系；
\item 如果模型把样品编号、数据来源等无材料意义字段排在前面，也请说明可能存在什么问题。
\end{enumerate}
\end{userbox}

\begin{aibox}
可这样理解这组结果：

\begin{enumerate}
\item “是否含硬质相”重要性最高，说明模型在预测硬度时较多依赖硬质相信息。若涂层中确实存在 TiC、碳化物或硼化物等硬质相，可与第二相强化联系起来解释。
\item “平均晶粒尺寸”重要性较高，且较大晶粒尺寸对应负向 SHAP 贡献，可能说明模型学习到了晶粒细化有利于提高硬度的趋势，可与细晶强化机制联系起来。
\item “原子半径差”有一定重要性，可与多主元合金中的晶格畸变和固溶强化联系起来，但需要结合相结构和成分分布判断。
\item “价电子浓度”可能与 FCC/BCC 相稳定性有关，但它对硬度的影响通常需要通过相结构、组织形貌和其他描述符共同解释。
\end{enumerate}

这些解释仍然只是模型层面的线索。若要形成可靠材料结论，需要用 XRD 确认物相，用 SEM/EDS 观察硬质相和成分分布，用晶粒尺寸统计验证细晶强化趋势，并检查硬度测试条件是否一致。如果无材料意义的字段在模型中排在前面，可能说明数据整理、特征选择或数据划分存在问题。
\end{aibox}
\end{samepage}

模型解释并不等于材料真理。特征重要性高，只说明模型在当前数据集中较多依赖该特征；SHAP 值为正，也只说明在当前模型中该特征对预测值有正向贡献。这些结果可能受到数据量、数据分布、特征相关性和模型类型影响。例如，如果所有含硬质相的样品刚好也具有较小晶粒尺寸，模型可能难以区分到底是硬质相起主要作用，还是晶粒细化起主要作用。因此，可解释性分析必须与实验表征和材料机理结合，不能孤立理解。


对于高熵合金涂层硬度预测，可按照以下思路解读模型解释结果。首先看特征重要性，判断模型主要关注哪些因素。其次看 SHAP 分析，判断这些因素对预测硬度的正负影响。然后把模型结果与材料强化机制对应起来：硬质相对应第二相强化，晶粒尺寸对应细晶强化，原子半径差对应晶格畸变和固溶强化，价电子浓度可能与相结构稳定性有关。最后，再回到实验数据中检查这些解释是否被组织形貌、物相分析和硬度测试支持。


例如，如果模型认为“是否含硬质相”最重要，而 SEM 和 XRD 也显示高硬度样品中确实存在较多 TiC 等硬质相，那么模型解释与实验观察相互支持。相反，如果模型认为某个表面上无关的字段最重要，例如样品编号或数据来源，那么就需要警惕模型可能学到了错误关联。这种情况并不是模型真正理解了材料规律，而可能是数据整理或特征设计存在问题。


除了特征重要性和 SHAP，还有一些更偏研究生层次的方法可用于材料模型解释。例如，符号回归和 SISSO\footnote{Sure Independence Screening and Sparsifying Operator，一种用于筛选描述符并构建稀疏表达式的方法。} 试图从数据中寻找简洁的数学表达式或组合描述符，使模型结果更接近材料专业人员能够理解的公式关系。对于入门阶段的研究者来说，可先掌握特征重要性和 SHAP 的基本思想；对于研究生，可进一步了解这些能够自动发现简洁规律的可解释建模方法。


本节的重点不是让研究者记住某一个解释算法，而是理解材料人工智能中的一个基本原则：模型预测之后，必须回到材料问题本身去解释。一个好的材料机器学习模型，不仅应该预测误差较小，还应该能够帮助研究者理解哪些成分、结构或组织因素可能控制性能。只有当模型结果能够与材料机制相互印证时，它才更可能成为材料研究和材料设计的有效工具。


通过高熵合金涂层硬度预测这个例子可看到，可解释性分析把机器学习结果和材料科学知识连接起来。模型指出哪些特征重要，材料科学知识帮助研究者解释这些特征为什么重要。但最终是否形成可靠规律，仍然需要实验表征、物理机理和专业判断共同验证。


下一节将进一步讨论物理知识嵌入。可解释性更多是在模型训练之后分析“模型为什么这样预测”，而物理知识嵌入则是在建模之前或建模过程中，就主动把材料科学规律加入数据选择、特征设计和模型约束中，使模型尽量避免学习错误关系。


\subsection{物理知识与模型约束}

上一节讨论了模型可解释性，即模型训练完成之后，研究者如何分析模型主要依据哪些特征做出预测，并将这些特征与材料机理联系起来。可解释性强调的是“事后解释”：模型已经训练好了，研究者再分析它为什么这样预测。


但是，在材料科学中，仅仅依靠事后解释还不够。更理想的做法是，在数据整理、特征选择、模型训练和结果判断的过程中，就主动引入材料科学知识，使模型尽量在合理的物理和化学范围内学习规律。这就是物理知识嵌入的基本思想。


物理知识嵌入并不一定是复杂的公式或高级算法。对于入门阶段，可把它理解为：不要把所有能收集到的字段都机械地丢给模型，而是根据材料机理选择有意义的数据和特征；不要让模型学习明显不合理的关系，而是用物理、化学和实验常识约束建模过程。这样得到的模型，通常比完全黑箱的模型更可靠，也更容易被研究者接受。


以前面的高熵合金涂层硬度预测任务为例，研究者已经用原子半径差、平均电负性、电负性差、价电子浓度、是否含硬质相和平均晶粒尺寸等特征预测涂层硬度。接下来需要判断的是：这些特征为什么可作为输入？哪些字段不应该作为输入？数据应该如何筛选，才能让模型学习到更接近材料规律的关系？


首先，物理知识可用于指导特征选择。对于高熵合金涂层硬度来说，硬度提高通常可能来自固溶强化、细晶强化、第二相强化和组织致密化等机制。因此，输入特征应该尽量围绕这些机制设计。


例如，原子半径差可与晶格畸变和固溶强化联系起来。多种元素固溶在同一晶格中时，元素原子半径差异会导致晶格畸变，从而可能阻碍位错运动，提高材料硬度。平均晶粒尺寸可与细晶强化联系起来。晶粒越细，晶界数量越多，位错运动越容易受到阻碍，硬度可能提高。是否含硬质相可与第二相强化联系起来。如果涂层中生成 TiC、碳化物或硼化物等硬质相，这些硬质相可能作为强化相提高涂层硬度。


因此，这些特征不是随便选出来的数字，而是和材料强化机制存在对应关系。模型使用这些特征预测硬度，才有可能学习到具有材料意义的规律。


这种关系可整理为表 \ref{tab:mechanism-feature-mapping}。


\begin{formalTable}[htbp]
\caption{材料机制与候选特征对应关系}
\label{tab:mechanism-feature-mapping}
\begin{tabularx}{\textwidth}{@{}YYY@{}}
\toprule
\textbf{材料机制} & \textbf{可转化的模型特征} & \textbf{对硬度的可能影响} \\
\midrule
固溶强化 & 原子半径差、电负性差、价电子浓度 & 晶格畸变和相稳定性变化可能影响硬度 \\
细晶强化 & 平均晶粒尺寸 & 晶粒细化可能提高硬度 \\
第二相强化 & 是否含硬质相、硬质相含量 & 硬质相可能提高涂层硬度 \\
组织致密化 & 孔隙率、裂纹数量、涂层缺陷 & 孔隙和裂纹可能降低硬度和服役性能 \\
\bottomrule
\end{tabularx}
\end{formalTable}

这个表格体现了物理知识嵌入的第一层含义：把材料机理转化为模型特征。模型输入不是任意字段，而是由材料知识筛选出来的、有明确含义的特征。


在实际建模中，这一步应先由材料机理提出候选特征，再结合已有实验记录判断哪些特征可获得、哪些特征需要补充测试。特征设计不应只看字段是否存在，还要看字段是否能对应到明确的材料机制。


其次，物理知识可帮助排除不合理特征。并不是所有出现在数据表中的字段都适合作为模型输入。例如，样品编号、测试日期、数据来源、实验人员姓名等字段，通常不应该作为材料性能预测特征。它们可能与硬度在数据表中出现某种偶然相关，但这种相关并不具有材料意义。如果模型根据样品编号预测硬度，就说明它可能学到了数据排序或实验批次的偶然关系，而不是材料规律。


前文已经讨论过目标泄漏，这里只强调它与物理知识嵌入的关系：输入字段必须在预测之前可获得，且不能由目标结果推导而来。例如，“硬度等级”“是否高硬度样品”“硬度排序”这类字段不能作为硬度预测输入，因为它们已经包含了目标硬度的信息。研究者需要检查每一个输入字段：它是否具有材料意义？是否直接包含目标结果？如果无法回答这些问题，就不应该轻易放入模型。


第三，物理知识可指导数据筛选和标准化。材料实验数据往往受到测试条件影响。以硬度为例，不同载荷、不同保载时间、不同测试位置可能得到不同硬度值。对于高熵合金涂层，表面硬度、截面硬度和不同层位硬度也可能存在差异。如果把不同测试条件下的硬度值直接混合到一个数据集中，模型可能学到混乱关系。


在建立硬度预测模型时，应该尽量保证硬度数据具有可比性。例如，尽量使用相同测试方法、相近测试载荷、相同单位和相同样品位置的数据。如果不同数据来源不可避免，就应在数据表中记录测试条件，必要时把测试条件作为附加字段，或者只选择条件一致的数据进行建模。


这种数据筛选本身就是物理知识嵌入的一部分。它告诉模型：哪些数据可放在一起学习，哪些数据不宜简单混用。对于材料科学来说，数据标准化不是形式工作，而是保证模型学习到真实规律的前提。


检查数据可比性时，应重点核验测试载荷、测试单位、热处理条件、样品位置和数据来源等字段。最终是否保留某条数据，仍然需要研究者根据实验条件和研究目标判断。


第四，物理知识可作为模型结果的合理性检查。模型训练完成后，如果预测结果明显违背材料常识，就需要警惕。例如，如果模型预测孔隙率越高硬度越高，或者晶粒尺寸越大硬度越高，就需要进一步检查数据是否存在偏差、特征是否相关、样本数量是否太少，或者模型是否受异常样品影响。当然，材料规律并不总是简单单调关系，但当模型结果与常识明显冲突时，必须进行核查，而不是直接接受。


对于高熵合金涂层硬度预测，可设定一些基本的合理性检查。例如，含硬质相的样品不一定总是硬度最高，但如果模型完全认为硬质相无关，就需要检查硬质相字段是否准确；晶粒细化通常有利于硬度提高，如果模型得出相反趋势，也要检查晶粒尺寸统计是否一致；孔隙率增加通常不利于力学性能，如果模型认为孔隙率越高越好，就需要检查数据是否存在偶然相关或测量问题。


第五，物理知识还可用于更高级的模型约束。在入门阶段，研究者主要通过特征选择、数据筛选和结果检查嵌入材料知识。对于更深入的研究，还可把物理规律直接写进模型或损失函数中。例如，在原子尺度模型中，可要求模型满足平移、旋转和置换对称性；在热力学相关任务中，可引入稳定性约束；在力学模型中，可要求某些预测量满足基本物理关系。这类方法通常称为物理约束机器学习或物理知识嵌入机器学习。


对于本章而言，不需要展开复杂数学形式，但要让研究者理解一个基本思想：材料模型不应只追求拟合数据，还要尽量尊重材料科学规律。一个模型如果预测误差很小，但违背基本物理常识，或者依赖无法解释的虚假相关字段，那么它在材料研究中的价值是有限的。


可用如下流程概括物理知识嵌入在高熵合金涂层硬度预测中的作用：


\begin{lightquote}
\centering 研究目标 → 材料机理分析 → 特征选择 → 数据筛选 → 模型训练 → 结果解释 → 实验和机理验证
\end{lightquote}

其中，材料机理分析不是最后才出现，而是从一开始就参与建模。研究者首先根据固溶强化、细晶强化、第二相强化等机制选择特征；再根据测试条件筛选可比数据；然后训练模型；最后再用实验表征和材料机理检查模型解释是否合理。


为了帮助研究者理解这个流程，可设计一个简单的操作任务。假设已经有一张高熵合金涂层硬度数据表，其中包含以下字段：样品编号、化学成分、原子半径差、电负性差、价电子浓度、是否含硬质相、晶粒尺寸、孔隙率、硬度等级、硬度/HV、测试载荷、数据来源。研究者需要判断哪些字段可作为输入特征，哪些字段应该保留但不作为模型输入，哪些字段必须删除或谨慎使用。


按照物理知识嵌入的原则，可这样处理：原子半径差、电负性差、价电子浓度、是否含硬质相、晶粒尺寸和孔隙率可作为候选输入特征，因为它们具有材料意义；测试载荷和数据来源应保留，用于判断数据是否可比，但不一定直接作为性能特征；样品编号通常不应作为输入；硬度等级不能作为输入，因为它可能由硬度结果推导而来，存在目标泄漏风险；硬度/HV 是预测目标 y，不应同时作为输入特征 X。


可整理为表 \ref{tab:physics-informed-field-treatment}。


\begin{formalTable}[htbp]
\caption{物理知识嵌入建模过程中的字段处理}
\label{tab:physics-informed-field-treatment}
\begin{tabularx}{\textwidth}{@{}YYY@{}}
\toprule
\textbf{字段} & \textbf{是否作为模型输入} & \textbf{原因} \\
\midrule
原子半径差 & 可 & 与晶格畸变和固溶强化有关 \\
电负性差 & 可 & 与元素化学差异和相形成有关 \\
价电子浓度 & 可 & 可能影响相结构稳定性 \\
是否含硬质相 & 可 & 对应第二相强化 \\
晶粒尺寸 & 可 & 对应细晶强化 \\
孔隙率 & 可 & 反映组织致密性 \\
测试载荷 & 谨慎使用或用于筛选 & 影响硬度可比性 \\
数据来源 & 不作为材料特征 & 用于追溯和分组检查 \\
样品编号 & 不使用 & 无材料意义 \\
硬度等级 & 不使用 & 可能造成目标泄漏 \\
硬度/HV & 目标值 & 模型要预测的 y \\
\bottomrule
\end{tabularx}
\end{formalTable}

这类字段筛选任务适合用提示词形式进行辅助检查。提示词应要求同时判断材料意义、数据可比性和目标泄漏风险，而不是只问“哪些列可输入模型”。

\begin{samepage}
\begin{userbox}
我有一张高熵合金涂层硬度数据表，字段包括：样品编号、化学成分、原子半径差、电负性差、价电子浓度、是否含硬质相、晶粒尺寸、孔隙率、硬度等级、硬度/HV、测试载荷、数据来源。请帮我判断哪些字段适合作为模型输入，哪些字段应保留但不作为输入，哪些字段可能造成目标泄漏或虚假相关。

要求如下：
\begin{enumerate}
\item 按固溶强化、细晶强化、第二相强化和组织致密化解释可用特征；
\item 指出哪些字段用于追溯或筛选数据，但不宜作为材料性能特征；
\item 明确指出硬度/HV 和硬度等级为什么不能作为输入特征；
\item 提醒哪些测试条件会影响硬度数据可比性；
\item 不要只给机器学习角度判断，要结合材料机理说明理由。
\end{enumerate}
\end{userbox}

\begin{aibox}
可按三类处理这些字段：

\begin{enumerate}
\item 可作为候选输入特征的字段：原子半径差、电负性差、价电子浓度、是否含硬质相、晶粒尺寸、孔隙率。这些字段分别可与固溶强化、相结构稳定性、第二相强化、细晶强化和组织致密化联系起来。
\item 应保留但不宜直接作为材料性能特征的字段：测试载荷、数据来源、化学成分和样品编号。测试载荷和数据来源主要用于判断数据可比性和可追溯性；样品编号通常没有材料意义；化学成分可保留，但进入模型前通常需要转化为成分描述符。
\item 不能作为输入特征的字段：硬度/HV 是预测目标，硬度等级往往由硬度结果划分而来，二者如果进入输入特征，会造成目标泄漏。
\end{enumerate}

此外，硬度数据是否可比还取决于测试载荷、保载时间、测试位置、样品表面状态和测试方法。如果这些条件不一致，应先进行筛选、分组或记录为元数据，再决定是否合并建模。
\end{aibox}
\end{samepage}

需要强调的是，物理知识嵌入并不意味着模型一定更复杂。很多时候，它体现为更谨慎的数据选择、更合理的特征设计和更严格的结果检查。对于材料人工智能入门阶段的研究者来说，最重要的不是马上构建复杂的物理约束模型，而是养成一种意识：建模前先问材料问题，建模中保留物理意义，建模后回到实验和机理验证。


本节说明了物理知识嵌入的基本思想。可解释性帮助研究者在模型训练后理解模型，物理知识嵌入则帮助研究者在建模前和建模过程中约束模型。二者共同说明：人工智能用于材料科学，并不是把数据机械地交给模型，而是让数据驱动方法与材料科学知识相互配合。至此，第\ref{chap:materials-cognition-prediction}章完成了从材料数据、材料表示、性能预测、原子尺度模拟到模型解释的完整链条。下一章将进一步讨论逆向设计问题，即如何根据目标性能反向寻找和设计新材料。

\subsubsection*{\textbf{小结}}

本节说明模型解释和物理知识嵌入共同服务于材料判断。
特征重要性和 SHAP 可提示模型依据，但不能直接证明因果机理，仍需 XRD、SEM、EDS 等证据支持。
物理知识应贯穿数据筛选、特征选择、模型训练和结果检查，避免模型学习虚假相关。

\subsection*{习题}

\begin{enumerate}[leftmargin=*]
\item SHAP 为什么不能直接证明材料机理？还需哪些表征？
\item 可解释性分析如何辅助材料筛选？请举例说明。
\item 只有 FCC 数据时，能否预测 BCC 合金？请说明风险和验证方案。
\end{enumerate}

\section*{本章小结}

本章围绕材料认知与性能预测展开，形成了从数据到模型、再从模型回到材料问题的基本链条。材料数据整理强调字段完整、来源可追溯和条件可比较；材料表示强调把化学式、结构和局域环境转化为描述符或图结构；性能预测部分说明了描述符模型、图神经网络和机器学习势在不同尺度任务中的作用；模型解释与物理知识嵌入则进一步说明，模型预测结果必须经过材料机理、实验表征和物理约束的共同检验。由此可见，人工智能在材料科学中的价值不在于替代实验和理论，而在于提高数据利用效率、缩小候选范围并辅助形成可验证的材料判断。

\chapter{材料设计与智能发现}
\label{chap:materials-design-discovery}


第\ref{chap:materials-cognition-prediction}章讨论正向预测，即在材料成分、结构或实验条件已知的情况下，预测材料可能具有的性能。本章进一步讨论逆向设计与智能发现：当目标性能已经给定时，如何从大量候选材料中筛选出值得验证的方案。这个过程并不是简单地让模型“生成答案”，而是要把目标性能、稳定性、可合成性、成本、安全性和实验条件共同纳入考虑。

本章按照材料发现的一般流程展开，重点解决四个问题。第一节说明如何把应用需求转化为材料设计目标、候选变量和约束边界。第二节讨论候选筛选和主动学习如何帮助研究者逐层缩小搜索空间，并安排下一轮实验或计算。第三节介绍生成设计如何提出库外候选，以及生成候选为什么必须经过有效性、稳定性、新颖性和可合成性验证。第四节讨论可信发现、方法选择和闭环迭代，强调模型推荐必须与计算验证、实验反馈和应用约束共同使用。通过本章学习，研究者应理解人工智能如何参与材料设计决策，同时认识候选材料最终仍需要通过物理判断、计算验证和实验结果逐步确认。




\clearpage
\section{问题转换与设计边界}


\begin{sectionkeypoints}
\item 正向预测回答“已知材料可能具有什么性质”，逆向设计回答“为了获得目标性质应当寻找什么材料”。
\item 逆向设计本质上是受约束搜索，不是单纯追求最高预测分数，而是在稳定性、可合成性、成本、安全和设备条件内寻找可验证候选。
\item 材料设计对象可能位于成分、原子结构、微结构、工艺和服役环境等不同层级，层级选择必须与研究目标一致。
\item 候选筛选、主动学习和生成设计是本章三条核心路线，实际研究中常常组合使用。
\end{sectionkeypoints}

\subsection{正向预测与逆向设计}

第\ref{chap:materials-cognition-prediction}章讨论的是正向问题：给定材料的成分、结构、工艺或表征信息，预测形成能、带隙、硬度、折射率或熔点等性质。第\ref{chap:materials-design-discovery}章把问题方向反过来：当目标性质已经确定时，寻找哪些材料、配方、结构或工艺值得进一步尝试。前者回答“这种材料具有什么性质”，后者回答“具有目标性质的材料可能在哪里”。


材料科学长期依赖经验规则和实验试错，但多组元合金、复杂晶体和玻璃配方的候选数量远超实验能够逐一验证的范围。逆向设计的作用不是取消经验和实验，而是把大规模穷举转化为有方向的搜索：先明确目标和约束，再用模型提出或排序候选，最后通过计算和实验逐步确认。


从数学上看，逆向设计可理解为受约束优化。材料对象 x 可以是化学式、晶体结构、玻璃配方、微结构参数或制备条件；正向模型 f(x) 估计目标性能；约束 g(x) 描述稳定性、可合成性、成本、安全和设备边界。搜索的目的不是让某个预测分数无限增大，而是在满足约束的前提下找到一组可验证的候选。正向预测、逆向设计与闭环发现之间的关系见表 \ref{tab:forward-inverse-closed-loop}。


\begin{formalTable}[htbp]
\caption{正向预测、逆向设计与闭环发现}
\label{tab:forward-inverse-closed-loop}
\footnotesize
\begin{tabularx}{\textwidth}{@{}YYYYY@{}}
\toprule
\textbf{问题类型} & \textbf{输入} & \textbf{输出} & \textbf{材料例子} & \textbf{主要风险} \\
\midrule
正向预测 & 候选材料 x & 性质 f(x) & 由晶体结构预测带隙 & 训练集外误差被低估 \\
逆向设计 & 目标性质 y & 候选材料 x & 按目标折射率设计玻璃 & 候选不可合成或不稳定 \\
闭环发现 & 已有数据 + 目标 & 下一步实验 & 主动学习选择合金配比 & 实验噪声导致错误反馈 \\
\bottomrule
\end{tabularx}
\end{formalTable}

\subsubsection*{\textbf{1. 问题转换：从应用目标到材料任务}}

逆向设计的第一步，不是直接选择模型，而是把应用语言转化为材料科学语言。应用目标往往比较笼统，例如“希望玻璃折射率更高”“希望发光材料颜色更接近蓝光”“希望涂层在高温下更稳定”。这些表述还不宜直接进入模型，因为模型需要明确的输入变量、目标性能和约束条件。只有把应用目标拆解成可计算、可测量、可验证的材料任务，后续的筛选、优化和生成才有依据。

以 AR\footnote{Augmented Reality，增强现实。} 光波导玻璃为例，“高折射率”只是一个目标方向。真正进入设计流程时，还需要进一步明确折射率测量波长、透过率范围、色散要求、玻璃化转变温度（Tg\footnote{Glass transition temperature，玻璃化转变温度。}）、热膨胀系数、密度、原料安全性和熔制温度窗口。如果只追求折射率，模型可能推荐含高极化率重金属氧化物较多的配方，但这类配方可能带来密度升高、成本增加、透过率下降或环保风险。因此，逆向设计不是单目标求最大值，而是多目标、多约束条件下的候选选择。

其他材料任务也需要类似拆解。例如 microLED\footnote{Micro Light-Emitting Diode，微型发光二极管。} 发光材料可把目标颜色转化为发光波长或带隙范围，但仍需补充缺陷态、载流子复合效率、热稳定性和外延匹配等约束。这里不展开为新的主线案例，只用它说明：目标性能通常只是材料设计任务的一部分。

具体而言，问题转换通常包括四个问题：目标性能是什么？候选变量是什么？约束条件是什么？验证方式是什么？目标性能决定模型要预测什么；候选变量决定模型可搜索哪些材料空间；约束条件决定哪些候选必须排除；验证方式决定候选是否能从模型结果走向真实材料。表中的 DFT\footnote{Density Functional Theory，密度泛函理论。} 指密度泛函理论计算。典型应用目标到材料设计任务的转换见表 \ref{tab:application-to-design-task}。

\begin{formalTable}[htbp]
\caption{应用目标到材料设计任务的转换}
\label{tab:application-to-design-task}
\footnotesize
\begin{tabularx}{\textwidth}{@{}YYYYY@{}}
\toprule
\textbf{应用目标} & \textbf{材料任务} & \textbf{候选变量} & \textbf{关键约束} & \textbf{验证方式} \\
\midrule
高折射率玻璃 & 预测并优化折射率、Tg、热膨胀系数 & 氧化物组成比例 & 玻璃形成、透明性、熔制温度 & 熔制实验、折射率和透过率测试 \\
蓝光发光材料 & 筛选合适带隙和稳定结构 & 元素组成、晶体结构、掺杂 & 缺陷、外延兼容、热稳定 & DFT、发光测试、器件验证 \\
高硬度涂层 & 优化成分、硬质相和组织参数 & 合金成分、工艺、相组成 & 裂纹、孔隙、结合强度 & 显微硬度、SEM/XRD、服役测试 \\
超高温陶瓷 & 寻找高熔点和抗氧化候选 & 化学组成、晶体结构 & 高温稳定、氧化行为、可烧结性 & 高温实验、热分析、模拟验证 \\
\bottomrule
\end{tabularx}
\end{formalTable}

这一转换过程适合借助 AI 辅助梳理，但不能让 AI 直接给出最终材料结论。提示词应要求它把应用目标拆成材料变量、目标性能、约束条件和验证方式，而不是只列出几个候选材料名称。

\begin{samepage}
\begin{userbox}
我想设计一种用于 AR 光波导的高折射率玻璃。请帮我把这个应用目标转化为材料逆向设计任务。

要求如下：
\begin{enumerate}
\item 明确目标性能，例如折射率、透过率、色散、Tg、热膨胀系数等；
\item 说明候选变量可包括哪些氧化物组成或配方参数；
\item 区分必须满足的硬约束和可折中的软目标；
\item 给出候选筛选后需要进行哪些计算或实验验证；
\item 不要直接编造最终玻璃配方，只做任务拆解。
\end{enumerate}
\end{userbox}

\begin{aibox}
可将“高折射率 AR 光波导玻璃”拆解为以下设计任务：

\begin{enumerate}
\item \textbf{目标性能：}折射率是核心目标，但还应同时记录透过率、色散、Tg、热膨胀系数、密度和加工温度等指标。
\item \textbf{候选变量：}可把各氧化物组分比例作为设计变量，例如网络形成体、网络修饰体和提高折射率的氧化物比例。
\item \textbf{硬约束：}配方比例之和应满足组成约束；原料不能超出安全和环保边界；熔制温度应在设备可承受范围内；样品应能够形成玻璃而不是明显析晶。
\item \textbf{软目标：}折射率越高越好，但不能以明显牺牲透明性、热稳定性和加工性为代价。
\item 验证方式：模型筛选后，需要进行玻璃形成能力判断、熔制实验、折射率测试、透过率测试和热膨胀系数测试。
\end{enumerate}

这样处理后，逆向设计任务不再只是“寻找高折射率材料”，而是“在可熔制、可成玻、透明且热性质合适的约束下，寻找高折射率玻璃配方”。
\end{aibox}
\end{samepage}

问题转换应当在建模之前完成，而不是在模型给出候选之后再补充解释。如果目标一开始没有定义清楚，后面所有筛选和排序都会变得模糊。例如，若只写“设计高性能玻璃”，模型无法判断“高性能”指的是高折射率、高透过率、低热膨胀系数还是高热稳定性；若只写“设计蓝光材料”，模型也不知道重点是发光波长、发光效率、热稳定性还是外延兼容性。本章强调问题转换，是为了让研究者意识到：人工智能材料设计首先是材料问题定义，其次才是模型选择。

一个规范的逆向设计任务，宜用一句话概括为“在什么约束下，寻找什么类型的材料，使什么目标性能达到什么范围”。例如，“在不使用高风险元素、熔制温度不超过设备上限、能够形成透明玻璃的约束下，寻找折射率较高且热膨胀系数合适的氧化物玻璃配方”。这样的表述比“让模型设计高折射率玻璃”更清楚，也更便于后续建立数据表、选择描述符和设计验证实验。

\subsection{设计对象与约束边界}

\subsubsection*{\textbf{1. 设计边界：材料对象分层}}

逆向设计并不总是直接生成一个化学式。材料对象通常包含多个层级：成分决定由哪些元素或组分构成；原子结构决定空间群、占位、缺陷和局域配位；微结构包含晶粒、孔隙、第二相和界面；工艺包含熔融、退火、沉积、烧结和冷却条件。同一成分经过不同工艺可得到不同结构和性能，因此设计对象的层级必须与研究目标相匹配。材料逆向设计常见对象层级见表 \ref{tab:inverse-design-object-levels}。


\begin{formalTable}[htbp]
\caption{材料逆向设计的对象层级}
\label{tab:inverse-design-object-levels}
\begin{tabularx}{\textwidth}{@{}YYYY@{}}
\toprule
\textbf{设计层级} & \textbf{常见变量} & \textbf{适用任务} & \textbf{容易忽略的问题} \\
\midrule
成分 & 元素/氧化物比例、掺杂量 & 玻璃配方、合金筛选、催化剂组分 & 相形成与工艺窗口 \\
原子结构 & 晶格、坐标、缺陷、表面 & 带隙、离子迁移、稳定性 & 对称性和局域环境 \\
微结构 & 晶粒、孔隙、第二相、取向 & 力学、热学、疲劳与可靠性 & 制备历史强相关 \\
工艺 & 温度、气氛、压力、时间 & 可合成性与性能优化 & 元数据缺失会误导模型 \\
服役环境 & 光、电、热、湿度、载荷 & 器件可靠性和寿命 & 实验条件与应用条件不一致 \\
\bottomrule
\end{tabularx}
\end{formalTable}

除了设计层级，还需要区分约束的类型。材料设计中的约束可分为硬约束和软约束。硬约束是必须满足的条件，例如有毒元素不能使用、设备最高温度不能超过、配方比例必须归一化、候选结构不能明显违反价态和原子距离规则。软约束是希望尽量优化但可折中的条件，例如成本越低越好、密度越低越好、热膨胀系数越接近目标越好。硬约束用于排除候选，软约束用于排序候选。

如果不区分硬约束和软约束，逆向设计很容易变成一个不清楚的打分问题。例如，把折射率、成本、热稳定性和环保要求都压缩成一个总分，表面上方便排序，但不同目标之间的取舍会被隐藏。逆向设计应把硬约束、软目标和数据边界分开说明。

设计边界还包括数据边界。模型只能在训练数据覆盖的范围内比较可靠地工作。如果训练数据主要来自硅酸盐玻璃，那么把模型直接用于完全不同的硫系玻璃或氟化物玻璃，预测结果就可能不稳定。如果模型只学习了室温性质，就不宜直接推断高温服役行为。如果数据没有记录制备工艺，模型也很难判断同一成分在不同热处理条件下的性能差异。因此，逆向设计开始之前必须说明数据从哪里来、覆盖哪些材料体系、哪些变量没有记录。不同边界类型及其处理方式见表 \ref{tab:inverse-design-boundary-types}。

\begin{formalTable}[htbp]
\caption{材料逆向设计中的边界类型}
\label{tab:inverse-design-boundary-types}
\footnotesize
\begin{tabularx}{\textwidth}{@{}YYYY@{}}
\toprule
\textbf{边界类型} & \textbf{含义} & \textbf{材料例子} & \textbf{处理方式} \\
\midrule
硬约束 & 不满足就直接排除 & 有毒元素、不可成玻、设备温度超限 & 筛选前写入规则 \\
软目标 & 可折中但要排序 & 折射率、密度、成本、热膨胀系数 & 作为多目标评价项 \\
数据边界 & 模型训练数据覆盖范围 & 只含某类玻璃或某类晶体 & 明确适用范围和外推风险 \\
工艺边界 & 实际制备条件限制 & 熔制温度、气氛、烧结时间 & 作为实验可行性约束 \\
\bottomrule
\end{tabularx}
\end{formalTable}

在教学中，可让研究者首先写出设计边界，再选择模型。这样的顺序比“先选一个算法再想怎么用”更合理。算法只能在已经定义清楚的问题上发挥作用；如果目标、变量和边界本身模糊，模型输出再复杂也难以转化为真实材料方案。

这一类边界梳理同样可使用 AI 辅助，但提示词必须要求区分硬约束、软目标和数据边界。

\begin{samepage}
\begin{userbox}
我准备做一个材料逆向设计任务，目标是寻找候选高折射率玻璃。请帮我列出这个任务的设计边界。

要求如下：
\begin{enumerate}
\item 区分硬约束、软目标、数据边界和工艺边界；
\item 每一类边界都要说明它为什么重要；
\item 不要把所有条件合并成一个总分；
\item 说明哪些信息需要在实验记录或数据库中保留；
\item 不要直接给出候选配方。
\end{enumerate}
\end{userbox}

\begin{aibox}
可从四类边界整理：

\begin{enumerate}
\item \textbf{硬约束：}原料安全、环保限制、组成比例合法、可形成玻璃、设备温度可达。这类条件不满足时，候选应直接排除。
\item \textbf{软目标：}折射率、透过率、密度、热膨胀系数、成本和加工性能。这些目标之间可能存在折中，应分别记录和排序。
\item \textbf{数据边界：}需要说明训练数据主要来自哪些玻璃体系、性质是在什么条件下测得、是否包含失败样品和析晶样品。超出数据覆盖范围时，模型预测应谨慎使用。
\item \textbf{工艺边界：}熔制温度、保温时间、冷却方式、气氛、坩埚材料等都会影响玻璃形成和性能，应作为实验元数据保存。
\end{enumerate}

这些边界不是附加说明，而是逆向设计任务的一部分。只有先限定边界，模型推荐的候选才有可能进入真实实验。
\end{aibox}
\end{samepage}

边界条件还应随研究阶段动态更新。早期阶段可设置较宽的组成范围，用来观察模型是否能发现潜在趋势；进入实验阶段后，边界应逐渐收紧，例如限定可采购原料、设备温度、样品尺寸和测试条件；进入应用阶段后，还要加入成本、安全、长期稳定性和批量制备要求。换言之，边界不是一次写完就不变，而是随着数据积累和实验反馈不断修正。

对于研究者来说，最容易忽略的是“不可见边界”。有些限制并不直接写在数据表中，例如某些配方虽然数据库中有记录，但实验室没有相应原料；某些结构计算上稳定，但现有设备无法提供所需压力或气氛；某些材料性能很好，但含有不适合应用场景的元素。这些不可见边界如果不提前说明，模型推荐就会表面上合理、实际不可执行。因此，材料逆向设计的边界既包括数据中的字段，也包括实验室条件和应用场景中的真实限制。

本章讨论三条核心路线。候选筛选适合从已有数据库或枚举空间中逐层过滤；主动学习适合实验昂贵、每轮只能验证少量样品的任务；生成设计用于突破已有候选库，在目标条件下提出新结构或新配方。三者并非互相替代，实际研究常按“筛选--优化--生成--验证”的方式组合使用。

\subsubsection*{\textbf{小结}}

本节说明逆向设计不是直接生成材料，而是把应用需求转化为目标、变量和约束。
正向预测回答已知材料的性能，逆向设计回答为了目标性能应寻找哪些候选。
材料对象层级、硬约束、软目标、数据边界和工艺边界需要在建模前明确。

\subsection*{习题}

\begin{enumerate}[leftmargin=*]
\item 正向预测和逆向设计在高折射率玻璃任务中如何配合？
\item 逆向设计为什么必须先明确目标和约束？
\item 如何把高折射率玻璃需求转化为计算目标？
\end{enumerate}

\clearpage
\section{候选筛选与主动学习}


\begin{sectionkeypoints}
\item 候选筛选的基本逻辑是先建立候选库，再按照成本由低到高、可信度由粗到细逐层过滤。
\item 筛选漏斗通常包含规则过滤、机器学习代理模型、高精度计算和实验验证等层级，每一级都应记录通过和排除理由。
\item 主动学习适合实验昂贵、小样本和需要多轮优化的材料任务，核心是同时考虑预测性能和模型不确定性。
\item 实验噪声、批次效应和失败样品都应进入数据记录，否则模型可能把制备差异误认为材料本征规律。
\item 候选排序和主动学习评分只负责辅助决策，最终实验批次仍需要材料约束、工艺可行性和人工复核。
\end{sectionkeypoints}

\subsection{筛选漏斗与分层验证}

高通量虚拟筛选的基本思路，是先建立候选库，再按照成本由低到高、可信度由粗到细逐层过滤。最外层使用元素排除、价态平衡、成分范围、成本或毒性等规则；中间层使用第\ref{chap:materials-cognition-prediction}章建立的正向模型快速估计形成能、带隙、弹性模量或折射率；内层使用 DFT、分子动力学或热力学计算精算；最后才进入实验合成与表征。


\begin{figure}[H]
\centering
\BookFigureImage{images_11/3.png}
\caption{候选筛选与验证层级示意图}
\label{fig:candidate-screening-validation}
\end{figure}
\FloatBarrier

如图 \ref{fig:candidate-screening-validation} 所示，筛选漏斗利用了不同方法之间的成本差异。规则过滤几乎不需要计算，机器学习预测可快速处理大量候选，而高精度计算和真实实验只能用于较小范围。把便宜但粗略的方法放在前面，把昂贵但可靠的方法放在后面，能够避免把有限资源消耗在明显不合理的候选上。


筛选的第一个常见门槛是稳定性。无机晶体可参考形成能和凸包能量，合金和陶瓷还要考虑相分离、氧化和高温相变，玻璃则需要考虑玻璃形成能力、析晶倾向和熔制窗口。计算稳定不等于实验一定能够合成，热力学判断还要与动力学路径、前驱体和设备条件结合。


第二个门槛是多目标要求。AR 光波导玻璃不仅需要较高折射率，还要兼顾透过率、色散、玻璃化转变温度、热膨胀系数、密度和加工性能；microLED 发光材料除带隙外，还要考虑缺陷、载流子复合、热稳定和外延兼容性。相应地，候选排序不应只依据单一预测值，而应保留不同目标之间的折中关系。筛选漏斗中不同层级的作用见表 \ref{tab:screening-funnel-levels}。


\begin{formalTable}[htbp]
\caption{筛选漏斗中不同层级的作用}
\label{tab:screening-funnel-levels}
\begin{tabularx}{\textwidth}{@{}YYYY@{}}
\toprule
\textbf{筛选层级} & \textbf{常见工具} & \textbf{适合回答的问题} & \textbf{学习要点} \\
\midrule
规则层 & 元素价态、毒性、成本、工艺窗口 & 哪些候选明显不可取？ & 规则便宜但粗糙，适合第一轮排除 \\
代理模型层 & 回归/分类模型、图神经网络、集成模型 & 哪些候选可能性能较好？ & 模型用于排序，不等于最终证明 \\
物理计算层 & DFT、相图、分子动力学、热力学模型 & 候选是否稳定、机制是否可信？ & 精度更高但成本也更高 \\
实验层 & 合成、表征、服役测试 & 候选能否成为真实材料？ & 实验验证是材料发现的终点而非附属步骤 \\
\bottomrule
\end{tabularx}
\end{formalTable}

\subsubsection*{\textbf{1. 案例示范：高折射率玻璃候选筛选}}

以高折射率玻璃为例，首先把应用要求转化为可计算指标；然后限定各氧化物比例之和、原料安全和设备温度等边界；代理模型预测折射率、Tg、热膨胀系数和析晶倾向；物理或经验模型检查玻璃形成与黏度窗口；最后制备少量代表性配方并测量真实光学损耗。每一级都应记录排除理由，使候选清单可复核。

具体而言，第一层规则可先排除明显不合理的配方。例如，各氧化物比例之和必须为 100\%，单一组分含量不能超过玻璃形成的经验范围，含有明显不适合教学或工程应用的高风险元素时应直接排除。第二层可用代理模型预测折射率和热性质，把候选从几千个缩小到几十个。第三层再用更可靠的热力学判断或经验模型检查玻璃形成能力、析晶风险和黏度窗口。最后才选择少量配方进入熔制实验。

筛选漏斗的一个重要原则是“保留失败理由”。如果一个候选被排除，应记录它是因为折射率不够、Tg 太低、热膨胀系数过高、含有风险元素，还是因为预测不确定性太大。这样做有两个好处。第一，后续可复核筛选规则是否过于严格；第二，失败信息可反馈给下一轮模型或实验设计。如果只保留最后通过的候选，而删除中间淘汰信息，材料发现过程就会失去可追溯性。

在初学阶段，研究者可将筛选漏斗理解为一个逐层提问的过程：第一问，这个候选是否明显违反基本规则？第二问，它的目标性能是否接近要求？第三问，它是否可能稳定存在或可制备？第四问，它是否值得消耗实验资源？每一层都不是绝对正确的判断，而是在不同成本下尽量减少无效尝试。筛选过程中需要保留的判断信息见表 \ref{tab:screening-funnel-records}。

\begin{formalTable}[htbp]
\caption{筛选漏斗中需要保留的判断信息}
\label{tab:screening-funnel-records}
\footnotesize
\begin{tabularx}{\textwidth}{@{}YYYY@{}}
\toprule
\textbf{筛选问题} & \textbf{排除依据} & \textbf{保留信息} & \textbf{后续用途} \\
\midrule
是否违反基本规则 & 比例错误、风险元素、明显价态异常 & 排除原因和原始配方 & 修正规则或候选生成空间 \\
性能是否接近目标 & 折射率、带隙、硬度等预测值偏离 & 预测值和模型版本 & 比较不同代理模型 \\
是否可能稳定 & 凸包能、析晶倾向、相分离风险 & 稳定性指标和计算方法 & 决定是否进入高精度计算 \\
是否值得实验 & 成本、设备温度、样品数量限制 & 排序分数和实验批次 & 形成可执行实验计划 \\
\bottomrule
\end{tabularx}
\end{formalTable}

这一过程可借助 AI 辅助把筛选规则写成清单，但规则本身必须由材料任务决定。可用如下提示词把流程规范化。

\begin{samepage}
\begin{userbox}
我正在设计一个高折射率玻璃候选筛选漏斗。候选数据表中有氧化物比例、预测折射率、预测 Tg、热膨胀系数、是否含风险元素、预测析晶风险等字段。请帮我把筛选流程整理成逐层规则。

要求如下：
\begin{enumerate}
\item 先列出规则过滤层、代理模型层、物理/经验验证层和实验层；
\item 每一层说明主要排除什么候选；
\item 明确哪些排除理由需要记录；
\item 提醒候选排序不应只看折射率；
\item 不要编造真实配方数据。
\end{enumerate}
\end{userbox}

\begin{aibox}
可将筛选漏斗设计为四层：

\begin{enumerate}
\item \textbf{规则过滤层：}检查氧化物比例是否合理、是否含风险元素、熔制温度是否超出设备范围。被排除的候选应记录具体违反哪条规则。
\item \textbf{代理模型层：}预测折射率、Tg 和热膨胀系数。折射率高但 Tg 过低或热膨胀系数过高的候选不宜直接排在最前。
\item \textbf{物理/经验验证层：}进一步检查玻璃形成能力、析晶倾向、黏度窗口和配方稳定性。此处的目标是排除“预测性能好但难以形成玻璃”的候选。
\item \textbf{实验层：}选择少量代表性配方进行熔制、退火、折射率和透过率测试，并记录失败样品原因。
\end{enumerate}

筛选报告中应保留每一级候选数量、通过率和排除理由。这样可判断筛选规则是否过严，也便于后续主动学习利用失败信息。
\end{aibox}
\end{samepage}

筛选漏斗还有一个容易被忽视的问题：前一级筛选规则会影响后一级能看到什么候选。如果第一层规则过于严格，可能会把一些有潜力但不符合传统经验的候选提前排除；如果第一层规则过于宽松，又会把大量明显不可行的候选留给后续高成本计算。为此，筛选规则需要根据任务目标反复校准，而不是一次确定后不再修改。

在课堂练习中，可让研究者比较两种筛选策略。一种策略偏保守，优先排除风险候选，适合实验资源很有限的任务；另一种策略偏探索，保留更多边界候选，适合希望发现新体系的任务。两种策略没有绝对优劣，关键在于是否说明了筛选目的、资源条件和验证能力。这样可避免把筛选漏斗理解为机械流程，而是理解为材料决策工具。


以下给出一个最小候选排序代码示例。代码把目标带隙、凸包能量和元素风险写成显式评分规则，用于说明筛选假设如何转化为可复核的计算步骤。评分只负责初步排序，不宜代替稳定性复算和实验验证。


\begin{lstlisting}[language=Python,escapeinside={(*@}{@*)}]
import pandas as pd
# (*@\textnormal{gap 的单位为 eV，hull 的单位为 eV/atom}@*)
cands = pd.DataFrame({
    "id": ["mp-1", "mp-2", "mp-3", "mp-4"],
    "formula": ["M1O2", "M2O3", "M3S", "PbM4O"],
    "gap": [2.10, 1.55, 2.35, 2.00],
    "hull": [0.018, 0.006, 0.065, 0.012],
    "elements": ["M1,O", "M2,O", "M3,S", "Pb,M4,O"],
})
target_gap = 2.0
cands["gap_score"] = 1 - (cands["gap"] - target_gap).abs() / target_gap
cands["stability_score"] = (0.08 - cands["hull"]).clip(lower=0) / 0.08
cands["element_penalty"] = cands["elements"].str.contains("Pb|Cd|Hg").astype(int)
cands["priority"] = (
    0.45 * cands["gap_score"]
    + 0.45 * cands["stability_score"]
    - 0.50 * cands["element_penalty"]
)
cols = ["id", "formula", "gap", "hull", "priority"]
print(cands.sort_values("priority", ascending=False)[cols])
\end{lstlisting}

\manualfigcaption{\textbf{代码示例：}候选材料初步排序示例}

\subsection{主动学习与实验决策}

高通量筛选是在固定候选库中做过滤，主动学习则是在搜索过程中不断更新模型。研究者首先用少量实验或计算训练代理模型，模型同时给出候选性能和不确定性，再选择下一批最值得验证的候选；新结果返回后重新训练模型，形成序贯迭代。


\begin{figure}[H]
\centering
\BookFigureImage{images_11/4.png}
\caption{主动学习与生成设计闭环示意图}
\label{fig:active-learning-generative-loop}
\end{figure}
\FloatBarrier

图 \ref{fig:active-learning-generative-loop} 概括了主动学习与生成设计之间的闭环关系。贝叶斯优化是常见的主动学习方法，适合小样本和高成本实验。它通常包含代理模型和采集函数两个部分：代理模型近似真实实验或高精度计算，采集函数综合预测均值与不确定性，决定下一步优先验证哪些候选。其核心不是总选择当前分数最高的样品，而是在“利用”和“探索”之间取得平衡。


利用型候选靠近模型认为的高性能区域，有利于尽快得到改进材料；探索型候选位于模型不确定但物理上仍合理的区域，用于扩展知识边界；重复型或近重复型候选用于估计制备和测量噪声。真实实验批次通常同时包含这几类样品，而不是只给出一个单一的“最佳配方”。主动学习批次中的候选类型见表 \ref{tab:active-learning-batch-types}。


\begin{formalTable}[htbp]
\caption{主动学习批次中的候选类型}
\label{tab:active-learning-batch-types}
\begin{tabularx}{\textwidth}{@{}YYY@{}}
\toprule
\textbf{候选类型} & \textbf{选择依据} & \textbf{在实验轮次中的作用} \\
\midrule
利用型 & 预测性能高且约束满足 & 尽快逼近可用材料 \\
探索型 & 模型不确定但物理上可行 & 扩展模型对未知区域的认识 \\
边界型 & 靠近工艺或相稳定边界 & 识别可制备窗口 \\
重复型 & 已有样品或近似重复配方 & 估计实验噪声和批次误差 \\
\bottomrule
\end{tabularx}
\end{formalTable}

实验噪声和批次效应是主动学习能否落地的关键。同一成分可能因混料、升温速率、保温时间、气氛、坩埚污染或仪器校准而表现不同。如果模型把这些差异全部当成材料本征规律，就会错误更新搜索方向。这意味着，样品批次、设备、制备程序、测量条件和失败原因都应进入数据记录。


失败样品同样提供信息。某类成分若反复析晶、开裂或产生杂相，说明该区域可能超出可制备边界。例如，某轮玻璃熔制实验中如果发现 La\textsubscript{2}O\textsubscript{3} 含量高于 12\% 的配方全部析晶，下一轮模型排序就应把这一区域标记为高风险，并降低相邻候选的推荐优先级。主动学习不应只学习“哪些候选性能高”，还应学习“哪些区域无法稳定制备”。这使下一轮实验决策同时考虑性能、不确定性、工艺风险和已有负结果。


\subsubsection*{\textbf{1. 案例示范：AR 光波导玻璃主动学习}}

AR 光波导玻璃可作为连续案例。初始数据由已有氧化物配方及其折射率、Tg、热膨胀系数和析晶情况构成；代理模型预测候选性质和不确定性；每轮选择若干高性能候选、若干未知区域候选和少量重复样品；实验结果再用于修正模型和玻璃形成边界。这样，模型输出被转化为可执行的实验批次。

主动学习与普通筛选最大的不同，是它把“下一步做什么实验”作为核心问题。普通筛选常常一次性给出候选排序，而主动学习会根据已经完成的实验不断调整模型。第一轮实验后，模型可能发现某个组成区域折射率较高但容易析晶；第二轮就应减少该区域中高风险候选的比例，同时在相邻区域选择少量探索样品。这样，实验不是简单验证模型，而是在持续修正模型。

在真实实验中，一轮实验通常不应只选择预测性能最高的候选。如果全部选择利用型样品，模型可能很快陷入局部区域，对未知空间了解不足；如果全部选择探索型样品，短期内又可能得不到性能改进。为此，一个合理批次通常同时包含高性能候选、模型不确定候选、边界候选和重复样品。重复样品表面上浪费名额，但它可估计实验噪声和批次误差，对判断模型是否真的改进很重要。

主动学习还需要记录失败类型。对于玻璃配方来说，失败可能表现为不能熔清、冷却后析晶、样品开裂、气泡过多、折射率测试失败或热膨胀系数不满足要求。不同失败类型对应不同边界：不能熔清可能是工艺温度问题，析晶可能是玻璃形成能力不足，开裂可能与热膨胀系数或退火制度有关。把这些失败统一写成“失败”会损失大量材料信息。主动学习实验结果的记录方式见表 \ref{tab:active-learning-result-records}。

\begin{formalTable}[htbp]
\caption{主动学习实验结果的记录方式}
\label{tab:active-learning-result-records}
\footnotesize
\begin{tabularx}{\textwidth}{@{}YYYY@{}}
\toprule
\textbf{实验结果类型} & \textbf{记录内容} & \textbf{对模型的作用} & \textbf{下一轮决策} \\
\midrule
性能达标 & 折射率、透过率、Tg、热膨胀系数 & 提供高价值正样本 & 在附近细化搜索 \\
性能一般但可成玻 & 完整性质数据和配方 & 扩展模型覆盖范围 & 可作为训练样本保留 \\
析晶或开裂 & 失败现象、热处理条件 & 标记不可制备边界 & 减少相邻高风险候选 \\
重复样品 & 批次、设备、测试误差 & 估计噪声水平 & 修正不确定性判断 \\
\bottomrule
\end{tabularx}
\end{formalTable}

AI 可辅助把主动学习输出整理成实验批次说明，但提示词中必须要求它保留利用、探索、边界和重复样品的比例，并提醒不能自动执行模型给出的批次。

\begin{samepage}
\begin{userbox}
我正在做 AR 光波导玻璃配方的主动学习优化。模型已经给出候选配方的预测折射率、不确定性、Tg、热膨胀系数、距离已有样本的远近以及是否满足基本约束。请帮我整理下一轮实验批次选择原则。

要求如下：
\begin{enumerate}
\item 区分利用型、探索型、边界型和重复型样品；
\item 说明每类样品为什么需要进入实验批次；
\item 提醒哪些实验元数据必须记录；
\item 说明失败样品如何反馈到下一轮模型；
\item 不要直接把模型最高分样品全部选为下一轮实验。
\end{enumerate}
\end{userbox}

\begin{aibox}
下一轮实验批次可按四类候选组合：

\begin{enumerate}
\item \textbf{利用型样品：}预测折射率较高、Tg 和热膨胀系数满足约束，用于尽快获得性能改进。
\item \textbf{探索型样品：}预测不确定性较高但物理上仍可行，用于扩展模型对未知组成区域的认识。
\item \textbf{边界型样品：}靠近玻璃形成或工艺边界，用于判断哪些区域容易析晶、开裂或无法熔制。
\item \textbf{重复型样品：}选择已有样品或近似重复样品，用于估计实验噪声、批次差异和测试误差。
\end{enumerate}

每个样品都应记录配方、原料批次、熔制温度、保温时间、冷却制度、样品状态、测试设备和失败原因。失败样品不应删除，而应作为不可制备区域或高风险区域的信息反馈给下一轮模型。
\end{aibox}
\end{samepage}

以下代码示例保留了主动学习最核心的计算关系：预测值体现“利用”，不确定性体现“探索”，Tg 和热膨胀系数等惩罚项体现工艺约束。代码输出的是待人工复核的实验批次，而不是可自动执行的最终方案。


\begin{lstlisting}[language=Python,escapeinside={(*@}{@*)}]
import pandas as pd
cand = pd.DataFrame({
    "id": ["G01", "G02", "G03", "G04", "G05", "G06"],
    "n_pred": [1.86, 1.91, 1.88, 1.83, 1.94, 1.89],
    "sigma": [0.012, 0.010, 0.006, 0.030, 0.018, 0.026],
    "Tg": [565, 548, 570, 535, 505, 552],
    "cte": [8.1, 8.7, 8.4, 8.9, 9.8, 8.0],
    "dist": [0.25, 0.18, 0.35, 0.12, 0.08, 0.10],
    "ok": [True, True, True, True, False, True],
})
target_cte = 8.3
cand["penalty"] = (
    (cand["Tg"] < 520).astype(int) * 0.20
    + (cand["cte"] - target_cte).abs() * 0.03
    + (~cand["ok"]).astype(int) * 10
)
cand["score"] = cand["n_pred"] + 0.4 * cand["sigma"] - cand["penalty"]
feasible = cand[cand["ok"]].copy()
use = feasible.sort_values("n_pred", ascending=False).head(2)
unused = feasible.drop(use.index)
explore = unused.sort_values("sigma", ascending=False).head(2)
unused = unused.drop(explore.index)
boundary = unused.sort_values("dist").head(1)
# (*@\textnormal{保留 G03 作为重复样品，用于评估实验重复性}@*)
repeat = feasible[feasible["id"].eq("G03")].copy()
batch = pd.concat([
    use.assign(kind="use"),
    explore.assign(kind="explore"),
    boundary.assign(kind="boundary"),
    repeat.assign(kind="repeat"),
])
cols = ["id", "kind", "n_pred", "sigma", "Tg", "cte", "score"]
print(batch[cols])
\end{lstlisting}

\manualfigcaption{\textbf{代码示例：}主动学习实验批次评分示例}

\subsubsection*{\textbf{小结}}

本节说明候选筛选应按照由低成本到高成本、由粗判断到精验证的漏斗推进。
规则过滤、代理模型、物理计算和实验验证各有作用，每一级都应保留通过和排除理由。
主动学习把预测值、不确定性和实验反馈结合起来，失败样品和噪声记录同样是有价值的数据。

\subsection*{习题}

\begin{enumerate}[leftmargin=*]
\item 高通量筛选为什么常采用分层漏斗？
\item 比较规则过滤、代理模型、DFT 和实验验证的作用。
\item 主动学习中的“利用”和“探索”分别指什么？
\end{enumerate}

\clearpage
\section{生成设计与候选验证}


\begin{sectionkeypoints}
\item 生成设计改变的是候选产生方式，可在目标条件下提出新的组成、结构或配方，但不能取消材料约束和实验验证。
\item 材料表示决定生成模型能够处理的对象。晶体、玻璃、分子和聚合物需要不同表示方式，不宜简单套用同一种输入格式。
\item 条件生成需要把目标性能转化为可计算条件，例如带隙范围、形成能、元素体系、成本边界或供应链限制。
\item 生成候选必须经过有效性、稳定性、新颖性、可合成性和应用性等多层验证，不应只看生成数量。
\item 生成设计通常需要与筛选、主动学习、DFT 计算和实验验证串联使用，才能形成可信的材料发现流程。
\end{sectionkeypoints}

\subsection{条件生成与验证流程}

候选筛选只能在已经列出的空间中寻找材料。如果候选库没有包含某类新结构或新配方，筛选再快也无法发现它。生成设计尝试学习已有材料的分布，并在目标条件下提出新的组成、结构或配方。它改变的是候选产生方式，而不是取消物理约束和实验验证。


材料表示决定生成模型能够处理什么。分子可使用字符串或分子图，晶体需要同时表示元素种类、周期性晶格和原子坐标，玻璃与非晶材料常以成分和短程有序统计表示，聚合物还要考虑单体序列和拓扑。生成表示若不能表达周期性、价态和局域配位，模型容易产生形式正确但没有材料意义的候选。表中的 VAE\footnote{Variational Autoencoder，变分自编码器。} 指变分自编码器，GAN\footnote{Generative Adversarial Network，生成对抗网络。} 指生成对抗网络，自回归模型\footnote{Autoregressive Model，自回归模型，按序列顺序逐步生成数据。} 按步骤生成序列或结构片段，扩散模型\footnote{Diffusion Model，扩散模型，从噪声出发逐步去噪生成数据。} 从噪声中逐步生成结构，SMILES\footnote{Simplified Molecular Input Line Entry System，简化分子线性输入系统。} 是分子线性表示方法。不同生成模型路线与材料对象的对应关系见表 \ref{tab:generative-model-routes}。


\begin{formalTable}[htbp]
\caption{生成模型路线与材料对象}
\label{tab:generative-model-routes}
\begin{tabularx}{\textwidth}{@{}YYYY@{}}
\toprule
\textbf{模型路线} & \textbf{材料生成中的直觉} & \textbf{典型对象} & \textbf{需要特别注意} \\
\midrule
VAE & 把材料压缩到连续潜空间并在其中优化 & 分子、聚合物、部分晶体表示 & 潜空间是否有物理意义 \\
GAN & 生成器与判别器对抗，逼近真实样本分布 & 显微组织、分子、结构图像 & 训练稳定性与约束不足 \\
自回归模型 & 按序列或步骤逐步生成 & SMILES、合成路线、结构 token & 语法有效性和错误传播 \\
扩散模型 & 从噪声逐步去噪生成结构 & 晶体、分子三维结构 & 周期边界、对称性和坐标处理 \\
混合工作流 & 生成、筛选、DFT、实验闭环结合 & 无机晶体、玻璃、催化剂 & 验证流程比模型名称更重要 \\
\bottomrule
\end{tabularx}
\end{formalTable}

VAE、GAN、自回归模型和扩散模型可采用不同方式学习材料分布。VAE 更适合把分子或配方压缩到连续潜空间中再优化；GAN 常用于生成图像式显微组织或分子结构，但训练稳定性和约束控制需要特别检查；自回归模型适合 SMILES、合成路线或结构 token 这类有顺序的表示；扩散模型适合逐步生成分子或晶体三维结构，但必须处理周期性边界、坐标对称性和原子距离合理性。本章不展开各模型的数学细节，只保留共同逻辑：模型从已有材料中学习可行区域，目标条件用于引导生成方向，化学和结构约束用于排除无效样本，正向模型和高精度计算用于评价生成结果。


\subsubsection*{\textbf{1. 案例示范：发光材料条件生成}}

条件生成是逆向设计的关键。目标条件可以是带隙范围、形成能、体模量、磁性、元素体系或供应链限制。以发光材料为例，目标波长可转化为带隙范围，生成模型提出候选组成和结构，再由第\ref{chap:materials-cognition-prediction}章的带隙模型、DFT 和实验逐级验证。目标带隙只是入口，缺陷、掺杂、热稳定和器件兼容性仍需单独检查。

对于发光材料设计，条件生成可从两个层次理解。第一层是性能条件，例如目标发光波长、带隙范围、热稳定性和缺陷容忍度。第二层是材料边界，例如元素体系、晶体结构类型、掺杂元素、成本和安全要求。如果只给生成模型一个带隙条件，模型可能提出带隙合适但不稳定、含风险元素或难以合成的候选。因此，条件生成需要同时输入目标和边界。

发光材料还具有明显的“性质链条”。目标颜色对应发光波长，发光波长与带隙有关，但材料是否真正发光还取决于缺陷、能级、载流子复合和非辐射损失。换言之，带隙只是候选进入下一步验证的第一道门槛。一个候选材料即使带隙接近目标，如果存在深能级缺陷、热稳定性差或与器件结构不兼容，也不宜直接认为适合应用。

由此可见，生成设计更适合作为“候选扩展器”。它可帮助研究者跳出已有数据库或人工枚举空间，提出一些值得进一步检查的新组成或新结构。但候选是否有意义，要经过规则检查、正向预测、结构弛豫、高精度计算和实验验证。生成数量本身没有价值，真正有价值的是通过层层验证后仍然保留下来的候选。发光材料条件生成中的约束类型见表 \ref{tab:luminescent-material-constraints}。

\begin{formalTable}[htbp]
\caption{发光材料条件生成中的约束类型}
\label{tab:luminescent-material-constraints}
\footnotesize
\begin{tabularx}{\textwidth}{@{}YYYY@{}}
\toprule
\textbf{条件类型} & \textbf{发光材料中的含义} & \textbf{可能数据字段} & \textbf{验证方式} \\
\midrule
目标条件 & 希望获得的发光颜色或带隙范围 & 目标波长、带隙范围 & 光谱测试、带隙预测、DFT \\
结构条件 & 候选结构是否合理 & 空间群、原子坐标、局域配位 & 结构弛豫、距离检查、声子计算 \\
稳定条件 & 候选是否可能存在 & 形成能、凸包能、热稳定性 & 高精度计算、热分析、实验合成 \\
应用条件 & 是否适合器件环境 & 缺陷、热稳定、外延兼容 & 发光测试、器件测试、可靠性测试 \\
\bottomrule
\end{tabularx}
\end{formalTable}

在使用 AI 辅助生成设计时，提示词应强调“提出候选后的验证链条”，而不是让 AI 直接给出材料名称。示例提示词可这样组织。

\begin{samepage}
\begin{userbox}
我想设计一种目标发光波长在蓝光范围附近的无机发光材料。请帮我把这个目标转化为条件生成设计流程。

要求如下：
\begin{enumerate}
\item 说明目标波长如何转化为带隙或能级相关条件；
\item 说明生成模型可提出哪些类型的候选，例如组成、掺杂或晶体结构；
\item 列出生成后必须检查的化学、结构、稳定性和应用约束；
\item 说明为什么不应只根据带隙判断材料是否适合发光器件；
\item 不要直接虚构某个“最佳材料”。
\end{enumerate}
\end{userbox}

\begin{aibox}
可把任务拆成以下流程：

\begin{enumerate}
\item 目标条件设定：将目标发光波长转化为近似带隙范围，同时记录发光效率、热稳定性和器件兼容性等要求。
\item \textbf{候选生成：}生成模型可提出候选组成、掺杂元素、晶体结构或局域配位环境，但这些结果只是候选，不是最终材料。
\item \textbf{有效性检查：}检查化学式、价态、原子距离、晶格体积和结构格式是否合理，排除明显无意义候选。
\item 稳定性验证：用形成能、凸包能、结构弛豫和必要的声子或热稳定性分析判断候选是否可能存在。
\item \textbf{应用验证：}进一步检查缺陷、非辐射复合、热稳定、外延匹配和器件环境适配性。
\end{enumerate}

带隙合适只能说明候选可能满足发光颜色要求，不能说明它一定具有高发光效率或可制备性。生成模型的作用是扩展候选空间，最终判断仍然需要计算和实验验证。
\end{aibox}
\end{samepage}

条件生成还需要保留“负面条件”。很多时候，研究者不仅知道想要什么，也知道不想要什么。例如，不希望候选含有某些高风险元素，不希望晶体结构出现过短原子距离，不希望配方超出已知可制备范围，不希望候选与训练集中已有材料完全重复。这些负面条件如果不写入生成和筛选流程，模型可能不断提出形式上满足目标性能、但材料上明显不可接受的候选。

具体而言，条件生成的输入可分为三类：目标条件、边界条件和排除条件。目标条件推动模型朝目标性能靠近；边界条件保证候选处于可研究范围；排除条件提前去掉明显不可接受的结果。对于材料科学任务，三类条件同样重要。只写目标条件，容易得到“性能诱人但不可用”的候选；只写排除条件，又可能使搜索空间过窄，失去发现新材料的机会。


\subsection{候选验证与路线组合}

\subsubsection*{\textbf{1. 验证流程：从有效性到应用性}}

生成候选首先要检查有效性：化学式、价态、原子距离、晶格体积和周期性是否合理；其次检查稳定性：结构弛豫后是否保持、形成能和声子是否可接受；再次检查新颖性：候选是否只是训练集或数据库中的重复结构；最后检查可合成性和应用条件。任何一级不通过，都不应把候选直接称为材料发现。生成设计的主要验证维度见表 \ref{tab:generative-design-validation}。


\begin{formalTable}[htbp]
\caption{生成设计的验证维度}
\label{tab:generative-design-validation}
\begin{tabularx}{\textwidth}{@{}YYYY@{}}
\toprule
\textbf{验证指标} & \textbf{问题} & \textbf{常用判断} & \textbf{为什么重要} \\
\midrule
有效性 & 结构是否符合基本规则？ & 价态、距离、格式、晶格范围 & 避免无意义候选 \\
稳定性 & 候选是否可能存在？ & 形成能、凸包能、声子、动力学 & 减少纸上材料 \\
新颖性 & 是否只是训练集复述？ & 结构匹配、组成相似度、数据库检索 & 证明发现价值 \\
可合成性 & 实验能否做出来？ & 合成路线、温度窗口、前驱体、杂相 & 连接真实材料 \\
应用性 & 能否满足场景要求？ & 服役环境、成本、可靠性、规模化 & 决定工程价值 \\
\bottomrule
\end{tabularx}
\end{formalTable}

可靠的验证链条通常分为四层。第一层是格式、几何和重复性检查；第二层是快速正向模型打分；第三层是 DFT、分子动力学或热力学精算；第四层是实验合成、结构表征和性能测试。研究报告还应说明每一级候选的通过率，区分生成数量、有效数量、稳定数量和实验验证数量。


筛选、主动学习和生成设计可串联使用：生成模型扩展候选空间，规则和代理模型快速排除明显不合理候选，主动学习决定哪些候选最值得进入下一轮高精度计算或实验。最终判断仍然来自材料约束和验证证据，而不是模型名称。

可合成性是生成设计最容易被低估的环节。一个结构在计算上稳定，不代表实验上容易合成。实验可达性受到前驱体、温度、压力、气氛、反应路径、动力学障碍和杂相竞争影响。对于玻璃材料，候选配方还必须能熔制、澄清、冷却成玻璃并避免严重析晶；对于晶体材料，候选结构还要考虑相竞争和合成窗口。可合成性判断不足，是“纸上材料”大量出现的重要原因。生成候选在不同证据阶段的表述边界见表 \ref{tab:generated-candidate-evidence}。

\begin{formalTable}[htbp]
\caption{生成候选在不同证据阶段的表述边界}
\label{tab:generated-candidate-evidence}
\footnotesize
\begin{tabularx}{\textwidth}{@{}YYYY@{}}
\toprule
\textbf{候选阶段} & \textbf{可采用的表述} & \textbf{不能怎样表述} & \textbf{需要补充证据} \\
\midrule
模型生成 & 生成了一批候选结构或配方 & 发现了一批新材料 & 格式、价态和几何检查 \\
模型打分 & 获得高优先级候选 & 证明候选性能优异 & 正向模型误差和适用范围 \\
高精度计算 & 候选在计算上接近稳定 & 候选一定可合成 & 结构弛豫、凸包能、相竞争 \\
实验合成 & 候选已被制备并测试 & 可直接工程应用 & 服役、可靠性和规模化验证 \\
\bottomrule
\end{tabularx}
\end{formalTable}

这一部分也适合用 AI 辅助整理验证清单。关键是要求 AI 区分证据层级，避免把模型输出写成最终结论。

\begin{samepage}
\begin{userbox}
我用生成模型得到了一批候选无机晶体结构。请帮我设计一个候选验证清单，判断这些候选能否进入后续计算或实验。

要求如下：
\begin{enumerate}
\item 按有效性、稳定性、新颖性、可合成性和应用性分层；
\item 每一层说明需要检查哪些信息；
\item 区分“模型推荐候选”和“已验证材料”的表述；
\item 提醒如何记录每一级候选数量和通过率；
\item 不要把生成模型输出直接称为材料发现。
\end{enumerate}
\end{userbox}

\begin{aibox}
可建立五层验证清单：

\begin{enumerate}
\item \textbf{有效性：}检查化学式、价态、原子距离、晶格体积、周期性和文件格式是否合理。
\item 稳定性：进行结构弛豫，检查形成能、凸包能、声子稳定性或分子动力学稳定性。
\item \textbf{新颖性：}与训练集、材料数据库和文献记录比对，判断是否为已有材料或高度相似结构。
\item \textbf{可合成性：}分析前驱体、温度窗口、相竞争、动力学路径和实验设备条件。
\item \textbf{应用性：}结合目标场景检查性能、成本、安全、服役环境和长期可靠性。
\end{enumerate}

报告中应分别写出生成数量、有效数量、稳定数量、新颖数量、进入实验数量和实验验证数量。模型生成和模型打分只能说明候选值得进一步检查，不宜直接等同于材料发现。
\end{aibox}
\end{samepage}

在实际写作中，还应避免把验证流程写成一条单向直线。候选在任何一级验证失败，都可能需要返回前面的步骤。例如，结构弛豫后候选发生明显畸变，说明生成结构可能不合理；DFT 显示候选远离稳定区域，说明筛选规则需要加强稳定性约束；实验中反复出现杂相，说明可合成性判断不足。这些失败不是无效信息，而是帮助修正候选生成和筛选规则的重要反馈。

因此，候选验证更像是一个带反馈的流程。生成模型提出候选，验证流程排除或确认候选，失败原因再返回去修改约束、训练数据或生成条件。只有把反馈写清楚，生成设计才不是“模型输出一批材料”，而是“模型、计算和实验共同收敛到更可靠候选”的过程。

\subsubsection*{\textbf{小结}}

本节说明生成设计的作用是扩展候选空间，而不是直接给出最终材料。
条件生成需要同时设置目标条件、边界条件和排除条件，并选择适合材料对象的表示方式。
生成候选必须经过有效性、稳定性、新颖性、可合成性和应用性验证，失败结果也应回流流程。

\subsection*{习题}

\begin{enumerate}[leftmargin=*]
\item 生成候选为什么不能直接称为材料发现？
\item 为发光材料生成任务设计三类验证指标。
\end{enumerate}

\clearpage
\section{可信发现与路线选择}


\begin{sectionkeypoints}
\item 可信材料发现需要区分数据、模型、材料、验证和应用等不同证据层级，不应把模型推荐直接等同于真实材料发现。
\item 方法选择应由材料问题决定。候选库明确时优先筛选，实验昂贵时优先主动学习，候选库不足时再考虑生成模型。
\item 安全、环境、供应链、成本、可加工性和长期可靠性应作为并列约束保留，不能被压缩成无法解释的单一分数。
\item 第\ref{chap:materials-cognition-prediction}章的材料表示、性质预测、不确定性、机器学习势和可解释性，是第\ref{chap:materials-design-discovery}章候选筛选、生成验证和闭环优化的基础。
\item 全模块的完整逻辑是：定义目标、限定边界、提出候选、逐层验证、实验反馈，并在新数据基础上继续更新模型。
\end{sectionkeypoints}

\subsection{可靠评价与方法选择}

人工智能能够扩大材料搜索范围，也会放大数据偏差和模型外推风险。材料数据库中的计算参数、实验条件和样品质量并不完全一致，成功结果通常多于失败结果，热门体系的数据也远多于冷门体系。模型会继承这些差异，因此候选排名并不是客观真理，而是特定数据和模型边界下的判断。


材料性能还高度依赖结构、工艺和服役环境。同一化学式在不同缺陷、微结构和热处理条件下可能表现迥异，如果数据只记录成分和最终性能，模型就会把未记录的工艺差异当成噪声。逆向设计因此必须明确材料对象和数据适用范围，不应把材料简化为一个化学式或单一分数。


“纸上材料”指计算或模型预测较好，但缺少可合成性、规模化或应用验证的候选。真正的材料发现需要区分不同证据层级：模型预测提供线索，高精度计算提高物理可信度，实验合成证明材料可达，器件测试和长期稳定性才说明应用潜力。材料发现结果的分层评价见表 \ref{tab:materials-discovery-evaluation}。


\begin{formalTable}[htbp]
\caption{材料发现结果的分层评价}
\label{tab:materials-discovery-evaluation}
\begin{tabularx}{\textwidth}{@{}YYY@{}}
\toprule
\textbf{评价层面} & \textbf{核心判断} & \textbf{可靠表述示例} \\
\midrule
数据 & 来源、去重、划分是否清楚 & 在某数据库和划分方式下模型表现良好 \\
模型 & 是否有基线、不确定性和边界说明 & 模型推荐了若干高优先级候选 \\
材料 & 是否满足化学、结构和稳定性约束 & 候选在计算上接近稳定且结构合理 \\
验证 & 是否经过 DFT、实验或器件测试 & 某候选已合成并测得目标性质 \\
应用 & 是否考虑成本、安全和可靠性 & 候选具备进一步器件验证价值 \\
\bottomrule
\end{tabularx}
\end{formalTable}

可信评价的关键，是把“模型表现好”和“材料真的好”分开。模型表现好，通常指在某个数据集、某种划分方式和某个评价指标下，预测误差较小或排序效果较好。材料真的好，则意味着候选在真实制备、结构稳定、性能测试和服役环境中都能经受验证。二者之间存在距离。材料人工智能研究中常见的问题，是把模型层面的高分候选直接写成材料发现结果，忽略了验证证据不足。

评价结果时，还要说明模型适用范围。例如，模型是在无机晶体数据库上训练的，还是在某类玻璃配方上训练的；数据是计算数据、实验数据，还是二者混合；测试集划分是随机划分，还是按材料体系、时间或组成空间划分。如果训练集和测试集高度相似，模型指标可能看起来很好，但对真正未知材料的预测能力并不一定强。对于逆向设计来说，外推风险比普通插值预测更突出，因为候选往往位于已有数据边界附近。

不确定性也是可信评价的重要组成部分。一个候选预测性能高，但模型不确定性也高，说明它可能值得探索，但不宜直接排在确定性较高的候选之前。相反，一个候选预测性能略低但不确定性小、约束满足、可合成性强，也可能更适合作为近期实验对象。候选报告中应同时给出预测值、不确定性、约束是否满足和推荐理由。材料发现结果的可靠表述方式见表 \ref{tab:reliable-discovery-expression}。

\begin{formalTable}[htbp]
\caption{材料发现结果的可靠表述方式}
\label{tab:reliable-discovery-expression}
\footnotesize
\begin{tabularx}{\textwidth}{@{}YYYY@{}}
\toprule
\textbf{常见表述} & \textbf{风险} & \textbf{更稳妥的写法} & \textbf{还需补充} \\
\midrule
模型发现了某材料 & 夸大模型作用 & 模型推荐该候选进入验证 & 计算或实验验证 \\
该候选性能优异 & 忽略不确定性和边界 & 在当前模型和数据范围内预测值较高 & 不确定性、误差范围 \\
生成了大量新材料 & 混淆生成与发现 & 生成了若干待验证候选 & 去重、新颖性和稳定性检查 \\
候选可用于器件 & 跳过应用验证 & 候选具备进一步器件验证价值 & 器件测试和长期可靠性 \\
\bottomrule
\end{tabularx}
\end{formalTable}

候选报告还需要特别注意“数量表述”。生成模型可能一次输出几千个候选，但这并不意味着发现了几千种新材料。更准确的写法应区分：生成候选数量、格式有效数量、结构合理数量、计算稳定数量、去重后新颖数量、进入实验数量和实验验证成功数量。只有经过实验合成和性能测试的候选，才能接近“发现”的表述；只经过模型打分的候选，只能称为“推荐候选”或“高优先级候选”。

新颖性检查也属于可信评价的一部分。生成模型可能复述训练集中已经存在的材料，或者只对已有结构做很小改动。如果不做数据库检索和结构匹配，就可能把已有材料误写成新发现。对于晶体材料，可比较化学组成、空间群、晶格参数和结构相似度；对于玻璃或配方材料，可比较组分比例和已有专利、文献或实验记录。新颖性不是单靠模型输出判断的，而需要与数据库和文献记录对照。

\subsubsection*{\textbf{1. 路线选择：问题决定方法}}

方法选择应由问题决定，而不是由模型新旧决定。候选空间已经明确时，高通量筛选通常最直接；实验成本高且每轮样品有限时，主动学习更适合；需要突破已有候选库时，生成模型才体现优势。玻璃形成边界、相稳定、毒性和设备温度等硬约束，应尽可能在搜索前加入，而不是在模型给出结果后再补救。不同问题情形下的方法选择见表 \ref{tab:inverse-design-method-selection}。


\begin{formalTable}[htbp]
\caption{材料逆向设计的方法选择}
\label{tab:inverse-design-method-selection}
\begin{tabularx}{\textwidth}{@{}YYYY@{}}
\toprule
\textbf{问题情形} & \textbf{优先方法} & \textbf{理由} & \textbf{典型材料场景} \\
\midrule
候选库明确 & 高通量筛选 & 流程清楚、可解释、验证成本低 & 从数据库筛选稳定光电材料 \\
实验昂贵且小样本 & 贝叶斯优化/主动学习 & 每一轮实验都要最大化信息量 & 玻璃配方、合金热处理、催化条件 \\
候选库不足 & 生成模型 + 后筛选 & 需要提出库外候选 & 新晶体结构、功能分子、聚合物序列 \\
安全或法规约束强 & 规则约束 + 人工审核 & 先排除不可接受候选 & 含毒元素、稀缺元素或高风险反应 \\
\bottomrule
\end{tabularx}
\end{formalTable}

可靠评价还应考虑安全、环境和供应链。候选若含有有毒、稀缺或高风险元素，即使单项性能突出，也可能不适合真实应用。相应地，性能、成本、环境影响、可加工性和生命周期不应被压缩成一个无法解释的总分，而应作为并列约束保留在候选报告中。

方法选择还应考虑课程或项目阶段。对于入门学习，最合适的方法往往不是最新的生成模型，而是可解释、可复核的筛选流程。因为研究者需要先理解数据字段、模型输入、约束条件和验证证据之间的关系。对于研究生或真实科研项目，如果已有稳定的数据积累和实验反馈机制，再引入主动学习或生成模型才更合理。否则，复杂模型可能只是增加技术难度，并不能提高材料发现质量。

在写作时，可把方法选择写成一组判断问题：候选空间是否已经明确？每次实验成本是否很高？是否需要提出已有数据库之外的新候选？是否存在严格安全或工艺边界？是否有能力进行后续高精度计算和实验验证？这些问题比“选择哪个先进算法”更重要。

这一节中的方法选择和可信评价也可借助 AI 辅助整理，但提示词应要求它给出理由和证据边界，不应只推荐一个模型名称。

\begin{samepage}
\begin{userbox}
我正在写材料逆向设计方案，候选空间包括已有数据库中的无机晶体和少量可枚举组成。实验成本较高，每轮只能验证 5～10 个样品。请帮我判断优先使用高通量筛选、主动学习还是生成模型。

要求如下：
\begin{enumerate}
\item 根据候选空间、实验成本、数据量和验证能力进行判断；
\item 说明每种方法适合和不适合的情况；
\item 给出推荐路线，但必须说明理由；
\item 提醒结果报告中应如何避免夸大模型作用；
\item 不要只因为生成模型较新就优先推荐它。
\end{enumerate}
\end{userbox}

\begin{aibox}
在这个任务中，候选空间已有数据库和可枚举组成，说明不一定需要首先使用生成模型。更合理的路线可以是：

\begin{enumerate}
\item 第一阶段使用高通量筛选：先用规则、正向模型和稳定性判断从已有候选库中排除明显不合理候选。
\item \textbf{第二阶段使用主动学习：}由于每轮实验只能验证 5～10 个样品，应让模型同时考虑预测性能和不确定性，选择高价值实验批次。
\item 生成模型暂时作为补充：只有当已有候选库不能覆盖目标性能区域，或者需要提出库外结构时，再考虑生成模型。
\end{enumerate}

报告中应写成“模型推荐若干高优先级候选，并通过后续计算或实验验证”，而不是“模型发现了最终材料”。如果没有实验或高精度计算支持，应明确说明候选仍处于待验证阶段。
\end{aibox}
\end{samepage}

\subsection{能力复用与闭环迭代}

\subsubsection*{\textbf{1. 两章闭环：正向模型服务逆向设计}}

两章之间形成正反闭环。第\ref{chap:materials-cognition-prediction}章把材料数据转化为描述符或结构图，训练正向模型并解释模型边界；第\ref{chap:materials-design-discovery}章把这些模型作为候选筛选和生成验证的评价器，再用新计算和实验结果更新数据。没有正向预测，逆向搜索缺少可靠评价；没有逆向设计，正向模型容易停留在被动打分。第\ref{chap:materials-cognition-prediction}章能力在第\ref{chap:materials-design-discovery}章中的复用关系见表 \ref{tab:chapter-one-capability-reuse}。


\begin{formalTable}[htbp]
\caption{第\ref{chap:materials-cognition-prediction}章能力在第\ref{chap:materials-design-discovery}章中的复用}
\label{tab:chapter-one-capability-reuse}
\begin{tabularx}{\textwidth}{@{}YYY@{}}
\toprule
\textbf{第\ref{chap:materials-cognition-prediction}章能力} & \textbf{第\ref{chap:materials-design-discovery}章中的用途} & \textbf{材料例子} \\
\midrule
性质预测 & 作为筛选和生成后的快速打分器 & 带隙筛选 microLED 候选 \\
材料表示 & 决定候选能否被模型有效评价 & 晶体图、玻璃成分、分子图 \\
不确定性量化 & 指导主动学习下一步实验 & 高折射率玻璃小样本优化 \\
机器学习势 & 验证候选结构和高温行为 & 超高温陶瓷和高熵合金 \\
可解释性 & 把模型推荐转化为材料规律 & 形成能力、熔点或缺陷趋势分析 \\
\bottomrule
\end{tabularx}
\end{formalTable}

完整逻辑可概括为：目标性质定义问题，材料知识限定边界，筛选、主动学习或生成模型提出候选，正向模型和物理计算逐层评价，实验结果决定最终结论并反馈到数据集。人工智能的价值在于提高搜索与决策效率，材料科学的责任则是保证问题、约束和证据真实可靠。

闭环迭代要落实到可操作的数据管理上。每一轮推荐都应记录模型版本、训练数据版本、候选来源、预测值、不确定性、约束检查结果和推荐理由；每一轮验证后，还要把成功样品、失败样品、测试条件和失败原因按同一格式回写数据表。这样，下一轮模型更新才知道哪些变化来自材料本身，哪些变化来自数据版本、实验批次或测试条件。

例如，在高折射率玻璃任务中，第\ref{chap:materials-cognition-prediction}章中讨论的数据整理和描述符构建可帮助建立折射率、Tg 和热膨胀系数预测模型；第\ref{chap:materials-design-discovery}章则利用这些模型筛选候选配方，并通过主动学习选择下一轮熔制实验。在高熵合金涂层任务中，第\ref{chap:materials-cognition-prediction}章的硬度预测、图神经网络和机器学习势可提供性能估计与结构理解；第\ref{chap:materials-design-discovery}章则进一步用这些模型决定哪些成分或硬质相值得验证。这样的闭环，才是人工智能真正进入材料发现流程的方式。

闭环迭代还要求新数据及时回流。实验失败、性能一般、模型预测错误的样品，都应被记录下来。只把成功样品加入数据库，会使模型越来越偏向成功区域，无法学习失败边界；只记录最终性能而不记录工艺条件，也会使模型无法区分成分效应和工艺效应。闭环发现不仅是“模型推荐--实验验证”，还包括“实验记录--数据更新--模型修正”。材料智能发现闭环中的数据回流内容见表 \ref{tab:closed-loop-data-feedback}。

\begin{formalTable}[htbp]
\caption{材料智能发现闭环中的数据回流}
\label{tab:closed-loop-data-feedback}
\footnotesize
\begin{tabularx}{\textwidth}{@{}YYYY@{}}
\toprule
\textbf{闭环环节} & \textbf{主要任务} & \textbf{需要记录的信息} & \textbf{常见问题} \\
\midrule
模型推荐 & 根据目标和约束提出候选 & 预测值、不确定性、推荐理由 & 只给排名，不说明边界 \\
计算验证 & 用更高精度方法复核候选 & 结构、能量、稳定性、失败原因 & 只保留通过候选 \\
实验验证 & 合成、表征和性能测试 & 工艺、样品状态、测试条件 & 缺少失败样品记录 \\
数据回流 & 更新数据集并重训模型 & 新数据、元数据、版本信息 & 数据来源不可追溯 \\
\bottomrule
\end{tabularx}
\end{formalTable}

AI 可辅助整理闭环报告结构，尤其是检查是否遗漏数据来源、验证证据和失败记录。但同样需要明确：AI 只能帮助组织报告，不宜替代实验和模型审查。

\begin{samepage}
\begin{userbox}
我想把一个材料智能发现项目写成闭环流程报告。项目包括正向模型预测、候选筛选、实验验证和数据回流。请帮我设计报告结构。

要求如下：
\begin{enumerate}
\item 说明第\ref{chap:materials-cognition-prediction}章中的材料表示、性能预测和模型解释如何服务第\ref{chap:materials-design-discovery}章的逆向设计；
\item 按模型推荐、计算验证、实验验证和数据回流四个环节组织；
\item 提醒每个环节需要记录哪些信息；
\item 特别说明失败样品和预测错误样品为什么也要记录；
\item 不要把闭环写成模型自动发现材料。
\end{enumerate}
\end{userbox}

\begin{aibox}
闭环流程报告可包括四部分：

\begin{enumerate}
\item \textbf{模型推荐：}说明使用了哪些材料表示和正向预测模型，给出候选的预测值、不确定性和推荐理由。
\item \textbf{计算验证：}记录结构弛豫、稳定性判断、形成能或其他高精度计算结果，并保留未通过候选的原因。
\item \textbf{实验验证：}记录制备工艺、样品状态、表征结果、性能测试条件和失败现象，不应只记录成功样品。
\item \textbf{数据回流：}说明新数据如何加入数据库，如何更新训练集，模型是否重新训练，以及新模型与旧模型相比是否改进。
\end{enumerate}

失败样品和预测错误样品同样重要，因为它们可帮助模型学习不可制备区域、实验噪声和模型外推边界。材料智能发现不是模型自动给出答案，而是模型、计算、实验和数据更新共同形成的迭代过程。
\end{aibox}
\end{samepage}

如果把两章放在一个完整模块中理解，第\ref{chap:materials-cognition-prediction}章更像是“建立评价能力”，第\ref{chap:materials-design-discovery}章更像是“利用评价能力做决策”。没有第\ref{chap:materials-cognition-prediction}章的数据整理、表示学习和模型评估，第\ref{chap:materials-design-discovery}章的筛选和生成就缺少可靠判断依据；没有第\ref{chap:materials-design-discovery}章的候选选择和实验反馈，第\ref{chap:materials-cognition-prediction}章的模型又容易停留在已有数据上的被动预测。二者结合起来，才能形成从材料认知到材料发现的完整教学闭环。

最后一节不只是总结第\ref{chap:materials-design-discovery}章，也是在说明整个模块的逻辑边界：人工智能可帮助研究者更快地整理材料数据、构建候选、预测性能、安排验证和更新模型，但不应跳过材料科学中的基本问题。材料是否真实存在，能否制备出来，是否满足应用场景，仍然要靠实验、计算和工程约束共同判断。

\subsubsection*{\textbf{小结}}

本节强调可信材料发现需要区分模型推荐、计算证据、实验验证和应用潜力。
方法选择应由问题决定，候选库明确时优先筛选，实验昂贵时优先主动学习，候选不足时再考虑生成。
第\ref{chap:materials-cognition-prediction}章建立评价能力，第\ref{chap:materials-design-discovery}章利用评价能力做决策，数据回流使整个模块形成闭环。

\subsection*{习题}

\begin{enumerate}[leftmargin=*]
\item 数据回流为什么是材料智能发现闭环的关键？
\item 如何根据候选库、实验成本和数据量选择材料发现路线？
\end{enumerate}

\section*{本章小结}

本章讨论了材料设计与智能发现的基本思路。逆向设计首先要求将应用需求转化为可计算目标、变量和约束；候选筛选通过规则、代理模型、高精度计算和实验验证逐层缩小范围；主动学习利用预测值和不确定性安排下一批实验或计算任务；生成设计能够提出新的结构或配方，但必须接受有效性、稳定性、新颖性、可合成性和应用条件等多层验证。全章强调，智能发现不是单次模型输出，而是目标定义、候选构建、分层验证和数据回流共同组成的闭环过程。

% 文件中不加载任何宏包，所有样式均调用项目现有的 package.tex。
\part*{微模块5\quad 人工智能在材料科学中的应用}
\addcontentsline{toc}{part}{微模块5\texorpdfstring{\quad}{ }人工智能在材料科学中的应用}

\chapter{材料认知与性能预测}
\label{chap:materials-cognition-prediction}


理解材料性能的来源，是材料科学研究的核心任务。材料的力学、热学、电学、磁学和光学性能，往往不是由单一因素决定的，而是受到成分组成、微观结构和制备工艺的共同影响。随着研究对象从简单合金、陶瓷和半导体扩展到高熵合金、复合材料、低维材料和复杂功能材料，材料候选空间迅速增大，单纯依靠实验试错已经难以系统揭示其中的规律。在这一背景下，人工智能方法的价值，主要体现在对分散实验数据、计算数据和文献知识的组织与学习上，使模型能够从中提取“成分--结构--工艺--性能”之间的关联，从而辅助材料性能预测和后续实验设计。

本章围绕“材料认知与性能预测”展开，重点解决四个问题。第一节说明材料人工智能为什么需要可靠数据基础，以及实验、计算、文献和数据库数据应如何整理。第二节讨论材料成分和晶体结构如何转化为描述符或图表示，使模型能够读取材料信息。第三节介绍描述符模型、图神经网络和机器学习原子间势（下文简称“机器学习势”）如何服务硬度预测、性质预测和熔点模拟。第四节进一步说明模型解释与物理知识嵌入的作用，帮助研究者判断模型结果是否符合材料机理和实验事实。通过本章学习，研究者需要建立一个基本认识：人工智能并不是替代材料理论和实验验证，而是与数据、计算和材料机理共同构成辅助材料认知的新工具。




\clearpage
\section{范式变革与数据基础}


\begin{sectionkeypoints}
\item 材料人工智能中的数据并不只是单个性能数值，而是包含材料组成、结构信息、制备工艺、测试条件、数据来源和目标性能的完整记录。
\item 材料数据主要来自实验、计算、文献和数据库四类渠道。不同来源的数据各有优势，也各有局限，使用时需要结合具体研究任务进行判断。
\item 材料数据库为模型训练、性能预测和候选筛选提供了重要的数据基础，但数据库中的数据不宜直接等同于最终模型输入，还需要经过字段筛选、来源核验和格式整理。
\item 在正式建模之前，应先整理原始数据清单，明确记录材料是什么、结构是什么、性能数值是多少、数据来自哪里，以及采用了什么计算方法或实验条件。
\item AI 可辅助解释数据库字段、设计表格框架和检查数据规范性，但不宜替代数据库原始记录、实验记录和论文原文核验。材料数据整理的核心要求是清楚、可追溯、可检查。
\end{sectionkeypoints}

\subsection{实验试错与数据驱动}

材料科学的核心任务之一，是理解和建立材料成分、结构、工艺与性能之间的关联。研究者希望理解：为什么某种成分的合金具有更高强度？为什么某种陶瓷能够承受更高温度？为什么某种半导体材料具有合适的带隙并能用于发光器件？这些问题表面上属于不同材料体系，但本质上都指向同一个核心：材料的内部组成与结构如何决定其外部性能。


在传统材料研发中，研究者通常依靠已有理论、文献经验和实验试错来提出材料方案，再通过制备、表征和性能测试不断调整成分与工艺。例如，为了提高合金强度，可尝试加入不同合金元素；为了改善陶瓷耐高温性能，可调节烧结温度、晶粒尺寸或相组成；为了优化光电材料性能，可改变材料的元素组成和晶体结构。这种方式直接可靠，也符合材料实验研究的基本规律。它的优势在于结论来自真实样品和真实测试条件，缺点是往往需要大量重复实验，研发周期较长，成本也较高。


当材料体系从简单单组元或二元材料扩展到多组元合金、复合材料、异质结构、低维材料和复杂光电材料时，实验试错面对的组合空间会迅速扩大。以高熵合金为例，四种、五种甚至更多主要元素的组合、比例和热处理工艺都可能改变组织结构和性能；同时，强度、硬度、韧性、耐腐蚀性、导电性或发光性能又受到晶体结构、缺陷、晶粒尺寸、相组成、界面状态和服役环境共同影响。传统经验仍然重要，但单纯依靠人工逐一尝试，已经难以在有限时间内完成系统筛选。


在这一背景下，材料研发逐渐从单纯依赖实验试错，转向实验、计算、数据和人工智能相结合的新方式。实验能够提供真实可靠的性能结果，计算能够在实验之前预测部分结构和能量信息，数据库能够积累已有材料的成分、结构和性能数据，而人工智能则能够从这些数据中学习复杂的材料规律。换言之，人工智能并不是脱离实验和理论单独发挥作用，而是建立在已有材料知识、实验数据和计算数据基础上的辅助工具。


数据驱动并不是否定实验和理论，而是把已有实验数据、计算数据、文献数据和数据库资源转化为可分析、可学习和可预测的信息。人工智能模型可从大量材料数据中学习成分、结构与性能之间的关系，并用于预测未知材料的性质。这样，研究者不再需要在庞大的候选空间中完全盲目试错，而可借助数据和模型缩小研究范围，把实验资源集中到更有希望的材料体系上。


\subsubsection*{\textbf{1. 人工智能用于材料科学的基本流程}}

在人工智能用于材料科学的过程中，研究者通常需要依次解决几个关键环节。第一个环节是材料数据整理。需要收集和规范实验数据、计算数据、文献数据和数据库资源，并记录材料的成分、结构、工艺、性能以及数据来源等信息。只有数据来源清楚、含义明确，后续模型学习才具有可靠基础。

第二个环节是材料表示。机器学习模型不宜直接理解“合金”“晶体结构”或“热处理工艺”等材料概念，因此需要把材料转化为机器可处理的形式。例如，可用描述符表示化学成分和结构信息，也可用图结构表示原子之间的连接关系。

第三个环节是模型预测。在材料被合理表示之后，可建立机器学习模型，预测材料的形成能、带隙、弹性模量、熔点、硬度等性能。对于涉及原子运动和微观演化的问题，还可进一步构建机器学习势（利用机器学习模型学习原子结构与能量、力之间关系的势函数），用于模拟扩散、相变和熔化等原子尺度行为。

最后，模型预测结果还需要回到材料科学问题本身。通过可解释性方法，可分析模型为什么做出某种预测，识别哪些成分、结构或工艺因素对性能影响更大，从而把模型结果进一步转化为可理解的材料规律。


在材料研究中还必须明确，人工智能不宜替代材料专业知识、实验验证和物理判断。人工智能的作用更像是一种增强手段：它可帮助研究者整理数据、发现规律、提出候选、预测性能和辅助决策，但最终结论仍然需要通过材料理论分析和实验结果进行验证。尤其是在材料科学中，模型给出的预测结果并不等同于真实材料性能。材料是否能够被实际制备出来，结构是否稳定，以及能否在实际服役环境中长期可靠工作，仍然需要进一步验证。


围绕材料性能预测和材料设计两类问题，人工智能用于材料科学可概括为两条基本主线。第一条是正向预测，即在材料成分、结构或实验条件已经给定的情况下，预测材料可能具有的性能和行为。例如，已知一种晶体材料的结构，可预测它的形成能、带隙、弹性模量或熔点。第二条是逆向设计，即先给出目标性能，再反过来寻找可能满足要求的材料。从概念上看，正向预测关注已知材料可能表现出什么性能，逆向设计关注为了获得目标性能应当寻找什么材料。


本章主要讨论第一条主线，即材料科学中的正向预测。本章将从材料数据出发，进一步讨论材料如何被表示成机器可理解的形式，机器学习如何预测材料性能，机器学习势如何模拟原子尺度行为，以及可解释人工智能如何帮助研究者从模型结果中提炼材料规律。


\subsection{数据来源与材料数据库}

人工智能要预测材料性能，首先需要数据。对于材料科学而言，数据并不只是一个性能数值，而是对材料组成、制备过程、测试条件和性能结果的完整记录。一个可靠的材料数据条目，至少应该说明材料的化学组成、结构信息、制备工艺、测试条件、数据来源和目标性能。缺少这些基本信息，模型就难以建立清晰的输入与输出对应关系，所得预测结果也难以作为可靠依据。


材料数据主要来自实验、计算、文献和数据库四类渠道。


第一类是实验数据。实验数据来自材料制备、显微表征和性能测试过程，例如硬度、拉伸强度、腐蚀电流密度、熔点、热导率、带隙、发光波长等。实验数据最接近真实材料和真实服役环境，但获取成本较高，也容易受到制备工艺、测试条件和仪器参数的影响。例如，同一种合金在不同热处理温度下可能得到不同硬度，同一种涂层在不同腐蚀介质中可能表现出不同耐蚀性能。因此，实验数据不应只记录结果，还要尽量记录测试条件。


第二类是计算数据。计算数据主要来自第一性原理计算、分子动力学模拟、相场模拟等方法。例如，第一性原理计算可得到形成能、能带结构、弹性常数、态密度、晶体结构稳定性等信息；分子动力学可模拟扩散、相变和高温行为。计算数据的优点是容易标准化，也适合大规模生成；但计算数据依赖具体计算方法和参数，不宜简单等同于实验结果。


第三类是文献数据。大量材料知识存在于论文、专利、实验报告和技术资料中。例如，一篇论文可能同时记录材料成分、制备工艺、热处理条件、显微组织和性能结果。文献数据的优点是信息丰富，缺点是格式不统一，很多信息隐藏在正文、表格、图片或补充材料中，需要人工整理或借助文本处理工具提取。


第四类是材料数据库。材料数据库可理解为经过整理的材料信息资源库，它把大量材料的化学组成、晶体结构、计算性质、实验性质和文献信息集中到一个可检索的平台中。对于材料人工智能而言，数据库是重要数据来源，因为它可为模型训练、性能预测和材料筛选提供相对规范的数据基础。


常见材料数据库包括 Materials Project、OQMD\footnote{Open Quantum Materials Database，开放量子材料数据库。}、AFLOW\footnote{Automatic Flow Framework for Materials Discovery，自动化材料发现计算框架。}、NOMAD\footnote{Novel Materials Discovery，材料数据管理与共享平台。}、ICSD\footnote{Inorganic Crystal Structure Database，无机晶体结构数据库。}、JARVIS\footnote{Joint Automated Repository for Various Integrated Simulations，材料计算数据库与工具平台。} 等。不同数据库的侧重点并不完全相同：有的偏向无机晶体结构和稳定性数据，有的偏向大规模热力学计算结果，有的强调计算数据管理和共享，还有的主要提供实验晶体结构信息。教材中列举数据库名称，目的不是要求逐一掌握平台功能，而是说明不同数据来源对应不同研究任务。

对于入门阶段的研究者而言，不必一开始深入掌握每个数据库的所有功能，但需要知道不同数据库主要提供哪类材料数据，以及这些数据适合用于哪些研究任务。只有明确数据来源和适用范围，才能更合理地将数据库用于材料性能预测、候选筛选和模型训练。


使用材料数据库时，首先要明确自己的研究任务。不同任务需要的数据字段不同。例如，如果任务是预测无机光电材料的带隙，就需要关注化学式、晶体结构、空间群、带隙、形成能、数据来源和计算方法等字段；如果任务是预测合金硬度，就需要关注成分、相结构、制备工艺、热处理条件、硬度值和测试方法等字段；如果任务是筛选超高温陶瓷，就需要关注化学组成、晶体结构、熔点、密度、热稳定性和抗氧化相关数据。


因此，数据库只是提供了材料信息的来源，真正进入模型训练之前，还需要把这些信息整理成结构清楚、字段明确、来源可追溯的原始数据清单。下面以无机光电材料带隙数据整理为例，说明从数据库原始记录到可建模数据清单的具体做法。

\subsubsection*{\textbf{1. 案例示范：无机光电材料带隙数据整理}}

以“无机光电材料带隙预测”为例，数据准备的第一步不是训练模型，而是从数据库或文献中整理出一张原始数据清单。这个清单应该回答几个基本问题：材料是什么？结构是什么？带隙是多少？这个带隙来自哪里？是实验值还是计算值？如果是计算值，采用了什么计算方法？如果是实验值，是否记录了测试条件？


一张初步整理后的原始数据清单可按表 \ref{tab:raw-bandgap-data} 设计。表中的 eV\footnote{Electron volt，电子伏特，常用于表示电子能量和带隙大小。} 是带隙常用单位。


\begin{formalTable}[htbp]
\caption{原始带隙数据清单示例}
\label{tab:raw-bandgap-data}
\scriptsize
\setlength{\tabcolsep}{2.8pt}
\begin{tabularx}{\textwidth}{@{}YYYYYYYYY@{}}
\toprule
\textbf{材料编号} & \textbf{化学式} & \textbf{晶体结构} & \textbf{空间群} & \textbf{带隙/eV} & \textbf{数据类型} & \textbf{数据来源} & \textbf{方法或条件} & \textbf{备注} \\
\midrule
001 & GaN & 纤锌矿 & P6\textsubscript{3}mc & 3.40 & 实验 & 文献 & 室温测试 & 直接带隙 \\
002 & ZnO & 纤锌矿 & P6\textsubscript{3}mc & 3.37 & 实验 & 文献 & 室温测试 & 直接带隙 \\
003 & SrTiO\textsubscript{3} & 钙钛矿 & Pm-3m & 3.20 & 数据库 & 数据库记录 & 需注明方法 & 需核对来源 \\
\bottomrule
\end{tabularx}
\end{formalTable}

表 \ref{tab:raw-bandgap-data} 的作用是保存原始材料信息，而不是直接代表最终模型输入。它的重点在于把分散的信息整理成清楚、可追溯、可检查的记录。后续是否把化学式转化为成分特征，是否把晶体结构转化为结构描述符，属于下一节要讨论的材料表示问题。


在整理这类原始数据清单时，可借助 AI 辅助完成部分工作。此处的重点不是让 AI 代替研究者判断数据是否可靠，而是让 AI 帮助理解字段含义、设计表格框架和检查数据规范性。以下以无机光电材料带隙数据整理为例给出提示词示范。

\begin{samepage}
\begin{userbox}
我正在整理无机光电材料的带隙数据，准备用于后续机器学习建模。数据可能来自 Materials Project、JARVIS、文献或实验记录。请帮我设计一张原始数据清单的表头，并说明每一列的作用。

要求如下：
\begin{enumerate}
\item 表格要保留材料编号、化学式、晶体结构、空间群、带隙数值、带隙单位、数据类型、数据来源、计算方法或实验条件等信息；
\item 区分实验值和计算值；
\item 提醒哪些字段是后续建模前必须核验的；
\item 不要虚构具体材料数据，只需要给出表头设计和填写说明。
\end{enumerate}
\end{userbox}

\begin{aibox}
可将原始数据清单设计为以下字段：

\begin{itemize}
\item 材料编号：用于唯一标识每一条材料记录，便于后续追踪和检查。
\item 化学式：记录材料的元素组成，例如 GaN、ZnO、SrTiO$_3$ 等。
\item \textbf{晶体结构：}记录材料所属结构类型，例如纤锌矿、钙钛矿等。
\item \textbf{空间群：}记录晶体结构的对称性信息。
\item 带隙数值：记录材料的带隙大小。
\item 带隙单位：通常为 eV，避免不同来源数据混用时产生单位错误。
\item \textbf{数据类型：}说明该带隙是实验值还是计算值。
\item \textbf{数据来源：}记录数据来自数据库、论文、实验报告还是其他资料。
\item \textbf{方法或条件：}如果是计算值，应注明计算方法；如果是实验值，应注明测试温度、测试方法等条件。
\item \textbf{备注：}记录是否为直接带隙、是否需要进一步核验等补充信息。
\end{itemize}

在正式建模前，尤其需要核验带隙数值、数据类型、数据来源和方法或条件等字段。同一种材料在不同晶体结构、不同计算方法或不同测试条件下可能具有不同带隙，因此这些信息不应省略。
\end{aibox}
\end{samepage}

上面的提示词示范说明，AI 可帮助研究者快速形成数据表框架，但 AI 输出并不宜直接作为最终数据结果。研究者还需要回到数据库原始页面或论文原文，核对具体数值、计算方法、实验条件和材料结构是否准确。

因此，在材料数据整理阶段，AI 更适合充当数据整理助手，而不是数据来源本身。它可帮助设计表格、解释字段和生成检查规则，但具体数据是否准确，仍然需要根据数据库原始记录、实验记录或论文原文进行确认。尤其在材料科学中，同一种化学式可能对应不同晶体结构，不同计算方法也可能得到不同带隙值，因此数据来源、计算方法和实验条件必须保留。


为了提高数据质量，原始数据整理时应注意几个问题。第一，要统一单位。例如带隙通常用 eV，硬度可能用 HV\footnote{Vickers Hardness，维氏硬度单位。} 或 GPa\footnote{Gigapascal，吉帕，常用于表示压力、应力或弹性模量。}，能量可能用 eV/atom 或 kJ/mol。第二，要区分实验数据和计算数据。二者可同时记录，但不能在没有说明的情况下混用。第三，要保留数据来源。每一条数据宜能追溯到数据库编号、论文或实验记录。第四，要检查缺失值和重复值。同一个材料如果重复出现，需要判断它们是否对应不同结构、不同测试条件或不同计算方法。第五，要记录必要条件。对于实验数据，应尽量记录温度、压力、测试环境和样品状态；对于计算数据，应记录计算方法和参数来源。

\subsubsection*{\textbf{小结}}

本节说明材料研发从实验试错走向数据、计算和实验协同的基本背景。
材料数据应包含成分、结构、工艺、性能、来源和测试条件，而不只是单一性能数值。
数据库和 AI 可提高整理效率，但数据核验、条件记录和材料判断仍需由研究者完成。

\subsection*{习题}

\begin{enumerate}[leftmargin=*]
\item 材料数据为什么不能只记录性能值？还应补充哪些字段？
\item 比较四类材料数据来源，并说明选择依据。
\end{enumerate}

\clearpage
\section{特征构建与结构表示}


\begin{sectionkeypoints}
\item 材料特征构建的目的，是把化学式、晶体结构、空间群和性能数值等原始记录转化为机器学习模型能够读取的数值特征。
\item 成分描述符可由元素组成和元素属性计算得到，例如平均原子序数、平均电负性、电负性差和平均原子半径等。
\item 结构描述符用于补充原子排列信息，常见字段包括空间群编号、晶胞体积、密度和晶格常数等。
\item 图表示把晶体看成原子节点和原子间边组成的结构，可更直接地保留局域配位、原子间距离和周期性信息。
\item 描述符和图表示都需要材料判断。特征不是越多越好，应避免目标泄漏，并检查所选特征是否具有明确的材料意义。
\end{sectionkeypoints}

\subsection{成分与结构描述符}

上一节介绍了材料数据的来源和材料数据库的作用。通过实验、计算、文献或数据库，研究者可获得材料的化学式、晶体结构、空间群、性能数值和数据来源等信息。但是，这些信息只是原始材料记录，还不宜直接被大多数机器学习模型使用。机器学习模型真正能够处理的，通常是一组数值特征。因此，在进行性能预测之前，需要把材料信息进一步转化为模型可读取的数值表达。


这一过程称为材料的特征化，也常称为描述符构建。描述符指用数字表达材料中可能影响性能的关键信息。描述符不是随意选择的一组数字，而应该尽量反映材料的成分、结构和物理化学特征。例如，对于光电材料带隙预测，元素电负性、原子半径、晶体结构、空间群和晶胞参数等信息都可能与电子结构有关，因此可作为候选描述符。


\subsubsection*{\textbf{1. 案例示范：无机光电材料特征构建}}

本节围绕一个任务展开：无机光电材料带隙预测。假设研究者已经从材料数据库或文献中整理出表 \ref{tab:bandgap-raw-records} 所示的原始数据清单。


\begin{formalTable}[htbp]
\caption{无机光电材料原始记录示例}
\label{tab:bandgap-raw-records}
\footnotesize
\begin{tabularx}{\textwidth}{@{}YYYYY@{}}
\toprule
\textbf{材料} & \textbf{化学式} & \textbf{晶体结构} & \textbf{空间群} & \textbf{带隙/eV} \\
\midrule
氮化镓 & GaN & 纤锌矿 & P6\textsubscript{3}mc & 3.40 \\
氧化锌 & ZnO & 纤锌矿 & P6\textsubscript{3}mc & 3.37 \\
钛酸锶 & SrTiO\textsubscript{3} & 钙钛矿 & Pm-3m & 3.20 \\
\bottomrule
\end{tabularx}
\end{formalTable}

这个表格适合人阅读，但模型不宜直接理解“GaN”“纤锌矿”或“P6\textsubscript{3}mc”这些符号和文字。因此，第一步是把化学式转化为元素组成。


以 GaN 为例，它由 Ga 和 N 两种元素组成，原子比例为 1:1。为了便于模型处理，可把这个比例归一化为原子分数，即 Ga 为 0.50，N 为 0.50。ZnO 同理，Zn 为 0.50，O 为 0.50。SrTiO\textsubscript{3} 中 Sr、Ti 和 O 的比例为 1:1:3，归一化后分别为 0.20、0.20 和 0.60。


经过这一步，原始化学式可转化为表 \ref{tab:formula-composition-example} 所示形式。


\begin{formalTable}[htbp]
\caption{化学式到元素组成的转换示例}
\label{tab:formula-composition-example}
\begin{tabularx}{\textwidth}{@{}YYY@{}}
\toprule
\textbf{化学式} & \textbf{元素组成} & \textbf{原子分数} \\
\midrule
GaN & Ga, N & Ga: 0.50；N: 0.50 \\
ZnO & Zn, O & Zn: 0.50；O: 0.50 \\
SrTiO\textsubscript{3} & Sr, Ti, O & Sr: 0.20；Ti: 0.20；O: 0.60 \\
\bottomrule
\end{tabularx}
\end{formalTable}

这一步表面上简单，但当数据表中有几百或几千个化学式时，人工逐个拆解就很低效。实际建模时，通常需要用程序批量解析化学式，并检查程序是否能够正确处理下标、括号、掺杂和复杂化学式。只有化学式解析正确，后续成分描述符才有意义。


得到元素组成后，第二步是引入元素属性。机器学习模型并不知道 Ga、N、Zn、O、Sr、Ti 这些元素各自有什么物理化学特征，因此需要查找或建立元素属性表。这个属性表可包括原子序数、原子半径、电负性、价电子数、元素熔点等信息。


例如，GaN 的两个组成元素是 Ga 和 N。如果需要计算平均电负性，就需要知道 Ga 和 N 各自的电负性，再按照原子分数加权平均。一般来说，若第 i 种元素的原子分数为 xi，对应元素属性为 Pi，则该材料的平均属性可表示为：


\begin{equation*}\text{平均属性}=\sum_i x_i P_i\end{equation*}

这个公式可用于计算平均原子序数、平均电负性、平均原子半径、平均价电子数等特征。除了平均值，还可计算最大值、最小值、极差和方差。例如，电负性差可用组成元素中最大电负性与最小电负性的差值表示，用来描述元素之间化学性质差异。


对于带隙预测来说，电负性和电负性差具有直观意义。带隙与材料的电子结构和成键特征有关，而元素之间电负性差异会影响电子分布和键合方式。因此，平均电负性、电负性差、平均原子半径等都可作为初步描述符。当然，这些描述符并不能完全决定带隙，但它们为模型提供了可学习的材料信息。


经过这一步，数据表可进一步转化为表 \ref{tab:composition-descriptor-example} 所示的成分描述符表。


\begin{formalTable}[htbp]
\caption{成分描述符表构建示例}
\label{tab:composition-descriptor-example}
\footnotesize
\setlength{\tabcolsep}{2.8pt}
\begin{tabularx}{\textwidth}{@{}YYYYYY@{}}
\toprule
\textbf{化学式} & \textbf{平均原子序数} & \textbf{平均电负性} & \textbf{电负性差} & \textbf{平均原子半径} & \textbf{带隙/eV} \\
\midrule
GaN & 数值 & 数值 & 数值 & 数值 & 3.40 \\
ZnO & 数值 & 数值 & 数值 & 数值 & 3.37 \\
SrTiO\textsubscript{3} & 数值 & 数值 & 数值 & 数值 & 3.20 \\
\bottomrule
\end{tabularx}
\end{formalTable}

此处的“数值”需要根据元素属性表计算得到。实际操作时，可先准备一个元素属性 CSV\footnote{Comma-Separated Values，逗号分隔值文件。} 文件，其中包含元素符号、原子序数、电负性、原子半径等字段。然后将化学式解析结果与元素属性表关联起来，自动计算每个材料的平均值和差值。这样可减少重复计算，但计算结果仍然需要通过抽样检查来确认。


第三步，是加入结构描述符。仅有成分描述符还不够，因为材料性能不仅取决于由哪些元素组成，还取决于这些原子如何排列。GaN 和 ZnO 都是纤锌矿结构，而 SrTiO\textsubscript{3} 是钙钛矿结构。不同晶体结构对应不同的原子排列、局域环境和电子能带特征，因此结构信息对于带隙预测也很重要。


对于入门阶段的研究者，可先使用比较容易获得的结构描述符，例如空间群编号、晶胞体积、密度、晶格常数等。以空间群为例，P6\textsubscript{3}mc 可对应空间群编号 186，Pm-3m 可对应空间群编号 221。这样，文字形式的空间群就可转化为数值字段。晶胞体积和密度则可从晶体结构文件或数据库记录中获得。


加入结构描述符后，表格可进一步扩展为表 \ref{tab:structure-descriptor-example}。


\begin{formalTable}[htbp]
\caption{成分与结构描述符表构建示例}
\label{tab:structure-descriptor-example}
\scriptsize
\setlength{\tabcolsep}{2.8pt}
\begin{tabularx}{\textwidth}{@{}YYYYYYYY@{}}
\toprule
\textbf{化学式} & \textbf{平均原子序数} & \textbf{平均电负性} & \textbf{电负性差} & \textbf{平均原子半径} & \textbf{空间群编号} & \textbf{晶胞体积} & \textbf{带隙/eV} \\
\midrule
GaN & 数值 & 数值 & 数值 & 数值 & 186 & 数值 & 3.40 \\
ZnO & 数值 & 数值 & 数值 & 数值 & 186 & 数值 & 3.37 \\
SrTiO\textsubscript{3} & 数值 & 数值 & 数值 & 数值 & 221 & 数值 & 3.20 \\
\bottomrule
\end{tabularx}
\end{formalTable}

在这个表中，平均原子序数、平均电负性、电负性差、平均原子半径、空间群编号和晶胞体积等列就是模型输入特征，带隙/eV 是模型要预测的目标值。这样，原始的材料记录就被转化为了机器学习可处理的特征表。


该流程可概括为下列步骤，并可结合图 \ref{fig:materials-data-feature} 理解：


\begin{lightquote}
\centering 原始材料记录 → 化学式解析 → 元素属性计算 → 结构字段数值化 → 特征表
\end{lightquote}

\begin{figure}[H]
\centering
\BookFigureImage{images_11/1.png}
\caption{材料数据整理与特征构建示意图}
\label{fig:materials-data-feature}
\end{figure}
\FloatBarrier

在这一环节，如果借助 AI 辅助处理，应像数据整理阶段一样把任务、已有字段、期望输出和限制条件说明清楚。可按下面方式提出任务。

\begin{samepage}
\begin{userbox}
我已经整理了一张无机光电材料带隙数据表，字段包括化学式、晶体结构、空间群、带隙/eV。现在准备用于机器学习建模，请帮我设计从原始表格到特征表的处理流程。

要求如下：
\begin{enumerate}
\item 说明如何把化学式解析为元素组成和原子分数；
\item 说明如何利用元素属性表计算平均原子序数、平均电负性、电负性差和平均原子半径；
\item 指出晶体结构、空间群等文字字段哪些需要进一步数值化；
\item 提醒哪些字段不宜直接作为模型输入，以及如何避免目标泄漏；
\item 不要虚构具体数值，只给出处理步骤和检查要点。
\end{enumerate}
\end{userbox}

\begin{aibox}
可把处理流程分为五步：

\begin{enumerate}
\item 化学式解析：将 GaN、ZnO、SrTiO\textsubscript{3} 等化学式拆分为元素种类和原子比例，再归一化为原子分数。
\item \textbf{元素属性匹配：}准备元素属性表，至少包含元素符号、原子序数、电负性、原子半径等字段，并与解析后的元素组成对应。
\item \textbf{成分描述符计算：}按原子分数加权计算平均原子序数、平均电负性、平均原子半径，并计算电负性差等差异类特征。
\item \textbf{结构字段转换：}晶体结构名称和空间群符号不宜直接输入多数模型，应转换为空间群编号、晶胞体积、密度或其他可计算结构描述符。
\item \textbf{建模前检查：}带隙/eV 是预测目标，不能作为输入特征；发光颜色等由带隙直接决定的字段也要谨慎使用，避免目标泄漏。
\end{enumerate}

最终得到的特征表应做到字段清楚、单位一致、输入特征和预测目标分离，并保留原始数据来源，便于后续核验。
\end{aibox}
\end{samepage}

在实践中，描述符不是越多越好。如果数据样本较少，而输入特征过多，模型可能会过拟合。对于入门任务，可先选择少量具有明确材料意义、且容易从数据库获得的描述符，例如平均原子序数、平均电负性、电负性差、平均原子半径、空间群编号和晶胞体积。等研究者理解基本流程后，再逐步加入更复杂的局域环境描述符。


还要避免目标泄漏。目标泄漏指的是把与预测目标直接相关、甚至由目标结果推导出来的信息放入输入特征中。例如，如果目标是预测带隙，就不应把“发光颜色”作为输入特征，因为发光颜色本身往往已经反映了带隙大小。这样的模型即使看起来预测准确，也没有真正学习材料成分和结构对带隙的影响。


\subsection{图表示与物理对称性}

上一节介绍了成分描述符和结构描述符。通过这些描述符，研究者可将化学式、元素属性、空间群和晶胞体积等信息转化为数值特征，用于机器学习模型预测材料性能。这种方法直观、容易理解，也适合入门学习。但它仍然有一个明显问题：很多描述符是人工设计的，模型看到的是人工预先整理好的特征，而不是材料结构本身。


成分描述符依赖人工设计，适合快速形成特征表，但难以完整保留原子排列、局域配位和相互作用距离等结构信息。带隙、形成能、弹性性质、离子扩散能力等，都与原子之间的连接关系、距离和局域环境密切相关。如果模型只使用平均电负性、平均原子半径、空间群编号等手工描述符，可能会丢失一部分精细结构信息。


图表示为这一问题提供了另一种思路：把材料看作原子节点与原子间边组成的图结构。在这张图中，原子是节点，原子之间的邻近关系或相互作用是边。这样，模型就不再只是读取一组人工设计的表格特征，而是可直接从材料的原子结构中学习信息。


继续以无机光电材料带隙预测为例，假设研究者要预测 GaN 的带隙。上一节中，研究者可使用平均电负性、电负性差、平均原子半径、空间群编号等描述符表示 GaN。但如果采用图表示，问题就变为：Ga 原子和 N 原子在晶体中如何排列，它们之间的距离是多少，每个原子周围有多少邻居。


在晶体图中，节点通常代表原子。对于 GaN 来说，图中的节点可以是 Ga 原子和 N 原子。每个节点需要带有节点特征，用来描述这个原子的基本信息。例如，一个 Ga 节点可包含 Ga 的原子序数、元素类型、电负性、原子半径、价电子数等信息；一个 N 节点也可包含 N 的对应元素属性。这样，模型就知道每个节点代表什么元素。


边通常代表原子之间的邻近关系。晶体中原子不是孤立存在的，材料性能往往取决于原子周围的局域环境。因此，需要根据一定规则判断哪些原子之间建立边。最常见的方法是设置一个截断半径：如果两个原子之间的距离小于某个阈值，就认为它们在图中相连。例如，可规定距离小于 4 \AA{} 的原子之间建立边。这样，Ga 原子周围一定距离内的 N 原子会与它相连，形成局域原子环境。


每条边也可带有边特征。边特征可包括原子间距离、相对位置、键类型或距离展开后的数值表示。对于带隙预测来说，原子间距离和局域配位环境会影响轨道重叠、成键方式和电子能带结构。因此，图表示比简单成分描述符更接近材料真实结构。


\subsubsection*{\textbf{1. 案例示范：GaN 晶体图构建}}

以 GaN 为例，构建晶体图可按照以下步骤进行。


第一步，获得 GaN 的晶体结构文件。这个文件可来自材料数据库，常见格式包括 CIF\footnote{Crystallographic Information File，晶体学信息文件。} 文件或其他晶体结构格式。结构文件中记录了晶格参数、原子种类和原子坐标。对于图表示来说，晶体结构文件是关键输入，因为它提供了原子在空间中的排列信息。


第二步，读取原子信息。程序需要从结构文件中读取每个原子的元素类型和坐标。例如，一个晶胞中可能包含 Ga 原子和 N 原子，每个原子都有对应的空间位置。此时，每个原子就可被看作图中的一个节点。


第三步，为节点添加特征。对于每个原子节点，可用元素属性作为节点特征。例如，Ga 节点可用原子序数、电负性、原子半径、价电子数等表示；N 节点也用同样类型的属性表示。这样，不同元素虽然都是“节点”，但节点特征不同，模型可区分它们的化学性质。


第四步，根据原子间距离建立边。设定一个截断半径，例如 4 \AA{}。程序计算所有原子之间的距离，如果两个原子距离小于截断半径，就在它们之间建立一条边。由于晶体具有周期性，计算距离时还要考虑相邻晶胞中的原子，而不是只看一个孤立晶胞。否则，边界附近原子的邻居关系会被错误截断。


第五步，为边添加特征。最简单的边特征是两个原子之间的距离。例如，Ga--N 距离可作为边特征输入模型。更复杂的方法会把距离展开成一组函数值，使模型更容易学习不同距离对材料性质的影响。


第六步，将整张图与目标性能对应起来。对于带隙预测任务，GaN 的晶体图作为模型输入，GaN 的带隙数值作为模型输出。模型学习的目标，就是从晶体图中预测带隙。换言之，模型不再只学习“平均电负性--带隙”的关系，而是学习“原子种类、局域连接、原子距离和晶体结构--带隙”的关系。


用表格形式表示，图数据可概括为表 \ref{tab:gan-crystal-graph-data-components} 所示内容。


\begin{formalTable}[htbp]
\caption{GaN 晶体图数据组成示例}
\label{tab:gan-crystal-graph-data-components}
\begin{tabularx}{\textwidth}{@{}YYY@{}}
\toprule
\textbf{图的组成} & \textbf{在 GaN 中对应什么} & \textbf{作用} \\
\midrule
节点 & Ga 原子、N 原子 & 表示晶体中的原子 \\
节点特征 & 原子序数、电负性、原子半径等 & 描述每个原子的元素属性 \\
边 & 截断半径内相邻原子之间的连接 & 表示局域原子关系 \\
边特征 & 原子间距离、相对位置等 & 描述原子之间的空间关系 \\
图标签 & GaN 的带隙值 & 模型需要预测的目标性能 \\
\bottomrule
\end{tabularx}
\end{formalTable}

晶体图构建涉及结构文件读取、周期性边界条件和邻接关系检查，适合借助 AI 辅助梳理代码流程。但提示词不应只写“构建晶体图”，而要明确输入文件、截断半径、输出内容和检查要求。

\begin{samepage}
\begin{userbox}
我有一个 GaN 的 CIF 晶体结构文件，准备用它构建用于带隙预测的晶体图。请帮我设计一段 Python 处理流程，用 pymatgen 读取结构，并根据 4 \AA{} 截断半径建立原子邻接关系。

要求如下：
\begin{enumerate}
\item 输出每个原子的元素种类和坐标；
\item 将原子作为节点，并说明可使用哪些节点特征；
\item 根据 4 \AA{} 截断半径生成边列表，并说明是否考虑周期性边界条件；
\item 输出节点数量、边数量和部分邻居关系，便于检查图是否合理；
\item 不要假设 CIF 文件一定正确，提醒需要核对哪些结构信息。
\end{enumerate}
\end{userbox}

\begin{aibox}
可按以下流程处理：

\begin{enumerate}
\item 用 pymatgen 读取 GaN.cif，先检查晶格参数、原子种类、原子坐标和空间群信息是否完整。
\item 将每个原子作为图中的一个节点。节点特征可从元素属性中选取，例如原子序数、电负性、原子半径和价电子数。
\item 使用结构对象的邻居搜索功能，在 4 \AA{} 截断半径内寻找相邻原子。对于晶体材料，应使用考虑周期性边界条件的方法，而不是只计算单个晶胞内部距离。
\item 将邻居关系保存为边列表，边特征可先使用原子间距离。后续如果用于图神经网络，还可把距离展开为多个函数值。
\item 检查节点数、边数和每类原子的邻居数量。如果某些原子没有邻居，或者边数量明显异常，应重新检查截断半径、CIF 文件和周期性设置。
\end{enumerate}

这段流程只能帮助完成结构到图的转换。截断半径是否合理、节点特征是否足够、图结构是否符合 GaN 的晶体环境，仍然需要结合晶体结构知识判断。
\end{aibox}
\end{samepage}

图表示之所以适合材料结构，还因为它能够更自然地满足材料中的一些物理要求。材料性质不应该因为坐标系移动、整体旋转或原子编号改变而发生变化。例如，同一个 GaN 晶体，如果研究者把整个晶体平移一段距离，它的带隙不应该改变；如果研究者重新编号原子，材料性质也不应该改变；如果只是改变观察角度，材料本身也没有发生变化。


这些要求可概括为物理对称性。对于材料机器学习，常见的物理对称性包括平移不变性、旋转不变性和置换不变性。


平移不变性是指材料整体移动后，性质不应改变。换言之，原子坐标整体加上同一个位移，材料的带隙、形成能等性质不应变化。因此，模型不应只依赖原子的绝对坐标，而应该更多使用原子之间的相对位置、距离或局域环境。


旋转不变性是指材料整体旋转后，性质不应改变。对于带隙、形成能、总能量等标量性质，材料旋转一个角度后，预测值应该保持不变。因此，很多模型会使用原子间距离这类旋转不变的量作为边特征。


置换不变性是指相同类型原子的编号顺序改变后，性质不应改变。在晶体结构文件中，原子编号只是人为排列顺序，不代表真实物理差异。例如，同一个晶体中两个等价 Ga 原子交换编号后，化学组成、局域环境和带隙都没有改变，模型输出也不应改变。图神经网络通过对邻居信息进行聚合，可在一定程度上满足这种要求。


除了不变性，还有一个更进一步的概念叫等变性。从概念上看，不变性要求某些输出在变换后保持不变，而等变性要求输出随着输入变换而按照物理规律相应变化。例如，能量是标量，材料整体旋转后能量不变；力是矢量，如果结构旋转，力的方向也应该随之旋转。可以把它理解为：同一个原子环境换了观察方向，能量数值不变，但每个原子受力的方向要跟着坐标系一起转动。对于带隙预测这类标量任务，理解不变性已经足够；对于后续机器学习势、力预测和分子动力学模拟，等变性会更加重要。


因此，图表示不仅是一种数据格式，更是一种更符合材料物理特征的表示方式。它把材料看作原子节点和相互作用边组成的结构，同时尽量尊重材料中的平移、旋转和置换对称性。相比只使用手工描述符，图表示能够保留更多局域结构信息，也为后续图神经网络模型奠定基础。


图表示并不意味着完全不需要人工判断。构建晶体图时，截断半径、节点特征、边特征表示和周期性边界条件处理都会影响模型结果。即使借助工具生成代码并检查图结构，上述关键设置是否合理，仍然需要结合材料结构和预测任务来判断。

\subsubsection*{\textbf{小结}}

本节说明材料信息需要转化为描述符、结构字段或晶体图，才能进入机器学习模型。
成分和结构描述符适合表格建模，图表示则更适合保留原子连接、距离和局域环境。
特征构建必须避免目标泄漏，并结合物理对称性、截断半径和边特征检查输入是否合理。

\subsection*{习题}

\begin{enumerate}[leftmargin=*]
\item 比较成分描述符和结构描述符的作用及适用场景。
\item 图表示如何保留比成分描述符更多的结构信息？
\item 成分描述符和图表示各有什么优势与局限？请结合带隙预测举例说明。
\end{enumerate}

\clearpage
\section{性能预测与尺度模拟}


\begin{sectionkeypoints}
\item 材料性能预测的核心，是让模型学习材料特征与目标性能之间的对应关系。硬度、带隙、形成能、弹性模量和熔点都可成为预测目标。
\item 描述符模型适合表格数据和入门任务，关键步骤包括整理特征表、划分训练测试集、训练模型、评价误差和检查特征重要性。
\item 图神经网络适合处理晶体结构数据，能够从原子节点、原子间边和局域环境中学习材料性质。
\item 机器学习势学习的是原子结构与能量、力之间的关系，可连接第一性原理计算和分子动力学模拟。
\item 模型预测和尺度模拟都必须进行可靠性检查。训练集外推、过拟合、数据泄漏、结构文件错误和模拟参数不合理都会影响结论。
\end{sectionkeypoints}

\subsection{描述符模型与硬度预测}
\label{subsec:descriptor-hardness-prediction}

有了材料描述符特征表之后，下一步是建立特征与性能之间的映射关系。也就是说，模型需要根据成分、结构、工艺或显微组织等输入特征，预测硬度、带隙、形成能、熔点等目标性能。


材料性能预测的基本思想，是让模型学习“材料特征”和“目标性能”之间的对应关系。材料特征可来自成分、结构、工艺和显微组织，目标性能可以是硬度、强度、带隙、形成能、熔点或耐腐蚀性能等。对于传统机器学习方法而言，材料首先要被转化成一张特征表，每一行代表一个材料样品，每一列代表一个可计算的特征，最后一列通常是需要预测的性能。


\subsubsection*{\textbf{1. 案例示范：高熵合金涂层硬度预测}}

本节只围绕一个具体任务展开：高熵合金涂层硬度预测。假设研究者希望根据高熵合金涂层的成分和组织信息，预测涂层硬度。这个例子适合说明传统机器学习预测流程，因为硬度是材料实验中常见的数值型性能，既受成分影响，也受相结构、晶粒尺寸和硬质相等组织因素影响。


对于一个高熵合金涂层样品，原始实验记录可能包括名义成分、制备工艺、XRD\footnote{X-ray Diffraction，X 射线衍射。} 相分析、显微组织观察和显微硬度测试结果。经过前面的描述符构建过程后，这些信息可整理为表 \ref{tab:hea-hardness-features} 所示形式。


\begin{formalTable}[htbp]
\caption{高熵合金涂层硬度预测特征表示示例}
\label{tab:hea-hardness-features}
\scriptsize
\setlength{\tabcolsep}{2.8pt}
\begin{tabularx}{\textwidth}{@{}YYYYYYYY@{}}
\toprule
\textbf{样品编号} & \textbf{原子半径差} & \textbf{平均电负性} & \textbf{电负性差} & \textbf{价电子浓度} & \textbf{是否含硬质相} & \textbf{平均晶粒尺寸/\(\mu\)m} & \textbf{硬度/HV} \\
\midrule
S1 & 数值 & 数值 & 数值 & 数值 & 0 & 数值 & 数值 \\
S2 & 数值 & 数值 & 数值 & 数值 & 1 & 数值 & 数值 \\
S3 & 数值 & 数值 & 数值 & 数值 & 1 & 数值 & 数值 \\
S4 & 数值 & 数值 & 数值 & 数值 & 1 & 数值 & 数值 \\
\bottomrule
\end{tabularx}
\end{formalTable}

表 \ref{tab:hea-hardness-features} 中，“原子半径差、平均电负性、电负性差、价电子浓度、是否含硬质相、平均晶粒尺寸”等列是模型输入特征，“硬度/HV”是模型需要预测的目标值。换成机器学习语言，就是把输入特征组成矩阵 X，把硬度组成目标向量 y。模型训练的目的，就是从 X 中学习到 y 的变化规律。


在这个任务中，硬度是一个连续数值，因此它属于回归任务。回归任务的目标不是判断材料属于哪一类，而是预测一个具体数值。例如，模型输入某个涂层样品的成分和组织特征后，输出预测硬度是多少 HV。与之相对，如果任务是判断材料是否耐腐蚀、是否稳定、是否属于某种相结构，那就更接近分类任务。


为了让这个过程更加具体，以下给出一个精简的机器学习代码框架。假设研究者已经整理好了一个高熵合金涂层硬度数据表，文件名为 hea\_\allowbreak{}coating\_\allowbreak{}hardness.csv。表格中包含若干材料描述符和目标性能“硬度/HV”。代码的目标是读取数据、划分训练集和测试集、训练随机森林回归模型，并评价预测误差。代码采用平均绝对误差（Mean Absolute Error，MAE）、均方根误差（Root Mean Squared Error，RMSE）和决定系数（Coefficient of Determination，\(R^2\)）评价预测误差与拟合程度。


\begin{lstlisting}[language=Python,escapeinside={(*@}{@*)}]
# (*@\textnormal{1. 导入常用库}@*)
import pandas as pd
from sklearn.model_selection import train_test_split
from sklearn.ensemble import RandomForestRegressor
from sklearn.metrics import mean_absolute_error, mean_squared_error, r2_score
# (*@\textnormal{2. 读取高熵合金涂层硬度数据表}@*)
data = pd.read_csv("hea_coating_hardness.csv")
# (*@\textnormal{3. 指定输入特征和预测目标}@*)
# (*@\textnormal{这些特征来自前面已经构建好的材料描述符}@*)
feature_cols = [
    "atomic_radius_diff",             # (*@\textnormal{原子半径差}@*)
    "mean_electronegativity",         # (*@\textnormal{平均电负性}@*)
    "electronegativity_diff",         # (*@\textnormal{电负性差}@*)
    "valence_electron_concentration", # (*@\textnormal{价电子浓度}@*)
    "has_hard_phase",                 # (*@\textnormal{是否含硬质相：0 表示否，1 表示是}@*)
    "grain_size"                      # (*@\textnormal{平均晶粒尺寸}@*)
]
X = data[feature_cols]
y = data["hardness_HV"]
# (*@\textnormal{4. 划分训练集和测试集}@*)
X_train, X_test, y_train, y_test = train_test_split(
    X,
    y,
    test_size=0.2,
    random_state=42
)
# (*@\textnormal{5. 建立并训练随机森林回归模型}@*)
model = RandomForestRegressor(
    n_estimators=200,
    random_state=42
)
model.fit(X_train, y_train)
# (*@\textnormal{6. 在测试集上进行硬度预测}@*)
y_pred = model.predict(X_test)
# (*@\textnormal{7. 评价模型预测效果}@*)
mae = mean_absolute_error(y_test, y_pred)
rmse = mean_squared_error(y_test, y_pred) ** 0.5
r2 = r2_score(y_test, y_pred)
print("MAE:", mae)
print("RMSE:", rmse)
print("R2:", r2)
# (*@\textnormal{8. 输出特征重要性}@*)
importance = pd.DataFrame({
    "feature": feature_cols,
    "importance": model.feature_importances_
}).sort_values(by="importance", ascending=False)
print(importance)
\end{lstlisting}

这段代码对应了材料性能预测的基本流程。首先，pd.read\_\allowbreak{}csv() 用于读取已经整理好的材料数据表。此处的数据表不是原始实验记录，而是已经经过描述符构建后的特征表。换言之，代码中的输入不是“FeCoCrNiTi”这样的化学式，而是“原子半径差、平均电负性、电负性差、价电子浓度、是否含硬质相、平均晶粒尺寸”等数值特征。


其中，X = data[feature\_\allowbreak{}cols] 表示模型输入，y = data["hardness\_\allowbreak{}HV"] 表示模型输出。对于本例来说，X 是涂层样品的成分和组织特征，y 是实验测得的硬度值。模型训练时要学习的，就是这些特征与硬度之间的关系。


train\_\allowbreak{}test\_\allowbreak{}split() 用于划分训练集和测试集。训练集用于让模型学习规律，测试集用于检验模型对新样品的预测能力。此处设置 test\_\allowbreak{}size=0.2，表示 80\% 的样品用于训练，20\% 的样品用于测试。对于材料数据来说，样本数量往往不大，因此测试集必须严格独立，不能参与模型训练。否则，模型评估结果会偏高，不能代表真实预测能力。


本例使用的模型是随机森林回归模型。随机森林由多棵决策树组成，能够处理非线性关系，对小规模表格数据比较友好，也可输出特征重要性。对于高熵合金涂层硬度预测来说，成分描述符、硬质相和晶粒尺寸与硬度之间通常不是简单线性关系，因此随机森林适合作为入门模型。此处选择随机森林并不意味着它一定最优，而是因为它便于研究者理解完整建模流程。


模型训练完成后，model.predict(X\_\allowbreak{}test) 会对测试集样品进行硬度预测。此时模型只看到测试样品的输入特征，不知道真实硬度。随后，研究者用真实硬度和预测硬度进行比较，计算 MAE、RMSE 和 R\textsuperscript{2}。


MAE 是平均绝对误差，表示模型预测值与真实值平均相差多少。例如，如果 MAE 为 30 HV，可理解为模型对测试样品的硬度平均预测误差约为 30 HV。RMSE 是均方根误差，对较大的预测误差更加敏感。如果少数样品预测偏差很大，RMSE 会明显增大。R\textsuperscript{2} 是决定系数，用来衡量预测值与真实值变化趋势的一致程度，越接近 1，说明模型预测效果越好。


对于材料研究来说，仅有误差指标还不够。宜进一步绘制“真实硬度--预测硬度”散点图。如果模型预测较好，散点应大致分布在 y = x 参考线附近。图中的坐标轴、单位、图例和参考线都要清楚，否则模型结果不便于与实验报告或论文表达衔接。


除了预测误差，随机森林模型还可输出特征重要性。代码中的 model.feature\_\allowbreak{}importances\_\allowbreak{} 可显示不同输入特征对硬度预测的贡献大小。如果结果显示“是否含硬质相”“平均晶粒尺寸”“原子半径差”等特征重要性较高，就可结合材料强化机制进行分析。例如，硬质相可能带来第二相强化，晶粒细化可能带来细晶强化，原子半径差可能与固溶强化和晶格畸变有关。


在代码实现阶段，AI 更适合承担“字段核对、代码改写、报错解释和绘图补充”等辅助工作，而不是替研究者判断哪些特征一定正确。可用下面的提示词约束任务边界。

\begin{samepage}
\begin{userbox}
我正在用随机森林预测高熵合金涂层硬度，CSV 表头包括原子半径差、平均电负性、电负性差、价电子浓度、是否含硬质相、晶粒尺寸和硬度/HV。请帮我检查这段建模流程是否清楚，并补充测试集真实硬度与预测硬度散点图的绘图步骤。

要求如下：
\begin{enumerate}
\item 判断哪些列可作为输入特征，哪一列是预测目标；
\item 提醒哪些字段需要确认数据来源或实验条件；
\item 说明如果出现列名不存在、缺失值或字符串字段无法输入模型，应如何检查；
\item 补充绘制真实值--预测值散点图和 y=x 参考线的思路；
\item 不要替我判断模型一定可靠，只给出代码检查和结果呈现建议。
\end{enumerate}
\end{userbox}

\begin{aibox}
表 \ref{tab:hea-hardness-features} 中，原子半径差、平均电负性、电负性差、价电子浓度、是否含硬质相和晶粒尺寸可作为候选输入特征，硬度/HV 是预测目标，不应同时放入输入特征。

建模前需要重点检查三类问题：
\begin{enumerate}
\item \textbf{字段是否存在：}代码中指定的输入特征名称必须与 CSV 表头完全一致，否则会出现列名不存在的错误。
\item \textbf{数据是否可输入模型：}输入特征应为数值型；如果存在空值、文字标签或单位混用，需要先清洗或编码。
\item \textbf{材料含义是否可靠：}是否含硬质相应来自 XRD、SEM\footnote{Scanning Electron Microscopy，扫描电子显微镜。} 或其他表征结果，晶粒尺寸应说明统计方法，硬度值应记录测试载荷和保载时间。
\end{enumerate}

模型评价时，除 MAE、RMSE 和 R\textsuperscript{2} 外，可绘制测试集真实硬度与预测硬度散点图，并添加 y=x 参考线。如果散点明显偏离参考线，说明模型在部分样品上的预测误差较大，需要回到数据质量、样本数量和特征设计中继续检查。
\end{aibox}
\end{samepage}

在实际建模时，还应特别注意过拟合。材料数据往往样本少、特征多，如果只有十几个样品却加入几十个描述符，模型可能在训练集上表现很好，但在测试集上预测很差。这说明模型可能只是记住了已有样品，而没有真正学到可推广的规律。为了降低过拟合风险，可减少无关特征，采用交叉验证，增加样本数量，或者选择更简单、更稳定的模型。是否属于过拟合，仍然要结合训练误差、测试误差、数据规模和材料任务共同判断。


数据泄漏同样需要警惕。预测硬度时，不应把由硬度结果直接推导出的字段作为输入特征。例如，“是否高硬度样品”“硬度等级”这类字段不能作为输入，因为它们本身已经包含了目标信息。否则模型看起来预测很准，实际上并没有从成分和组织中学习硬度规律。


\subsection{图神经网络与性质预测}

上一节以高熵合金涂层硬度预测为例，介绍了传统机器学习的基本流程。那种方法的特点是：研究者首先根据材料知识人工构造描述符，例如原子半径差、电负性差、价电子浓度、是否含硬质相和晶粒尺寸，然后再把这些描述符输入随机森林等模型中预测硬度。这种方法直观、容易操作，也适合材料数据量较少的入门任务。


但是，传统描述符方法也有局限。材料性能不仅取决于几个平均特征，还与原子在晶体中的排列方式、局域配位环境、原子间距离和结构稳定性有关。尤其对于高熵合金涂层这类复杂体系，材料中可能同时存在 FCC\footnote{Face-Centered Cubic，面心立方结构。} 固溶体、BCC\footnote{Body-Centered Cubic，体心立方结构。} 相、碳化物或硼化物等不同结构。仅用“是否含硬质相”“平均晶粒尺寸”这样的表格特征，很难完整表达原子尺度结构对性能的影响。


图神经网络提供了另一种建模思路。它不再把材料完全压缩成一行人工描述符，而是把材料看作由原子节点和原子间关系组成的图结构。在晶体图中，每个原子是一个节点，原子之间的邻近关系是边，节点和边都可带有特征。模型通过在图上进行信息传递，学习局域原子环境与材料性质之间的关系。


这里要限定问题边界：高熵合金涂层是一个多尺度、多相组织体系，真实涂层中包含晶粒、晶界、第二相、孔隙、界面和残余应力等复杂因素。图神经网络不宜简单地把一张 SEM 图或一个涂层样品直接转化成完整晶体图，然后准确预测整体硬度。更合理的入门方式是：选取涂层中的代表性晶体相或候选晶体结构，例如 FCC 固溶体、BCC 固溶体或 TiC 等硬质相，把它们转化为晶体图，学习结构与形成能、弹性模量或稳定性等性质之间的关系。然后再把这些原子尺度信息作为理解涂层硬度和强化机制的补充。


这里的重点不是用图神经网络直接预测整个涂层硬度，而是说明如何把涂层中的代表性晶体结构转化为图，并用图神经网络预测与性能相关的材料性质。


假设研究者研究的是一种 FeCoCrNiTi 高熵合金涂层。实验表明，涂层中可能存在 FCC 固溶体、BCC 相和 TiC 等硬质相。上一节中，研究者用人工描述符预测涂层硬度；现在研究者从更微观的角度思考：如果想理解某个相为什么可能更稳定、为什么某种硬质相会提高硬度，就需要进一步关注原子尺度结构。图神经网络正适合处理这类晶体结构信息。


对于一个代表性晶体结构，原始数据通常来自 CIF 文件或数据库结构文件。一个结构文件中记录了晶格参数、原子种类和原子坐标。图神经网络的第一步，就是把这个晶体结构转化为晶体图。


\subsubsection*{\textbf{1. 案例示范：TiC 晶体图性质预测}}

以 TiC 硬质相为例，它在高熵合金涂层中常被视为提高硬度的重要相之一。构建 TiC 晶体图时，可按照以下步骤进行。


第一步，读取晶体结构文件。结构文件中包含 Ti 和 C 原子的种类、坐标以及晶格参数。对于模型来说，这些信息是图结构的基础。


第二步，把原子转化为节点。TiC 中的 Ti 原子和 C 原子都可看作图中的节点。每个节点需要包含节点特征，例如原子序数、电负性、原子半径、价电子数等。这样，模型能够区分 Ti 节点和 C 节点，并理解它们具有不同化学属性。


第三步，根据原子间距离建立边。在晶体中，一个原子的性质受到周围邻居原子的影响。因此，可设置一个截断半径，例如 5 \AA{}。如果两个原子之间的距离小于这个截断半径，就在它们之间建立一条边。这样，模型就可知道哪些原子彼此相邻。


第四步，为边添加特征。最常见的边特征是原子间距离。距离越近，原子之间相互作用通常越强；距离不同，键合环境也不同。因此，边特征可帮助模型学习局域结构对性质的影响。


第五步，把整个晶体图与目标性质对应起来。如果目标是预测形成能，那么每一个晶体图对应一个形成能数值；如果目标是预测弹性模量，那么每一个晶体图对应一个弹性模量数值。模型训练的目标，就是学习“晶体图结构”与“目标性质”之间的关系。


该流程可概括为下列步骤，并可结合图 \ref{fig:crystal-graph-prediction} 理解：


\begin{lightquote}
\centering 晶体结构文件 → 原子节点 → 原子间边 → 节点特征和边特征 → 晶体图 → 性质预测
\end{lightquote}

\begin{figure}[H]
\centering
\BookFigureImage{images_11/2.png}
\caption{晶体图表示与性能预测示意图}
\label{fig:crystal-graph-prediction}
\end{figure}
\FloatBarrier

图 \ref{fig:crystal-graph-prediction} 所示流程中，图数据的组成可进一步概括为表 \ref{tab:crystal-graph-data-components}。


\begin{formalTable}[htbp]
\caption{晶体图数据组成示例}
\label{tab:crystal-graph-data-components}
\begin{tabularx}{\textwidth}{@{}YY@{}}
\toprule
\textbf{图数据组成} & \textbf{在 TiC 或高熵合金相关晶体结构中表示什么} \\
\midrule
节点 & Ti、C 或 Fe、Co、Cr、Ni、Ti 等原子 \\
节点特征 & 原子序数、电负性、原子半径、价电子数等 \\
边 & 截断半径内的原子邻近关系 \\
边特征 & 原子间距离或距离展开特征 \\
图标签 & 形成能、弹性模量、体模量或其他目标性质 \\
\bottomrule
\end{tabularx}
\end{formalTable}

图神经网络的核心是信息传递。每个原子节点在初始时只包含自己的元素属性，例如 Ti 节点知道自己是 Ti，C 节点知道自己是 C。经过一轮信息传递后，一个节点会接收周围邻居节点的信息。经过多轮信息传递后，节点就不仅包含自身信息，还包含局域原子环境信息。最后，模型把所有节点的信息汇总成整个晶体的表示，再输出目标性质。


这个过程可用材料语言理解：模型不只是看“这个材料含有哪些元素”，而是进一步看“这些元素在空间中如何排列、每个原子周围是什么邻居、原子间距离是多少、局域环境是否稳定”。这正是图神经网络相比传统描述符方法的重要优势。


为了让这个过程更具体，以下给出一个教学用的精简代码框架。它的目的不是提供完整可直接运行的项目，而是帮助研究者理解图神经网络建模的基本步骤。


\begin{lstlisting}[language=Python,escapeinside={(*@}{@*)}]
# (*@\textnormal{教学用简化框架：从晶体结构构建图并预测性质}@*)
# (*@\textnormal{实际运行时需要安装 pymatgen、torch、torch\_\allowbreak{}geometric 等库}@*)
from pymatgen.core import Structure
import torch
# (*@\textnormal{1. 读取晶体结构文件，例如 TiC.cif}@*)
structure = Structure.from_file("TiC.cif")
# (*@\textnormal{2. 设置截断半径}@*)
cutoff = 5.0
# (*@\textnormal{3. 构建节点特征：此处仅用原子序数作为示例}@*)
node_features = []
for site in structure:
    atomic_number = site.specie.Z
    node_features.append([atomic_number])
x = torch.tensor(node_features, dtype=torch.float)
# (*@\textnormal{4. 构建边：寻找截断半径内的邻居原子}@*)
edge_index = []
edge_attr = []
for i, site in enumerate(structure):
    neighbors = structure.get_neighbors(site, cutoff)
    for neighbor in neighbors:
        j = neighbor.index
        distance = neighbor.nn_distance
        edge_index.append([i, j])
        edge_attr.append([distance])
edge_index = torch.tensor(edge_index, dtype=torch.long).t().contiguous()
edge_attr = torch.tensor(edge_attr, dtype=torch.float)
# (*@\textnormal{5. 图标签：例如该结构的形成能或弹性模量}@*)
# (*@\textnormal{此处用一个占位数值表示，实际应来自数据库或计算结果}@*)
y = torch.tensor([target_property], dtype=torch.float)
# (*@\textnormal{至此，一个晶体结构已经被转化为图数据：}@*)
# (*@\textnormal{x 表示节点特征}@*)
# (*@\textnormal{edge\_\allowbreak{}index 表示边连接关系}@*)
# (*@\textnormal{edge\_\allowbreak{}attr 表示边特征}@*)
# (*@\textnormal{y 表示目标性质}@*)
\end{lstlisting}

这段代码展示了晶体图构建的基本思想。Structure.from\_\allowbreak{}file("TiC.cif") 用来读取晶体结构文件；node\_\allowbreak{}features 保存每个原子的节点特征；edge\_\allowbreak{}index 保存哪些原子之间存在边；edge\_\allowbreak{}attr 保存边特征，例如原子间距离；y 是模型要预测的目标性质。


实际使用时，节点特征不应只包含原子序数，还可加入电负性、原子半径、价电子数等元素属性。边特征也不一定只用一个距离数值，可把距离展开成一组函数，以便模型更好学习不同距离范围内的相互作用。对于不同材料任务，目标 y 也可不同。如果关注结构稳定性，可预测形成能；如果关注力学性能，可预测弹性模量或体模量；如果进一步结合实验数据，也可尝试预测与硬度相关的性质。


如果借助 AI 辅助构建 TiC 晶体图，提示词应突出“代表性晶体相”和“目标性质”两个信息，避免把图神经网络误写成直接预测整个涂层硬度的工具。示例任务可写成如下形式。

\begin{samepage}
\begin{userbox}
我想以高熵合金涂层中的 TiC 硬质相为例，构建一个用于性质预测的晶体图。请帮我设计从 TiC.cif 到图数据的处理流程，目标性质可以是形成能、弹性模量或体模量。

要求如下：
\begin{enumerate}
\item 说明如何读取 TiC.cif，并把 Ti 和 C 原子转化为节点；
\item 说明节点特征可选择哪些元素属性；
\item 根据 5 \AA{} 截断半径建立边，并说明为什么要考虑周期性边界条件；
\item 说明边特征和图标签分别是什么；
\item 提醒这个晶体图只能表示 TiC 这一代表性相，不宜直接代表整个高熵合金涂层硬度。
\end{enumerate}
\end{userbox}

\begin{aibox}
可按以下流程构建图数据：

\begin{enumerate}
\item \textbf{读取结构：}使用 pymatgen 读取 TiC.cif，检查晶格参数、Ti/C 原子坐标和结构完整性。
\item \textbf{构建节点：}把每个 Ti 原子和 C 原子作为节点。节点特征可包括原子序数、电负性、原子半径和价电子数。
\item \textbf{构建边：}使用 5 \AA{} 截断半径寻找邻近原子。由于 TiC 是周期晶体，邻居搜索必须考虑相邻晶胞，否则会漏掉晶胞边界附近的邻居。
\item \textbf{添加边特征：}可先使用原子间距离作为边特征，后续再根据模型需要使用距离展开特征。
\item \textbf{设置图标签：}如果训练目标是形成能，图标签就是形成能；如果目标是弹性模量或体模量，图标签应来自数据库、第一性原理计算或可靠实验/计算资料。
\end{enumerate}

这里需要专门说明，TiC 晶体图只能帮助理解涂层中某个硬质相的结构性质。真实高熵合金涂层硬度还受到晶粒尺寸、相分布、界面、孔隙和残余应力影响，不能由单个 TiC 晶体图直接决定。
\end{aibox}
\end{samepage}

与传统描述符模型相比，图神经网络的优势在于它能更直接地保留原子结构信息。传统模型需要人工预先设计“平均电负性”“原子半径差”等特征，而图神经网络可从原子种类、原子间距离和局域连接关系中自动学习结构信息。对于晶体材料来说，这种表示方式更接近材料本身。


但是，图神经网络也不是万能的。首先，它需要可靠的晶体结构输入。如果 CIF 文件错误、原子占位不清楚或结构不完整，构建出的图也会有问题。其次，它通常需要较多训练数据。如果只有少量结构样本，图神经网络也可能过拟合。再次，对于真实高熵合金涂层来说，整体硬度还受到晶粒尺寸、相分布、孔隙率、界面结合和残余应力等显微组织因素影响，单个晶体图不能完全代表整个涂层。


因此，在高熵合金涂层研究中，图神经网络更适合作为结构层面的补充工具。它可帮助研究者理解某个候选相、硬质相或代表性晶体结构的稳定性和力学相关性质；而整体涂层硬度仍然需要结合上一节中的表格特征模型、显微组织数据和实验测试结果共同分析。


\subsection{原子间势与熔点模拟}

前两节讨论的是材料性能预测。无论是传统机器学习模型，还是图神经网络模型，它们通常都是输入一个材料样品或晶体结构，然后输出一个性能数值，例如硬度、形成能、带隙或弹性模量。这类模型回答的是“这个材料可能具有什么性质”。


但是，材料科学中还有另一类重要问题：材料在原子尺度上如何运动、如何扩散、如何相变、如何熔化。例如，超高温陶瓷在高温下能否保持结构稳定？原子在升温过程中是否发生剧烈扩散？材料在什么温度附近开始熔化？这些问题不是简单预测一个静态数值就能完全回答的，而需要模拟原子随时间的运动过程。


分子动力学可用来研究这类原子尺度动态行为。在分子动力学中，研究者根据原子之间的相互作用计算每个原子的受力，再根据力更新原子位置，从而模拟材料随时间演化的过程。问题在于，原子间相互作用如何计算。第一性原理分子动力学精度较高，但计算成本高，难以处理大体系和长时间尺度；传统经验势速度快，但对复杂材料体系的适用性有限。机器学习势正是在这种背景下发展起来的。


机器学习势的核心思想，是用机器学习模型学习第一性原理计算得到的势能面。从概念上看，第一性原理计算可给出某一组原子结构的能量、原子受力和应力；机器学习势则利用这些数据学习“原子结构--能量/力”之间的关系。训练完成后，机器学习势可像传统势函数一样快速计算能量和力，从而支持更大规模、更长时间的分子动力学模拟。


\subsubsection*{\textbf{1. 案例示范：TiC 熔化行为模拟}}

本节只围绕一个任务展开：用机器学习势辅助模拟 TiC 超高温陶瓷的熔化行为。TiC 是典型的超高温陶瓷之一，具有高熔点、高硬度和良好的高温稳定性，也常作为高熵合金涂层或复合涂层中的硬质增强相。通过这个例子，可说明机器学习势如何从第一性原理数据出发，进一步服务于高温分子动力学模拟。


首先需要明确，机器学习势与前面性能预测模型的目标不同。前面随机森林模型预测的是硬度，图神经网络可预测形成能或弹性模量；而机器学习势要预测的是原子结构对应的总能量和每个原子的受力。它不是直接告诉研究者“TiC 的熔点是多少”，而是先学会在不同原子构型下计算能量和力，然后通过分子动力学模拟升温过程，再从模拟结果中判断熔化行为。


对于 TiC 熔点模拟，可按照以下流程理解。


第一步，准备 TiC 的初始晶体结构。研究者需要获得 TiC 的晶体结构文件，例如 CIF 文件。这个结构文件包含 Ti 和 C 原子的排列方式、晶格参数和空间结构。初始结构通常可来自数据库、文献或已有计算结果。


第二步，生成不同构型的训练数据。机器学习势不应只看一个理想晶体结构，因为真实高温模拟中原子会振动、偏移，甚至出现缺陷和局部无序。因此，需要构造一批不同的 TiC 原子构型，例如平衡晶体结构、轻微扰动结构、不同体积结构、不同温度下短时间分子动力学得到的结构等。这样，训练数据才能覆盖模型后续可能遇到的原子环境。


第三步，用第一性原理计算这些构型的能量和力。对于每一个 TiC 构型，需要通过第一性原理计算得到总能量、每个原子的力以及必要时的应力。这一步计算成本较高，但只需要对训练集中的有限构型进行。可理解为：第一性原理计算负责提供高质量“标准答案”，机器学习势负责学习这些标准答案中的规律。


第四步，训练机器学习势。训练时，模型输入原子种类和原子坐标，输出预测能量和原子力。训练目标是让模型预测的能量和力尽可能接近第一性原理计算结果。常见机器学习势包括 Behler--Parrinello 神经网络势、高斯近似势、DeepMD\footnote{Deep Potential Molecular Dynamics，深度势分子动力学。}、MACE\footnote{Message Passing Atomic Cluster Expansion，消息传递原子簇展开模型。}、NequIP\footnote{Neural Equivariant Interatomic Potentials，神经等变原子间势。}、M3GNet\footnote{Materials 3-body Graph Network，材料三体图网络。}、CHGNet\footnote{Crystal Hamiltonian Graph Neural Network，晶体哈密顿量图神经网络。} 等。教材入门阶段不需要详细展开每种模型的数学细节，重点是理解它们都在学习原子结构与能量/力之间的关系。


第五步，验证机器学习势是否可靠。训练完成后，不宜直接用于高温模拟，而要先在测试构型上验证预测误差。常见做法是比较模型预测能量与第一性原理能量之间的误差，以及模型预测力与第一性原理力之间的误差。如果误差较大，说明模型还没有学好 TiC 的势能面，需要增加训练数据、调整模型参数或补充高温构型。


第六步，用训练好的机器学习势进行分子动力学模拟。此时，机器学习势可快速计算每一步原子的能量和力，进而模拟 TiC 在不同温度下的原子运动。为了观察熔化行为，可逐步升高温度，例如从较低温度开始，逐步升温到高温区间，记录能量、体积、径向分布函数、均方位移等随温度变化的情况。


第七步，根据模拟结果判断熔化行为。材料熔化时，原子排列会从有序晶体逐渐转变为无序液态结构。模拟中可通过多个指标辅助判断：能量--温度曲线是否出现突变，体积是否发生明显变化，径向分布函数的长程有序峰是否消失，均方位移是否显著增加。通过这些现象，可估计 TiC 的熔化温度范围。


这个过程可概括为：


\begin{lightquote}
\centering TiC 晶体结构 → 多种原子构型 → 第一性原理能量和力 → 训练机器学习势 → 验证误差 → 分子动力学升温模拟 → 分析熔化行为
\end{lightquote}

为了让研究者更清楚地理解流程，以下给出一个教学用的简化代码框架。它不是完整可直接运行的工程代码，因为真实机器学习势训练通常需要专门软件和大量计算资源。此处的目的，是帮助理解每一步在做什么。


\begin{lstlisting}[language=Python,escapeinside={(*@}{@*)}]
# (*@\textnormal{教学用简化框架：机器学习势辅助 TiC 熔化模拟}@*)
# (*@\textnormal{实际研究中可使用 DeepMD、MACE、NequIP、CHGNet、M3GNet 等工具}@*)
# (*@\textnormal{1. 读取 TiC 初始晶体结构}@*)
structure = read_structure("TiC.cif")
# (*@\textnormal{2. 生成训练构型}@*)
# (*@\textnormal{包括平衡结构、扰动结构、不同体积结构和高温采样结构}@*)
train_structures = generate_configurations(
    structure,
    perturb=True,
    volume_changes=True,
    temperature_sampling=True
)
# (*@\textnormal{3. 对训练构型进行第一性原理计算}@*)
# (*@\textnormal{得到每个构型的能量、原子力和应力}@*)
dft_dataset = []
for config in train_structures:
    energy, forces, stress = run_dft_calculation(config)
    dft_dataset.append({
        "structure": config,
        "energy": energy,
        "forces": forces,
        "stress": stress
    })
# (*@\textnormal{4. 训练机器学习势}@*)
ml_potential = train_machine_learning_potential(
    dataset=dft_dataset,
    target_properties=["energy", "forces", "stress"]
)
# (*@\textnormal{5. 在测试构型上验证模型误差}@*)
test_results = evaluate_potential(
    model=ml_potential,
    test_dataset="test_configurations"
)
print(test_results)
# (*@\textnormal{6. 使用训练好的机器学习势进行分子动力学升温模拟}@*)
md_results = run_molecular_dynamics(
    structure=structure,
    potential=ml_potential,
    temperature_schedule=[1000, 1500, 2000, 2500, 3000, 3500, 4000],
    simulation_time="several ps or longer"
)
# (*@\textnormal{7. 分析熔化行为}@*)
plot_energy_temperature(md_results)
plot_radial_distribution_function(md_results)
plot_mean_square_displacement(md_results)
\end{lstlisting}

这段框架中，read\_\allowbreak{}structure() 表示读取 TiC 的晶体结构；generate\_\allowbreak{}configurations() 表示生成不同原子构型；run\_\allowbreak{}dft\_\allowbreak{}calculation() 表示用第一性原理计算能量和力；train\_\allowbreak{}machine\_\allowbreak{}learning\_\allowbreak{}potential() 表示训练机器学习势；run\_\allowbreak{}molecular\_\allowbreak{}dynamics() 表示用训练好的势进行分子动力学模拟；最后通过能量曲线、径向分布函数和均方位移分析熔化行为。


与前两节代码不同，此处给出的是流程框架，而不是几十行就能完成的入门脚本。原因是机器学习势涉及第一性原理计算、模型训练和分子动力学模拟，真实流程需要较多软件和计算资源。但这个框架可帮助研究者理解机器学习势在材料研究中的位置：它不是直接替代实验，也不是简单回归一个性能值，而是连接第一性原理计算和大尺度原子模拟的桥梁。


在机器学习势和分子动力学任务中，AI 辅助更适合用于解释文件字段、检查流程是否完整和整理输出指标，而不宜直接判断熔点结论。下面的写法更适合用于流程核查。

\begin{samepage}
\begin{userbox}
我想用机器学习势辅助模拟 TiC 的熔化行为。已有 TiC 初始结构，计划生成扰动构型和高温采样构型，用第一性原理计算能量与力，再训练机器学习势并进行升温分子动力学模拟。请帮我检查这个流程是否完整，并说明每类输出指标的物理意义。

要求如下：
\begin{enumerate}
\item 说明训练数据应包含哪些 TiC 原子构型；
\item 说明第一性原理计算需要记录哪些目标量；
\item 检查机器学习势训练后应如何验证能量和力误差；
\item 说明升温分子动力学中能量、体积、径向分布函数和均方位移分别反映什么；
\item 明确提醒不应只根据单一指标判断熔点。
\end{enumerate}
\end{userbox}

\begin{aibox}
这个流程可分为四个检查重点：

\begin{enumerate}
\item \textbf{训练构型覆盖范围：}除平衡 TiC 晶体结构外，还应包含轻微扰动结构、不同体积结构、不同温度采样结构，以及接近高温无序状态的构型。否则模型可能只适用于低温平衡结构。
\item \textbf{第一性原理标注目标：}每个构型至少应记录总能量和每个原子的力；如果后续模拟涉及体积变化，还应考虑应力信息。
\item \textbf{机器学习势验证：}训练后应在独立测试构型上比较预测能量、预测力与第一性原理结果的误差，不应只看训练集误差。
\item \textbf{熔化行为分析：}能量--温度曲线突变可能反映相变，体积突变可能反映结构膨胀，径向分布函数长程有序峰消失说明晶体有序性降低，均方位移显著增大说明原子扩散增强。
\end{enumerate}

熔点判断应综合多个指标，并结合体系大小、升温速率、模拟时间和训练数据覆盖范围共同分析。单独依靠 AI 对曲线的文字解释，不能作为可靠的材料结论。
\end{aibox}
\end{samepage}

需要特别强调的是，机器学习势的可靠性高度依赖训练数据。如果训练集中只包含低温平衡结构，而模拟时却用于高温熔化过程，模型很可能外推到训练分布之外，预测结果就不可靠。因此，在 TiC 熔点模拟中，训练数据应尽量包含高温扰动构型、不同体积构型和接近熔化状态的构型。否则，模型可能在低温结构上误差很小，却在高温模拟中产生错误行为。


此外，机器学习势不能完全替代第一性原理计算和实验验证。它的优势是加速模拟，但前提是训练数据可靠、验证误差合理、模拟条件设置正确。对于超高温陶瓷这类材料，熔点模拟还可能受到模拟体系大小、升温速率、时间尺度和缺陷状态影响。因此，模拟得到的熔化温度应被理解为计算预测结果，需要与实验数据或更高精度计算进行对照。

\subsubsection*{\textbf{小结}}

本节围绕硬度预测、图神经网络性质预测和机器学习势展开。
描述符模型用于表格化的宏观性能预测，图神经网络用于原子结构关联的性质预测，机器学习势用于原子尺度动力学模拟。
性能指标、数据泄漏、过拟合、结构文件质量和训练构型覆盖都会影响模型结论的可靠性。

\subsection*{习题}

\begin{enumerate}[leftmargin=*]
\item 比较描述符模型、图神经网络和机器学习势的适用任务。
\item 机器学习势适合哪些模拟场景？训练覆盖为何重要？
\item 为高熵合金涂层硬度预测设计输入特征，并说明依据。
\end{enumerate}

\clearpage
\section{模型解释与知识嵌入}


\begin{sectionkeypoints}
\item 可解释性分析关注模型训练之后的问题，即模型主要依赖哪些特征、这些特征如何影响预测结果，以及解释是否符合材料机理。
\item 特征重要性和 SHAP\footnote{SHapley Additive exPlanations，一种基于博弈论的可解释性分析方法。} 可帮助把机器学习输出转化为材料语言，但它们反映的是当前数据和当前模型中的关系，不等同于材料真理。
\item 物理知识嵌入强调在建模前和建模过程中主动引入材料规律，包括特征选择、数据筛选、目标泄漏检查和结果合理性判断。
\item 对高熵合金涂层硬度预测而言，原子半径差、晶粒尺寸、硬质相和孔隙率等特征应与固溶强化、细晶强化、第二相强化和组织致密化联系起来理解。
\item 一个可靠的材料人工智能模型，不仅要有较小预测误差，还要能被实验表征、物理机理和材料常识共同支持。
\end{sectionkeypoints}

\subsection{模型解释与特征归因}

前面几节已经介绍了材料性能预测的基本方法。通过描述符和机器学习模型，研究者可根据材料成分、结构和组织信息预测硬度、带隙、形成能、熔点等性能。对于工程应用来说，预测准确当然重要；但对于材料科学研究来说，仅仅得到一个预测数值还不够。研究者还需要进一步知道：模型为什么给出这样的预测？它主要依据了哪些材料特征？这些特征是否符合已有材料机理？


这正是模型可解释性要解决的问题。可解释性关注的是模型在预测时主要依赖哪些输入特征，以及这些特征如何影响预测结果。对于材料科学而言，可解释性尤其重要，因为材料研究的目标不仅是预测某个样品性能，更是提炼可理解、可验证、可迁移的材料规律。


本节继续围绕前面的高熵合金涂层硬度预测任务展开。假设研究者已经训练了一个随机森林模型，用来预测不同高熵合金涂层的硬度。模型的输入特征包括原子半径差、平均电负性、电负性差、价电子浓度、是否含硬质相和平均晶粒尺寸，输出目标是硬度/HV。现在的问题不再是“模型能不能预测硬度”，而是“模型主要根据什么判断某个涂层硬度较高”。


如果只看预测误差，例如 MAE、RMSE 和 R\textsuperscript{2}，研究者只能知道模型预测效果大致如何，却不知道模型内部学到了什么。例如，一个模型可能在测试集上误差较小，但它可能主要依赖了某个偶然相关的字段，而不是真正反映材料强化机制的特征。这样的模型即使暂时预测准确，也很难让研究者放心。因此，在材料人工智能中，模型解释往往和模型评价同样重要。


对于随机森林这类模型，最常见的解释方法是特征重要性分析。特征重要性可粗略告诉研究者：模型在预测硬度时，哪些输入特征贡献更大。假设训练后的模型输出表 \ref{tab:hardness-feature-importance} 所示结果。


\begin{formalTable}[htbp]
\caption{随机森林硬度预测特征重要性示例}
\label{tab:hardness-feature-importance}
\begin{tabularx}{\textwidth}{@{}YY@{}}
\toprule
\textbf{特征} & \textbf{重要性} \\
\midrule
是否含硬质相 & 0.34 \\
平均晶粒尺寸 & 0.26 \\
原子半径差 & 0.18 \\
价电子浓度 & 0.13 \\
电负性差 & 0.06 \\
平均电负性 & 0.03 \\
\bottomrule
\end{tabularx}
\end{formalTable}

这个结果说明，在当前数据集和模型中，“是否含硬质相”和“平均晶粒尺寸”对硬度预测影响最大。接下来，研究者需要把这个模型结果转化为材料语言。含硬质相的涂层往往可能通过第二相强化提高硬度；晶粒尺寸越小，晶界数量越多，可能带来细晶强化；原子半径差较大时，晶格畸变增强，也可能产生固溶强化。因此，如果模型认为这些特征重要，并且这种排序与材料强化机制基本一致，就说明模型结果具有一定可解释性。


但是，特征重要性只能告诉研究者“哪些特征重要”，不宜直接告诉研究者“这个特征让硬度升高还是降低”。例如，平均晶粒尺寸重要，并不意味着晶粒尺寸越大硬度越高，也不意味着晶粒尺寸越小一定硬度越高。它只能说明模型在预测时经常用到这个特征。为了进一步分析特征对预测结果的正负影响，可引入 SHAP 方法。


SHAP 可理解为一种更细致的解释工具。它不仅能分析某个特征整体上是否重要，还可分析在某一个样品中，某个特征是把预测硬度往高处推，还是往低处拉。例如，对于一个含 TiC 硬质相、晶粒较细的高熵合金涂层样品，SHAP 分析可能显示：“是否含硬质相”对预测硬度有正贡献，“平均晶粒尺寸较小”也对预测硬度有正贡献。这种结果就可和第二相强化、细晶强化联系起来解释。


为了更清楚地说明可解释性分析，以下给出一个精简代码框架。假设研究者已经在第 \ref{subsec:descriptor-hardness-prediction} 节中训练好了随机森林硬度预测模型 model，并且已经得到输入特征 X\_\allowbreak{}train、X\_\allowbreak{}test 和特征名称 feature\_\allowbreak{}cols。以下代码用于输出随机森林特征重要性，并进一步使用 SHAP 分析模型预测结果。


\begin{lstlisting}[language=Python,escapeinside={(*@}{@*)}]
# (*@\textnormal{1. 随机森林特征重要性}@*)
importance = pd.DataFrame({
    "feature": feature_cols,
    "importance": model.feature_importances_
}).sort_values(by="importance", ascending=False)
print(importance)
# (*@\textnormal{2. 绘制特征重要性柱状图}@*)
plt.figure(figsize=(6, 4))
plt.barh(importance["feature"], importance["importance"])
plt.xlabel("Feature importance")
plt.ylabel("Feature")
plt.title("Feature importance for hardness prediction")
plt.gca().invert_yaxis()
plt.tight_layout()
plt.show()
# (*@\textnormal{3. 使用 SHAP 解释模型}@*)
# (*@\textnormal{需要先安装 shap: pip install shap}@*)
import shap
explainer = shap.TreeExplainer(model)
shap_values = explainer.shap_values(X_test)
# (*@\textnormal{4. 绘制 SHAP 总结图}@*)
shap.summary_plot(
    shap_values,
    X_test,
    feature_names=feature_cols
)
\end{lstlisting}

这段代码分为两个部分。第一部分使用 model.feature\_\allowbreak{}importances\_\allowbreak{} 输出随机森林自身的特征重要性，并画出柱状图。这个图可帮助研究者快速判断模型最关注哪些特征。第二部分使用 SHAP 对模型进行解释，shap.summary\_\allowbreak{}plot() 可显示每个特征对预测结果的整体影响。


可解释性分析中的 AI 辅助，应服务于“把模型输出转化为材料语言”，而不是替代机理判断。提示词中需要同时给出模型结果和要求它说明证据限制。可按下面方式组织问题。

\begin{samepage}
\begin{userbox}
我用随机森林预测高熵合金涂层硬度，模型输出的特征重要性排序为：是否含硬质相 0.34，平均晶粒尺寸 0.26，原子半径差 0.18，价电子浓度 0.13，电负性差 0.06，平均电负性 0.03。SHAP 图显示，含硬质相通常对预测硬度有正贡献，晶粒尺寸较大时常出现负贡献。

请帮我把这些模型解释结果转化为材料语言。

要求如下：
\begin{enumerate}
\item 分别联系第二相强化、细晶强化、固溶强化或相稳定性解释这些特征；
\item 说明这些解释需要哪些实验表征支持，例如 XRD、SEM、EDS\footnote{Energy Dispersive Spectroscopy，能谱分析。} 或硬度测试；
\item 提醒特征重要性和 SHAP 只能说明当前模型中的关系，不宜直接证明因果关系；
\item 如果模型把样品编号、数据来源等无材料意义字段排在前面，也请说明可能存在什么问题。
\end{enumerate}
\end{userbox}

\begin{aibox}
可这样理解这组结果：

\begin{enumerate}
\item “是否含硬质相”重要性最高，说明模型在预测硬度时较多依赖硬质相信息。若涂层中确实存在 TiC、碳化物或硼化物等硬质相，可与第二相强化联系起来解释。
\item “平均晶粒尺寸”重要性较高，且较大晶粒尺寸对应负向 SHAP 贡献，可能说明模型学习到了晶粒细化有利于提高硬度的趋势，可与细晶强化机制联系起来。
\item “原子半径差”有一定重要性，可与多主元合金中的晶格畸变和固溶强化联系起来，但需要结合相结构和成分分布判断。
\item “价电子浓度”可能与 FCC/BCC 相稳定性有关，但它对硬度的影响通常需要通过相结构、组织形貌和其他描述符共同解释。
\end{enumerate}

这些解释仍然只是模型层面的线索。若要形成可靠材料结论，需要用 XRD 确认物相，用 SEM/EDS 观察硬质相和成分分布，用晶粒尺寸统计验证细晶强化趋势，并检查硬度测试条件是否一致。如果无材料意义的字段在模型中排在前面，可能说明数据整理、特征选择或数据划分存在问题。
\end{aibox}
\end{samepage}

模型解释并不等于材料真理。特征重要性高，只说明模型在当前数据集中较多依赖该特征；SHAP 值为正，也只说明在当前模型中该特征对预测值有正向贡献。这些结果可能受到数据量、数据分布、特征相关性和模型类型影响。例如，如果所有含硬质相的样品刚好也具有较小晶粒尺寸，模型可能难以区分到底是硬质相起主要作用，还是晶粒细化起主要作用。因此，可解释性分析必须与实验表征和材料机理结合，不能孤立理解。


对于高熵合金涂层硬度预测，可按照以下思路解读模型解释结果。首先看特征重要性，判断模型主要关注哪些因素。其次看 SHAP 分析，判断这些因素对预测硬度的正负影响。然后把模型结果与材料强化机制对应起来：硬质相对应第二相强化，晶粒尺寸对应细晶强化，原子半径差对应晶格畸变和固溶强化，价电子浓度可能与相结构稳定性有关。最后，再回到实验数据中检查这些解释是否被组织形貌、物相分析和硬度测试支持。


例如，如果模型认为“是否含硬质相”最重要，而 SEM 和 XRD 也显示高硬度样品中确实存在较多 TiC 等硬质相，那么模型解释与实验观察相互支持。相反，如果模型认为某个表面上无关的字段最重要，例如样品编号或数据来源，那么就需要警惕模型可能学到了错误关联。这种情况并不是模型真正理解了材料规律，而可能是数据整理或特征设计存在问题。


除了特征重要性和 SHAP，还有一些更偏研究生层次的方法可用于材料模型解释。例如，符号回归和 SISSO\footnote{Sure Independence Screening and Sparsifying Operator，一种用于筛选描述符并构建稀疏表达式的方法。} 试图从数据中寻找简洁的数学表达式或组合描述符，使模型结果更接近材料专业人员能够理解的公式关系。对于入门阶段的研究者来说，可先掌握特征重要性和 SHAP 的基本思想；对于研究生，可进一步了解这些能够自动发现简洁规律的可解释建模方法。


本节的重点不是让研究者记住某一个解释算法，而是理解材料人工智能中的一个基本原则：模型预测之后，必须回到材料问题本身去解释。一个好的材料机器学习模型，不仅应该预测误差较小，还应该能够帮助研究者理解哪些成分、结构或组织因素可能控制性能。只有当模型结果能够与材料机制相互印证时，它才更可能成为材料研究和材料设计的有效工具。


通过高熵合金涂层硬度预测这个例子可看到，可解释性分析把机器学习结果和材料科学知识连接起来。模型指出哪些特征重要，材料科学知识帮助研究者解释这些特征为什么重要。但最终是否形成可靠规律，仍然需要实验表征、物理机理和专业判断共同验证。


下一节将进一步讨论物理知识嵌入。可解释性更多是在模型训练之后分析“模型为什么这样预测”，而物理知识嵌入则是在建模之前或建模过程中，就主动把材料科学规律加入数据选择、特征设计和模型约束中，使模型尽量避免学习错误关系。


\subsection{物理知识与模型约束}

上一节讨论了模型可解释性，即模型训练完成之后，研究者如何分析模型主要依据哪些特征做出预测，并将这些特征与材料机理联系起来。可解释性强调的是“事后解释”：模型已经训练好了，研究者再分析它为什么这样预测。


但是，在材料科学中，仅仅依靠事后解释还不够。更理想的做法是，在数据整理、特征选择、模型训练和结果判断的过程中，就主动引入材料科学知识，使模型尽量在合理的物理和化学范围内学习规律。这就是物理知识嵌入的基本思想。


物理知识嵌入并不一定是复杂的公式或高级算法。对于入门阶段，可把它理解为：不要把所有能收集到的字段都机械地丢给模型，而是根据材料机理选择有意义的数据和特征；不要让模型学习明显不合理的关系，而是用物理、化学和实验常识约束建模过程。这样得到的模型，通常比完全黑箱的模型更可靠，也更容易被研究者接受。


以前面的高熵合金涂层硬度预测任务为例，研究者已经用原子半径差、平均电负性、电负性差、价电子浓度、是否含硬质相和平均晶粒尺寸等特征预测涂层硬度。接下来需要判断的是：这些特征为什么可作为输入？哪些字段不应该作为输入？数据应该如何筛选，才能让模型学习到更接近材料规律的关系？


首先，物理知识可用于指导特征选择。对于高熵合金涂层硬度来说，硬度提高通常可能来自固溶强化、细晶强化、第二相强化和组织致密化等机制。因此，输入特征应该尽量围绕这些机制设计。


例如，原子半径差可与晶格畸变和固溶强化联系起来。多种元素固溶在同一晶格中时，元素原子半径差异会导致晶格畸变，从而可能阻碍位错运动，提高材料硬度。平均晶粒尺寸可与细晶强化联系起来。晶粒越细，晶界数量越多，位错运动越容易受到阻碍，硬度可能提高。是否含硬质相可与第二相强化联系起来。如果涂层中生成 TiC、碳化物或硼化物等硬质相，这些硬质相可能作为强化相提高涂层硬度。


因此，这些特征不是随便选出来的数字，而是和材料强化机制存在对应关系。模型使用这些特征预测硬度，才有可能学习到具有材料意义的规律。


这种关系可整理为表 \ref{tab:mechanism-feature-mapping}。


\begin{formalTable}[htbp]
\caption{材料机制与候选特征对应关系}
\label{tab:mechanism-feature-mapping}
\begin{tabularx}{\textwidth}{@{}YYY@{}}
\toprule
\textbf{材料机制} & \textbf{可转化的模型特征} & \textbf{对硬度的可能影响} \\
\midrule
固溶强化 & 原子半径差、电负性差、价电子浓度 & 晶格畸变和相稳定性变化可能影响硬度 \\
细晶强化 & 平均晶粒尺寸 & 晶粒细化可能提高硬度 \\
第二相强化 & 是否含硬质相、硬质相含量 & 硬质相可能提高涂层硬度 \\
组织致密化 & 孔隙率、裂纹数量、涂层缺陷 & 孔隙和裂纹可能降低硬度和服役性能 \\
\bottomrule
\end{tabularx}
\end{formalTable}

这个表格体现了物理知识嵌入的第一层含义：把材料机理转化为模型特征。模型输入不是任意字段，而是由材料知识筛选出来的、有明确含义的特征。


在实际建模中，这一步应先由材料机理提出候选特征，再结合已有实验记录判断哪些特征可获得、哪些特征需要补充测试。特征设计不应只看字段是否存在，还要看字段是否能对应到明确的材料机制。


其次，物理知识可帮助排除不合理特征。并不是所有出现在数据表中的字段都适合作为模型输入。例如，样品编号、测试日期、数据来源、实验人员姓名等字段，通常不应该作为材料性能预测特征。它们可能与硬度在数据表中出现某种偶然相关，但这种相关并不具有材料意义。如果模型根据样品编号预测硬度，就说明它可能学到了数据排序或实验批次的偶然关系，而不是材料规律。


前文已经讨论过目标泄漏，这里只强调它与物理知识嵌入的关系：输入字段必须在预测之前可获得，且不能由目标结果推导而来。例如，“硬度等级”“是否高硬度样品”“硬度排序”这类字段不能作为硬度预测输入，因为它们已经包含了目标硬度的信息。研究者需要检查每一个输入字段：它是否具有材料意义？是否直接包含目标结果？如果无法回答这些问题，就不应该轻易放入模型。


第三，物理知识可指导数据筛选和标准化。材料实验数据往往受到测试条件影响。以硬度为例，不同载荷、不同保载时间、不同测试位置可能得到不同硬度值。对于高熵合金涂层，表面硬度、截面硬度和不同层位硬度也可能存在差异。如果把不同测试条件下的硬度值直接混合到一个数据集中，模型可能学到混乱关系。


在建立硬度预测模型时，应该尽量保证硬度数据具有可比性。例如，尽量使用相同测试方法、相近测试载荷、相同单位和相同样品位置的数据。如果不同数据来源不可避免，就应在数据表中记录测试条件，必要时把测试条件作为附加字段，或者只选择条件一致的数据进行建模。


这种数据筛选本身就是物理知识嵌入的一部分。它告诉模型：哪些数据可放在一起学习，哪些数据不宜简单混用。对于材料科学来说，数据标准化不是形式工作，而是保证模型学习到真实规律的前提。


检查数据可比性时，应重点核验测试载荷、测试单位、热处理条件、样品位置和数据来源等字段。最终是否保留某条数据，仍然需要研究者根据实验条件和研究目标判断。


第四，物理知识可作为模型结果的合理性检查。模型训练完成后，如果预测结果明显违背材料常识，就需要警惕。例如，如果模型预测孔隙率越高硬度越高，或者晶粒尺寸越大硬度越高，就需要进一步检查数据是否存在偏差、特征是否相关、样本数量是否太少，或者模型是否受异常样品影响。当然，材料规律并不总是简单单调关系，但当模型结果与常识明显冲突时，必须进行核查，而不是直接接受。


对于高熵合金涂层硬度预测，可设定一些基本的合理性检查。例如，含硬质相的样品不一定总是硬度最高，但如果模型完全认为硬质相无关，就需要检查硬质相字段是否准确；晶粒细化通常有利于硬度提高，如果模型得出相反趋势，也要检查晶粒尺寸统计是否一致；孔隙率增加通常不利于力学性能，如果模型认为孔隙率越高越好，就需要检查数据是否存在偶然相关或测量问题。


第五，物理知识还可用于更高级的模型约束。在入门阶段，研究者主要通过特征选择、数据筛选和结果检查嵌入材料知识。对于更深入的研究，还可把物理规律直接写进模型或损失函数中。例如，在原子尺度模型中，可要求模型满足平移、旋转和置换对称性；在热力学相关任务中，可引入稳定性约束；在力学模型中，可要求某些预测量满足基本物理关系。这类方法通常称为物理约束机器学习或物理知识嵌入机器学习。


对于本章而言，不需要展开复杂数学形式，但要让研究者理解一个基本思想：材料模型不应只追求拟合数据，还要尽量尊重材料科学规律。一个模型如果预测误差很小，但违背基本物理常识，或者依赖无法解释的虚假相关字段，那么它在材料研究中的价值是有限的。


可用如下流程概括物理知识嵌入在高熵合金涂层硬度预测中的作用：


\begin{lightquote}
\centering 研究目标 → 材料机理分析 → 特征选择 → 数据筛选 → 模型训练 → 结果解释 → 实验和机理验证
\end{lightquote}

其中，材料机理分析不是最后才出现，而是从一开始就参与建模。研究者首先根据固溶强化、细晶强化、第二相强化等机制选择特征；再根据测试条件筛选可比数据；然后训练模型；最后再用实验表征和材料机理检查模型解释是否合理。


为了帮助研究者理解这个流程，可设计一个简单的操作任务。假设已经有一张高熵合金涂层硬度数据表，其中包含以下字段：样品编号、化学成分、原子半径差、电负性差、价电子浓度、是否含硬质相、晶粒尺寸、孔隙率、硬度等级、硬度/HV、测试载荷、数据来源。研究者需要判断哪些字段可作为输入特征，哪些字段应该保留但不作为模型输入，哪些字段必须删除或谨慎使用。


按照物理知识嵌入的原则，可这样处理：原子半径差、电负性差、价电子浓度、是否含硬质相、晶粒尺寸和孔隙率可作为候选输入特征，因为它们具有材料意义；测试载荷和数据来源应保留，用于判断数据是否可比，但不一定直接作为性能特征；样品编号通常不应作为输入；硬度等级不能作为输入，因为它可能由硬度结果推导而来，存在目标泄漏风险；硬度/HV 是预测目标 y，不应同时作为输入特征 X。


可整理为表 \ref{tab:physics-informed-field-treatment}。


\begin{formalTable}[htbp]
\caption{物理知识嵌入建模过程中的字段处理}
\label{tab:physics-informed-field-treatment}
\begin{tabularx}{\textwidth}{@{}YYY@{}}
\toprule
\textbf{字段} & \textbf{是否作为模型输入} & \textbf{原因} \\
\midrule
原子半径差 & 可 & 与晶格畸变和固溶强化有关 \\
电负性差 & 可 & 与元素化学差异和相形成有关 \\
价电子浓度 & 可 & 可能影响相结构稳定性 \\
是否含硬质相 & 可 & 对应第二相强化 \\
晶粒尺寸 & 可 & 对应细晶强化 \\
孔隙率 & 可 & 反映组织致密性 \\
测试载荷 & 谨慎使用或用于筛选 & 影响硬度可比性 \\
数据来源 & 不作为材料特征 & 用于追溯和分组检查 \\
样品编号 & 不使用 & 无材料意义 \\
硬度等级 & 不使用 & 可能造成目标泄漏 \\
硬度/HV & 目标值 & 模型要预测的 y \\
\bottomrule
\end{tabularx}
\end{formalTable}

这类字段筛选任务适合用提示词形式进行辅助检查。提示词应要求同时判断材料意义、数据可比性和目标泄漏风险，而不是只问“哪些列可输入模型”。

\begin{samepage}
\begin{userbox}
我有一张高熵合金涂层硬度数据表，字段包括：样品编号、化学成分、原子半径差、电负性差、价电子浓度、是否含硬质相、晶粒尺寸、孔隙率、硬度等级、硬度/HV、测试载荷、数据来源。请帮我判断哪些字段适合作为模型输入，哪些字段应保留但不作为输入，哪些字段可能造成目标泄漏或虚假相关。

要求如下：
\begin{enumerate}
\item 按固溶强化、细晶强化、第二相强化和组织致密化解释可用特征；
\item 指出哪些字段用于追溯或筛选数据，但不宜作为材料性能特征；
\item 明确指出硬度/HV 和硬度等级为什么不能作为输入特征；
\item 提醒哪些测试条件会影响硬度数据可比性；
\item 不要只给机器学习角度判断，要结合材料机理说明理由。
\end{enumerate}
\end{userbox}

\begin{aibox}
可按三类处理这些字段：

\begin{enumerate}
\item 可作为候选输入特征的字段：原子半径差、电负性差、价电子浓度、是否含硬质相、晶粒尺寸、孔隙率。这些字段分别可与固溶强化、相结构稳定性、第二相强化、细晶强化和组织致密化联系起来。
\item 应保留但不宜直接作为材料性能特征的字段：测试载荷、数据来源、化学成分和样品编号。测试载荷和数据来源主要用于判断数据可比性和可追溯性；样品编号通常没有材料意义；化学成分可保留，但进入模型前通常需要转化为成分描述符。
\item 不能作为输入特征的字段：硬度/HV 是预测目标，硬度等级往往由硬度结果划分而来，二者如果进入输入特征，会造成目标泄漏。
\end{enumerate}

此外，硬度数据是否可比还取决于测试载荷、保载时间、测试位置、样品表面状态和测试方法。如果这些条件不一致，应先进行筛选、分组或记录为元数据，再决定是否合并建模。
\end{aibox}
\end{samepage}

需要强调的是，物理知识嵌入并不意味着模型一定更复杂。很多时候，它体现为更谨慎的数据选择、更合理的特征设计和更严格的结果检查。对于材料人工智能入门阶段的研究者来说，最重要的不是马上构建复杂的物理约束模型，而是养成一种意识：建模前先问材料问题，建模中保留物理意义，建模后回到实验和机理验证。


本节说明了物理知识嵌入的基本思想。可解释性帮助研究者在模型训练后理解模型，物理知识嵌入则帮助研究者在建模前和建模过程中约束模型。二者共同说明：人工智能用于材料科学，并不是把数据机械地交给模型，而是让数据驱动方法与材料科学知识相互配合。至此，第\ref{chap:materials-cognition-prediction}章完成了从材料数据、材料表示、性能预测、原子尺度模拟到模型解释的完整链条。下一章将进一步讨论逆向设计问题，即如何根据目标性能反向寻找和设计新材料。

\subsubsection*{\textbf{小结}}

本节说明模型解释和物理知识嵌入共同服务于材料判断。
特征重要性和 SHAP 可提示模型依据，但不能直接证明因果机理，仍需 XRD、SEM、EDS 等证据支持。
物理知识应贯穿数据筛选、特征选择、模型训练和结果检查，避免模型学习虚假相关。

\subsection*{习题}

\begin{enumerate}[leftmargin=*]
\item SHAP 为什么不能直接证明材料机理？还需哪些表征？
\item 可解释性分析如何辅助材料筛选？请举例说明。
\item 只有 FCC 数据时，能否预测 BCC 合金？请说明风险和验证方案。
\end{enumerate}

\section*{本章小结}

本章围绕材料认知与性能预测展开，形成了从数据到模型、再从模型回到材料问题的基本链条。材料数据整理强调字段完整、来源可追溯和条件可比较；材料表示强调把化学式、结构和局域环境转化为描述符或图结构；性能预测部分说明了描述符模型、图神经网络和机器学习势在不同尺度任务中的作用；模型解释与物理知识嵌入则进一步说明，模型预测结果必须经过材料机理、实验表征和物理约束的共同检验。由此可见，人工智能在材料科学中的价值不在于替代实验和理论，而在于提高数据利用效率、缩小候选范围并辅助形成可验证的材料判断。

\chapter{材料设计与智能发现}
\label{chap:materials-design-discovery}


第\ref{chap:materials-cognition-prediction}章讨论正向预测，即在材料成分、结构或实验条件已知的情况下，预测材料可能具有的性能。本章进一步讨论逆向设计与智能发现：当目标性能已经给定时，如何从大量候选材料中筛选出值得验证的方案。这个过程并不是简单地让模型“生成答案”，而是要把目标性能、稳定性、可合成性、成本、安全性和实验条件共同纳入考虑。

本章按照材料发现的一般流程展开，重点解决四个问题。第一节说明如何把应用需求转化为材料设计目标、候选变量和约束边界。第二节讨论候选筛选和主动学习如何帮助研究者逐层缩小搜索空间，并安排下一轮实验或计算。第三节介绍生成设计如何提出库外候选，以及生成候选为什么必须经过有效性、稳定性、新颖性和可合成性验证。第四节讨论可信发现、方法选择和闭环迭代，强调模型推荐必须与计算验证、实验反馈和应用约束共同使用。通过本章学习，研究者应理解人工智能如何参与材料设计决策，同时认识候选材料最终仍需要通过物理判断、计算验证和实验结果逐步确认。




\clearpage
\section{问题转换与设计边界}


\begin{sectionkeypoints}
\item 正向预测回答“已知材料可能具有什么性质”，逆向设计回答“为了获得目标性质应当寻找什么材料”。
\item 逆向设计本质上是受约束搜索，不是单纯追求最高预测分数，而是在稳定性、可合成性、成本、安全和设备条件内寻找可验证候选。
\item 材料设计对象可能位于成分、原子结构、微结构、工艺和服役环境等不同层级，层级选择必须与研究目标一致。
\item 候选筛选、主动学习和生成设计是本章三条核心路线，实际研究中常常组合使用。
\end{sectionkeypoints}

\subsection{正向预测与逆向设计}

第\ref{chap:materials-cognition-prediction}章讨论的是正向问题：给定材料的成分、结构、工艺或表征信息，预测形成能、带隙、硬度、折射率或熔点等性质。第\ref{chap:materials-design-discovery}章把问题方向反过来：当目标性质已经确定时，寻找哪些材料、配方、结构或工艺值得进一步尝试。前者回答“这种材料具有什么性质”，后者回答“具有目标性质的材料可能在哪里”。


材料科学长期依赖经验规则和实验试错，但多组元合金、复杂晶体和玻璃配方的候选数量远超实验能够逐一验证的范围。逆向设计的作用不是取消经验和实验，而是把大规模穷举转化为有方向的搜索：先明确目标和约束，再用模型提出或排序候选，最后通过计算和实验逐步确认。


从数学上看，逆向设计可理解为受约束优化。材料对象 x 可以是化学式、晶体结构、玻璃配方、微结构参数或制备条件；正向模型 f(x) 估计目标性能；约束 g(x) 描述稳定性、可合成性、成本、安全和设备边界。搜索的目的不是让某个预测分数无限增大，而是在满足约束的前提下找到一组可验证的候选。正向预测、逆向设计与闭环发现之间的关系见表 \ref{tab:forward-inverse-closed-loop}。


\begin{formalTable}[htbp]
\caption{正向预测、逆向设计与闭环发现}
\label{tab:forward-inverse-closed-loop}
\footnotesize
\begin{tabularx}{\textwidth}{@{}YYYYY@{}}
\toprule
\textbf{问题类型} & \textbf{输入} & \textbf{输出} & \textbf{材料例子} & \textbf{主要风险} \\
\midrule
正向预测 & 候选材料 x & 性质 f(x) & 由晶体结构预测带隙 & 训练集外误差被低估 \\
逆向设计 & 目标性质 y & 候选材料 x & 按目标折射率设计玻璃 & 候选不可合成或不稳定 \\
闭环发现 & 已有数据 + 目标 & 下一步实验 & 主动学习选择合金配比 & 实验噪声导致错误反馈 \\
\bottomrule
\end{tabularx}
\end{formalTable}

\subsubsection*{\textbf{1. 问题转换：从应用目标到材料任务}}

逆向设计的第一步，不是直接选择模型，而是把应用语言转化为材料科学语言。应用目标往往比较笼统，例如“希望玻璃折射率更高”“希望发光材料颜色更接近蓝光”“希望涂层在高温下更稳定”。这些表述还不宜直接进入模型，因为模型需要明确的输入变量、目标性能和约束条件。只有把应用目标拆解成可计算、可测量、可验证的材料任务，后续的筛选、优化和生成才有依据。

以 AR\footnote{Augmented Reality，增强现实。} 光波导玻璃为例，“高折射率”只是一个目标方向。真正进入设计流程时，还需要进一步明确折射率测量波长、透过率范围、色散要求、玻璃化转变温度（Tg\footnote{Glass transition temperature，玻璃化转变温度。}）、热膨胀系数、密度、原料安全性和熔制温度窗口。如果只追求折射率，模型可能推荐含高极化率重金属氧化物较多的配方，但这类配方可能带来密度升高、成本增加、透过率下降或环保风险。因此，逆向设计不是单目标求最大值，而是多目标、多约束条件下的候选选择。

其他材料任务也需要类似拆解。例如 microLED\footnote{Micro Light-Emitting Diode，微型发光二极管。} 发光材料可把目标颜色转化为发光波长或带隙范围，但仍需补充缺陷态、载流子复合效率、热稳定性和外延匹配等约束。这里不展开为新的主线案例，只用它说明：目标性能通常只是材料设计任务的一部分。

具体而言，问题转换通常包括四个问题：目标性能是什么？候选变量是什么？约束条件是什么？验证方式是什么？目标性能决定模型要预测什么；候选变量决定模型可搜索哪些材料空间；约束条件决定哪些候选必须排除；验证方式决定候选是否能从模型结果走向真实材料。表中的 DFT\footnote{Density Functional Theory，密度泛函理论。} 指密度泛函理论计算。典型应用目标到材料设计任务的转换见表 \ref{tab:application-to-design-task}。

\begin{formalTable}[htbp]
\caption{应用目标到材料设计任务的转换}
\label{tab:application-to-design-task}
\footnotesize
\begin{tabularx}{\textwidth}{@{}YYYYY@{}}
\toprule
\textbf{应用目标} & \textbf{材料任务} & \textbf{候选变量} & \textbf{关键约束} & \textbf{验证方式} \\
\midrule
高折射率玻璃 & 预测并优化折射率、Tg、热膨胀系数 & 氧化物组成比例 & 玻璃形成、透明性、熔制温度 & 熔制实验、折射率和透过率测试 \\
蓝光发光材料 & 筛选合适带隙和稳定结构 & 元素组成、晶体结构、掺杂 & 缺陷、外延兼容、热稳定 & DFT、发光测试、器件验证 \\
高硬度涂层 & 优化成分、硬质相和组织参数 & 合金成分、工艺、相组成 & 裂纹、孔隙、结合强度 & 显微硬度、SEM/XRD、服役测试 \\
超高温陶瓷 & 寻找高熔点和抗氧化候选 & 化学组成、晶体结构 & 高温稳定、氧化行为、可烧结性 & 高温实验、热分析、模拟验证 \\
\bottomrule
\end{tabularx}
\end{formalTable}

这一转换过程适合借助 AI 辅助梳理，但不能让 AI 直接给出最终材料结论。提示词应要求它把应用目标拆成材料变量、目标性能、约束条件和验证方式，而不是只列出几个候选材料名称。

\begin{samepage}
\begin{userbox}
我想设计一种用于 AR 光波导的高折射率玻璃。请帮我把这个应用目标转化为材料逆向设计任务。

要求如下：
\begin{enumerate}
\item 明确目标性能，例如折射率、透过率、色散、Tg、热膨胀系数等；
\item 说明候选变量可包括哪些氧化物组成或配方参数；
\item 区分必须满足的硬约束和可折中的软目标；
\item 给出候选筛选后需要进行哪些计算或实验验证；
\item 不要直接编造最终玻璃配方，只做任务拆解。
\end{enumerate}
\end{userbox}

\begin{aibox}
可将“高折射率 AR 光波导玻璃”拆解为以下设计任务：

\begin{enumerate}
\item \textbf{目标性能：}折射率是核心目标，但还应同时记录透过率、色散、Tg、热膨胀系数、密度和加工温度等指标。
\item \textbf{候选变量：}可把各氧化物组分比例作为设计变量，例如网络形成体、网络修饰体和提高折射率的氧化物比例。
\item \textbf{硬约束：}配方比例之和应满足组成约束；原料不能超出安全和环保边界；熔制温度应在设备可承受范围内；样品应能够形成玻璃而不是明显析晶。
\item \textbf{软目标：}折射率越高越好，但不能以明显牺牲透明性、热稳定性和加工性为代价。
\item 验证方式：模型筛选后，需要进行玻璃形成能力判断、熔制实验、折射率测试、透过率测试和热膨胀系数测试。
\end{enumerate}

这样处理后，逆向设计任务不再只是“寻找高折射率材料”，而是“在可熔制、可成玻、透明且热性质合适的约束下，寻找高折射率玻璃配方”。
\end{aibox}
\end{samepage}

问题转换应当在建模之前完成，而不是在模型给出候选之后再补充解释。如果目标一开始没有定义清楚，后面所有筛选和排序都会变得模糊。例如，若只写“设计高性能玻璃”，模型无法判断“高性能”指的是高折射率、高透过率、低热膨胀系数还是高热稳定性；若只写“设计蓝光材料”，模型也不知道重点是发光波长、发光效率、热稳定性还是外延兼容性。本章强调问题转换，是为了让研究者意识到：人工智能材料设计首先是材料问题定义，其次才是模型选择。

一个规范的逆向设计任务，宜用一句话概括为“在什么约束下，寻找什么类型的材料，使什么目标性能达到什么范围”。例如，“在不使用高风险元素、熔制温度不超过设备上限、能够形成透明玻璃的约束下，寻找折射率较高且热膨胀系数合适的氧化物玻璃配方”。这样的表述比“让模型设计高折射率玻璃”更清楚，也更便于后续建立数据表、选择描述符和设计验证实验。

\subsection{设计对象与约束边界}

\subsubsection*{\textbf{1. 设计边界：材料对象分层}}

逆向设计并不总是直接生成一个化学式。材料对象通常包含多个层级：成分决定由哪些元素或组分构成；原子结构决定空间群、占位、缺陷和局域配位；微结构包含晶粒、孔隙、第二相和界面；工艺包含熔融、退火、沉积、烧结和冷却条件。同一成分经过不同工艺可得到不同结构和性能，因此设计对象的层级必须与研究目标相匹配。材料逆向设计常见对象层级见表 \ref{tab:inverse-design-object-levels}。


\begin{formalTable}[htbp]
\caption{材料逆向设计的对象层级}
\label{tab:inverse-design-object-levels}
\begin{tabularx}{\textwidth}{@{}YYYY@{}}
\toprule
\textbf{设计层级} & \textbf{常见变量} & \textbf{适用任务} & \textbf{容易忽略的问题} \\
\midrule
成分 & 元素/氧化物比例、掺杂量 & 玻璃配方、合金筛选、催化剂组分 & 相形成与工艺窗口 \\
原子结构 & 晶格、坐标、缺陷、表面 & 带隙、离子迁移、稳定性 & 对称性和局域环境 \\
微结构 & 晶粒、孔隙、第二相、取向 & 力学、热学、疲劳与可靠性 & 制备历史强相关 \\
工艺 & 温度、气氛、压力、时间 & 可合成性与性能优化 & 元数据缺失会误导模型 \\
服役环境 & 光、电、热、湿度、载荷 & 器件可靠性和寿命 & 实验条件与应用条件不一致 \\
\bottomrule
\end{tabularx}
\end{formalTable}

除了设计层级，还需要区分约束的类型。材料设计中的约束可分为硬约束和软约束。硬约束是必须满足的条件，例如有毒元素不能使用、设备最高温度不能超过、配方比例必须归一化、候选结构不能明显违反价态和原子距离规则。软约束是希望尽量优化但可折中的条件，例如成本越低越好、密度越低越好、热膨胀系数越接近目标越好。硬约束用于排除候选，软约束用于排序候选。

如果不区分硬约束和软约束，逆向设计很容易变成一个不清楚的打分问题。例如，把折射率、成本、热稳定性和环保要求都压缩成一个总分，表面上方便排序，但不同目标之间的取舍会被隐藏。逆向设计应把硬约束、软目标和数据边界分开说明。

设计边界还包括数据边界。模型只能在训练数据覆盖的范围内比较可靠地工作。如果训练数据主要来自硅酸盐玻璃，那么把模型直接用于完全不同的硫系玻璃或氟化物玻璃，预测结果就可能不稳定。如果模型只学习了室温性质，就不宜直接推断高温服役行为。如果数据没有记录制备工艺，模型也很难判断同一成分在不同热处理条件下的性能差异。因此，逆向设计开始之前必须说明数据从哪里来、覆盖哪些材料体系、哪些变量没有记录。不同边界类型及其处理方式见表 \ref{tab:inverse-design-boundary-types}。

\begin{formalTable}[htbp]
\caption{材料逆向设计中的边界类型}
\label{tab:inverse-design-boundary-types}
\footnotesize
\begin{tabularx}{\textwidth}{@{}YYYY@{}}
\toprule
\textbf{边界类型} & \textbf{含义} & \textbf{材料例子} & \textbf{处理方式} \\
\midrule
硬约束 & 不满足就直接排除 & 有毒元素、不可成玻、设备温度超限 & 筛选前写入规则 \\
软目标 & 可折中但要排序 & 折射率、密度、成本、热膨胀系数 & 作为多目标评价项 \\
数据边界 & 模型训练数据覆盖范围 & 只含某类玻璃或某类晶体 & 明确适用范围和外推风险 \\
工艺边界 & 实际制备条件限制 & 熔制温度、气氛、烧结时间 & 作为实验可行性约束 \\
\bottomrule
\end{tabularx}
\end{formalTable}

在教学中，可让研究者首先写出设计边界，再选择模型。这样的顺序比“先选一个算法再想怎么用”更合理。算法只能在已经定义清楚的问题上发挥作用；如果目标、变量和边界本身模糊，模型输出再复杂也难以转化为真实材料方案。

这一类边界梳理同样可使用 AI 辅助，但提示词必须要求区分硬约束、软目标和数据边界。

\begin{samepage}
\begin{userbox}
我准备做一个材料逆向设计任务，目标是寻找候选高折射率玻璃。请帮我列出这个任务的设计边界。

要求如下：
\begin{enumerate}
\item 区分硬约束、软目标、数据边界和工艺边界；
\item 每一类边界都要说明它为什么重要；
\item 不要把所有条件合并成一个总分；
\item 说明哪些信息需要在实验记录或数据库中保留；
\item 不要直接给出候选配方。
\end{enumerate}
\end{userbox}

\begin{aibox}
可从四类边界整理：

\begin{enumerate}
\item \textbf{硬约束：}原料安全、环保限制、组成比例合法、可形成玻璃、设备温度可达。这类条件不满足时，候选应直接排除。
\item \textbf{软目标：}折射率、透过率、密度、热膨胀系数、成本和加工性能。这些目标之间可能存在折中，应分别记录和排序。
\item \textbf{数据边界：}需要说明训练数据主要来自哪些玻璃体系、性质是在什么条件下测得、是否包含失败样品和析晶样品。超出数据覆盖范围时，模型预测应谨慎使用。
\item \textbf{工艺边界：}熔制温度、保温时间、冷却方式、气氛、坩埚材料等都会影响玻璃形成和性能，应作为实验元数据保存。
\end{enumerate}

这些边界不是附加说明，而是逆向设计任务的一部分。只有先限定边界，模型推荐的候选才有可能进入真实实验。
\end{aibox}
\end{samepage}

边界条件还应随研究阶段动态更新。早期阶段可设置较宽的组成范围，用来观察模型是否能发现潜在趋势；进入实验阶段后，边界应逐渐收紧，例如限定可采购原料、设备温度、样品尺寸和测试条件；进入应用阶段后，还要加入成本、安全、长期稳定性和批量制备要求。换言之，边界不是一次写完就不变，而是随着数据积累和实验反馈不断修正。

对于研究者来说，最容易忽略的是“不可见边界”。有些限制并不直接写在数据表中，例如某些配方虽然数据库中有记录，但实验室没有相应原料；某些结构计算上稳定，但现有设备无法提供所需压力或气氛；某些材料性能很好，但含有不适合应用场景的元素。这些不可见边界如果不提前说明，模型推荐就会表面上合理、实际不可执行。因此，材料逆向设计的边界既包括数据中的字段，也包括实验室条件和应用场景中的真实限制。

本章讨论三条核心路线。候选筛选适合从已有数据库或枚举空间中逐层过滤；主动学习适合实验昂贵、每轮只能验证少量样品的任务；生成设计用于突破已有候选库，在目标条件下提出新结构或新配方。三者并非互相替代，实际研究常按“筛选--优化--生成--验证”的方式组合使用。

\subsubsection*{\textbf{小结}}

本节说明逆向设计不是直接生成材料，而是把应用需求转化为目标、变量和约束。
正向预测回答已知材料的性能，逆向设计回答为了目标性能应寻找哪些候选。
材料对象层级、硬约束、软目标、数据边界和工艺边界需要在建模前明确。

\subsection*{习题}

\begin{enumerate}[leftmargin=*]
\item 正向预测和逆向设计在高折射率玻璃任务中如何配合？
\item 逆向设计为什么必须先明确目标和约束？
\item 如何把高折射率玻璃需求转化为计算目标？
\end{enumerate}

\clearpage
\section{候选筛选与主动学习}


\begin{sectionkeypoints}
\item 候选筛选的基本逻辑是先建立候选库，再按照成本由低到高、可信度由粗到细逐层过滤。
\item 筛选漏斗通常包含规则过滤、机器学习代理模型、高精度计算和实验验证等层级，每一级都应记录通过和排除理由。
\item 主动学习适合实验昂贵、小样本和需要多轮优化的材料任务，核心是同时考虑预测性能和模型不确定性。
\item 实验噪声、批次效应和失败样品都应进入数据记录，否则模型可能把制备差异误认为材料本征规律。
\item 候选排序和主动学习评分只负责辅助决策，最终实验批次仍需要材料约束、工艺可行性和人工复核。
\end{sectionkeypoints}

\subsection{筛选漏斗与分层验证}

高通量虚拟筛选的基本思路，是先建立候选库，再按照成本由低到高、可信度由粗到细逐层过滤。最外层使用元素排除、价态平衡、成分范围、成本或毒性等规则；中间层使用第\ref{chap:materials-cognition-prediction}章建立的正向模型快速估计形成能、带隙、弹性模量或折射率；内层使用 DFT、分子动力学或热力学计算精算；最后才进入实验合成与表征。


\begin{figure}[H]
\centering
\BookFigureImage{images_11/3.png}
\caption{候选筛选与验证层级示意图}
\label{fig:candidate-screening-validation}
\end{figure}
\FloatBarrier

如图 \ref{fig:candidate-screening-validation} 所示，筛选漏斗利用了不同方法之间的成本差异。规则过滤几乎不需要计算，机器学习预测可快速处理大量候选，而高精度计算和真实实验只能用于较小范围。把便宜但粗略的方法放在前面，把昂贵但可靠的方法放在后面，能够避免把有限资源消耗在明显不合理的候选上。


筛选的第一个常见门槛是稳定性。无机晶体可参考形成能和凸包能量，合金和陶瓷还要考虑相分离、氧化和高温相变，玻璃则需要考虑玻璃形成能力、析晶倾向和熔制窗口。计算稳定不等于实验一定能够合成，热力学判断还要与动力学路径、前驱体和设备条件结合。


第二个门槛是多目标要求。AR 光波导玻璃不仅需要较高折射率，还要兼顾透过率、色散、玻璃化转变温度、热膨胀系数、密度和加工性能；microLED 发光材料除带隙外，还要考虑缺陷、载流子复合、热稳定和外延兼容性。相应地，候选排序不应只依据单一预测值，而应保留不同目标之间的折中关系。筛选漏斗中不同层级的作用见表 \ref{tab:screening-funnel-levels}。


\begin{formalTable}[htbp]
\caption{筛选漏斗中不同层级的作用}
\label{tab:screening-funnel-levels}
\begin{tabularx}{\textwidth}{@{}YYYY@{}}
\toprule
\textbf{筛选层级} & \textbf{常见工具} & \textbf{适合回答的问题} & \textbf{学习要点} \\
\midrule
规则层 & 元素价态、毒性、成本、工艺窗口 & 哪些候选明显不可取？ & 规则便宜但粗糙，适合第一轮排除 \\
代理模型层 & 回归/分类模型、图神经网络、集成模型 & 哪些候选可能性能较好？ & 模型用于排序，不等于最终证明 \\
物理计算层 & DFT、相图、分子动力学、热力学模型 & 候选是否稳定、机制是否可信？ & 精度更高但成本也更高 \\
实验层 & 合成、表征、服役测试 & 候选能否成为真实材料？ & 实验验证是材料发现的终点而非附属步骤 \\
\bottomrule
\end{tabularx}
\end{formalTable}

\subsubsection*{\textbf{1. 案例示范：高折射率玻璃候选筛选}}

以高折射率玻璃为例，首先把应用要求转化为可计算指标；然后限定各氧化物比例之和、原料安全和设备温度等边界；代理模型预测折射率、Tg、热膨胀系数和析晶倾向；物理或经验模型检查玻璃形成与黏度窗口；最后制备少量代表性配方并测量真实光学损耗。每一级都应记录排除理由，使候选清单可复核。

具体而言，第一层规则可先排除明显不合理的配方。例如，各氧化物比例之和必须为 100\%，单一组分含量不能超过玻璃形成的经验范围，含有明显不适合教学或工程应用的高风险元素时应直接排除。第二层可用代理模型预测折射率和热性质，把候选从几千个缩小到几十个。第三层再用更可靠的热力学判断或经验模型检查玻璃形成能力、析晶风险和黏度窗口。最后才选择少量配方进入熔制实验。

筛选漏斗的一个重要原则是“保留失败理由”。如果一个候选被排除，应记录它是因为折射率不够、Tg 太低、热膨胀系数过高、含有风险元素，还是因为预测不确定性太大。这样做有两个好处。第一，后续可复核筛选规则是否过于严格；第二，失败信息可反馈给下一轮模型或实验设计。如果只保留最后通过的候选，而删除中间淘汰信息，材料发现过程就会失去可追溯性。

在初学阶段，研究者可将筛选漏斗理解为一个逐层提问的过程：第一问，这个候选是否明显违反基本规则？第二问，它的目标性能是否接近要求？第三问，它是否可能稳定存在或可制备？第四问，它是否值得消耗实验资源？每一层都不是绝对正确的判断，而是在不同成本下尽量减少无效尝试。筛选过程中需要保留的判断信息见表 \ref{tab:screening-funnel-records}。

\begin{formalTable}[htbp]
\caption{筛选漏斗中需要保留的判断信息}
\label{tab:screening-funnel-records}
\footnotesize
\begin{tabularx}{\textwidth}{@{}YYYY@{}}
\toprule
\textbf{筛选问题} & \textbf{排除依据} & \textbf{保留信息} & \textbf{后续用途} \\
\midrule
是否违反基本规则 & 比例错误、风险元素、明显价态异常 & 排除原因和原始配方 & 修正规则或候选生成空间 \\
性能是否接近目标 & 折射率、带隙、硬度等预测值偏离 & 预测值和模型版本 & 比较不同代理模型 \\
是否可能稳定 & 凸包能、析晶倾向、相分离风险 & 稳定性指标和计算方法 & 决定是否进入高精度计算 \\
是否值得实验 & 成本、设备温度、样品数量限制 & 排序分数和实验批次 & 形成可执行实验计划 \\
\bottomrule
\end{tabularx}
\end{formalTable}

这一过程可借助 AI 辅助把筛选规则写成清单，但规则本身必须由材料任务决定。可用如下提示词把流程规范化。

\begin{samepage}
\begin{userbox}
我正在设计一个高折射率玻璃候选筛选漏斗。候选数据表中有氧化物比例、预测折射率、预测 Tg、热膨胀系数、是否含风险元素、预测析晶风险等字段。请帮我把筛选流程整理成逐层规则。

要求如下：
\begin{enumerate}
\item 先列出规则过滤层、代理模型层、物理/经验验证层和实验层；
\item 每一层说明主要排除什么候选；
\item 明确哪些排除理由需要记录；
\item 提醒候选排序不应只看折射率；
\item 不要编造真实配方数据。
\end{enumerate}
\end{userbox}

\begin{aibox}
可将筛选漏斗设计为四层：

\begin{enumerate}
\item \textbf{规则过滤层：}检查氧化物比例是否合理、是否含风险元素、熔制温度是否超出设备范围。被排除的候选应记录具体违反哪条规则。
\item \textbf{代理模型层：}预测折射率、Tg 和热膨胀系数。折射率高但 Tg 过低或热膨胀系数过高的候选不宜直接排在最前。
\item \textbf{物理/经验验证层：}进一步检查玻璃形成能力、析晶倾向、黏度窗口和配方稳定性。此处的目标是排除“预测性能好但难以形成玻璃”的候选。
\item \textbf{实验层：}选择少量代表性配方进行熔制、退火、折射率和透过率测试，并记录失败样品原因。
\end{enumerate}

筛选报告中应保留每一级候选数量、通过率和排除理由。这样可判断筛选规则是否过严，也便于后续主动学习利用失败信息。
\end{aibox}
\end{samepage}

筛选漏斗还有一个容易被忽视的问题：前一级筛选规则会影响后一级能看到什么候选。如果第一层规则过于严格，可能会把一些有潜力但不符合传统经验的候选提前排除；如果第一层规则过于宽松，又会把大量明显不可行的候选留给后续高成本计算。为此，筛选规则需要根据任务目标反复校准，而不是一次确定后不再修改。

在课堂练习中，可让研究者比较两种筛选策略。一种策略偏保守，优先排除风险候选，适合实验资源很有限的任务；另一种策略偏探索，保留更多边界候选，适合希望发现新体系的任务。两种策略没有绝对优劣，关键在于是否说明了筛选目的、资源条件和验证能力。这样可避免把筛选漏斗理解为机械流程，而是理解为材料决策工具。


以下给出一个最小候选排序代码示例。代码把目标带隙、凸包能量和元素风险写成显式评分规则，用于说明筛选假设如何转化为可复核的计算步骤。评分只负责初步排序，不宜代替稳定性复算和实验验证。


\begin{lstlisting}[language=Python,escapeinside={(*@}{@*)}]
import pandas as pd
# (*@\textnormal{gap 的单位为 eV，hull 的单位为 eV/atom}@*)
cands = pd.DataFrame({
    "id": ["mp-1", "mp-2", "mp-3", "mp-4"],
    "formula": ["M1O2", "M2O3", "M3S", "PbM4O"],
    "gap": [2.10, 1.55, 2.35, 2.00],
    "hull": [0.018, 0.006, 0.065, 0.012],
    "elements": ["M1,O", "M2,O", "M3,S", "Pb,M4,O"],
})
target_gap = 2.0
cands["gap_score"] = 1 - (cands["gap"] - target_gap).abs() / target_gap
cands["stability_score"] = (0.08 - cands["hull"]).clip(lower=0) / 0.08
cands["element_penalty"] = cands["elements"].str.contains("Pb|Cd|Hg").astype(int)
cands["priority"] = (
    0.45 * cands["gap_score"]
    + 0.45 * cands["stability_score"]
    - 0.50 * cands["element_penalty"]
)
cols = ["id", "formula", "gap", "hull", "priority"]
print(cands.sort_values("priority", ascending=False)[cols])
\end{lstlisting}

\manualfigcaption{\textbf{代码示例：}候选材料初步排序示例}

\subsection{主动学习与实验决策}

高通量筛选是在固定候选库中做过滤，主动学习则是在搜索过程中不断更新模型。研究者首先用少量实验或计算训练代理模型，模型同时给出候选性能和不确定性，再选择下一批最值得验证的候选；新结果返回后重新训练模型，形成序贯迭代。


\begin{figure}[H]
\centering
\BookFigureImage{images_11/4.png}
\caption{主动学习与生成设计闭环示意图}
\label{fig:active-learning-generative-loop}
\end{figure}
\FloatBarrier

图 \ref{fig:active-learning-generative-loop} 概括了主动学习与生成设计之间的闭环关系。贝叶斯优化是常见的主动学习方法，适合小样本和高成本实验。它通常包含代理模型和采集函数两个部分：代理模型近似真实实验或高精度计算，采集函数综合预测均值与不确定性，决定下一步优先验证哪些候选。其核心不是总选择当前分数最高的样品，而是在“利用”和“探索”之间取得平衡。


利用型候选靠近模型认为的高性能区域，有利于尽快得到改进材料；探索型候选位于模型不确定但物理上仍合理的区域，用于扩展知识边界；重复型或近重复型候选用于估计制备和测量噪声。真实实验批次通常同时包含这几类样品，而不是只给出一个单一的“最佳配方”。主动学习批次中的候选类型见表 \ref{tab:active-learning-batch-types}。


\begin{formalTable}[htbp]
\caption{主动学习批次中的候选类型}
\label{tab:active-learning-batch-types}
\begin{tabularx}{\textwidth}{@{}YYY@{}}
\toprule
\textbf{候选类型} & \textbf{选择依据} & \textbf{在实验轮次中的作用} \\
\midrule
利用型 & 预测性能高且约束满足 & 尽快逼近可用材料 \\
探索型 & 模型不确定但物理上可行 & 扩展模型对未知区域的认识 \\
边界型 & 靠近工艺或相稳定边界 & 识别可制备窗口 \\
重复型 & 已有样品或近似重复配方 & 估计实验噪声和批次误差 \\
\bottomrule
\end{tabularx}
\end{formalTable}

实验噪声和批次效应是主动学习能否落地的关键。同一成分可能因混料、升温速率、保温时间、气氛、坩埚污染或仪器校准而表现不同。如果模型把这些差异全部当成材料本征规律，就会错误更新搜索方向。这意味着，样品批次、设备、制备程序、测量条件和失败原因都应进入数据记录。


失败样品同样提供信息。某类成分若反复析晶、开裂或产生杂相，说明该区域可能超出可制备边界。例如，某轮玻璃熔制实验中如果发现 La\textsubscript{2}O\textsubscript{3} 含量高于 12\% 的配方全部析晶，下一轮模型排序就应把这一区域标记为高风险，并降低相邻候选的推荐优先级。主动学习不应只学习“哪些候选性能高”，还应学习“哪些区域无法稳定制备”。这使下一轮实验决策同时考虑性能、不确定性、工艺风险和已有负结果。


\subsubsection*{\textbf{1. 案例示范：AR 光波导玻璃主动学习}}

AR 光波导玻璃可作为连续案例。初始数据由已有氧化物配方及其折射率、Tg、热膨胀系数和析晶情况构成；代理模型预测候选性质和不确定性；每轮选择若干高性能候选、若干未知区域候选和少量重复样品；实验结果再用于修正模型和玻璃形成边界。这样，模型输出被转化为可执行的实验批次。

主动学习与普通筛选最大的不同，是它把“下一步做什么实验”作为核心问题。普通筛选常常一次性给出候选排序，而主动学习会根据已经完成的实验不断调整模型。第一轮实验后，模型可能发现某个组成区域折射率较高但容易析晶；第二轮就应减少该区域中高风险候选的比例，同时在相邻区域选择少量探索样品。这样，实验不是简单验证模型，而是在持续修正模型。

在真实实验中，一轮实验通常不应只选择预测性能最高的候选。如果全部选择利用型样品，模型可能很快陷入局部区域，对未知空间了解不足；如果全部选择探索型样品，短期内又可能得不到性能改进。为此，一个合理批次通常同时包含高性能候选、模型不确定候选、边界候选和重复样品。重复样品表面上浪费名额，但它可估计实验噪声和批次误差，对判断模型是否真的改进很重要。

主动学习还需要记录失败类型。对于玻璃配方来说，失败可能表现为不能熔清、冷却后析晶、样品开裂、气泡过多、折射率测试失败或热膨胀系数不满足要求。不同失败类型对应不同边界：不能熔清可能是工艺温度问题，析晶可能是玻璃形成能力不足，开裂可能与热膨胀系数或退火制度有关。把这些失败统一写成“失败”会损失大量材料信息。主动学习实验结果的记录方式见表 \ref{tab:active-learning-result-records}。

\begin{formalTable}[htbp]
\caption{主动学习实验结果的记录方式}
\label{tab:active-learning-result-records}
\footnotesize
\begin{tabularx}{\textwidth}{@{}YYYY@{}}
\toprule
\textbf{实验结果类型} & \textbf{记录内容} & \textbf{对模型的作用} & \textbf{下一轮决策} \\
\midrule
性能达标 & 折射率、透过率、Tg、热膨胀系数 & 提供高价值正样本 & 在附近细化搜索 \\
性能一般但可成玻 & 完整性质数据和配方 & 扩展模型覆盖范围 & 可作为训练样本保留 \\
析晶或开裂 & 失败现象、热处理条件 & 标记不可制备边界 & 减少相邻高风险候选 \\
重复样品 & 批次、设备、测试误差 & 估计噪声水平 & 修正不确定性判断 \\
\bottomrule
\end{tabularx}
\end{formalTable}

AI 可辅助把主动学习输出整理成实验批次说明，但提示词中必须要求它保留利用、探索、边界和重复样品的比例，并提醒不能自动执行模型给出的批次。

\begin{samepage}
\begin{userbox}
我正在做 AR 光波导玻璃配方的主动学习优化。模型已经给出候选配方的预测折射率、不确定性、Tg、热膨胀系数、距离已有样本的远近以及是否满足基本约束。请帮我整理下一轮实验批次选择原则。

要求如下：
\begin{enumerate}
\item 区分利用型、探索型、边界型和重复型样品；
\item 说明每类样品为什么需要进入实验批次；
\item 提醒哪些实验元数据必须记录；
\item 说明失败样品如何反馈到下一轮模型；
\item 不要直接把模型最高分样品全部选为下一轮实验。
\end{enumerate}
\end{userbox}

\begin{aibox}
下一轮实验批次可按四类候选组合：

\begin{enumerate}
\item \textbf{利用型样品：}预测折射率较高、Tg 和热膨胀系数满足约束，用于尽快获得性能改进。
\item \textbf{探索型样品：}预测不确定性较高但物理上仍可行，用于扩展模型对未知组成区域的认识。
\item \textbf{边界型样品：}靠近玻璃形成或工艺边界，用于判断哪些区域容易析晶、开裂或无法熔制。
\item \textbf{重复型样品：}选择已有样品或近似重复样品，用于估计实验噪声、批次差异和测试误差。
\end{enumerate}

每个样品都应记录配方、原料批次、熔制温度、保温时间、冷却制度、样品状态、测试设备和失败原因。失败样品不应删除，而应作为不可制备区域或高风险区域的信息反馈给下一轮模型。
\end{aibox}
\end{samepage}

以下代码示例保留了主动学习最核心的计算关系：预测值体现“利用”，不确定性体现“探索”，Tg 和热膨胀系数等惩罚项体现工艺约束。代码输出的是待人工复核的实验批次，而不是可自动执行的最终方案。


\begin{lstlisting}[language=Python,escapeinside={(*@}{@*)}]
import pandas as pd
cand = pd.DataFrame({
    "id": ["G01", "G02", "G03", "G04", "G05", "G06"],
    "n_pred": [1.86, 1.91, 1.88, 1.83, 1.94, 1.89],
    "sigma": [0.012, 0.010, 0.006, 0.030, 0.018, 0.026],
    "Tg": [565, 548, 570, 535, 505, 552],
    "cte": [8.1, 8.7, 8.4, 8.9, 9.8, 8.0],
    "dist": [0.25, 0.18, 0.35, 0.12, 0.08, 0.10],
    "ok": [True, True, True, True, False, True],
})
target_cte = 8.3
cand["penalty"] = (
    (cand["Tg"] < 520).astype(int) * 0.20
    + (cand["cte"] - target_cte).abs() * 0.03
    + (~cand["ok"]).astype(int) * 10
)
cand["score"] = cand["n_pred"] + 0.4 * cand["sigma"] - cand["penalty"]
feasible = cand[cand["ok"]].copy()
use = feasible.sort_values("n_pred", ascending=False).head(2)
unused = feasible.drop(use.index)
explore = unused.sort_values("sigma", ascending=False).head(2)
unused = unused.drop(explore.index)
boundary = unused.sort_values("dist").head(1)
# (*@\textnormal{保留 G03 作为重复样品，用于评估实验重复性}@*)
repeat = feasible[feasible["id"].eq("G03")].copy()
batch = pd.concat([
    use.assign(kind="use"),
    explore.assign(kind="explore"),
    boundary.assign(kind="boundary"),
    repeat.assign(kind="repeat"),
])
cols = ["id", "kind", "n_pred", "sigma", "Tg", "cte", "score"]
print(batch[cols])
\end{lstlisting}

\manualfigcaption{\textbf{代码示例：}主动学习实验批次评分示例}

\subsubsection*{\textbf{小结}}

本节说明候选筛选应按照由低成本到高成本、由粗判断到精验证的漏斗推进。
规则过滤、代理模型、物理计算和实验验证各有作用，每一级都应保留通过和排除理由。
主动学习把预测值、不确定性和实验反馈结合起来，失败样品和噪声记录同样是有价值的数据。

\subsection*{习题}

\begin{enumerate}[leftmargin=*]
\item 高通量筛选为什么常采用分层漏斗？
\item 比较规则过滤、代理模型、DFT 和实验验证的作用。
\item 主动学习中的“利用”和“探索”分别指什么？
\end{enumerate}

\clearpage
\section{生成设计与候选验证}


\begin{sectionkeypoints}
\item 生成设计改变的是候选产生方式，可在目标条件下提出新的组成、结构或配方，但不能取消材料约束和实验验证。
\item 材料表示决定生成模型能够处理的对象。晶体、玻璃、分子和聚合物需要不同表示方式，不宜简单套用同一种输入格式。
\item 条件生成需要把目标性能转化为可计算条件，例如带隙范围、形成能、元素体系、成本边界或供应链限制。
\item 生成候选必须经过有效性、稳定性、新颖性、可合成性和应用性等多层验证，不应只看生成数量。
\item 生成设计通常需要与筛选、主动学习、DFT 计算和实验验证串联使用，才能形成可信的材料发现流程。
\end{sectionkeypoints}

\subsection{条件生成与验证流程}

候选筛选只能在已经列出的空间中寻找材料。如果候选库没有包含某类新结构或新配方，筛选再快也无法发现它。生成设计尝试学习已有材料的分布，并在目标条件下提出新的组成、结构或配方。它改变的是候选产生方式，而不是取消物理约束和实验验证。


材料表示决定生成模型能够处理什么。分子可使用字符串或分子图，晶体需要同时表示元素种类、周期性晶格和原子坐标，玻璃与非晶材料常以成分和短程有序统计表示，聚合物还要考虑单体序列和拓扑。生成表示若不能表达周期性、价态和局域配位，模型容易产生形式正确但没有材料意义的候选。表中的 VAE\footnote{Variational Autoencoder，变分自编码器。} 指变分自编码器，GAN\footnote{Generative Adversarial Network，生成对抗网络。} 指生成对抗网络，自回归模型\footnote{Autoregressive Model，自回归模型，按序列顺序逐步生成数据。} 按步骤生成序列或结构片段，扩散模型\footnote{Diffusion Model，扩散模型，从噪声出发逐步去噪生成数据。} 从噪声中逐步生成结构，SMILES\footnote{Simplified Molecular Input Line Entry System，简化分子线性输入系统。} 是分子线性表示方法。不同生成模型路线与材料对象的对应关系见表 \ref{tab:generative-model-routes}。


\begin{formalTable}[htbp]
\caption{生成模型路线与材料对象}
\label{tab:generative-model-routes}
\begin{tabularx}{\textwidth}{@{}YYYY@{}}
\toprule
\textbf{模型路线} & \textbf{材料生成中的直觉} & \textbf{典型对象} & \textbf{需要特别注意} \\
\midrule
VAE & 把材料压缩到连续潜空间并在其中优化 & 分子、聚合物、部分晶体表示 & 潜空间是否有物理意义 \\
GAN & 生成器与判别器对抗，逼近真实样本分布 & 显微组织、分子、结构图像 & 训练稳定性与约束不足 \\
自回归模型 & 按序列或步骤逐步生成 & SMILES、合成路线、结构 token & 语法有效性和错误传播 \\
扩散模型 & 从噪声逐步去噪生成结构 & 晶体、分子三维结构 & 周期边界、对称性和坐标处理 \\
混合工作流 & 生成、筛选、DFT、实验闭环结合 & 无机晶体、玻璃、催化剂 & 验证流程比模型名称更重要 \\
\bottomrule
\end{tabularx}
\end{formalTable}

VAE、GAN、自回归模型和扩散模型可采用不同方式学习材料分布。VAE 更适合把分子或配方压缩到连续潜空间中再优化；GAN 常用于生成图像式显微组织或分子结构，但训练稳定性和约束控制需要特别检查；自回归模型适合 SMILES、合成路线或结构 token 这类有顺序的表示；扩散模型适合逐步生成分子或晶体三维结构，但必须处理周期性边界、坐标对称性和原子距离合理性。本章不展开各模型的数学细节，只保留共同逻辑：模型从已有材料中学习可行区域，目标条件用于引导生成方向，化学和结构约束用于排除无效样本，正向模型和高精度计算用于评价生成结果。


\subsubsection*{\textbf{1. 案例示范：发光材料条件生成}}

条件生成是逆向设计的关键。目标条件可以是带隙范围、形成能、体模量、磁性、元素体系或供应链限制。以发光材料为例，目标波长可转化为带隙范围，生成模型提出候选组成和结构，再由第\ref{chap:materials-cognition-prediction}章的带隙模型、DFT 和实验逐级验证。目标带隙只是入口，缺陷、掺杂、热稳定和器件兼容性仍需单独检查。

对于发光材料设计，条件生成可从两个层次理解。第一层是性能条件，例如目标发光波长、带隙范围、热稳定性和缺陷容忍度。第二层是材料边界，例如元素体系、晶体结构类型、掺杂元素、成本和安全要求。如果只给生成模型一个带隙条件，模型可能提出带隙合适但不稳定、含风险元素或难以合成的候选。因此，条件生成需要同时输入目标和边界。

发光材料还具有明显的“性质链条”。目标颜色对应发光波长，发光波长与带隙有关，但材料是否真正发光还取决于缺陷、能级、载流子复合和非辐射损失。换言之，带隙只是候选进入下一步验证的第一道门槛。一个候选材料即使带隙接近目标，如果存在深能级缺陷、热稳定性差或与器件结构不兼容，也不宜直接认为适合应用。

由此可见，生成设计更适合作为“候选扩展器”。它可帮助研究者跳出已有数据库或人工枚举空间，提出一些值得进一步检查的新组成或新结构。但候选是否有意义，要经过规则检查、正向预测、结构弛豫、高精度计算和实验验证。生成数量本身没有价值，真正有价值的是通过层层验证后仍然保留下来的候选。发光材料条件生成中的约束类型见表 \ref{tab:luminescent-material-constraints}。

\begin{formalTable}[htbp]
\caption{发光材料条件生成中的约束类型}
\label{tab:luminescent-material-constraints}
\footnotesize
\begin{tabularx}{\textwidth}{@{}YYYY@{}}
\toprule
\textbf{条件类型} & \textbf{发光材料中的含义} & \textbf{可能数据字段} & \textbf{验证方式} \\
\midrule
目标条件 & 希望获得的发光颜色或带隙范围 & 目标波长、带隙范围 & 光谱测试、带隙预测、DFT \\
结构条件 & 候选结构是否合理 & 空间群、原子坐标、局域配位 & 结构弛豫、距离检查、声子计算 \\
稳定条件 & 候选是否可能存在 & 形成能、凸包能、热稳定性 & 高精度计算、热分析、实验合成 \\
应用条件 & 是否适合器件环境 & 缺陷、热稳定、外延兼容 & 发光测试、器件测试、可靠性测试 \\
\bottomrule
\end{tabularx}
\end{formalTable}

在使用 AI 辅助生成设计时，提示词应强调“提出候选后的验证链条”，而不是让 AI 直接给出材料名称。示例提示词可这样组织。

\begin{samepage}
\begin{userbox}
我想设计一种目标发光波长在蓝光范围附近的无机发光材料。请帮我把这个目标转化为条件生成设计流程。

要求如下：
\begin{enumerate}
\item 说明目标波长如何转化为带隙或能级相关条件；
\item 说明生成模型可提出哪些类型的候选，例如组成、掺杂或晶体结构；
\item 列出生成后必须检查的化学、结构、稳定性和应用约束；
\item 说明为什么不应只根据带隙判断材料是否适合发光器件；
\item 不要直接虚构某个“最佳材料”。
\end{enumerate}
\end{userbox}

\begin{aibox}
可把任务拆成以下流程：

\begin{enumerate}
\item 目标条件设定：将目标发光波长转化为近似带隙范围，同时记录发光效率、热稳定性和器件兼容性等要求。
\item \textbf{候选生成：}生成模型可提出候选组成、掺杂元素、晶体结构或局域配位环境，但这些结果只是候选，不是最终材料。
\item \textbf{有效性检查：}检查化学式、价态、原子距离、晶格体积和结构格式是否合理，排除明显无意义候选。
\item 稳定性验证：用形成能、凸包能、结构弛豫和必要的声子或热稳定性分析判断候选是否可能存在。
\item \textbf{应用验证：}进一步检查缺陷、非辐射复合、热稳定、外延匹配和器件环境适配性。
\end{enumerate}

带隙合适只能说明候选可能满足发光颜色要求，不能说明它一定具有高发光效率或可制备性。生成模型的作用是扩展候选空间，最终判断仍然需要计算和实验验证。
\end{aibox}
\end{samepage}

条件生成还需要保留“负面条件”。很多时候，研究者不仅知道想要什么，也知道不想要什么。例如，不希望候选含有某些高风险元素，不希望晶体结构出现过短原子距离，不希望配方超出已知可制备范围，不希望候选与训练集中已有材料完全重复。这些负面条件如果不写入生成和筛选流程，模型可能不断提出形式上满足目标性能、但材料上明显不可接受的候选。

具体而言，条件生成的输入可分为三类：目标条件、边界条件和排除条件。目标条件推动模型朝目标性能靠近；边界条件保证候选处于可研究范围；排除条件提前去掉明显不可接受的结果。对于材料科学任务，三类条件同样重要。只写目标条件，容易得到“性能诱人但不可用”的候选；只写排除条件，又可能使搜索空间过窄，失去发现新材料的机会。


\subsection{候选验证与路线组合}

\subsubsection*{\textbf{1. 验证流程：从有效性到应用性}}

生成候选首先要检查有效性：化学式、价态、原子距离、晶格体积和周期性是否合理；其次检查稳定性：结构弛豫后是否保持、形成能和声子是否可接受；再次检查新颖性：候选是否只是训练集或数据库中的重复结构；最后检查可合成性和应用条件。任何一级不通过，都不应把候选直接称为材料发现。生成设计的主要验证维度见表 \ref{tab:generative-design-validation}。


\begin{formalTable}[htbp]
\caption{生成设计的验证维度}
\label{tab:generative-design-validation}
\begin{tabularx}{\textwidth}{@{}YYYY@{}}
\toprule
\textbf{验证指标} & \textbf{问题} & \textbf{常用判断} & \textbf{为什么重要} \\
\midrule
有效性 & 结构是否符合基本规则？ & 价态、距离、格式、晶格范围 & 避免无意义候选 \\
稳定性 & 候选是否可能存在？ & 形成能、凸包能、声子、动力学 & 减少纸上材料 \\
新颖性 & 是否只是训练集复述？ & 结构匹配、组成相似度、数据库检索 & 证明发现价值 \\
可合成性 & 实验能否做出来？ & 合成路线、温度窗口、前驱体、杂相 & 连接真实材料 \\
应用性 & 能否满足场景要求？ & 服役环境、成本、可靠性、规模化 & 决定工程价值 \\
\bottomrule
\end{tabularx}
\end{formalTable}

可靠的验证链条通常分为四层。第一层是格式、几何和重复性检查；第二层是快速正向模型打分；第三层是 DFT、分子动力学或热力学精算；第四层是实验合成、结构表征和性能测试。研究报告还应说明每一级候选的通过率，区分生成数量、有效数量、稳定数量和实验验证数量。


筛选、主动学习和生成设计可串联使用：生成模型扩展候选空间，规则和代理模型快速排除明显不合理候选，主动学习决定哪些候选最值得进入下一轮高精度计算或实验。最终判断仍然来自材料约束和验证证据，而不是模型名称。

可合成性是生成设计最容易被低估的环节。一个结构在计算上稳定，不代表实验上容易合成。实验可达性受到前驱体、温度、压力、气氛、反应路径、动力学障碍和杂相竞争影响。对于玻璃材料，候选配方还必须能熔制、澄清、冷却成玻璃并避免严重析晶；对于晶体材料，候选结构还要考虑相竞争和合成窗口。可合成性判断不足，是“纸上材料”大量出现的重要原因。生成候选在不同证据阶段的表述边界见表 \ref{tab:generated-candidate-evidence}。

\begin{formalTable}[htbp]
\caption{生成候选在不同证据阶段的表述边界}
\label{tab:generated-candidate-evidence}
\footnotesize
\begin{tabularx}{\textwidth}{@{}YYYY@{}}
\toprule
\textbf{候选阶段} & \textbf{可采用的表述} & \textbf{不能怎样表述} & \textbf{需要补充证据} \\
\midrule
模型生成 & 生成了一批候选结构或配方 & 发现了一批新材料 & 格式、价态和几何检查 \\
模型打分 & 获得高优先级候选 & 证明候选性能优异 & 正向模型误差和适用范围 \\
高精度计算 & 候选在计算上接近稳定 & 候选一定可合成 & 结构弛豫、凸包能、相竞争 \\
实验合成 & 候选已被制备并测试 & 可直接工程应用 & 服役、可靠性和规模化验证 \\
\bottomrule
\end{tabularx}
\end{formalTable}

这一部分也适合用 AI 辅助整理验证清单。关键是要求 AI 区分证据层级，避免把模型输出写成最终结论。

\begin{samepage}
\begin{userbox}
我用生成模型得到了一批候选无机晶体结构。请帮我设计一个候选验证清单，判断这些候选能否进入后续计算或实验。

要求如下：
\begin{enumerate}
\item 按有效性、稳定性、新颖性、可合成性和应用性分层；
\item 每一层说明需要检查哪些信息；
\item 区分“模型推荐候选”和“已验证材料”的表述；
\item 提醒如何记录每一级候选数量和通过率；
\item 不要把生成模型输出直接称为材料发现。
\end{enumerate}
\end{userbox}

\begin{aibox}
可建立五层验证清单：

\begin{enumerate}
\item \textbf{有效性：}检查化学式、价态、原子距离、晶格体积、周期性和文件格式是否合理。
\item 稳定性：进行结构弛豫，检查形成能、凸包能、声子稳定性或分子动力学稳定性。
\item \textbf{新颖性：}与训练集、材料数据库和文献记录比对，判断是否为已有材料或高度相似结构。
\item \textbf{可合成性：}分析前驱体、温度窗口、相竞争、动力学路径和实验设备条件。
\item \textbf{应用性：}结合目标场景检查性能、成本、安全、服役环境和长期可靠性。
\end{enumerate}

报告中应分别写出生成数量、有效数量、稳定数量、新颖数量、进入实验数量和实验验证数量。模型生成和模型打分只能说明候选值得进一步检查，不宜直接等同于材料发现。
\end{aibox}
\end{samepage}

在实际写作中，还应避免把验证流程写成一条单向直线。候选在任何一级验证失败，都可能需要返回前面的步骤。例如，结构弛豫后候选发生明显畸变，说明生成结构可能不合理；DFT 显示候选远离稳定区域，说明筛选规则需要加强稳定性约束；实验中反复出现杂相，说明可合成性判断不足。这些失败不是无效信息，而是帮助修正候选生成和筛选规则的重要反馈。

因此，候选验证更像是一个带反馈的流程。生成模型提出候选，验证流程排除或确认候选，失败原因再返回去修改约束、训练数据或生成条件。只有把反馈写清楚，生成设计才不是“模型输出一批材料”，而是“模型、计算和实验共同收敛到更可靠候选”的过程。

\subsubsection*{\textbf{小结}}

本节说明生成设计的作用是扩展候选空间，而不是直接给出最终材料。
条件生成需要同时设置目标条件、边界条件和排除条件，并选择适合材料对象的表示方式。
生成候选必须经过有效性、稳定性、新颖性、可合成性和应用性验证，失败结果也应回流流程。

\subsection*{习题}

\begin{enumerate}[leftmargin=*]
\item 生成候选为什么不能直接称为材料发现？
\item 为发光材料生成任务设计三类验证指标。
\end{enumerate}

\clearpage
\section{可信发现与路线选择}


\begin{sectionkeypoints}
\item 可信材料发现需要区分数据、模型、材料、验证和应用等不同证据层级，不应把模型推荐直接等同于真实材料发现。
\item 方法选择应由材料问题决定。候选库明确时优先筛选，实验昂贵时优先主动学习，候选库不足时再考虑生成模型。
\item 安全、环境、供应链、成本、可加工性和长期可靠性应作为并列约束保留，不能被压缩成无法解释的单一分数。
\item 第\ref{chap:materials-cognition-prediction}章的材料表示、性质预测、不确定性、机器学习势和可解释性，是第\ref{chap:materials-design-discovery}章候选筛选、生成验证和闭环优化的基础。
\item 全模块的完整逻辑是：定义目标、限定边界、提出候选、逐层验证、实验反馈，并在新数据基础上继续更新模型。
\end{sectionkeypoints}

\subsection{可靠评价与方法选择}

人工智能能够扩大材料搜索范围，也会放大数据偏差和模型外推风险。材料数据库中的计算参数、实验条件和样品质量并不完全一致，成功结果通常多于失败结果，热门体系的数据也远多于冷门体系。模型会继承这些差异，因此候选排名并不是客观真理，而是特定数据和模型边界下的判断。


材料性能还高度依赖结构、工艺和服役环境。同一化学式在不同缺陷、微结构和热处理条件下可能表现迥异，如果数据只记录成分和最终性能，模型就会把未记录的工艺差异当成噪声。逆向设计因此必须明确材料对象和数据适用范围，不应把材料简化为一个化学式或单一分数。


“纸上材料”指计算或模型预测较好，但缺少可合成性、规模化或应用验证的候选。真正的材料发现需要区分不同证据层级：模型预测提供线索，高精度计算提高物理可信度，实验合成证明材料可达，器件测试和长期稳定性才说明应用潜力。材料发现结果的分层评价见表 \ref{tab:materials-discovery-evaluation}。


\begin{formalTable}[htbp]
\caption{材料发现结果的分层评价}
\label{tab:materials-discovery-evaluation}
\begin{tabularx}{\textwidth}{@{}YYY@{}}
\toprule
\textbf{评价层面} & \textbf{核心判断} & \textbf{可靠表述示例} \\
\midrule
数据 & 来源、去重、划分是否清楚 & 在某数据库和划分方式下模型表现良好 \\
模型 & 是否有基线、不确定性和边界说明 & 模型推荐了若干高优先级候选 \\
材料 & 是否满足化学、结构和稳定性约束 & 候选在计算上接近稳定且结构合理 \\
验证 & 是否经过 DFT、实验或器件测试 & 某候选已合成并测得目标性质 \\
应用 & 是否考虑成本、安全和可靠性 & 候选具备进一步器件验证价值 \\
\bottomrule
\end{tabularx}
\end{formalTable}

可信评价的关键，是把“模型表现好”和“材料真的好”分开。模型表现好，通常指在某个数据集、某种划分方式和某个评价指标下，预测误差较小或排序效果较好。材料真的好，则意味着候选在真实制备、结构稳定、性能测试和服役环境中都能经受验证。二者之间存在距离。材料人工智能研究中常见的问题，是把模型层面的高分候选直接写成材料发现结果，忽略了验证证据不足。

评价结果时，还要说明模型适用范围。例如，模型是在无机晶体数据库上训练的，还是在某类玻璃配方上训练的；数据是计算数据、实验数据，还是二者混合；测试集划分是随机划分，还是按材料体系、时间或组成空间划分。如果训练集和测试集高度相似，模型指标可能看起来很好，但对真正未知材料的预测能力并不一定强。对于逆向设计来说，外推风险比普通插值预测更突出，因为候选往往位于已有数据边界附近。

不确定性也是可信评价的重要组成部分。一个候选预测性能高，但模型不确定性也高，说明它可能值得探索，但不宜直接排在确定性较高的候选之前。相反，一个候选预测性能略低但不确定性小、约束满足、可合成性强，也可能更适合作为近期实验对象。候选报告中应同时给出预测值、不确定性、约束是否满足和推荐理由。材料发现结果的可靠表述方式见表 \ref{tab:reliable-discovery-expression}。

\begin{formalTable}[htbp]
\caption{材料发现结果的可靠表述方式}
\label{tab:reliable-discovery-expression}
\footnotesize
\begin{tabularx}{\textwidth}{@{}YYYY@{}}
\toprule
\textbf{常见表述} & \textbf{风险} & \textbf{更稳妥的写法} & \textbf{还需补充} \\
\midrule
模型发现了某材料 & 夸大模型作用 & 模型推荐该候选进入验证 & 计算或实验验证 \\
该候选性能优异 & 忽略不确定性和边界 & 在当前模型和数据范围内预测值较高 & 不确定性、误差范围 \\
生成了大量新材料 & 混淆生成与发现 & 生成了若干待验证候选 & 去重、新颖性和稳定性检查 \\
候选可用于器件 & 跳过应用验证 & 候选具备进一步器件验证价值 & 器件测试和长期可靠性 \\
\bottomrule
\end{tabularx}
\end{formalTable}

候选报告还需要特别注意“数量表述”。生成模型可能一次输出几千个候选，但这并不意味着发现了几千种新材料。更准确的写法应区分：生成候选数量、格式有效数量、结构合理数量、计算稳定数量、去重后新颖数量、进入实验数量和实验验证成功数量。只有经过实验合成和性能测试的候选，才能接近“发现”的表述；只经过模型打分的候选，只能称为“推荐候选”或“高优先级候选”。

新颖性检查也属于可信评价的一部分。生成模型可能复述训练集中已经存在的材料，或者只对已有结构做很小改动。如果不做数据库检索和结构匹配，就可能把已有材料误写成新发现。对于晶体材料，可比较化学组成、空间群、晶格参数和结构相似度；对于玻璃或配方材料，可比较组分比例和已有专利、文献或实验记录。新颖性不是单靠模型输出判断的，而需要与数据库和文献记录对照。

\subsubsection*{\textbf{1. 路线选择：问题决定方法}}

方法选择应由问题决定，而不是由模型新旧决定。候选空间已经明确时，高通量筛选通常最直接；实验成本高且每轮样品有限时，主动学习更适合；需要突破已有候选库时，生成模型才体现优势。玻璃形成边界、相稳定、毒性和设备温度等硬约束，应尽可能在搜索前加入，而不是在模型给出结果后再补救。不同问题情形下的方法选择见表 \ref{tab:inverse-design-method-selection}。


\begin{formalTable}[htbp]
\caption{材料逆向设计的方法选择}
\label{tab:inverse-design-method-selection}
\begin{tabularx}{\textwidth}{@{}YYYY@{}}
\toprule
\textbf{问题情形} & \textbf{优先方法} & \textbf{理由} & \textbf{典型材料场景} \\
\midrule
候选库明确 & 高通量筛选 & 流程清楚、可解释、验证成本低 & 从数据库筛选稳定光电材料 \\
实验昂贵且小样本 & 贝叶斯优化/主动学习 & 每一轮实验都要最大化信息量 & 玻璃配方、合金热处理、催化条件 \\
候选库不足 & 生成模型 + 后筛选 & 需要提出库外候选 & 新晶体结构、功能分子、聚合物序列 \\
安全或法规约束强 & 规则约束 + 人工审核 & 先排除不可接受候选 & 含毒元素、稀缺元素或高风险反应 \\
\bottomrule
\end{tabularx}
\end{formalTable}

可靠评价还应考虑安全、环境和供应链。候选若含有有毒、稀缺或高风险元素，即使单项性能突出，也可能不适合真实应用。相应地，性能、成本、环境影响、可加工性和生命周期不应被压缩成一个无法解释的总分，而应作为并列约束保留在候选报告中。

方法选择还应考虑课程或项目阶段。对于入门学习，最合适的方法往往不是最新的生成模型，而是可解释、可复核的筛选流程。因为研究者需要先理解数据字段、模型输入、约束条件和验证证据之间的关系。对于研究生或真实科研项目，如果已有稳定的数据积累和实验反馈机制，再引入主动学习或生成模型才更合理。否则，复杂模型可能只是增加技术难度，并不能提高材料发现质量。

在写作时，可把方法选择写成一组判断问题：候选空间是否已经明确？每次实验成本是否很高？是否需要提出已有数据库之外的新候选？是否存在严格安全或工艺边界？是否有能力进行后续高精度计算和实验验证？这些问题比“选择哪个先进算法”更重要。

这一节中的方法选择和可信评价也可借助 AI 辅助整理，但提示词应要求它给出理由和证据边界，不应只推荐一个模型名称。

\begin{samepage}
\begin{userbox}
我正在写材料逆向设计方案，候选空间包括已有数据库中的无机晶体和少量可枚举组成。实验成本较高，每轮只能验证 5～10 个样品。请帮我判断优先使用高通量筛选、主动学习还是生成模型。

要求如下：
\begin{enumerate}
\item 根据候选空间、实验成本、数据量和验证能力进行判断；
\item 说明每种方法适合和不适合的情况；
\item 给出推荐路线，但必须说明理由；
\item 提醒结果报告中应如何避免夸大模型作用；
\item 不要只因为生成模型较新就优先推荐它。
\end{enumerate}
\end{userbox}

\begin{aibox}
在这个任务中，候选空间已有数据库和可枚举组成，说明不一定需要首先使用生成模型。更合理的路线可以是：

\begin{enumerate}
\item 第一阶段使用高通量筛选：先用规则、正向模型和稳定性判断从已有候选库中排除明显不合理候选。
\item \textbf{第二阶段使用主动学习：}由于每轮实验只能验证 5～10 个样品，应让模型同时考虑预测性能和不确定性，选择高价值实验批次。
\item 生成模型暂时作为补充：只有当已有候选库不能覆盖目标性能区域，或者需要提出库外结构时，再考虑生成模型。
\end{enumerate}

报告中应写成“模型推荐若干高优先级候选，并通过后续计算或实验验证”，而不是“模型发现了最终材料”。如果没有实验或高精度计算支持，应明确说明候选仍处于待验证阶段。
\end{aibox}
\end{samepage}

\subsection{能力复用与闭环迭代}

\subsubsection*{\textbf{1. 两章闭环：正向模型服务逆向设计}}

两章之间形成正反闭环。第\ref{chap:materials-cognition-prediction}章把材料数据转化为描述符或结构图，训练正向模型并解释模型边界；第\ref{chap:materials-design-discovery}章把这些模型作为候选筛选和生成验证的评价器，再用新计算和实验结果更新数据。没有正向预测，逆向搜索缺少可靠评价；没有逆向设计，正向模型容易停留在被动打分。第\ref{chap:materials-cognition-prediction}章能力在第\ref{chap:materials-design-discovery}章中的复用关系见表 \ref{tab:chapter-one-capability-reuse}。


\begin{formalTable}[htbp]
\caption{第\ref{chap:materials-cognition-prediction}章能力在第\ref{chap:materials-design-discovery}章中的复用}
\label{tab:chapter-one-capability-reuse}
\begin{tabularx}{\textwidth}{@{}YYY@{}}
\toprule
\textbf{第\ref{chap:materials-cognition-prediction}章能力} & \textbf{第\ref{chap:materials-design-discovery}章中的用途} & \textbf{材料例子} \\
\midrule
性质预测 & 作为筛选和生成后的快速打分器 & 带隙筛选 microLED 候选 \\
材料表示 & 决定候选能否被模型有效评价 & 晶体图、玻璃成分、分子图 \\
不确定性量化 & 指导主动学习下一步实验 & 高折射率玻璃小样本优化 \\
机器学习势 & 验证候选结构和高温行为 & 超高温陶瓷和高熵合金 \\
可解释性 & 把模型推荐转化为材料规律 & 形成能力、熔点或缺陷趋势分析 \\
\bottomrule
\end{tabularx}
\end{formalTable}

完整逻辑可概括为：目标性质定义问题，材料知识限定边界，筛选、主动学习或生成模型提出候选，正向模型和物理计算逐层评价，实验结果决定最终结论并反馈到数据集。人工智能的价值在于提高搜索与决策效率，材料科学的责任则是保证问题、约束和证据真实可靠。

闭环迭代要落实到可操作的数据管理上。每一轮推荐都应记录模型版本、训练数据版本、候选来源、预测值、不确定性、约束检查结果和推荐理由；每一轮验证后，还要把成功样品、失败样品、测试条件和失败原因按同一格式回写数据表。这样，下一轮模型更新才知道哪些变化来自材料本身，哪些变化来自数据版本、实验批次或测试条件。

例如，在高折射率玻璃任务中，第\ref{chap:materials-cognition-prediction}章中讨论的数据整理和描述符构建可帮助建立折射率、Tg 和热膨胀系数预测模型；第\ref{chap:materials-design-discovery}章则利用这些模型筛选候选配方，并通过主动学习选择下一轮熔制实验。在高熵合金涂层任务中，第\ref{chap:materials-cognition-prediction}章的硬度预测、图神经网络和机器学习势可提供性能估计与结构理解；第\ref{chap:materials-design-discovery}章则进一步用这些模型决定哪些成分或硬质相值得验证。这样的闭环，才是人工智能真正进入材料发现流程的方式。

闭环迭代还要求新数据及时回流。实验失败、性能一般、模型预测错误的样品，都应被记录下来。只把成功样品加入数据库，会使模型越来越偏向成功区域，无法学习失败边界；只记录最终性能而不记录工艺条件，也会使模型无法区分成分效应和工艺效应。闭环发现不仅是“模型推荐--实验验证”，还包括“实验记录--数据更新--模型修正”。材料智能发现闭环中的数据回流内容见表 \ref{tab:closed-loop-data-feedback}。

\begin{formalTable}[htbp]
\caption{材料智能发现闭环中的数据回流}
\label{tab:closed-loop-data-feedback}
\footnotesize
\begin{tabularx}{\textwidth}{@{}YYYY@{}}
\toprule
\textbf{闭环环节} & \textbf{主要任务} & \textbf{需要记录的信息} & \textbf{常见问题} \\
\midrule
模型推荐 & 根据目标和约束提出候选 & 预测值、不确定性、推荐理由 & 只给排名，不说明边界 \\
计算验证 & 用更高精度方法复核候选 & 结构、能量、稳定性、失败原因 & 只保留通过候选 \\
实验验证 & 合成、表征和性能测试 & 工艺、样品状态、测试条件 & 缺少失败样品记录 \\
数据回流 & 更新数据集并重训模型 & 新数据、元数据、版本信息 & 数据来源不可追溯 \\
\bottomrule
\end{tabularx}
\end{formalTable}

AI 可辅助整理闭环报告结构，尤其是检查是否遗漏数据来源、验证证据和失败记录。但同样需要明确：AI 只能帮助组织报告，不宜替代实验和模型审查。

\begin{samepage}
\begin{userbox}
我想把一个材料智能发现项目写成闭环流程报告。项目包括正向模型预测、候选筛选、实验验证和数据回流。请帮我设计报告结构。

要求如下：
\begin{enumerate}
\item 说明第\ref{chap:materials-cognition-prediction}章中的材料表示、性能预测和模型解释如何服务第\ref{chap:materials-design-discovery}章的逆向设计；
\item 按模型推荐、计算验证、实验验证和数据回流四个环节组织；
\item 提醒每个环节需要记录哪些信息；
\item 特别说明失败样品和预测错误样品为什么也要记录；
\item 不要把闭环写成模型自动发现材料。
\end{enumerate}
\end{userbox}

\begin{aibox}
闭环流程报告可包括四部分：

\begin{enumerate}
\item \textbf{模型推荐：}说明使用了哪些材料表示和正向预测模型，给出候选的预测值、不确定性和推荐理由。
\item \textbf{计算验证：}记录结构弛豫、稳定性判断、形成能或其他高精度计算结果，并保留未通过候选的原因。
\item \textbf{实验验证：}记录制备工艺、样品状态、表征结果、性能测试条件和失败现象，不应只记录成功样品。
\item \textbf{数据回流：}说明新数据如何加入数据库，如何更新训练集，模型是否重新训练，以及新模型与旧模型相比是否改进。
\end{enumerate}

失败样品和预测错误样品同样重要，因为它们可帮助模型学习不可制备区域、实验噪声和模型外推边界。材料智能发现不是模型自动给出答案，而是模型、计算、实验和数据更新共同形成的迭代过程。
\end{aibox}
\end{samepage}

如果把两章放在一个完整模块中理解，第\ref{chap:materials-cognition-prediction}章更像是“建立评价能力”，第\ref{chap:materials-design-discovery}章更像是“利用评价能力做决策”。没有第\ref{chap:materials-cognition-prediction}章的数据整理、表示学习和模型评估，第\ref{chap:materials-design-discovery}章的筛选和生成就缺少可靠判断依据；没有第\ref{chap:materials-design-discovery}章的候选选择和实验反馈，第\ref{chap:materials-cognition-prediction}章的模型又容易停留在已有数据上的被动预测。二者结合起来，才能形成从材料认知到材料发现的完整教学闭环。

最后一节不只是总结第\ref{chap:materials-design-discovery}章，也是在说明整个模块的逻辑边界：人工智能可帮助研究者更快地整理材料数据、构建候选、预测性能、安排验证和更新模型，但不应跳过材料科学中的基本问题。材料是否真实存在，能否制备出来，是否满足应用场景，仍然要靠实验、计算和工程约束共同判断。

\subsubsection*{\textbf{小结}}

本节强调可信材料发现需要区分模型推荐、计算证据、实验验证和应用潜力。
方法选择应由问题决定，候选库明确时优先筛选，实验昂贵时优先主动学习，候选不足时再考虑生成。
第\ref{chap:materials-cognition-prediction}章建立评价能力，第\ref{chap:materials-design-discovery}章利用评价能力做决策，数据回流使整个模块形成闭环。

\subsection*{习题}

\begin{enumerate}[leftmargin=*]
\item 数据回流为什么是材料智能发现闭环的关键？
\item 如何根据候选库、实验成本和数据量选择材料发现路线？
\end{enumerate}

\section*{本章小结}

本章讨论了材料设计与智能发现的基本思路。逆向设计首先要求将应用需求转化为可计算目标、变量和约束；候选筛选通过规则、代理模型、高精度计算和实验验证逐层缩小范围；主动学习利用预测值和不确定性安排下一批实验或计算任务；生成设计能够提出新的结构或配方，但必须接受有效性、稳定性、新颖性、可合成性和应用条件等多层验证。全章强调，智能发现不是单次模型输出，而是目标定义、候选构建、分层验证和数据回流共同组成的闭环过程。

% 文件中不加载任何宏包，所有样式均调用项目现有的 package.tex。
\part*{微模块5\quad 人工智能在材料科学中的应用}
\addcontentsline{toc}{part}{微模块5\texorpdfstring{\quad}{ }人工智能在材料科学中的应用}

\chapter{材料认知与性能预测}
\label{chap:materials-cognition-prediction}


理解材料性能的来源，是材料科学研究的核心任务。材料的力学、热学、电学、磁学和光学性能，往往不是由单一因素决定的，而是受到成分组成、微观结构和制备工艺的共同影响。随着研究对象从简单合金、陶瓷和半导体扩展到高熵合金、复合材料、低维材料和复杂功能材料，材料候选空间迅速增大，单纯依靠实验试错已经难以系统揭示其中的规律。在这一背景下，人工智能方法的价值，主要体现在对分散实验数据、计算数据和文献知识的组织与学习上，使模型能够从中提取“成分--结构--工艺--性能”之间的关联，从而辅助材料性能预测和后续实验设计。

本章围绕“材料认知与性能预测”展开，重点解决四个问题。第一节说明材料人工智能为什么需要可靠数据基础，以及实验、计算、文献和数据库数据应如何整理。第二节讨论材料成分和晶体结构如何转化为描述符或图表示，使模型能够读取材料信息。第三节介绍描述符模型、图神经网络和机器学习原子间势（下文简称“机器学习势”）如何服务硬度预测、性质预测和熔点模拟。第四节进一步说明模型解释与物理知识嵌入的作用，帮助研究者判断模型结果是否符合材料机理和实验事实。通过本章学习，研究者需要建立一个基本认识：人工智能并不是替代材料理论和实验验证，而是与数据、计算和材料机理共同构成辅助材料认知的新工具。




\clearpage
\section{范式变革与数据基础}


\begin{sectionkeypoints}
\item 材料人工智能中的数据并不只是单个性能数值，而是包含材料组成、结构信息、制备工艺、测试条件、数据来源和目标性能的完整记录。
\item 材料数据主要来自实验、计算、文献和数据库四类渠道。不同来源的数据各有优势，也各有局限，使用时需要结合具体研究任务进行判断。
\item 材料数据库为模型训练、性能预测和候选筛选提供了重要的数据基础，但数据库中的数据不宜直接等同于最终模型输入，还需要经过字段筛选、来源核验和格式整理。
\item 在正式建模之前，应先整理原始数据清单，明确记录材料是什么、结构是什么、性能数值是多少、数据来自哪里，以及采用了什么计算方法或实验条件。
\item AI 可辅助解释数据库字段、设计表格框架和检查数据规范性，但不宜替代数据库原始记录、实验记录和论文原文核验。材料数据整理的核心要求是清楚、可追溯、可检查。
\end{sectionkeypoints}

\subsection{实验试错与数据驱动}

材料科学的核心任务之一，是理解和建立材料成分、结构、工艺与性能之间的关联。研究者希望理解：为什么某种成分的合金具有更高强度？为什么某种陶瓷能够承受更高温度？为什么某种半导体材料具有合适的带隙并能用于发光器件？这些问题表面上属于不同材料体系，但本质上都指向同一个核心：材料的内部组成与结构如何决定其外部性能。


在传统材料研发中，研究者通常依靠已有理论、文献经验和实验试错来提出材料方案，再通过制备、表征和性能测试不断调整成分与工艺。例如，为了提高合金强度，可尝试加入不同合金元素；为了改善陶瓷耐高温性能，可调节烧结温度、晶粒尺寸或相组成；为了优化光电材料性能，可改变材料的元素组成和晶体结构。这种方式直接可靠，也符合材料实验研究的基本规律。它的优势在于结论来自真实样品和真实测试条件，缺点是往往需要大量重复实验，研发周期较长，成本也较高。


当材料体系从简单单组元或二元材料扩展到多组元合金、复合材料、异质结构、低维材料和复杂光电材料时，实验试错面对的组合空间会迅速扩大。以高熵合金为例，四种、五种甚至更多主要元素的组合、比例和热处理工艺都可能改变组织结构和性能；同时，强度、硬度、韧性、耐腐蚀性、导电性或发光性能又受到晶体结构、缺陷、晶粒尺寸、相组成、界面状态和服役环境共同影响。传统经验仍然重要，但单纯依靠人工逐一尝试，已经难以在有限时间内完成系统筛选。


在这一背景下，材料研发逐渐从单纯依赖实验试错，转向实验、计算、数据和人工智能相结合的新方式。实验能够提供真实可靠的性能结果，计算能够在实验之前预测部分结构和能量信息，数据库能够积累已有材料的成分、结构和性能数据，而人工智能则能够从这些数据中学习复杂的材料规律。换言之，人工智能并不是脱离实验和理论单独发挥作用，而是建立在已有材料知识、实验数据和计算数据基础上的辅助工具。


数据驱动并不是否定实验和理论，而是把已有实验数据、计算数据、文献数据和数据库资源转化为可分析、可学习和可预测的信息。人工智能模型可从大量材料数据中学习成分、结构与性能之间的关系，并用于预测未知材料的性质。这样，研究者不再需要在庞大的候选空间中完全盲目试错，而可借助数据和模型缩小研究范围，把实验资源集中到更有希望的材料体系上。


\subsubsection*{\textbf{1. 人工智能用于材料科学的基本流程}}

在人工智能用于材料科学的过程中，研究者通常需要依次解决几个关键环节。第一个环节是材料数据整理。需要收集和规范实验数据、计算数据、文献数据和数据库资源，并记录材料的成分、结构、工艺、性能以及数据来源等信息。只有数据来源清楚、含义明确，后续模型学习才具有可靠基础。

第二个环节是材料表示。机器学习模型不宜直接理解“合金”“晶体结构”或“热处理工艺”等材料概念，因此需要把材料转化为机器可处理的形式。例如，可用描述符表示化学成分和结构信息，也可用图结构表示原子之间的连接关系。

第三个环节是模型预测。在材料被合理表示之后，可建立机器学习模型，预测材料的形成能、带隙、弹性模量、熔点、硬度等性能。对于涉及原子运动和微观演化的问题，还可进一步构建机器学习势（利用机器学习模型学习原子结构与能量、力之间关系的势函数），用于模拟扩散、相变和熔化等原子尺度行为。

最后，模型预测结果还需要回到材料科学问题本身。通过可解释性方法，可分析模型为什么做出某种预测，识别哪些成分、结构或工艺因素对性能影响更大，从而把模型结果进一步转化为可理解的材料规律。


在材料研究中还必须明确，人工智能不宜替代材料专业知识、实验验证和物理判断。人工智能的作用更像是一种增强手段：它可帮助研究者整理数据、发现规律、提出候选、预测性能和辅助决策，但最终结论仍然需要通过材料理论分析和实验结果进行验证。尤其是在材料科学中，模型给出的预测结果并不等同于真实材料性能。材料是否能够被实际制备出来，结构是否稳定，以及能否在实际服役环境中长期可靠工作，仍然需要进一步验证。


围绕材料性能预测和材料设计两类问题，人工智能用于材料科学可概括为两条基本主线。第一条是正向预测，即在材料成分、结构或实验条件已经给定的情况下，预测材料可能具有的性能和行为。例如，已知一种晶体材料的结构，可预测它的形成能、带隙、弹性模量或熔点。第二条是逆向设计，即先给出目标性能，再反过来寻找可能满足要求的材料。从概念上看，正向预测关注已知材料可能表现出什么性能，逆向设计关注为了获得目标性能应当寻找什么材料。


本章主要讨论第一条主线，即材料科学中的正向预测。本章将从材料数据出发，进一步讨论材料如何被表示成机器可理解的形式，机器学习如何预测材料性能，机器学习势如何模拟原子尺度行为，以及可解释人工智能如何帮助研究者从模型结果中提炼材料规律。


\subsection{数据来源与材料数据库}

人工智能要预测材料性能，首先需要数据。对于材料科学而言，数据并不只是一个性能数值，而是对材料组成、制备过程、测试条件和性能结果的完整记录。一个可靠的材料数据条目，至少应该说明材料的化学组成、结构信息、制备工艺、测试条件、数据来源和目标性能。缺少这些基本信息，模型就难以建立清晰的输入与输出对应关系，所得预测结果也难以作为可靠依据。


材料数据主要来自实验、计算、文献和数据库四类渠道。


第一类是实验数据。实验数据来自材料制备、显微表征和性能测试过程，例如硬度、拉伸强度、腐蚀电流密度、熔点、热导率、带隙、发光波长等。实验数据最接近真实材料和真实服役环境，但获取成本较高，也容易受到制备工艺、测试条件和仪器参数的影响。例如，同一种合金在不同热处理温度下可能得到不同硬度，同一种涂层在不同腐蚀介质中可能表现出不同耐蚀性能。因此，实验数据不应只记录结果，还要尽量记录测试条件。


第二类是计算数据。计算数据主要来自第一性原理计算、分子动力学模拟、相场模拟等方法。例如，第一性原理计算可得到形成能、能带结构、弹性常数、态密度、晶体结构稳定性等信息；分子动力学可模拟扩散、相变和高温行为。计算数据的优点是容易标准化，也适合大规模生成；但计算数据依赖具体计算方法和参数，不宜简单等同于实验结果。


第三类是文献数据。大量材料知识存在于论文、专利、实验报告和技术资料中。例如，一篇论文可能同时记录材料成分、制备工艺、热处理条件、显微组织和性能结果。文献数据的优点是信息丰富，缺点是格式不统一，很多信息隐藏在正文、表格、图片或补充材料中，需要人工整理或借助文本处理工具提取。


第四类是材料数据库。材料数据库可理解为经过整理的材料信息资源库，它把大量材料的化学组成、晶体结构、计算性质、实验性质和文献信息集中到一个可检索的平台中。对于材料人工智能而言，数据库是重要数据来源，因为它可为模型训练、性能预测和材料筛选提供相对规范的数据基础。


常见材料数据库包括 Materials Project、OQMD\footnote{Open Quantum Materials Database，开放量子材料数据库。}、AFLOW\footnote{Automatic Flow Framework for Materials Discovery，自动化材料发现计算框架。}、NOMAD\footnote{Novel Materials Discovery，材料数据管理与共享平台。}、ICSD\footnote{Inorganic Crystal Structure Database，无机晶体结构数据库。}、JARVIS\footnote{Joint Automated Repository for Various Integrated Simulations，材料计算数据库与工具平台。} 等。不同数据库的侧重点并不完全相同：有的偏向无机晶体结构和稳定性数据，有的偏向大规模热力学计算结果，有的强调计算数据管理和共享，还有的主要提供实验晶体结构信息。教材中列举数据库名称，目的不是要求逐一掌握平台功能，而是说明不同数据来源对应不同研究任务。

对于入门阶段的研究者而言，不必一开始深入掌握每个数据库的所有功能，但需要知道不同数据库主要提供哪类材料数据，以及这些数据适合用于哪些研究任务。只有明确数据来源和适用范围，才能更合理地将数据库用于材料性能预测、候选筛选和模型训练。


使用材料数据库时，首先要明确自己的研究任务。不同任务需要的数据字段不同。例如，如果任务是预测无机光电材料的带隙，就需要关注化学式、晶体结构、空间群、带隙、形成能、数据来源和计算方法等字段；如果任务是预测合金硬度，就需要关注成分、相结构、制备工艺、热处理条件、硬度值和测试方法等字段；如果任务是筛选超高温陶瓷，就需要关注化学组成、晶体结构、熔点、密度、热稳定性和抗氧化相关数据。


因此，数据库只是提供了材料信息的来源，真正进入模型训练之前，还需要把这些信息整理成结构清楚、字段明确、来源可追溯的原始数据清单。下面以无机光电材料带隙数据整理为例，说明从数据库原始记录到可建模数据清单的具体做法。

\subsubsection*{\textbf{1. 案例示范：无机光电材料带隙数据整理}}

以“无机光电材料带隙预测”为例，数据准备的第一步不是训练模型，而是从数据库或文献中整理出一张原始数据清单。这个清单应该回答几个基本问题：材料是什么？结构是什么？带隙是多少？这个带隙来自哪里？是实验值还是计算值？如果是计算值，采用了什么计算方法？如果是实验值，是否记录了测试条件？


一张初步整理后的原始数据清单可按表 \ref{tab:raw-bandgap-data} 设计。表中的 eV\footnote{Electron volt，电子伏特，常用于表示电子能量和带隙大小。} 是带隙常用单位。


\begin{formalTable}[htbp]
\caption{原始带隙数据清单示例}
\label{tab:raw-bandgap-data}
\scriptsize
\setlength{\tabcolsep}{2.8pt}
\begin{tabularx}{\textwidth}{@{}YYYYYYYYY@{}}
\toprule
\textbf{材料编号} & \textbf{化学式} & \textbf{晶体结构} & \textbf{空间群} & \textbf{带隙/eV} & \textbf{数据类型} & \textbf{数据来源} & \textbf{方法或条件} & \textbf{备注} \\
\midrule
001 & GaN & 纤锌矿 & P6\textsubscript{3}mc & 3.40 & 实验 & 文献 & 室温测试 & 直接带隙 \\
002 & ZnO & 纤锌矿 & P6\textsubscript{3}mc & 3.37 & 实验 & 文献 & 室温测试 & 直接带隙 \\
003 & SrTiO\textsubscript{3} & 钙钛矿 & Pm-3m & 3.20 & 数据库 & 数据库记录 & 需注明方法 & 需核对来源 \\
\bottomrule
\end{tabularx}
\end{formalTable}

表 \ref{tab:raw-bandgap-data} 的作用是保存原始材料信息，而不是直接代表最终模型输入。它的重点在于把分散的信息整理成清楚、可追溯、可检查的记录。后续是否把化学式转化为成分特征，是否把晶体结构转化为结构描述符，属于下一节要讨论的材料表示问题。


在整理这类原始数据清单时，可借助 AI 辅助完成部分工作。此处的重点不是让 AI 代替研究者判断数据是否可靠，而是让 AI 帮助理解字段含义、设计表格框架和检查数据规范性。以下以无机光电材料带隙数据整理为例给出提示词示范。

\begin{samepage}
\begin{userbox}
我正在整理无机光电材料的带隙数据，准备用于后续机器学习建模。数据可能来自 Materials Project、JARVIS、文献或实验记录。请帮我设计一张原始数据清单的表头，并说明每一列的作用。

要求如下：
\begin{enumerate}
\item 表格要保留材料编号、化学式、晶体结构、空间群、带隙数值、带隙单位、数据类型、数据来源、计算方法或实验条件等信息；
\item 区分实验值和计算值；
\item 提醒哪些字段是后续建模前必须核验的；
\item 不要虚构具体材料数据，只需要给出表头设计和填写说明。
\end{enumerate}
\end{userbox}

\begin{aibox}
可将原始数据清单设计为以下字段：

\begin{itemize}
\item 材料编号：用于唯一标识每一条材料记录，便于后续追踪和检查。
\item 化学式：记录材料的元素组成，例如 GaN、ZnO、SrTiO$_3$ 等。
\item \textbf{晶体结构：}记录材料所属结构类型，例如纤锌矿、钙钛矿等。
\item \textbf{空间群：}记录晶体结构的对称性信息。
\item 带隙数值：记录材料的带隙大小。
\item 带隙单位：通常为 eV，避免不同来源数据混用时产生单位错误。
\item \textbf{数据类型：}说明该带隙是实验值还是计算值。
\item \textbf{数据来源：}记录数据来自数据库、论文、实验报告还是其他资料。
\item \textbf{方法或条件：}如果是计算值，应注明计算方法；如果是实验值，应注明测试温度、测试方法等条件。
\item \textbf{备注：}记录是否为直接带隙、是否需要进一步核验等补充信息。
\end{itemize}

在正式建模前，尤其需要核验带隙数值、数据类型、数据来源和方法或条件等字段。同一种材料在不同晶体结构、不同计算方法或不同测试条件下可能具有不同带隙，因此这些信息不应省略。
\end{aibox}
\end{samepage}

上面的提示词示范说明，AI 可帮助研究者快速形成数据表框架，但 AI 输出并不宜直接作为最终数据结果。研究者还需要回到数据库原始页面或论文原文，核对具体数值、计算方法、实验条件和材料结构是否准确。

因此，在材料数据整理阶段，AI 更适合充当数据整理助手，而不是数据来源本身。它可帮助设计表格、解释字段和生成检查规则，但具体数据是否准确，仍然需要根据数据库原始记录、实验记录或论文原文进行确认。尤其在材料科学中，同一种化学式可能对应不同晶体结构，不同计算方法也可能得到不同带隙值，因此数据来源、计算方法和实验条件必须保留。


为了提高数据质量，原始数据整理时应注意几个问题。第一，要统一单位。例如带隙通常用 eV，硬度可能用 HV\footnote{Vickers Hardness，维氏硬度单位。} 或 GPa\footnote{Gigapascal，吉帕，常用于表示压力、应力或弹性模量。}，能量可能用 eV/atom 或 kJ/mol。第二，要区分实验数据和计算数据。二者可同时记录，但不能在没有说明的情况下混用。第三，要保留数据来源。每一条数据宜能追溯到数据库编号、论文或实验记录。第四，要检查缺失值和重复值。同一个材料如果重复出现，需要判断它们是否对应不同结构、不同测试条件或不同计算方法。第五，要记录必要条件。对于实验数据，应尽量记录温度、压力、测试环境和样品状态；对于计算数据，应记录计算方法和参数来源。

\subsubsection*{\textbf{小结}}

本节说明材料研发从实验试错走向数据、计算和实验协同的基本背景。
材料数据应包含成分、结构、工艺、性能、来源和测试条件，而不只是单一性能数值。
数据库和 AI 可提高整理效率，但数据核验、条件记录和材料判断仍需由研究者完成。

\subsection*{习题}

\begin{enumerate}[leftmargin=*]
\item 材料数据为什么不能只记录性能值？还应补充哪些字段？
\item 比较四类材料数据来源，并说明选择依据。
\end{enumerate}

\clearpage
\section{特征构建与结构表示}


\begin{sectionkeypoints}
\item 材料特征构建的目的，是把化学式、晶体结构、空间群和性能数值等原始记录转化为机器学习模型能够读取的数值特征。
\item 成分描述符可由元素组成和元素属性计算得到，例如平均原子序数、平均电负性、电负性差和平均原子半径等。
\item 结构描述符用于补充原子排列信息，常见字段包括空间群编号、晶胞体积、密度和晶格常数等。
\item 图表示把晶体看成原子节点和原子间边组成的结构，可更直接地保留局域配位、原子间距离和周期性信息。
\item 描述符和图表示都需要材料判断。特征不是越多越好，应避免目标泄漏，并检查所选特征是否具有明确的材料意义。
\end{sectionkeypoints}

\subsection{成分与结构描述符}

上一节介绍了材料数据的来源和材料数据库的作用。通过实验、计算、文献或数据库，研究者可获得材料的化学式、晶体结构、空间群、性能数值和数据来源等信息。但是，这些信息只是原始材料记录，还不宜直接被大多数机器学习模型使用。机器学习模型真正能够处理的，通常是一组数值特征。因此，在进行性能预测之前，需要把材料信息进一步转化为模型可读取的数值表达。


这一过程称为材料的特征化，也常称为描述符构建。描述符指用数字表达材料中可能影响性能的关键信息。描述符不是随意选择的一组数字，而应该尽量反映材料的成分、结构和物理化学特征。例如，对于光电材料带隙预测，元素电负性、原子半径、晶体结构、空间群和晶胞参数等信息都可能与电子结构有关，因此可作为候选描述符。


\subsubsection*{\textbf{1. 案例示范：无机光电材料特征构建}}

本节围绕一个任务展开：无机光电材料带隙预测。假设研究者已经从材料数据库或文献中整理出表 \ref{tab:bandgap-raw-records} 所示的原始数据清单。


\begin{formalTable}[htbp]
\caption{无机光电材料原始记录示例}
\label{tab:bandgap-raw-records}
\footnotesize
\begin{tabularx}{\textwidth}{@{}YYYYY@{}}
\toprule
\textbf{材料} & \textbf{化学式} & \textbf{晶体结构} & \textbf{空间群} & \textbf{带隙/eV} \\
\midrule
氮化镓 & GaN & 纤锌矿 & P6\textsubscript{3}mc & 3.40 \\
氧化锌 & ZnO & 纤锌矿 & P6\textsubscript{3}mc & 3.37 \\
钛酸锶 & SrTiO\textsubscript{3} & 钙钛矿 & Pm-3m & 3.20 \\
\bottomrule
\end{tabularx}
\end{formalTable}

这个表格适合人阅读，但模型不宜直接理解“GaN”“纤锌矿”或“P6\textsubscript{3}mc”这些符号和文字。因此，第一步是把化学式转化为元素组成。


以 GaN 为例，它由 Ga 和 N 两种元素组成，原子比例为 1:1。为了便于模型处理，可把这个比例归一化为原子分数，即 Ga 为 0.50，N 为 0.50。ZnO 同理，Zn 为 0.50，O 为 0.50。SrTiO\textsubscript{3} 中 Sr、Ti 和 O 的比例为 1:1:3，归一化后分别为 0.20、0.20 和 0.60。


经过这一步，原始化学式可转化为表 \ref{tab:formula-composition-example} 所示形式。


\begin{formalTable}[htbp]
\caption{化学式到元素组成的转换示例}
\label{tab:formula-composition-example}
\begin{tabularx}{\textwidth}{@{}YYY@{}}
\toprule
\textbf{化学式} & \textbf{元素组成} & \textbf{原子分数} \\
\midrule
GaN & Ga, N & Ga: 0.50；N: 0.50 \\
ZnO & Zn, O & Zn: 0.50；O: 0.50 \\
SrTiO\textsubscript{3} & Sr, Ti, O & Sr: 0.20；Ti: 0.20；O: 0.60 \\
\bottomrule
\end{tabularx}
\end{formalTable}

这一步表面上简单，但当数据表中有几百或几千个化学式时，人工逐个拆解就很低效。实际建模时，通常需要用程序批量解析化学式，并检查程序是否能够正确处理下标、括号、掺杂和复杂化学式。只有化学式解析正确，后续成分描述符才有意义。


得到元素组成后，第二步是引入元素属性。机器学习模型并不知道 Ga、N、Zn、O、Sr、Ti 这些元素各自有什么物理化学特征，因此需要查找或建立元素属性表。这个属性表可包括原子序数、原子半径、电负性、价电子数、元素熔点等信息。


例如，GaN 的两个组成元素是 Ga 和 N。如果需要计算平均电负性，就需要知道 Ga 和 N 各自的电负性，再按照原子分数加权平均。一般来说，若第 i 种元素的原子分数为 xi，对应元素属性为 Pi，则该材料的平均属性可表示为：


\begin{equation*}\text{平均属性}=\sum_i x_i P_i\end{equation*}

这个公式可用于计算平均原子序数、平均电负性、平均原子半径、平均价电子数等特征。除了平均值，还可计算最大值、最小值、极差和方差。例如，电负性差可用组成元素中最大电负性与最小电负性的差值表示，用来描述元素之间化学性质差异。


对于带隙预测来说，电负性和电负性差具有直观意义。带隙与材料的电子结构和成键特征有关，而元素之间电负性差异会影响电子分布和键合方式。因此，平均电负性、电负性差、平均原子半径等都可作为初步描述符。当然，这些描述符并不能完全决定带隙，但它们为模型提供了可学习的材料信息。


经过这一步，数据表可进一步转化为表 \ref{tab:composition-descriptor-example} 所示的成分描述符表。


\begin{formalTable}[htbp]
\caption{成分描述符表构建示例}
\label{tab:composition-descriptor-example}
\footnotesize
\setlength{\tabcolsep}{2.8pt}
\begin{tabularx}{\textwidth}{@{}YYYYYY@{}}
\toprule
\textbf{化学式} & \textbf{平均原子序数} & \textbf{平均电负性} & \textbf{电负性差} & \textbf{平均原子半径} & \textbf{带隙/eV} \\
\midrule
GaN & 数值 & 数值 & 数值 & 数值 & 3.40 \\
ZnO & 数值 & 数值 & 数值 & 数值 & 3.37 \\
SrTiO\textsubscript{3} & 数值 & 数值 & 数值 & 数值 & 3.20 \\
\bottomrule
\end{tabularx}
\end{formalTable}

此处的“数值”需要根据元素属性表计算得到。实际操作时，可先准备一个元素属性 CSV\footnote{Comma-Separated Values，逗号分隔值文件。} 文件，其中包含元素符号、原子序数、电负性、原子半径等字段。然后将化学式解析结果与元素属性表关联起来，自动计算每个材料的平均值和差值。这样可减少重复计算，但计算结果仍然需要通过抽样检查来确认。


第三步，是加入结构描述符。仅有成分描述符还不够，因为材料性能不仅取决于由哪些元素组成，还取决于这些原子如何排列。GaN 和 ZnO 都是纤锌矿结构，而 SrTiO\textsubscript{3} 是钙钛矿结构。不同晶体结构对应不同的原子排列、局域环境和电子能带特征，因此结构信息对于带隙预测也很重要。


对于入门阶段的研究者，可先使用比较容易获得的结构描述符，例如空间群编号、晶胞体积、密度、晶格常数等。以空间群为例，P6\textsubscript{3}mc 可对应空间群编号 186，Pm-3m 可对应空间群编号 221。这样，文字形式的空间群就可转化为数值字段。晶胞体积和密度则可从晶体结构文件或数据库记录中获得。


加入结构描述符后，表格可进一步扩展为表 \ref{tab:structure-descriptor-example}。


\begin{formalTable}[htbp]
\caption{成分与结构描述符表构建示例}
\label{tab:structure-descriptor-example}
\scriptsize
\setlength{\tabcolsep}{2.8pt}
\begin{tabularx}{\textwidth}{@{}YYYYYYYY@{}}
\toprule
\textbf{化学式} & \textbf{平均原子序数} & \textbf{平均电负性} & \textbf{电负性差} & \textbf{平均原子半径} & \textbf{空间群编号} & \textbf{晶胞体积} & \textbf{带隙/eV} \\
\midrule
GaN & 数值 & 数值 & 数值 & 数值 & 186 & 数值 & 3.40 \\
ZnO & 数值 & 数值 & 数值 & 数值 & 186 & 数值 & 3.37 \\
SrTiO\textsubscript{3} & 数值 & 数值 & 数值 & 数值 & 221 & 数值 & 3.20 \\
\bottomrule
\end{tabularx}
\end{formalTable}

在这个表中，平均原子序数、平均电负性、电负性差、平均原子半径、空间群编号和晶胞体积等列就是模型输入特征，带隙/eV 是模型要预测的目标值。这样，原始的材料记录就被转化为了机器学习可处理的特征表。


该流程可概括为下列步骤，并可结合图 \ref{fig:materials-data-feature} 理解：


\begin{lightquote}
\centering 原始材料记录 → 化学式解析 → 元素属性计算 → 结构字段数值化 → 特征表
\end{lightquote}

\begin{figure}[H]
\centering
\BookFigureImage{images_11/1.png}
\caption{材料数据整理与特征构建示意图}
\label{fig:materials-data-feature}
\end{figure}
\FloatBarrier

在这一环节，如果借助 AI 辅助处理，应像数据整理阶段一样把任务、已有字段、期望输出和限制条件说明清楚。可按下面方式提出任务。

\begin{samepage}
\begin{userbox}
我已经整理了一张无机光电材料带隙数据表，字段包括化学式、晶体结构、空间群、带隙/eV。现在准备用于机器学习建模，请帮我设计从原始表格到特征表的处理流程。

要求如下：
\begin{enumerate}
\item 说明如何把化学式解析为元素组成和原子分数；
\item 说明如何利用元素属性表计算平均原子序数、平均电负性、电负性差和平均原子半径；
\item 指出晶体结构、空间群等文字字段哪些需要进一步数值化；
\item 提醒哪些字段不宜直接作为模型输入，以及如何避免目标泄漏；
\item 不要虚构具体数值，只给出处理步骤和检查要点。
\end{enumerate}
\end{userbox}

\begin{aibox}
可把处理流程分为五步：

\begin{enumerate}
\item 化学式解析：将 GaN、ZnO、SrTiO\textsubscript{3} 等化学式拆分为元素种类和原子比例，再归一化为原子分数。
\item \textbf{元素属性匹配：}准备元素属性表，至少包含元素符号、原子序数、电负性、原子半径等字段，并与解析后的元素组成对应。
\item \textbf{成分描述符计算：}按原子分数加权计算平均原子序数、平均电负性、平均原子半径，并计算电负性差等差异类特征。
\item \textbf{结构字段转换：}晶体结构名称和空间群符号不宜直接输入多数模型，应转换为空间群编号、晶胞体积、密度或其他可计算结构描述符。
\item \textbf{建模前检查：}带隙/eV 是预测目标，不能作为输入特征；发光颜色等由带隙直接决定的字段也要谨慎使用，避免目标泄漏。
\end{enumerate}

最终得到的特征表应做到字段清楚、单位一致、输入特征和预测目标分离，并保留原始数据来源，便于后续核验。
\end{aibox}
\end{samepage}

在实践中，描述符不是越多越好。如果数据样本较少，而输入特征过多，模型可能会过拟合。对于入门任务，可先选择少量具有明确材料意义、且容易从数据库获得的描述符，例如平均原子序数、平均电负性、电负性差、平均原子半径、空间群编号和晶胞体积。等研究者理解基本流程后，再逐步加入更复杂的局域环境描述符。


还要避免目标泄漏。目标泄漏指的是把与预测目标直接相关、甚至由目标结果推导出来的信息放入输入特征中。例如，如果目标是预测带隙，就不应把“发光颜色”作为输入特征，因为发光颜色本身往往已经反映了带隙大小。这样的模型即使看起来预测准确，也没有真正学习材料成分和结构对带隙的影响。


\subsection{图表示与物理对称性}

上一节介绍了成分描述符和结构描述符。通过这些描述符，研究者可将化学式、元素属性、空间群和晶胞体积等信息转化为数值特征，用于机器学习模型预测材料性能。这种方法直观、容易理解，也适合入门学习。但它仍然有一个明显问题：很多描述符是人工设计的，模型看到的是人工预先整理好的特征，而不是材料结构本身。


成分描述符依赖人工设计，适合快速形成特征表，但难以完整保留原子排列、局域配位和相互作用距离等结构信息。带隙、形成能、弹性性质、离子扩散能力等，都与原子之间的连接关系、距离和局域环境密切相关。如果模型只使用平均电负性、平均原子半径、空间群编号等手工描述符，可能会丢失一部分精细结构信息。


图表示为这一问题提供了另一种思路：把材料看作原子节点与原子间边组成的图结构。在这张图中，原子是节点，原子之间的邻近关系或相互作用是边。这样，模型就不再只是读取一组人工设计的表格特征，而是可直接从材料的原子结构中学习信息。


继续以无机光电材料带隙预测为例，假设研究者要预测 GaN 的带隙。上一节中，研究者可使用平均电负性、电负性差、平均原子半径、空间群编号等描述符表示 GaN。但如果采用图表示，问题就变为：Ga 原子和 N 原子在晶体中如何排列，它们之间的距离是多少，每个原子周围有多少邻居。


在晶体图中，节点通常代表原子。对于 GaN 来说，图中的节点可以是 Ga 原子和 N 原子。每个节点需要带有节点特征，用来描述这个原子的基本信息。例如，一个 Ga 节点可包含 Ga 的原子序数、元素类型、电负性、原子半径、价电子数等信息；一个 N 节点也可包含 N 的对应元素属性。这样，模型就知道每个节点代表什么元素。


边通常代表原子之间的邻近关系。晶体中原子不是孤立存在的，材料性能往往取决于原子周围的局域环境。因此，需要根据一定规则判断哪些原子之间建立边。最常见的方法是设置一个截断半径：如果两个原子之间的距离小于某个阈值，就认为它们在图中相连。例如，可规定距离小于 4 \AA{} 的原子之间建立边。这样，Ga 原子周围一定距离内的 N 原子会与它相连，形成局域原子环境。


每条边也可带有边特征。边特征可包括原子间距离、相对位置、键类型或距离展开后的数值表示。对于带隙预测来说，原子间距离和局域配位环境会影响轨道重叠、成键方式和电子能带结构。因此，图表示比简单成分描述符更接近材料真实结构。


\subsubsection*{\textbf{1. 案例示范：GaN 晶体图构建}}

以 GaN 为例，构建晶体图可按照以下步骤进行。


第一步，获得 GaN 的晶体结构文件。这个文件可来自材料数据库，常见格式包括 CIF\footnote{Crystallographic Information File，晶体学信息文件。} 文件或其他晶体结构格式。结构文件中记录了晶格参数、原子种类和原子坐标。对于图表示来说，晶体结构文件是关键输入，因为它提供了原子在空间中的排列信息。


第二步，读取原子信息。程序需要从结构文件中读取每个原子的元素类型和坐标。例如，一个晶胞中可能包含 Ga 原子和 N 原子，每个原子都有对应的空间位置。此时，每个原子就可被看作图中的一个节点。


第三步，为节点添加特征。对于每个原子节点，可用元素属性作为节点特征。例如，Ga 节点可用原子序数、电负性、原子半径、价电子数等表示；N 节点也用同样类型的属性表示。这样，不同元素虽然都是“节点”，但节点特征不同，模型可区分它们的化学性质。


第四步，根据原子间距离建立边。设定一个截断半径，例如 4 \AA{}。程序计算所有原子之间的距离，如果两个原子距离小于截断半径，就在它们之间建立一条边。由于晶体具有周期性，计算距离时还要考虑相邻晶胞中的原子，而不是只看一个孤立晶胞。否则，边界附近原子的邻居关系会被错误截断。


第五步，为边添加特征。最简单的边特征是两个原子之间的距离。例如，Ga--N 距离可作为边特征输入模型。更复杂的方法会把距离展开成一组函数值，使模型更容易学习不同距离对材料性质的影响。


第六步，将整张图与目标性能对应起来。对于带隙预测任务，GaN 的晶体图作为模型输入，GaN 的带隙数值作为模型输出。模型学习的目标，就是从晶体图中预测带隙。换言之，模型不再只学习“平均电负性--带隙”的关系，而是学习“原子种类、局域连接、原子距离和晶体结构--带隙”的关系。


用表格形式表示，图数据可概括为表 \ref{tab:gan-crystal-graph-data-components} 所示内容。


\begin{formalTable}[htbp]
\caption{GaN 晶体图数据组成示例}
\label{tab:gan-crystal-graph-data-components}
\begin{tabularx}{\textwidth}{@{}YYY@{}}
\toprule
\textbf{图的组成} & \textbf{在 GaN 中对应什么} & \textbf{作用} \\
\midrule
节点 & Ga 原子、N 原子 & 表示晶体中的原子 \\
节点特征 & 原子序数、电负性、原子半径等 & 描述每个原子的元素属性 \\
边 & 截断半径内相邻原子之间的连接 & 表示局域原子关系 \\
边特征 & 原子间距离、相对位置等 & 描述原子之间的空间关系 \\
图标签 & GaN 的带隙值 & 模型需要预测的目标性能 \\
\bottomrule
\end{tabularx}
\end{formalTable}

晶体图构建涉及结构文件读取、周期性边界条件和邻接关系检查，适合借助 AI 辅助梳理代码流程。但提示词不应只写“构建晶体图”，而要明确输入文件、截断半径、输出内容和检查要求。

\begin{samepage}
\begin{userbox}
我有一个 GaN 的 CIF 晶体结构文件，准备用它构建用于带隙预测的晶体图。请帮我设计一段 Python 处理流程，用 pymatgen 读取结构，并根据 4 \AA{} 截断半径建立原子邻接关系。

要求如下：
\begin{enumerate}
\item 输出每个原子的元素种类和坐标；
\item 将原子作为节点，并说明可使用哪些节点特征；
\item 根据 4 \AA{} 截断半径生成边列表，并说明是否考虑周期性边界条件；
\item 输出节点数量、边数量和部分邻居关系，便于检查图是否合理；
\item 不要假设 CIF 文件一定正确，提醒需要核对哪些结构信息。
\end{enumerate}
\end{userbox}

\begin{aibox}
可按以下流程处理：

\begin{enumerate}
\item 用 pymatgen 读取 GaN.cif，先检查晶格参数、原子种类、原子坐标和空间群信息是否完整。
\item 将每个原子作为图中的一个节点。节点特征可从元素属性中选取，例如原子序数、电负性、原子半径和价电子数。
\item 使用结构对象的邻居搜索功能，在 4 \AA{} 截断半径内寻找相邻原子。对于晶体材料，应使用考虑周期性边界条件的方法，而不是只计算单个晶胞内部距离。
\item 将邻居关系保存为边列表，边特征可先使用原子间距离。后续如果用于图神经网络，还可把距离展开为多个函数值。
\item 检查节点数、边数和每类原子的邻居数量。如果某些原子没有邻居，或者边数量明显异常，应重新检查截断半径、CIF 文件和周期性设置。
\end{enumerate}

这段流程只能帮助完成结构到图的转换。截断半径是否合理、节点特征是否足够、图结构是否符合 GaN 的晶体环境，仍然需要结合晶体结构知识判断。
\end{aibox}
\end{samepage}

图表示之所以适合材料结构，还因为它能够更自然地满足材料中的一些物理要求。材料性质不应该因为坐标系移动、整体旋转或原子编号改变而发生变化。例如，同一个 GaN 晶体，如果研究者把整个晶体平移一段距离，它的带隙不应该改变；如果研究者重新编号原子，材料性质也不应该改变；如果只是改变观察角度，材料本身也没有发生变化。


这些要求可概括为物理对称性。对于材料机器学习，常见的物理对称性包括平移不变性、旋转不变性和置换不变性。


平移不变性是指材料整体移动后，性质不应改变。换言之，原子坐标整体加上同一个位移，材料的带隙、形成能等性质不应变化。因此，模型不应只依赖原子的绝对坐标，而应该更多使用原子之间的相对位置、距离或局域环境。


旋转不变性是指材料整体旋转后，性质不应改变。对于带隙、形成能、总能量等标量性质，材料旋转一个角度后，预测值应该保持不变。因此，很多模型会使用原子间距离这类旋转不变的量作为边特征。


置换不变性是指相同类型原子的编号顺序改变后，性质不应改变。在晶体结构文件中，原子编号只是人为排列顺序，不代表真实物理差异。例如，同一个晶体中两个等价 Ga 原子交换编号后，化学组成、局域环境和带隙都没有改变，模型输出也不应改变。图神经网络通过对邻居信息进行聚合，可在一定程度上满足这种要求。


除了不变性，还有一个更进一步的概念叫等变性。从概念上看，不变性要求某些输出在变换后保持不变，而等变性要求输出随着输入变换而按照物理规律相应变化。例如，能量是标量，材料整体旋转后能量不变；力是矢量，如果结构旋转，力的方向也应该随之旋转。可以把它理解为：同一个原子环境换了观察方向，能量数值不变，但每个原子受力的方向要跟着坐标系一起转动。对于带隙预测这类标量任务，理解不变性已经足够；对于后续机器学习势、力预测和分子动力学模拟，等变性会更加重要。


因此，图表示不仅是一种数据格式，更是一种更符合材料物理特征的表示方式。它把材料看作原子节点和相互作用边组成的结构，同时尽量尊重材料中的平移、旋转和置换对称性。相比只使用手工描述符，图表示能够保留更多局域结构信息，也为后续图神经网络模型奠定基础。


图表示并不意味着完全不需要人工判断。构建晶体图时，截断半径、节点特征、边特征表示和周期性边界条件处理都会影响模型结果。即使借助工具生成代码并检查图结构，上述关键设置是否合理，仍然需要结合材料结构和预测任务来判断。

\subsubsection*{\textbf{小结}}

本节说明材料信息需要转化为描述符、结构字段或晶体图，才能进入机器学习模型。
成分和结构描述符适合表格建模，图表示则更适合保留原子连接、距离和局域环境。
特征构建必须避免目标泄漏，并结合物理对称性、截断半径和边特征检查输入是否合理。

\subsection*{习题}

\begin{enumerate}[leftmargin=*]
\item 比较成分描述符和结构描述符的作用及适用场景。
\item 图表示如何保留比成分描述符更多的结构信息？
\item 成分描述符和图表示各有什么优势与局限？请结合带隙预测举例说明。
\end{enumerate}

\clearpage
\section{性能预测与尺度模拟}


\begin{sectionkeypoints}
\item 材料性能预测的核心，是让模型学习材料特征与目标性能之间的对应关系。硬度、带隙、形成能、弹性模量和熔点都可成为预测目标。
\item 描述符模型适合表格数据和入门任务，关键步骤包括整理特征表、划分训练测试集、训练模型、评价误差和检查特征重要性。
\item 图神经网络适合处理晶体结构数据，能够从原子节点、原子间边和局域环境中学习材料性质。
\item 机器学习势学习的是原子结构与能量、力之间的关系，可连接第一性原理计算和分子动力学模拟。
\item 模型预测和尺度模拟都必须进行可靠性检查。训练集外推、过拟合、数据泄漏、结构文件错误和模拟参数不合理都会影响结论。
\end{sectionkeypoints}

\subsection{描述符模型与硬度预测}
\label{subsec:descriptor-hardness-prediction}

有了材料描述符特征表之后，下一步是建立特征与性能之间的映射关系。也就是说，模型需要根据成分、结构、工艺或显微组织等输入特征，预测硬度、带隙、形成能、熔点等目标性能。


材料性能预测的基本思想，是让模型学习“材料特征”和“目标性能”之间的对应关系。材料特征可来自成分、结构、工艺和显微组织，目标性能可以是硬度、强度、带隙、形成能、熔点或耐腐蚀性能等。对于传统机器学习方法而言，材料首先要被转化成一张特征表，每一行代表一个材料样品，每一列代表一个可计算的特征，最后一列通常是需要预测的性能。


\subsubsection*{\textbf{1. 案例示范：高熵合金涂层硬度预测}}

本节只围绕一个具体任务展开：高熵合金涂层硬度预测。假设研究者希望根据高熵合金涂层的成分和组织信息，预测涂层硬度。这个例子适合说明传统机器学习预测流程，因为硬度是材料实验中常见的数值型性能，既受成分影响，也受相结构、晶粒尺寸和硬质相等组织因素影响。


对于一个高熵合金涂层样品，原始实验记录可能包括名义成分、制备工艺、XRD\footnote{X-ray Diffraction，X 射线衍射。} 相分析、显微组织观察和显微硬度测试结果。经过前面的描述符构建过程后，这些信息可整理为表 \ref{tab:hea-hardness-features} 所示形式。


\begin{formalTable}[htbp]
\caption{高熵合金涂层硬度预测特征表示示例}
\label{tab:hea-hardness-features}
\scriptsize
\setlength{\tabcolsep}{2.8pt}
\begin{tabularx}{\textwidth}{@{}YYYYYYYY@{}}
\toprule
\textbf{样品编号} & \textbf{原子半径差} & \textbf{平均电负性} & \textbf{电负性差} & \textbf{价电子浓度} & \textbf{是否含硬质相} & \textbf{平均晶粒尺寸/\(\mu\)m} & \textbf{硬度/HV} \\
\midrule
S1 & 数值 & 数值 & 数值 & 数值 & 0 & 数值 & 数值 \\
S2 & 数值 & 数值 & 数值 & 数值 & 1 & 数值 & 数值 \\
S3 & 数值 & 数值 & 数值 & 数值 & 1 & 数值 & 数值 \\
S4 & 数值 & 数值 & 数值 & 数值 & 1 & 数值 & 数值 \\
\bottomrule
\end{tabularx}
\end{formalTable}

表 \ref{tab:hea-hardness-features} 中，“原子半径差、平均电负性、电负性差、价电子浓度、是否含硬质相、平均晶粒尺寸”等列是模型输入特征，“硬度/HV”是模型需要预测的目标值。换成机器学习语言，就是把输入特征组成矩阵 X，把硬度组成目标向量 y。模型训练的目的，就是从 X 中学习到 y 的变化规律。


在这个任务中，硬度是一个连续数值，因此它属于回归任务。回归任务的目标不是判断材料属于哪一类，而是预测一个具体数值。例如，模型输入某个涂层样品的成分和组织特征后，输出预测硬度是多少 HV。与之相对，如果任务是判断材料是否耐腐蚀、是否稳定、是否属于某种相结构，那就更接近分类任务。


为了让这个过程更加具体，以下给出一个精简的机器学习代码框架。假设研究者已经整理好了一个高熵合金涂层硬度数据表，文件名为 hea\_\allowbreak{}coating\_\allowbreak{}hardness.csv。表格中包含若干材料描述符和目标性能“硬度/HV”。代码的目标是读取数据、划分训练集和测试集、训练随机森林回归模型，并评价预测误差。代码采用平均绝对误差（Mean Absolute Error，MAE）、均方根误差（Root Mean Squared Error，RMSE）和决定系数（Coefficient of Determination，\(R^2\)）评价预测误差与拟合程度。


\begin{lstlisting}[language=Python,escapeinside={(*@}{@*)}]
# (*@\textnormal{1. 导入常用库}@*)
import pandas as pd
from sklearn.model_selection import train_test_split
from sklearn.ensemble import RandomForestRegressor
from sklearn.metrics import mean_absolute_error, mean_squared_error, r2_score
# (*@\textnormal{2. 读取高熵合金涂层硬度数据表}@*)
data = pd.read_csv("hea_coating_hardness.csv")
# (*@\textnormal{3. 指定输入特征和预测目标}@*)
# (*@\textnormal{这些特征来自前面已经构建好的材料描述符}@*)
feature_cols = [
    "atomic_radius_diff",             # (*@\textnormal{原子半径差}@*)
    "mean_electronegativity",         # (*@\textnormal{平均电负性}@*)
    "electronegativity_diff",         # (*@\textnormal{电负性差}@*)
    "valence_electron_concentration", # (*@\textnormal{价电子浓度}@*)
    "has_hard_phase",                 # (*@\textnormal{是否含硬质相：0 表示否，1 表示是}@*)
    "grain_size"                      # (*@\textnormal{平均晶粒尺寸}@*)
]
X = data[feature_cols]
y = data["hardness_HV"]
# (*@\textnormal{4. 划分训练集和测试集}@*)
X_train, X_test, y_train, y_test = train_test_split(
    X,
    y,
    test_size=0.2,
    random_state=42
)
# (*@\textnormal{5. 建立并训练随机森林回归模型}@*)
model = RandomForestRegressor(
    n_estimators=200,
    random_state=42
)
model.fit(X_train, y_train)
# (*@\textnormal{6. 在测试集上进行硬度预测}@*)
y_pred = model.predict(X_test)
# (*@\textnormal{7. 评价模型预测效果}@*)
mae = mean_absolute_error(y_test, y_pred)
rmse = mean_squared_error(y_test, y_pred) ** 0.5
r2 = r2_score(y_test, y_pred)
print("MAE:", mae)
print("RMSE:", rmse)
print("R2:", r2)
# (*@\textnormal{8. 输出特征重要性}@*)
importance = pd.DataFrame({
    "feature": feature_cols,
    "importance": model.feature_importances_
}).sort_values(by="importance", ascending=False)
print(importance)
\end{lstlisting}

这段代码对应了材料性能预测的基本流程。首先，pd.read\_\allowbreak{}csv() 用于读取已经整理好的材料数据表。此处的数据表不是原始实验记录，而是已经经过描述符构建后的特征表。换言之，代码中的输入不是“FeCoCrNiTi”这样的化学式，而是“原子半径差、平均电负性、电负性差、价电子浓度、是否含硬质相、平均晶粒尺寸”等数值特征。


其中，X = data[feature\_\allowbreak{}cols] 表示模型输入，y = data["hardness\_\allowbreak{}HV"] 表示模型输出。对于本例来说，X 是涂层样品的成分和组织特征，y 是实验测得的硬度值。模型训练时要学习的，就是这些特征与硬度之间的关系。


train\_\allowbreak{}test\_\allowbreak{}split() 用于划分训练集和测试集。训练集用于让模型学习规律，测试集用于检验模型对新样品的预测能力。此处设置 test\_\allowbreak{}size=0.2，表示 80\% 的样品用于训练，20\% 的样品用于测试。对于材料数据来说，样本数量往往不大，因此测试集必须严格独立，不能参与模型训练。否则，模型评估结果会偏高，不能代表真实预测能力。


本例使用的模型是随机森林回归模型。随机森林由多棵决策树组成，能够处理非线性关系，对小规模表格数据比较友好，也可输出特征重要性。对于高熵合金涂层硬度预测来说，成分描述符、硬质相和晶粒尺寸与硬度之间通常不是简单线性关系，因此随机森林适合作为入门模型。此处选择随机森林并不意味着它一定最优，而是因为它便于研究者理解完整建模流程。


模型训练完成后，model.predict(X\_\allowbreak{}test) 会对测试集样品进行硬度预测。此时模型只看到测试样品的输入特征，不知道真实硬度。随后，研究者用真实硬度和预测硬度进行比较，计算 MAE、RMSE 和 R\textsuperscript{2}。


MAE 是平均绝对误差，表示模型预测值与真实值平均相差多少。例如，如果 MAE 为 30 HV，可理解为模型对测试样品的硬度平均预测误差约为 30 HV。RMSE 是均方根误差，对较大的预测误差更加敏感。如果少数样品预测偏差很大，RMSE 会明显增大。R\textsuperscript{2} 是决定系数，用来衡量预测值与真实值变化趋势的一致程度，越接近 1，说明模型预测效果越好。


对于材料研究来说，仅有误差指标还不够。宜进一步绘制“真实硬度--预测硬度”散点图。如果模型预测较好，散点应大致分布在 y = x 参考线附近。图中的坐标轴、单位、图例和参考线都要清楚，否则模型结果不便于与实验报告或论文表达衔接。


除了预测误差，随机森林模型还可输出特征重要性。代码中的 model.feature\_\allowbreak{}importances\_\allowbreak{} 可显示不同输入特征对硬度预测的贡献大小。如果结果显示“是否含硬质相”“平均晶粒尺寸”“原子半径差”等特征重要性较高，就可结合材料强化机制进行分析。例如，硬质相可能带来第二相强化，晶粒细化可能带来细晶强化，原子半径差可能与固溶强化和晶格畸变有关。


在代码实现阶段，AI 更适合承担“字段核对、代码改写、报错解释和绘图补充”等辅助工作，而不是替研究者判断哪些特征一定正确。可用下面的提示词约束任务边界。

\begin{samepage}
\begin{userbox}
我正在用随机森林预测高熵合金涂层硬度，CSV 表头包括原子半径差、平均电负性、电负性差、价电子浓度、是否含硬质相、晶粒尺寸和硬度/HV。请帮我检查这段建模流程是否清楚，并补充测试集真实硬度与预测硬度散点图的绘图步骤。

要求如下：
\begin{enumerate}
\item 判断哪些列可作为输入特征，哪一列是预测目标；
\item 提醒哪些字段需要确认数据来源或实验条件；
\item 说明如果出现列名不存在、缺失值或字符串字段无法输入模型，应如何检查；
\item 补充绘制真实值--预测值散点图和 y=x 参考线的思路；
\item 不要替我判断模型一定可靠，只给出代码检查和结果呈现建议。
\end{enumerate}
\end{userbox}

\begin{aibox}
表 \ref{tab:hea-hardness-features} 中，原子半径差、平均电负性、电负性差、价电子浓度、是否含硬质相和晶粒尺寸可作为候选输入特征，硬度/HV 是预测目标，不应同时放入输入特征。

建模前需要重点检查三类问题：
\begin{enumerate}
\item \textbf{字段是否存在：}代码中指定的输入特征名称必须与 CSV 表头完全一致，否则会出现列名不存在的错误。
\item \textbf{数据是否可输入模型：}输入特征应为数值型；如果存在空值、文字标签或单位混用，需要先清洗或编码。
\item \textbf{材料含义是否可靠：}是否含硬质相应来自 XRD、SEM\footnote{Scanning Electron Microscopy，扫描电子显微镜。} 或其他表征结果，晶粒尺寸应说明统计方法，硬度值应记录测试载荷和保载时间。
\end{enumerate}

模型评价时，除 MAE、RMSE 和 R\textsuperscript{2} 外，可绘制测试集真实硬度与预测硬度散点图，并添加 y=x 参考线。如果散点明显偏离参考线，说明模型在部分样品上的预测误差较大，需要回到数据质量、样本数量和特征设计中继续检查。
\end{aibox}
\end{samepage}

在实际建模时，还应特别注意过拟合。材料数据往往样本少、特征多，如果只有十几个样品却加入几十个描述符，模型可能在训练集上表现很好，但在测试集上预测很差。这说明模型可能只是记住了已有样品，而没有真正学到可推广的规律。为了降低过拟合风险，可减少无关特征，采用交叉验证，增加样本数量，或者选择更简单、更稳定的模型。是否属于过拟合，仍然要结合训练误差、测试误差、数据规模和材料任务共同判断。


数据泄漏同样需要警惕。预测硬度时，不应把由硬度结果直接推导出的字段作为输入特征。例如，“是否高硬度样品”“硬度等级”这类字段不能作为输入，因为它们本身已经包含了目标信息。否则模型看起来预测很准，实际上并没有从成分和组织中学习硬度规律。


\subsection{图神经网络与性质预测}

上一节以高熵合金涂层硬度预测为例，介绍了传统机器学习的基本流程。那种方法的特点是：研究者首先根据材料知识人工构造描述符，例如原子半径差、电负性差、价电子浓度、是否含硬质相和晶粒尺寸，然后再把这些描述符输入随机森林等模型中预测硬度。这种方法直观、容易操作，也适合材料数据量较少的入门任务。


但是，传统描述符方法也有局限。材料性能不仅取决于几个平均特征，还与原子在晶体中的排列方式、局域配位环境、原子间距离和结构稳定性有关。尤其对于高熵合金涂层这类复杂体系，材料中可能同时存在 FCC\footnote{Face-Centered Cubic，面心立方结构。} 固溶体、BCC\footnote{Body-Centered Cubic，体心立方结构。} 相、碳化物或硼化物等不同结构。仅用“是否含硬质相”“平均晶粒尺寸”这样的表格特征，很难完整表达原子尺度结构对性能的影响。


图神经网络提供了另一种建模思路。它不再把材料完全压缩成一行人工描述符，而是把材料看作由原子节点和原子间关系组成的图结构。在晶体图中，每个原子是一个节点，原子之间的邻近关系是边，节点和边都可带有特征。模型通过在图上进行信息传递，学习局域原子环境与材料性质之间的关系。


这里要限定问题边界：高熵合金涂层是一个多尺度、多相组织体系，真实涂层中包含晶粒、晶界、第二相、孔隙、界面和残余应力等复杂因素。图神经网络不宜简单地把一张 SEM 图或一个涂层样品直接转化成完整晶体图，然后准确预测整体硬度。更合理的入门方式是：选取涂层中的代表性晶体相或候选晶体结构，例如 FCC 固溶体、BCC 固溶体或 TiC 等硬质相，把它们转化为晶体图，学习结构与形成能、弹性模量或稳定性等性质之间的关系。然后再把这些原子尺度信息作为理解涂层硬度和强化机制的补充。


这里的重点不是用图神经网络直接预测整个涂层硬度，而是说明如何把涂层中的代表性晶体结构转化为图，并用图神经网络预测与性能相关的材料性质。


假设研究者研究的是一种 FeCoCrNiTi 高熵合金涂层。实验表明，涂层中可能存在 FCC 固溶体、BCC 相和 TiC 等硬质相。上一节中，研究者用人工描述符预测涂层硬度；现在研究者从更微观的角度思考：如果想理解某个相为什么可能更稳定、为什么某种硬质相会提高硬度，就需要进一步关注原子尺度结构。图神经网络正适合处理这类晶体结构信息。


对于一个代表性晶体结构，原始数据通常来自 CIF 文件或数据库结构文件。一个结构文件中记录了晶格参数、原子种类和原子坐标。图神经网络的第一步，就是把这个晶体结构转化为晶体图。


\subsubsection*{\textbf{1. 案例示范：TiC 晶体图性质预测}}

以 TiC 硬质相为例，它在高熵合金涂层中常被视为提高硬度的重要相之一。构建 TiC 晶体图时，可按照以下步骤进行。


第一步，读取晶体结构文件。结构文件中包含 Ti 和 C 原子的种类、坐标以及晶格参数。对于模型来说，这些信息是图结构的基础。


第二步，把原子转化为节点。TiC 中的 Ti 原子和 C 原子都可看作图中的节点。每个节点需要包含节点特征，例如原子序数、电负性、原子半径、价电子数等。这样，模型能够区分 Ti 节点和 C 节点，并理解它们具有不同化学属性。


第三步，根据原子间距离建立边。在晶体中，一个原子的性质受到周围邻居原子的影响。因此，可设置一个截断半径，例如 5 \AA{}。如果两个原子之间的距离小于这个截断半径，就在它们之间建立一条边。这样，模型就可知道哪些原子彼此相邻。


第四步，为边添加特征。最常见的边特征是原子间距离。距离越近，原子之间相互作用通常越强；距离不同，键合环境也不同。因此，边特征可帮助模型学习局域结构对性质的影响。


第五步，把整个晶体图与目标性质对应起来。如果目标是预测形成能，那么每一个晶体图对应一个形成能数值；如果目标是预测弹性模量，那么每一个晶体图对应一个弹性模量数值。模型训练的目标，就是学习“晶体图结构”与“目标性质”之间的关系。


该流程可概括为下列步骤，并可结合图 \ref{fig:crystal-graph-prediction} 理解：


\begin{lightquote}
\centering 晶体结构文件 → 原子节点 → 原子间边 → 节点特征和边特征 → 晶体图 → 性质预测
\end{lightquote}

\begin{figure}[H]
\centering
\BookFigureImage{images_11/2.png}
\caption{晶体图表示与性能预测示意图}
\label{fig:crystal-graph-prediction}
\end{figure}
\FloatBarrier

图 \ref{fig:crystal-graph-prediction} 所示流程中，图数据的组成可进一步概括为表 \ref{tab:crystal-graph-data-components}。


\begin{formalTable}[htbp]
\caption{晶体图数据组成示例}
\label{tab:crystal-graph-data-components}
\begin{tabularx}{\textwidth}{@{}YY@{}}
\toprule
\textbf{图数据组成} & \textbf{在 TiC 或高熵合金相关晶体结构中表示什么} \\
\midrule
节点 & Ti、C 或 Fe、Co、Cr、Ni、Ti 等原子 \\
节点特征 & 原子序数、电负性、原子半径、价电子数等 \\
边 & 截断半径内的原子邻近关系 \\
边特征 & 原子间距离或距离展开特征 \\
图标签 & 形成能、弹性模量、体模量或其他目标性质 \\
\bottomrule
\end{tabularx}
\end{formalTable}

图神经网络的核心是信息传递。每个原子节点在初始时只包含自己的元素属性，例如 Ti 节点知道自己是 Ti，C 节点知道自己是 C。经过一轮信息传递后，一个节点会接收周围邻居节点的信息。经过多轮信息传递后，节点就不仅包含自身信息，还包含局域原子环境信息。最后，模型把所有节点的信息汇总成整个晶体的表示，再输出目标性质。


这个过程可用材料语言理解：模型不只是看“这个材料含有哪些元素”，而是进一步看“这些元素在空间中如何排列、每个原子周围是什么邻居、原子间距离是多少、局域环境是否稳定”。这正是图神经网络相比传统描述符方法的重要优势。


为了让这个过程更具体，以下给出一个教学用的精简代码框架。它的目的不是提供完整可直接运行的项目，而是帮助研究者理解图神经网络建模的基本步骤。


\begin{lstlisting}[language=Python,escapeinside={(*@}{@*)}]
# (*@\textnormal{教学用简化框架：从晶体结构构建图并预测性质}@*)
# (*@\textnormal{实际运行时需要安装 pymatgen、torch、torch\_\allowbreak{}geometric 等库}@*)
from pymatgen.core import Structure
import torch
# (*@\textnormal{1. 读取晶体结构文件，例如 TiC.cif}@*)
structure = Structure.from_file("TiC.cif")
# (*@\textnormal{2. 设置截断半径}@*)
cutoff = 5.0
# (*@\textnormal{3. 构建节点特征：此处仅用原子序数作为示例}@*)
node_features = []
for site in structure:
    atomic_number = site.specie.Z
    node_features.append([atomic_number])
x = torch.tensor(node_features, dtype=torch.float)
# (*@\textnormal{4. 构建边：寻找截断半径内的邻居原子}@*)
edge_index = []
edge_attr = []
for i, site in enumerate(structure):
    neighbors = structure.get_neighbors(site, cutoff)
    for neighbor in neighbors:
        j = neighbor.index
        distance = neighbor.nn_distance
        edge_index.append([i, j])
        edge_attr.append([distance])
edge_index = torch.tensor(edge_index, dtype=torch.long).t().contiguous()
edge_attr = torch.tensor(edge_attr, dtype=torch.float)
# (*@\textnormal{5. 图标签：例如该结构的形成能或弹性模量}@*)
# (*@\textnormal{此处用一个占位数值表示，实际应来自数据库或计算结果}@*)
y = torch.tensor([target_property], dtype=torch.float)
# (*@\textnormal{至此，一个晶体结构已经被转化为图数据：}@*)
# (*@\textnormal{x 表示节点特征}@*)
# (*@\textnormal{edge\_\allowbreak{}index 表示边连接关系}@*)
# (*@\textnormal{edge\_\allowbreak{}attr 表示边特征}@*)
# (*@\textnormal{y 表示目标性质}@*)
\end{lstlisting}

这段代码展示了晶体图构建的基本思想。Structure.from\_\allowbreak{}file("TiC.cif") 用来读取晶体结构文件；node\_\allowbreak{}features 保存每个原子的节点特征；edge\_\allowbreak{}index 保存哪些原子之间存在边；edge\_\allowbreak{}attr 保存边特征，例如原子间距离；y 是模型要预测的目标性质。


实际使用时，节点特征不应只包含原子序数，还可加入电负性、原子半径、价电子数等元素属性。边特征也不一定只用一个距离数值，可把距离展开成一组函数，以便模型更好学习不同距离范围内的相互作用。对于不同材料任务，目标 y 也可不同。如果关注结构稳定性，可预测形成能；如果关注力学性能，可预测弹性模量或体模量；如果进一步结合实验数据，也可尝试预测与硬度相关的性质。


如果借助 AI 辅助构建 TiC 晶体图，提示词应突出“代表性晶体相”和“目标性质”两个信息，避免把图神经网络误写成直接预测整个涂层硬度的工具。示例任务可写成如下形式。

\begin{samepage}
\begin{userbox}
我想以高熵合金涂层中的 TiC 硬质相为例，构建一个用于性质预测的晶体图。请帮我设计从 TiC.cif 到图数据的处理流程，目标性质可以是形成能、弹性模量或体模量。

要求如下：
\begin{enumerate}
\item 说明如何读取 TiC.cif，并把 Ti 和 C 原子转化为节点；
\item 说明节点特征可选择哪些元素属性；
\item 根据 5 \AA{} 截断半径建立边，并说明为什么要考虑周期性边界条件；
\item 说明边特征和图标签分别是什么；
\item 提醒这个晶体图只能表示 TiC 这一代表性相，不宜直接代表整个高熵合金涂层硬度。
\end{enumerate}
\end{userbox}

\begin{aibox}
可按以下流程构建图数据：

\begin{enumerate}
\item \textbf{读取结构：}使用 pymatgen 读取 TiC.cif，检查晶格参数、Ti/C 原子坐标和结构完整性。
\item \textbf{构建节点：}把每个 Ti 原子和 C 原子作为节点。节点特征可包括原子序数、电负性、原子半径和价电子数。
\item \textbf{构建边：}使用 5 \AA{} 截断半径寻找邻近原子。由于 TiC 是周期晶体，邻居搜索必须考虑相邻晶胞，否则会漏掉晶胞边界附近的邻居。
\item \textbf{添加边特征：}可先使用原子间距离作为边特征，后续再根据模型需要使用距离展开特征。
\item \textbf{设置图标签：}如果训练目标是形成能，图标签就是形成能；如果目标是弹性模量或体模量，图标签应来自数据库、第一性原理计算或可靠实验/计算资料。
\end{enumerate}

这里需要专门说明，TiC 晶体图只能帮助理解涂层中某个硬质相的结构性质。真实高熵合金涂层硬度还受到晶粒尺寸、相分布、界面、孔隙和残余应力影响，不能由单个 TiC 晶体图直接决定。
\end{aibox}
\end{samepage}

与传统描述符模型相比，图神经网络的优势在于它能更直接地保留原子结构信息。传统模型需要人工预先设计“平均电负性”“原子半径差”等特征，而图神经网络可从原子种类、原子间距离和局域连接关系中自动学习结构信息。对于晶体材料来说，这种表示方式更接近材料本身。


但是，图神经网络也不是万能的。首先，它需要可靠的晶体结构输入。如果 CIF 文件错误、原子占位不清楚或结构不完整，构建出的图也会有问题。其次，它通常需要较多训练数据。如果只有少量结构样本，图神经网络也可能过拟合。再次，对于真实高熵合金涂层来说，整体硬度还受到晶粒尺寸、相分布、孔隙率、界面结合和残余应力等显微组织因素影响，单个晶体图不能完全代表整个涂层。


因此，在高熵合金涂层研究中，图神经网络更适合作为结构层面的补充工具。它可帮助研究者理解某个候选相、硬质相或代表性晶体结构的稳定性和力学相关性质；而整体涂层硬度仍然需要结合上一节中的表格特征模型、显微组织数据和实验测试结果共同分析。


\subsection{原子间势与熔点模拟}

前两节讨论的是材料性能预测。无论是传统机器学习模型，还是图神经网络模型，它们通常都是输入一个材料样品或晶体结构，然后输出一个性能数值，例如硬度、形成能、带隙或弹性模量。这类模型回答的是“这个材料可能具有什么性质”。


但是，材料科学中还有另一类重要问题：材料在原子尺度上如何运动、如何扩散、如何相变、如何熔化。例如，超高温陶瓷在高温下能否保持结构稳定？原子在升温过程中是否发生剧烈扩散？材料在什么温度附近开始熔化？这些问题不是简单预测一个静态数值就能完全回答的，而需要模拟原子随时间的运动过程。


分子动力学可用来研究这类原子尺度动态行为。在分子动力学中，研究者根据原子之间的相互作用计算每个原子的受力，再根据力更新原子位置，从而模拟材料随时间演化的过程。问题在于，原子间相互作用如何计算。第一性原理分子动力学精度较高，但计算成本高，难以处理大体系和长时间尺度；传统经验势速度快，但对复杂材料体系的适用性有限。机器学习势正是在这种背景下发展起来的。


机器学习势的核心思想，是用机器学习模型学习第一性原理计算得到的势能面。从概念上看，第一性原理计算可给出某一组原子结构的能量、原子受力和应力；机器学习势则利用这些数据学习“原子结构--能量/力”之间的关系。训练完成后，机器学习势可像传统势函数一样快速计算能量和力，从而支持更大规模、更长时间的分子动力学模拟。


\subsubsection*{\textbf{1. 案例示范：TiC 熔化行为模拟}}

本节只围绕一个任务展开：用机器学习势辅助模拟 TiC 超高温陶瓷的熔化行为。TiC 是典型的超高温陶瓷之一，具有高熔点、高硬度和良好的高温稳定性，也常作为高熵合金涂层或复合涂层中的硬质增强相。通过这个例子，可说明机器学习势如何从第一性原理数据出发，进一步服务于高温分子动力学模拟。


首先需要明确，机器学习势与前面性能预测模型的目标不同。前面随机森林模型预测的是硬度，图神经网络可预测形成能或弹性模量；而机器学习势要预测的是原子结构对应的总能量和每个原子的受力。它不是直接告诉研究者“TiC 的熔点是多少”，而是先学会在不同原子构型下计算能量和力，然后通过分子动力学模拟升温过程，再从模拟结果中判断熔化行为。


对于 TiC 熔点模拟，可按照以下流程理解。


第一步，准备 TiC 的初始晶体结构。研究者需要获得 TiC 的晶体结构文件，例如 CIF 文件。这个结构文件包含 Ti 和 C 原子的排列方式、晶格参数和空间结构。初始结构通常可来自数据库、文献或已有计算结果。


第二步，生成不同构型的训练数据。机器学习势不应只看一个理想晶体结构，因为真实高温模拟中原子会振动、偏移，甚至出现缺陷和局部无序。因此，需要构造一批不同的 TiC 原子构型，例如平衡晶体结构、轻微扰动结构、不同体积结构、不同温度下短时间分子动力学得到的结构等。这样，训练数据才能覆盖模型后续可能遇到的原子环境。


第三步，用第一性原理计算这些构型的能量和力。对于每一个 TiC 构型，需要通过第一性原理计算得到总能量、每个原子的力以及必要时的应力。这一步计算成本较高，但只需要对训练集中的有限构型进行。可理解为：第一性原理计算负责提供高质量“标准答案”，机器学习势负责学习这些标准答案中的规律。


第四步，训练机器学习势。训练时，模型输入原子种类和原子坐标，输出预测能量和原子力。训练目标是让模型预测的能量和力尽可能接近第一性原理计算结果。常见机器学习势包括 Behler--Parrinello 神经网络势、高斯近似势、DeepMD\footnote{Deep Potential Molecular Dynamics，深度势分子动力学。}、MACE\footnote{Message Passing Atomic Cluster Expansion，消息传递原子簇展开模型。}、NequIP\footnote{Neural Equivariant Interatomic Potentials，神经等变原子间势。}、M3GNet\footnote{Materials 3-body Graph Network，材料三体图网络。}、CHGNet\footnote{Crystal Hamiltonian Graph Neural Network，晶体哈密顿量图神经网络。} 等。教材入门阶段不需要详细展开每种模型的数学细节，重点是理解它们都在学习原子结构与能量/力之间的关系。


第五步，验证机器学习势是否可靠。训练完成后，不宜直接用于高温模拟，而要先在测试构型上验证预测误差。常见做法是比较模型预测能量与第一性原理能量之间的误差，以及模型预测力与第一性原理力之间的误差。如果误差较大，说明模型还没有学好 TiC 的势能面，需要增加训练数据、调整模型参数或补充高温构型。


第六步，用训练好的机器学习势进行分子动力学模拟。此时，机器学习势可快速计算每一步原子的能量和力，进而模拟 TiC 在不同温度下的原子运动。为了观察熔化行为，可逐步升高温度，例如从较低温度开始，逐步升温到高温区间，记录能量、体积、径向分布函数、均方位移等随温度变化的情况。


第七步，根据模拟结果判断熔化行为。材料熔化时，原子排列会从有序晶体逐渐转变为无序液态结构。模拟中可通过多个指标辅助判断：能量--温度曲线是否出现突变，体积是否发生明显变化，径向分布函数的长程有序峰是否消失，均方位移是否显著增加。通过这些现象，可估计 TiC 的熔化温度范围。


这个过程可概括为：


\begin{lightquote}
\centering TiC 晶体结构 → 多种原子构型 → 第一性原理能量和力 → 训练机器学习势 → 验证误差 → 分子动力学升温模拟 → 分析熔化行为
\end{lightquote}

为了让研究者更清楚地理解流程，以下给出一个教学用的简化代码框架。它不是完整可直接运行的工程代码，因为真实机器学习势训练通常需要专门软件和大量计算资源。此处的目的，是帮助理解每一步在做什么。


\begin{lstlisting}[language=Python,escapeinside={(*@}{@*)}]
# (*@\textnormal{教学用简化框架：机器学习势辅助 TiC 熔化模拟}@*)
# (*@\textnormal{实际研究中可使用 DeepMD、MACE、NequIP、CHGNet、M3GNet 等工具}@*)
# (*@\textnormal{1. 读取 TiC 初始晶体结构}@*)
structure = read_structure("TiC.cif")
# (*@\textnormal{2. 生成训练构型}@*)
# (*@\textnormal{包括平衡结构、扰动结构、不同体积结构和高温采样结构}@*)
train_structures = generate_configurations(
    structure,
    perturb=True,
    volume_changes=True,
    temperature_sampling=True
)
# (*@\textnormal{3. 对训练构型进行第一性原理计算}@*)
# (*@\textnormal{得到每个构型的能量、原子力和应力}@*)
dft_dataset = []
for config in train_structures:
    energy, forces, stress = run_dft_calculation(config)
    dft_dataset.append({
        "structure": config,
        "energy": energy,
        "forces": forces,
        "stress": stress
    })
# (*@\textnormal{4. 训练机器学习势}@*)
ml_potential = train_machine_learning_potential(
    dataset=dft_dataset,
    target_properties=["energy", "forces", "stress"]
)
# (*@\textnormal{5. 在测试构型上验证模型误差}@*)
test_results = evaluate_potential(
    model=ml_potential,
    test_dataset="test_configurations"
)
print(test_results)
# (*@\textnormal{6. 使用训练好的机器学习势进行分子动力学升温模拟}@*)
md_results = run_molecular_dynamics(
    structure=structure,
    potential=ml_potential,
    temperature_schedule=[1000, 1500, 2000, 2500, 3000, 3500, 4000],
    simulation_time="several ps or longer"
)
# (*@\textnormal{7. 分析熔化行为}@*)
plot_energy_temperature(md_results)
plot_radial_distribution_function(md_results)
plot_mean_square_displacement(md_results)
\end{lstlisting}

这段框架中，read\_\allowbreak{}structure() 表示读取 TiC 的晶体结构；generate\_\allowbreak{}configurations() 表示生成不同原子构型；run\_\allowbreak{}dft\_\allowbreak{}calculation() 表示用第一性原理计算能量和力；train\_\allowbreak{}machine\_\allowbreak{}learning\_\allowbreak{}potential() 表示训练机器学习势；run\_\allowbreak{}molecular\_\allowbreak{}dynamics() 表示用训练好的势进行分子动力学模拟；最后通过能量曲线、径向分布函数和均方位移分析熔化行为。


与前两节代码不同，此处给出的是流程框架，而不是几十行就能完成的入门脚本。原因是机器学习势涉及第一性原理计算、模型训练和分子动力学模拟，真实流程需要较多软件和计算资源。但这个框架可帮助研究者理解机器学习势在材料研究中的位置：它不是直接替代实验，也不是简单回归一个性能值，而是连接第一性原理计算和大尺度原子模拟的桥梁。


在机器学习势和分子动力学任务中，AI 辅助更适合用于解释文件字段、检查流程是否完整和整理输出指标，而不宜直接判断熔点结论。下面的写法更适合用于流程核查。

\begin{samepage}
\begin{userbox}
我想用机器学习势辅助模拟 TiC 的熔化行为。已有 TiC 初始结构，计划生成扰动构型和高温采样构型，用第一性原理计算能量与力，再训练机器学习势并进行升温分子动力学模拟。请帮我检查这个流程是否完整，并说明每类输出指标的物理意义。

要求如下：
\begin{enumerate}
\item 说明训练数据应包含哪些 TiC 原子构型；
\item 说明第一性原理计算需要记录哪些目标量；
\item 检查机器学习势训练后应如何验证能量和力误差；
\item 说明升温分子动力学中能量、体积、径向分布函数和均方位移分别反映什么；
\item 明确提醒不应只根据单一指标判断熔点。
\end{enumerate}
\end{userbox}

\begin{aibox}
这个流程可分为四个检查重点：

\begin{enumerate}
\item \textbf{训练构型覆盖范围：}除平衡 TiC 晶体结构外，还应包含轻微扰动结构、不同体积结构、不同温度采样结构，以及接近高温无序状态的构型。否则模型可能只适用于低温平衡结构。
\item \textbf{第一性原理标注目标：}每个构型至少应记录总能量和每个原子的力；如果后续模拟涉及体积变化，还应考虑应力信息。
\item \textbf{机器学习势验证：}训练后应在独立测试构型上比较预测能量、预测力与第一性原理结果的误差，不应只看训练集误差。
\item \textbf{熔化行为分析：}能量--温度曲线突变可能反映相变，体积突变可能反映结构膨胀，径向分布函数长程有序峰消失说明晶体有序性降低，均方位移显著增大说明原子扩散增强。
\end{enumerate}

熔点判断应综合多个指标，并结合体系大小、升温速率、模拟时间和训练数据覆盖范围共同分析。单独依靠 AI 对曲线的文字解释，不能作为可靠的材料结论。
\end{aibox}
\end{samepage}

需要特别强调的是，机器学习势的可靠性高度依赖训练数据。如果训练集中只包含低温平衡结构，而模拟时却用于高温熔化过程，模型很可能外推到训练分布之外，预测结果就不可靠。因此，在 TiC 熔点模拟中，训练数据应尽量包含高温扰动构型、不同体积构型和接近熔化状态的构型。否则，模型可能在低温结构上误差很小，却在高温模拟中产生错误行为。


此外，机器学习势不能完全替代第一性原理计算和实验验证。它的优势是加速模拟，但前提是训练数据可靠、验证误差合理、模拟条件设置正确。对于超高温陶瓷这类材料，熔点模拟还可能受到模拟体系大小、升温速率、时间尺度和缺陷状态影响。因此，模拟得到的熔化温度应被理解为计算预测结果，需要与实验数据或更高精度计算进行对照。

\subsubsection*{\textbf{小结}}

本节围绕硬度预测、图神经网络性质预测和机器学习势展开。
描述符模型用于表格化的宏观性能预测，图神经网络用于原子结构关联的性质预测，机器学习势用于原子尺度动力学模拟。
性能指标、数据泄漏、过拟合、结构文件质量和训练构型覆盖都会影响模型结论的可靠性。

\subsection*{习题}

\begin{enumerate}[leftmargin=*]
\item 比较描述符模型、图神经网络和机器学习势的适用任务。
\item 机器学习势适合哪些模拟场景？训练覆盖为何重要？
\item 为高熵合金涂层硬度预测设计输入特征，并说明依据。
\end{enumerate}

\clearpage
\section{模型解释与知识嵌入}


\begin{sectionkeypoints}
\item 可解释性分析关注模型训练之后的问题，即模型主要依赖哪些特征、这些特征如何影响预测结果，以及解释是否符合材料机理。
\item 特征重要性和 SHAP\footnote{SHapley Additive exPlanations，一种基于博弈论的可解释性分析方法。} 可帮助把机器学习输出转化为材料语言，但它们反映的是当前数据和当前模型中的关系，不等同于材料真理。
\item 物理知识嵌入强调在建模前和建模过程中主动引入材料规律，包括特征选择、数据筛选、目标泄漏检查和结果合理性判断。
\item 对高熵合金涂层硬度预测而言，原子半径差、晶粒尺寸、硬质相和孔隙率等特征应与固溶强化、细晶强化、第二相强化和组织致密化联系起来理解。
\item 一个可靠的材料人工智能模型，不仅要有较小预测误差，还要能被实验表征、物理机理和材料常识共同支持。
\end{sectionkeypoints}

\subsection{模型解释与特征归因}

前面几节已经介绍了材料性能预测的基本方法。通过描述符和机器学习模型，研究者可根据材料成分、结构和组织信息预测硬度、带隙、形成能、熔点等性能。对于工程应用来说，预测准确当然重要；但对于材料科学研究来说，仅仅得到一个预测数值还不够。研究者还需要进一步知道：模型为什么给出这样的预测？它主要依据了哪些材料特征？这些特征是否符合已有材料机理？


这正是模型可解释性要解决的问题。可解释性关注的是模型在预测时主要依赖哪些输入特征，以及这些特征如何影响预测结果。对于材料科学而言，可解释性尤其重要，因为材料研究的目标不仅是预测某个样品性能，更是提炼可理解、可验证、可迁移的材料规律。


本节继续围绕前面的高熵合金涂层硬度预测任务展开。假设研究者已经训练了一个随机森林模型，用来预测不同高熵合金涂层的硬度。模型的输入特征包括原子半径差、平均电负性、电负性差、价电子浓度、是否含硬质相和平均晶粒尺寸，输出目标是硬度/HV。现在的问题不再是“模型能不能预测硬度”，而是“模型主要根据什么判断某个涂层硬度较高”。


如果只看预测误差，例如 MAE、RMSE 和 R\textsuperscript{2}，研究者只能知道模型预测效果大致如何，却不知道模型内部学到了什么。例如，一个模型可能在测试集上误差较小，但它可能主要依赖了某个偶然相关的字段，而不是真正反映材料强化机制的特征。这样的模型即使暂时预测准确，也很难让研究者放心。因此，在材料人工智能中，模型解释往往和模型评价同样重要。


对于随机森林这类模型，最常见的解释方法是特征重要性分析。特征重要性可粗略告诉研究者：模型在预测硬度时，哪些输入特征贡献更大。假设训练后的模型输出表 \ref{tab:hardness-feature-importance} 所示结果。


\begin{formalTable}[htbp]
\caption{随机森林硬度预测特征重要性示例}
\label{tab:hardness-feature-importance}
\begin{tabularx}{\textwidth}{@{}YY@{}}
\toprule
\textbf{特征} & \textbf{重要性} \\
\midrule
是否含硬质相 & 0.34 \\
平均晶粒尺寸 & 0.26 \\
原子半径差 & 0.18 \\
价电子浓度 & 0.13 \\
电负性差 & 0.06 \\
平均电负性 & 0.03 \\
\bottomrule
\end{tabularx}
\end{formalTable}

这个结果说明，在当前数据集和模型中，“是否含硬质相”和“平均晶粒尺寸”对硬度预测影响最大。接下来，研究者需要把这个模型结果转化为材料语言。含硬质相的涂层往往可能通过第二相强化提高硬度；晶粒尺寸越小，晶界数量越多，可能带来细晶强化；原子半径差较大时，晶格畸变增强，也可能产生固溶强化。因此，如果模型认为这些特征重要，并且这种排序与材料强化机制基本一致，就说明模型结果具有一定可解释性。


但是，特征重要性只能告诉研究者“哪些特征重要”，不宜直接告诉研究者“这个特征让硬度升高还是降低”。例如，平均晶粒尺寸重要，并不意味着晶粒尺寸越大硬度越高，也不意味着晶粒尺寸越小一定硬度越高。它只能说明模型在预测时经常用到这个特征。为了进一步分析特征对预测结果的正负影响，可引入 SHAP 方法。


SHAP 可理解为一种更细致的解释工具。它不仅能分析某个特征整体上是否重要，还可分析在某一个样品中，某个特征是把预测硬度往高处推，还是往低处拉。例如，对于一个含 TiC 硬质相、晶粒较细的高熵合金涂层样品，SHAP 分析可能显示：“是否含硬质相”对预测硬度有正贡献，“平均晶粒尺寸较小”也对预测硬度有正贡献。这种结果就可和第二相强化、细晶强化联系起来解释。


为了更清楚地说明可解释性分析，以下给出一个精简代码框架。假设研究者已经在第 \ref{subsec:descriptor-hardness-prediction} 节中训练好了随机森林硬度预测模型 model，并且已经得到输入特征 X\_\allowbreak{}train、X\_\allowbreak{}test 和特征名称 feature\_\allowbreak{}cols。以下代码用于输出随机森林特征重要性，并进一步使用 SHAP 分析模型预测结果。


\begin{lstlisting}[language=Python,escapeinside={(*@}{@*)}]
# (*@\textnormal{1. 随机森林特征重要性}@*)
importance = pd.DataFrame({
    "feature": feature_cols,
    "importance": model.feature_importances_
}).sort_values(by="importance", ascending=False)
print(importance)
# (*@\textnormal{2. 绘制特征重要性柱状图}@*)
plt.figure(figsize=(6, 4))
plt.barh(importance["feature"], importance["importance"])
plt.xlabel("Feature importance")
plt.ylabel("Feature")
plt.title("Feature importance for hardness prediction")
plt.gca().invert_yaxis()
plt.tight_layout()
plt.show()
# (*@\textnormal{3. 使用 SHAP 解释模型}@*)
# (*@\textnormal{需要先安装 shap: pip install shap}@*)
import shap
explainer = shap.TreeExplainer(model)
shap_values = explainer.shap_values(X_test)
# (*@\textnormal{4. 绘制 SHAP 总结图}@*)
shap.summary_plot(
    shap_values,
    X_test,
    feature_names=feature_cols
)
\end{lstlisting}

这段代码分为两个部分。第一部分使用 model.feature\_\allowbreak{}importances\_\allowbreak{} 输出随机森林自身的特征重要性，并画出柱状图。这个图可帮助研究者快速判断模型最关注哪些特征。第二部分使用 SHAP 对模型进行解释，shap.summary\_\allowbreak{}plot() 可显示每个特征对预测结果的整体影响。


可解释性分析中的 AI 辅助，应服务于“把模型输出转化为材料语言”，而不是替代机理判断。提示词中需要同时给出模型结果和要求它说明证据限制。可按下面方式组织问题。

\begin{samepage}
\begin{userbox}
我用随机森林预测高熵合金涂层硬度，模型输出的特征重要性排序为：是否含硬质相 0.34，平均晶粒尺寸 0.26，原子半径差 0.18，价电子浓度 0.13，电负性差 0.06，平均电负性 0.03。SHAP 图显示，含硬质相通常对预测硬度有正贡献，晶粒尺寸较大时常出现负贡献。

请帮我把这些模型解释结果转化为材料语言。

要求如下：
\begin{enumerate}
\item 分别联系第二相强化、细晶强化、固溶强化或相稳定性解释这些特征；
\item 说明这些解释需要哪些实验表征支持，例如 XRD、SEM、EDS\footnote{Energy Dispersive Spectroscopy，能谱分析。} 或硬度测试；
\item 提醒特征重要性和 SHAP 只能说明当前模型中的关系，不宜直接证明因果关系；
\item 如果模型把样品编号、数据来源等无材料意义字段排在前面，也请说明可能存在什么问题。
\end{enumerate}
\end{userbox}

\begin{aibox}
可这样理解这组结果：

\begin{enumerate}
\item “是否含硬质相”重要性最高，说明模型在预测硬度时较多依赖硬质相信息。若涂层中确实存在 TiC、碳化物或硼化物等硬质相，可与第二相强化联系起来解释。
\item “平均晶粒尺寸”重要性较高，且较大晶粒尺寸对应负向 SHAP 贡献，可能说明模型学习到了晶粒细化有利于提高硬度的趋势，可与细晶强化机制联系起来。
\item “原子半径差”有一定重要性，可与多主元合金中的晶格畸变和固溶强化联系起来，但需要结合相结构和成分分布判断。
\item “价电子浓度”可能与 FCC/BCC 相稳定性有关，但它对硬度的影响通常需要通过相结构、组织形貌和其他描述符共同解释。
\end{enumerate}

这些解释仍然只是模型层面的线索。若要形成可靠材料结论，需要用 XRD 确认物相，用 SEM/EDS 观察硬质相和成分分布，用晶粒尺寸统计验证细晶强化趋势，并检查硬度测试条件是否一致。如果无材料意义的字段在模型中排在前面，可能说明数据整理、特征选择或数据划分存在问题。
\end{aibox}
\end{samepage}

模型解释并不等于材料真理。特征重要性高，只说明模型在当前数据集中较多依赖该特征；SHAP 值为正，也只说明在当前模型中该特征对预测值有正向贡献。这些结果可能受到数据量、数据分布、特征相关性和模型类型影响。例如，如果所有含硬质相的样品刚好也具有较小晶粒尺寸，模型可能难以区分到底是硬质相起主要作用，还是晶粒细化起主要作用。因此，可解释性分析必须与实验表征和材料机理结合，不能孤立理解。


对于高熵合金涂层硬度预测，可按照以下思路解读模型解释结果。首先看特征重要性，判断模型主要关注哪些因素。其次看 SHAP 分析，判断这些因素对预测硬度的正负影响。然后把模型结果与材料强化机制对应起来：硬质相对应第二相强化，晶粒尺寸对应细晶强化，原子半径差对应晶格畸变和固溶强化，价电子浓度可能与相结构稳定性有关。最后，再回到实验数据中检查这些解释是否被组织形貌、物相分析和硬度测试支持。


例如，如果模型认为“是否含硬质相”最重要，而 SEM 和 XRD 也显示高硬度样品中确实存在较多 TiC 等硬质相，那么模型解释与实验观察相互支持。相反，如果模型认为某个表面上无关的字段最重要，例如样品编号或数据来源，那么就需要警惕模型可能学到了错误关联。这种情况并不是模型真正理解了材料规律，而可能是数据整理或特征设计存在问题。


除了特征重要性和 SHAP，还有一些更偏研究生层次的方法可用于材料模型解释。例如，符号回归和 SISSO\footnote{Sure Independence Screening and Sparsifying Operator，一种用于筛选描述符并构建稀疏表达式的方法。} 试图从数据中寻找简洁的数学表达式或组合描述符，使模型结果更接近材料专业人员能够理解的公式关系。对于入门阶段的研究者来说，可先掌握特征重要性和 SHAP 的基本思想；对于研究生，可进一步了解这些能够自动发现简洁规律的可解释建模方法。


本节的重点不是让研究者记住某一个解释算法，而是理解材料人工智能中的一个基本原则：模型预测之后，必须回到材料问题本身去解释。一个好的材料机器学习模型，不仅应该预测误差较小，还应该能够帮助研究者理解哪些成分、结构或组织因素可能控制性能。只有当模型结果能够与材料机制相互印证时，它才更可能成为材料研究和材料设计的有效工具。


通过高熵合金涂层硬度预测这个例子可看到，可解释性分析把机器学习结果和材料科学知识连接起来。模型指出哪些特征重要，材料科学知识帮助研究者解释这些特征为什么重要。但最终是否形成可靠规律，仍然需要实验表征、物理机理和专业判断共同验证。


下一节将进一步讨论物理知识嵌入。可解释性更多是在模型训练之后分析“模型为什么这样预测”，而物理知识嵌入则是在建模之前或建模过程中，就主动把材料科学规律加入数据选择、特征设计和模型约束中，使模型尽量避免学习错误关系。


\subsection{物理知识与模型约束}

上一节讨论了模型可解释性，即模型训练完成之后，研究者如何分析模型主要依据哪些特征做出预测，并将这些特征与材料机理联系起来。可解释性强调的是“事后解释”：模型已经训练好了，研究者再分析它为什么这样预测。


但是，在材料科学中，仅仅依靠事后解释还不够。更理想的做法是，在数据整理、特征选择、模型训练和结果判断的过程中，就主动引入材料科学知识，使模型尽量在合理的物理和化学范围内学习规律。这就是物理知识嵌入的基本思想。


物理知识嵌入并不一定是复杂的公式或高级算法。对于入门阶段，可把它理解为：不要把所有能收集到的字段都机械地丢给模型，而是根据材料机理选择有意义的数据和特征；不要让模型学习明显不合理的关系，而是用物理、化学和实验常识约束建模过程。这样得到的模型，通常比完全黑箱的模型更可靠，也更容易被研究者接受。


以前面的高熵合金涂层硬度预测任务为例，研究者已经用原子半径差、平均电负性、电负性差、价电子浓度、是否含硬质相和平均晶粒尺寸等特征预测涂层硬度。接下来需要判断的是：这些特征为什么可作为输入？哪些字段不应该作为输入？数据应该如何筛选，才能让模型学习到更接近材料规律的关系？


首先，物理知识可用于指导特征选择。对于高熵合金涂层硬度来说，硬度提高通常可能来自固溶强化、细晶强化、第二相强化和组织致密化等机制。因此，输入特征应该尽量围绕这些机制设计。


例如，原子半径差可与晶格畸变和固溶强化联系起来。多种元素固溶在同一晶格中时，元素原子半径差异会导致晶格畸变，从而可能阻碍位错运动，提高材料硬度。平均晶粒尺寸可与细晶强化联系起来。晶粒越细，晶界数量越多，位错运动越容易受到阻碍，硬度可能提高。是否含硬质相可与第二相强化联系起来。如果涂层中生成 TiC、碳化物或硼化物等硬质相，这些硬质相可能作为强化相提高涂层硬度。


因此，这些特征不是随便选出来的数字，而是和材料强化机制存在对应关系。模型使用这些特征预测硬度，才有可能学习到具有材料意义的规律。


这种关系可整理为表 \ref{tab:mechanism-feature-mapping}。


\begin{formalTable}[htbp]
\caption{材料机制与候选特征对应关系}
\label{tab:mechanism-feature-mapping}
\begin{tabularx}{\textwidth}{@{}YYY@{}}
\toprule
\textbf{材料机制} & \textbf{可转化的模型特征} & \textbf{对硬度的可能影响} \\
\midrule
固溶强化 & 原子半径差、电负性差、价电子浓度 & 晶格畸变和相稳定性变化可能影响硬度 \\
细晶强化 & 平均晶粒尺寸 & 晶粒细化可能提高硬度 \\
第二相强化 & 是否含硬质相、硬质相含量 & 硬质相可能提高涂层硬度 \\
组织致密化 & 孔隙率、裂纹数量、涂层缺陷 & 孔隙和裂纹可能降低硬度和服役性能 \\
\bottomrule
\end{tabularx}
\end{formalTable}

这个表格体现了物理知识嵌入的第一层含义：把材料机理转化为模型特征。模型输入不是任意字段，而是由材料知识筛选出来的、有明确含义的特征。


在实际建模中，这一步应先由材料机理提出候选特征，再结合已有实验记录判断哪些特征可获得、哪些特征需要补充测试。特征设计不应只看字段是否存在，还要看字段是否能对应到明确的材料机制。


其次，物理知识可帮助排除不合理特征。并不是所有出现在数据表中的字段都适合作为模型输入。例如，样品编号、测试日期、数据来源、实验人员姓名等字段，通常不应该作为材料性能预测特征。它们可能与硬度在数据表中出现某种偶然相关，但这种相关并不具有材料意义。如果模型根据样品编号预测硬度，就说明它可能学到了数据排序或实验批次的偶然关系，而不是材料规律。


前文已经讨论过目标泄漏，这里只强调它与物理知识嵌入的关系：输入字段必须在预测之前可获得，且不能由目标结果推导而来。例如，“硬度等级”“是否高硬度样品”“硬度排序”这类字段不能作为硬度预测输入，因为它们已经包含了目标硬度的信息。研究者需要检查每一个输入字段：它是否具有材料意义？是否直接包含目标结果？如果无法回答这些问题，就不应该轻易放入模型。


第三，物理知识可指导数据筛选和标准化。材料实验数据往往受到测试条件影响。以硬度为例，不同载荷、不同保载时间、不同测试位置可能得到不同硬度值。对于高熵合金涂层，表面硬度、截面硬度和不同层位硬度也可能存在差异。如果把不同测试条件下的硬度值直接混合到一个数据集中，模型可能学到混乱关系。


在建立硬度预测模型时，应该尽量保证硬度数据具有可比性。例如，尽量使用相同测试方法、相近测试载荷、相同单位和相同样品位置的数据。如果不同数据来源不可避免，就应在数据表中记录测试条件，必要时把测试条件作为附加字段，或者只选择条件一致的数据进行建模。


这种数据筛选本身就是物理知识嵌入的一部分。它告诉模型：哪些数据可放在一起学习，哪些数据不宜简单混用。对于材料科学来说，数据标准化不是形式工作，而是保证模型学习到真实规律的前提。


检查数据可比性时，应重点核验测试载荷、测试单位、热处理条件、样品位置和数据来源等字段。最终是否保留某条数据，仍然需要研究者根据实验条件和研究目标判断。


第四，物理知识可作为模型结果的合理性检查。模型训练完成后，如果预测结果明显违背材料常识，就需要警惕。例如，如果模型预测孔隙率越高硬度越高，或者晶粒尺寸越大硬度越高，就需要进一步检查数据是否存在偏差、特征是否相关、样本数量是否太少，或者模型是否受异常样品影响。当然，材料规律并不总是简单单调关系，但当模型结果与常识明显冲突时，必须进行核查，而不是直接接受。


对于高熵合金涂层硬度预测，可设定一些基本的合理性检查。例如，含硬质相的样品不一定总是硬度最高，但如果模型完全认为硬质相无关，就需要检查硬质相字段是否准确；晶粒细化通常有利于硬度提高，如果模型得出相反趋势，也要检查晶粒尺寸统计是否一致；孔隙率增加通常不利于力学性能，如果模型认为孔隙率越高越好，就需要检查数据是否存在偶然相关或测量问题。


第五，物理知识还可用于更高级的模型约束。在入门阶段，研究者主要通过特征选择、数据筛选和结果检查嵌入材料知识。对于更深入的研究，还可把物理规律直接写进模型或损失函数中。例如，在原子尺度模型中，可要求模型满足平移、旋转和置换对称性；在热力学相关任务中，可引入稳定性约束；在力学模型中，可要求某些预测量满足基本物理关系。这类方法通常称为物理约束机器学习或物理知识嵌入机器学习。


对于本章而言，不需要展开复杂数学形式，但要让研究者理解一个基本思想：材料模型不应只追求拟合数据，还要尽量尊重材料科学规律。一个模型如果预测误差很小，但违背基本物理常识，或者依赖无法解释的虚假相关字段，那么它在材料研究中的价值是有限的。


可用如下流程概括物理知识嵌入在高熵合金涂层硬度预测中的作用：


\begin{lightquote}
\centering 研究目标 → 材料机理分析 → 特征选择 → 数据筛选 → 模型训练 → 结果解释 → 实验和机理验证
\end{lightquote}

其中，材料机理分析不是最后才出现，而是从一开始就参与建模。研究者首先根据固溶强化、细晶强化、第二相强化等机制选择特征；再根据测试条件筛选可比数据；然后训练模型；最后再用实验表征和材料机理检查模型解释是否合理。


为了帮助研究者理解这个流程，可设计一个简单的操作任务。假设已经有一张高熵合金涂层硬度数据表，其中包含以下字段：样品编号、化学成分、原子半径差、电负性差、价电子浓度、是否含硬质相、晶粒尺寸、孔隙率、硬度等级、硬度/HV、测试载荷、数据来源。研究者需要判断哪些字段可作为输入特征，哪些字段应该保留但不作为模型输入，哪些字段必须删除或谨慎使用。


按照物理知识嵌入的原则，可这样处理：原子半径差、电负性差、价电子浓度、是否含硬质相、晶粒尺寸和孔隙率可作为候选输入特征，因为它们具有材料意义；测试载荷和数据来源应保留，用于判断数据是否可比，但不一定直接作为性能特征；样品编号通常不应作为输入；硬度等级不能作为输入，因为它可能由硬度结果推导而来，存在目标泄漏风险；硬度/HV 是预测目标 y，不应同时作为输入特征 X。


可整理为表 \ref{tab:physics-informed-field-treatment}。


\begin{formalTable}[htbp]
\caption{物理知识嵌入建模过程中的字段处理}
\label{tab:physics-informed-field-treatment}
\begin{tabularx}{\textwidth}{@{}YYY@{}}
\toprule
\textbf{字段} & \textbf{是否作为模型输入} & \textbf{原因} \\
\midrule
原子半径差 & 可 & 与晶格畸变和固溶强化有关 \\
电负性差 & 可 & 与元素化学差异和相形成有关 \\
价电子浓度 & 可 & 可能影响相结构稳定性 \\
是否含硬质相 & 可 & 对应第二相强化 \\
晶粒尺寸 & 可 & 对应细晶强化 \\
孔隙率 & 可 & 反映组织致密性 \\
测试载荷 & 谨慎使用或用于筛选 & 影响硬度可比性 \\
数据来源 & 不作为材料特征 & 用于追溯和分组检查 \\
样品编号 & 不使用 & 无材料意义 \\
硬度等级 & 不使用 & 可能造成目标泄漏 \\
硬度/HV & 目标值 & 模型要预测的 y \\
\bottomrule
\end{tabularx}
\end{formalTable}

这类字段筛选任务适合用提示词形式进行辅助检查。提示词应要求同时判断材料意义、数据可比性和目标泄漏风险，而不是只问“哪些列可输入模型”。

\begin{samepage}
\begin{userbox}
我有一张高熵合金涂层硬度数据表，字段包括：样品编号、化学成分、原子半径差、电负性差、价电子浓度、是否含硬质相、晶粒尺寸、孔隙率、硬度等级、硬度/HV、测试载荷、数据来源。请帮我判断哪些字段适合作为模型输入，哪些字段应保留但不作为输入，哪些字段可能造成目标泄漏或虚假相关。

要求如下：
\begin{enumerate}
\item 按固溶强化、细晶强化、第二相强化和组织致密化解释可用特征；
\item 指出哪些字段用于追溯或筛选数据，但不宜作为材料性能特征；
\item 明确指出硬度/HV 和硬度等级为什么不能作为输入特征；
\item 提醒哪些测试条件会影响硬度数据可比性；
\item 不要只给机器学习角度判断，要结合材料机理说明理由。
\end{enumerate}
\end{userbox}

\begin{aibox}
可按三类处理这些字段：

\begin{enumerate}
\item 可作为候选输入特征的字段：原子半径差、电负性差、价电子浓度、是否含硬质相、晶粒尺寸、孔隙率。这些字段分别可与固溶强化、相结构稳定性、第二相强化、细晶强化和组织致密化联系起来。
\item 应保留但不宜直接作为材料性能特征的字段：测试载荷、数据来源、化学成分和样品编号。测试载荷和数据来源主要用于判断数据可比性和可追溯性；样品编号通常没有材料意义；化学成分可保留，但进入模型前通常需要转化为成分描述符。
\item 不能作为输入特征的字段：硬度/HV 是预测目标，硬度等级往往由硬度结果划分而来，二者如果进入输入特征，会造成目标泄漏。
\end{enumerate}

此外，硬度数据是否可比还取决于测试载荷、保载时间、测试位置、样品表面状态和测试方法。如果这些条件不一致，应先进行筛选、分组或记录为元数据，再决定是否合并建模。
\end{aibox}
\end{samepage}

需要强调的是，物理知识嵌入并不意味着模型一定更复杂。很多时候，它体现为更谨慎的数据选择、更合理的特征设计和更严格的结果检查。对于材料人工智能入门阶段的研究者来说，最重要的不是马上构建复杂的物理约束模型，而是养成一种意识：建模前先问材料问题，建模中保留物理意义，建模后回到实验和机理验证。


本节说明了物理知识嵌入的基本思想。可解释性帮助研究者在模型训练后理解模型，物理知识嵌入则帮助研究者在建模前和建模过程中约束模型。二者共同说明：人工智能用于材料科学，并不是把数据机械地交给模型，而是让数据驱动方法与材料科学知识相互配合。至此，第\ref{chap:materials-cognition-prediction}章完成了从材料数据、材料表示、性能预测、原子尺度模拟到模型解释的完整链条。下一章将进一步讨论逆向设计问题，即如何根据目标性能反向寻找和设计新材料。

\subsubsection*{\textbf{小结}}

本节说明模型解释和物理知识嵌入共同服务于材料判断。
特征重要性和 SHAP 可提示模型依据，但不能直接证明因果机理，仍需 XRD、SEM、EDS 等证据支持。
物理知识应贯穿数据筛选、特征选择、模型训练和结果检查，避免模型学习虚假相关。

\subsection*{习题}

\begin{enumerate}[leftmargin=*]
\item SHAP 为什么不能直接证明材料机理？还需哪些表征？
\item 可解释性分析如何辅助材料筛选？请举例说明。
\item 只有 FCC 数据时，能否预测 BCC 合金？请说明风险和验证方案。
\end{enumerate}

\section*{本章小结}

本章围绕材料认知与性能预测展开，形成了从数据到模型、再从模型回到材料问题的基本链条。材料数据整理强调字段完整、来源可追溯和条件可比较；材料表示强调把化学式、结构和局域环境转化为描述符或图结构；性能预测部分说明了描述符模型、图神经网络和机器学习势在不同尺度任务中的作用；模型解释与物理知识嵌入则进一步说明，模型预测结果必须经过材料机理、实验表征和物理约束的共同检验。由此可见，人工智能在材料科学中的价值不在于替代实验和理论，而在于提高数据利用效率、缩小候选范围并辅助形成可验证的材料判断。

\chapter{材料设计与智能发现}
\label{chap:materials-design-discovery}


第\ref{chap:materials-cognition-prediction}章讨论正向预测，即在材料成分、结构或实验条件已知的情况下，预测材料可能具有的性能。本章进一步讨论逆向设计与智能发现：当目标性能已经给定时，如何从大量候选材料中筛选出值得验证的方案。这个过程并不是简单地让模型“生成答案”，而是要把目标性能、稳定性、可合成性、成本、安全性和实验条件共同纳入考虑。

本章按照材料发现的一般流程展开，重点解决四个问题。第一节说明如何把应用需求转化为材料设计目标、候选变量和约束边界。第二节讨论候选筛选和主动学习如何帮助研究者逐层缩小搜索空间，并安排下一轮实验或计算。第三节介绍生成设计如何提出库外候选，以及生成候选为什么必须经过有效性、稳定性、新颖性和可合成性验证。第四节讨论可信发现、方法选择和闭环迭代，强调模型推荐必须与计算验证、实验反馈和应用约束共同使用。通过本章学习，研究者应理解人工智能如何参与材料设计决策，同时认识候选材料最终仍需要通过物理判断、计算验证和实验结果逐步确认。




\clearpage
\section{问题转换与设计边界}


\begin{sectionkeypoints}
\item 正向预测回答“已知材料可能具有什么性质”，逆向设计回答“为了获得目标性质应当寻找什么材料”。
\item 逆向设计本质上是受约束搜索，不是单纯追求最高预测分数，而是在稳定性、可合成性、成本、安全和设备条件内寻找可验证候选。
\item 材料设计对象可能位于成分、原子结构、微结构、工艺和服役环境等不同层级，层级选择必须与研究目标一致。
\item 候选筛选、主动学习和生成设计是本章三条核心路线，实际研究中常常组合使用。
\end{sectionkeypoints}

\subsection{正向预测与逆向设计}

第\ref{chap:materials-cognition-prediction}章讨论的是正向问题：给定材料的成分、结构、工艺或表征信息，预测形成能、带隙、硬度、折射率或熔点等性质。第\ref{chap:materials-design-discovery}章把问题方向反过来：当目标性质已经确定时，寻找哪些材料、配方、结构或工艺值得进一步尝试。前者回答“这种材料具有什么性质”，后者回答“具有目标性质的材料可能在哪里”。


材料科学长期依赖经验规则和实验试错，但多组元合金、复杂晶体和玻璃配方的候选数量远超实验能够逐一验证的范围。逆向设计的作用不是取消经验和实验，而是把大规模穷举转化为有方向的搜索：先明确目标和约束，再用模型提出或排序候选，最后通过计算和实验逐步确认。


从数学上看，逆向设计可理解为受约束优化。材料对象 x 可以是化学式、晶体结构、玻璃配方、微结构参数或制备条件；正向模型 f(x) 估计目标性能；约束 g(x) 描述稳定性、可合成性、成本、安全和设备边界。搜索的目的不是让某个预测分数无限增大，而是在满足约束的前提下找到一组可验证的候选。正向预测、逆向设计与闭环发现之间的关系见表 \ref{tab:forward-inverse-closed-loop}。


\begin{formalTable}[htbp]
\caption{正向预测、逆向设计与闭环发现}
\label{tab:forward-inverse-closed-loop}
\footnotesize
\begin{tabularx}{\textwidth}{@{}YYYYY@{}}
\toprule
\textbf{问题类型} & \textbf{输入} & \textbf{输出} & \textbf{材料例子} & \textbf{主要风险} \\
\midrule
正向预测 & 候选材料 x & 性质 f(x) & 由晶体结构预测带隙 & 训练集外误差被低估 \\
逆向设计 & 目标性质 y & 候选材料 x & 按目标折射率设计玻璃 & 候选不可合成或不稳定 \\
闭环发现 & 已有数据 + 目标 & 下一步实验 & 主动学习选择合金配比 & 实验噪声导致错误反馈 \\
\bottomrule
\end{tabularx}
\end{formalTable}

\subsubsection*{\textbf{1. 问题转换：从应用目标到材料任务}}

逆向设计的第一步，不是直接选择模型，而是把应用语言转化为材料科学语言。应用目标往往比较笼统，例如“希望玻璃折射率更高”“希望发光材料颜色更接近蓝光”“希望涂层在高温下更稳定”。这些表述还不宜直接进入模型，因为模型需要明确的输入变量、目标性能和约束条件。只有把应用目标拆解成可计算、可测量、可验证的材料任务，后续的筛选、优化和生成才有依据。

以 AR\footnote{Augmented Reality，增强现实。} 光波导玻璃为例，“高折射率”只是一个目标方向。真正进入设计流程时，还需要进一步明确折射率测量波长、透过率范围、色散要求、玻璃化转变温度（Tg\footnote{Glass transition temperature，玻璃化转变温度。}）、热膨胀系数、密度、原料安全性和熔制温度窗口。如果只追求折射率，模型可能推荐含高极化率重金属氧化物较多的配方，但这类配方可能带来密度升高、成本增加、透过率下降或环保风险。因此，逆向设计不是单目标求最大值，而是多目标、多约束条件下的候选选择。

其他材料任务也需要类似拆解。例如 microLED\footnote{Micro Light-Emitting Diode，微型发光二极管。} 发光材料可把目标颜色转化为发光波长或带隙范围，但仍需补充缺陷态、载流子复合效率、热稳定性和外延匹配等约束。这里不展开为新的主线案例，只用它说明：目标性能通常只是材料设计任务的一部分。

具体而言，问题转换通常包括四个问题：目标性能是什么？候选变量是什么？约束条件是什么？验证方式是什么？目标性能决定模型要预测什么；候选变量决定模型可搜索哪些材料空间；约束条件决定哪些候选必须排除；验证方式决定候选是否能从模型结果走向真实材料。表中的 DFT\footnote{Density Functional Theory，密度泛函理论。} 指密度泛函理论计算。典型应用目标到材料设计任务的转换见表 \ref{tab:application-to-design-task}。

\begin{formalTable}[htbp]
\caption{应用目标到材料设计任务的转换}
\label{tab:application-to-design-task}
\footnotesize
\begin{tabularx}{\textwidth}{@{}YYYYY@{}}
\toprule
\textbf{应用目标} & \textbf{材料任务} & \textbf{候选变量} & \textbf{关键约束} & \textbf{验证方式} \\
\midrule
高折射率玻璃 & 预测并优化折射率、Tg、热膨胀系数 & 氧化物组成比例 & 玻璃形成、透明性、熔制温度 & 熔制实验、折射率和透过率测试 \\
蓝光发光材料 & 筛选合适带隙和稳定结构 & 元素组成、晶体结构、掺杂 & 缺陷、外延兼容、热稳定 & DFT、发光测试、器件验证 \\
高硬度涂层 & 优化成分、硬质相和组织参数 & 合金成分、工艺、相组成 & 裂纹、孔隙、结合强度 & 显微硬度、SEM/XRD、服役测试 \\
超高温陶瓷 & 寻找高熔点和抗氧化候选 & 化学组成、晶体结构 & 高温稳定、氧化行为、可烧结性 & 高温实验、热分析、模拟验证 \\
\bottomrule
\end{tabularx}
\end{formalTable}

这一转换过程适合借助 AI 辅助梳理，但不能让 AI 直接给出最终材料结论。提示词应要求它把应用目标拆成材料变量、目标性能、约束条件和验证方式，而不是只列出几个候选材料名称。

\begin{samepage}
\begin{userbox}
我想设计一种用于 AR 光波导的高折射率玻璃。请帮我把这个应用目标转化为材料逆向设计任务。

要求如下：
\begin{enumerate}
\item 明确目标性能，例如折射率、透过率、色散、Tg、热膨胀系数等；
\item 说明候选变量可包括哪些氧化物组成或配方参数；
\item 区分必须满足的硬约束和可折中的软目标；
\item 给出候选筛选后需要进行哪些计算或实验验证；
\item 不要直接编造最终玻璃配方，只做任务拆解。
\end{enumerate}
\end{userbox}

\begin{aibox}
可将“高折射率 AR 光波导玻璃”拆解为以下设计任务：

\begin{enumerate}
\item \textbf{目标性能：}折射率是核心目标，但还应同时记录透过率、色散、Tg、热膨胀系数、密度和加工温度等指标。
\item \textbf{候选变量：}可把各氧化物组分比例作为设计变量，例如网络形成体、网络修饰体和提高折射率的氧化物比例。
\item \textbf{硬约束：}配方比例之和应满足组成约束；原料不能超出安全和环保边界；熔制温度应在设备可承受范围内；样品应能够形成玻璃而不是明显析晶。
\item \textbf{软目标：}折射率越高越好，但不能以明显牺牲透明性、热稳定性和加工性为代价。
\item 验证方式：模型筛选后，需要进行玻璃形成能力判断、熔制实验、折射率测试、透过率测试和热膨胀系数测试。
\end{enumerate}

这样处理后，逆向设计任务不再只是“寻找高折射率材料”，而是“在可熔制、可成玻、透明且热性质合适的约束下，寻找高折射率玻璃配方”。
\end{aibox}
\end{samepage}

问题转换应当在建模之前完成，而不是在模型给出候选之后再补充解释。如果目标一开始没有定义清楚，后面所有筛选和排序都会变得模糊。例如，若只写“设计高性能玻璃”，模型无法判断“高性能”指的是高折射率、高透过率、低热膨胀系数还是高热稳定性；若只写“设计蓝光材料”，模型也不知道重点是发光波长、发光效率、热稳定性还是外延兼容性。本章强调问题转换，是为了让研究者意识到：人工智能材料设计首先是材料问题定义，其次才是模型选择。

一个规范的逆向设计任务，宜用一句话概括为“在什么约束下，寻找什么类型的材料，使什么目标性能达到什么范围”。例如，“在不使用高风险元素、熔制温度不超过设备上限、能够形成透明玻璃的约束下，寻找折射率较高且热膨胀系数合适的氧化物玻璃配方”。这样的表述比“让模型设计高折射率玻璃”更清楚，也更便于后续建立数据表、选择描述符和设计验证实验。

\subsection{设计对象与约束边界}

\subsubsection*{\textbf{1. 设计边界：材料对象分层}}

逆向设计并不总是直接生成一个化学式。材料对象通常包含多个层级：成分决定由哪些元素或组分构成；原子结构决定空间群、占位、缺陷和局域配位；微结构包含晶粒、孔隙、第二相和界面；工艺包含熔融、退火、沉积、烧结和冷却条件。同一成分经过不同工艺可得到不同结构和性能，因此设计对象的层级必须与研究目标相匹配。材料逆向设计常见对象层级见表 \ref{tab:inverse-design-object-levels}。


\begin{formalTable}[htbp]
\caption{材料逆向设计的对象层级}
\label{tab:inverse-design-object-levels}
\begin{tabularx}{\textwidth}{@{}YYYY@{}}
\toprule
\textbf{设计层级} & \textbf{常见变量} & \textbf{适用任务} & \textbf{容易忽略的问题} \\
\midrule
成分 & 元素/氧化物比例、掺杂量 & 玻璃配方、合金筛选、催化剂组分 & 相形成与工艺窗口 \\
原子结构 & 晶格、坐标、缺陷、表面 & 带隙、离子迁移、稳定性 & 对称性和局域环境 \\
微结构 & 晶粒、孔隙、第二相、取向 & 力学、热学、疲劳与可靠性 & 制备历史强相关 \\
工艺 & 温度、气氛、压力、时间 & 可合成性与性能优化 & 元数据缺失会误导模型 \\
服役环境 & 光、电、热、湿度、载荷 & 器件可靠性和寿命 & 实验条件与应用条件不一致 \\
\bottomrule
\end{tabularx}
\end{formalTable}

除了设计层级，还需要区分约束的类型。材料设计中的约束可分为硬约束和软约束。硬约束是必须满足的条件，例如有毒元素不能使用、设备最高温度不能超过、配方比例必须归一化、候选结构不能明显违反价态和原子距离规则。软约束是希望尽量优化但可折中的条件，例如成本越低越好、密度越低越好、热膨胀系数越接近目标越好。硬约束用于排除候选，软约束用于排序候选。

如果不区分硬约束和软约束，逆向设计很容易变成一个不清楚的打分问题。例如，把折射率、成本、热稳定性和环保要求都压缩成一个总分，表面上方便排序，但不同目标之间的取舍会被隐藏。逆向设计应把硬约束、软目标和数据边界分开说明。

设计边界还包括数据边界。模型只能在训练数据覆盖的范围内比较可靠地工作。如果训练数据主要来自硅酸盐玻璃，那么把模型直接用于完全不同的硫系玻璃或氟化物玻璃，预测结果就可能不稳定。如果模型只学习了室温性质，就不宜直接推断高温服役行为。如果数据没有记录制备工艺，模型也很难判断同一成分在不同热处理条件下的性能差异。因此，逆向设计开始之前必须说明数据从哪里来、覆盖哪些材料体系、哪些变量没有记录。不同边界类型及其处理方式见表 \ref{tab:inverse-design-boundary-types}。

\begin{formalTable}[htbp]
\caption{材料逆向设计中的边界类型}
\label{tab:inverse-design-boundary-types}
\footnotesize
\begin{tabularx}{\textwidth}{@{}YYYY@{}}
\toprule
\textbf{边界类型} & \textbf{含义} & \textbf{材料例子} & \textbf{处理方式} \\
\midrule
硬约束 & 不满足就直接排除 & 有毒元素、不可成玻、设备温度超限 & 筛选前写入规则 \\
软目标 & 可折中但要排序 & 折射率、密度、成本、热膨胀系数 & 作为多目标评价项 \\
数据边界 & 模型训练数据覆盖范围 & 只含某类玻璃或某类晶体 & 明确适用范围和外推风险 \\
工艺边界 & 实际制备条件限制 & 熔制温度、气氛、烧结时间 & 作为实验可行性约束 \\
\bottomrule
\end{tabularx}
\end{formalTable}

在教学中，可让研究者首先写出设计边界，再选择模型。这样的顺序比“先选一个算法再想怎么用”更合理。算法只能在已经定义清楚的问题上发挥作用；如果目标、变量和边界本身模糊，模型输出再复杂也难以转化为真实材料方案。

这一类边界梳理同样可使用 AI 辅助，但提示词必须要求区分硬约束、软目标和数据边界。

\begin{samepage}
\begin{userbox}
我准备做一个材料逆向设计任务，目标是寻找候选高折射率玻璃。请帮我列出这个任务的设计边界。

要求如下：
\begin{enumerate}
\item 区分硬约束、软目标、数据边界和工艺边界；
\item 每一类边界都要说明它为什么重要；
\item 不要把所有条件合并成一个总分；
\item 说明哪些信息需要在实验记录或数据库中保留；
\item 不要直接给出候选配方。
\end{enumerate}
\end{userbox}

\begin{aibox}
可从四类边界整理：

\begin{enumerate}
\item \textbf{硬约束：}原料安全、环保限制、组成比例合法、可形成玻璃、设备温度可达。这类条件不满足时，候选应直接排除。
\item \textbf{软目标：}折射率、透过率、密度、热膨胀系数、成本和加工性能。这些目标之间可能存在折中，应分别记录和排序。
\item \textbf{数据边界：}需要说明训练数据主要来自哪些玻璃体系、性质是在什么条件下测得、是否包含失败样品和析晶样品。超出数据覆盖范围时，模型预测应谨慎使用。
\item \textbf{工艺边界：}熔制温度、保温时间、冷却方式、气氛、坩埚材料等都会影响玻璃形成和性能，应作为实验元数据保存。
\end{enumerate}

这些边界不是附加说明，而是逆向设计任务的一部分。只有先限定边界，模型推荐的候选才有可能进入真实实验。
\end{aibox}
\end{samepage}

边界条件还应随研究阶段动态更新。早期阶段可设置较宽的组成范围，用来观察模型是否能发现潜在趋势；进入实验阶段后，边界应逐渐收紧，例如限定可采购原料、设备温度、样品尺寸和测试条件；进入应用阶段后，还要加入成本、安全、长期稳定性和批量制备要求。换言之，边界不是一次写完就不变，而是随着数据积累和实验反馈不断修正。

对于研究者来说，最容易忽略的是“不可见边界”。有些限制并不直接写在数据表中，例如某些配方虽然数据库中有记录，但实验室没有相应原料；某些结构计算上稳定，但现有设备无法提供所需压力或气氛；某些材料性能很好，但含有不适合应用场景的元素。这些不可见边界如果不提前说明，模型推荐就会表面上合理、实际不可执行。因此，材料逆向设计的边界既包括数据中的字段，也包括实验室条件和应用场景中的真实限制。

本章讨论三条核心路线。候选筛选适合从已有数据库或枚举空间中逐层过滤；主动学习适合实验昂贵、每轮只能验证少量样品的任务；生成设计用于突破已有候选库，在目标条件下提出新结构或新配方。三者并非互相替代，实际研究常按“筛选--优化--生成--验证”的方式组合使用。

\subsubsection*{\textbf{小结}}

本节说明逆向设计不是直接生成材料，而是把应用需求转化为目标、变量和约束。
正向预测回答已知材料的性能，逆向设计回答为了目标性能应寻找哪些候选。
材料对象层级、硬约束、软目标、数据边界和工艺边界需要在建模前明确。

\subsection*{习题}

\begin{enumerate}[leftmargin=*]
\item 正向预测和逆向设计在高折射率玻璃任务中如何配合？
\item 逆向设计为什么必须先明确目标和约束？
\item 如何把高折射率玻璃需求转化为计算目标？
\end{enumerate}

\clearpage
\section{候选筛选与主动学习}


\begin{sectionkeypoints}
\item 候选筛选的基本逻辑是先建立候选库，再按照成本由低到高、可信度由粗到细逐层过滤。
\item 筛选漏斗通常包含规则过滤、机器学习代理模型、高精度计算和实验验证等层级，每一级都应记录通过和排除理由。
\item 主动学习适合实验昂贵、小样本和需要多轮优化的材料任务，核心是同时考虑预测性能和模型不确定性。
\item 实验噪声、批次效应和失败样品都应进入数据记录，否则模型可能把制备差异误认为材料本征规律。
\item 候选排序和主动学习评分只负责辅助决策，最终实验批次仍需要材料约束、工艺可行性和人工复核。
\end{sectionkeypoints}

\subsection{筛选漏斗与分层验证}

高通量虚拟筛选的基本思路，是先建立候选库，再按照成本由低到高、可信度由粗到细逐层过滤。最外层使用元素排除、价态平衡、成分范围、成本或毒性等规则；中间层使用第\ref{chap:materials-cognition-prediction}章建立的正向模型快速估计形成能、带隙、弹性模量或折射率；内层使用 DFT、分子动力学或热力学计算精算；最后才进入实验合成与表征。


\begin{figure}[H]
\centering
\BookFigureImage{images_11/3.png}
\caption{候选筛选与验证层级示意图}
\label{fig:candidate-screening-validation}
\end{figure}
\FloatBarrier

如图 \ref{fig:candidate-screening-validation} 所示，筛选漏斗利用了不同方法之间的成本差异。规则过滤几乎不需要计算，机器学习预测可快速处理大量候选，而高精度计算和真实实验只能用于较小范围。把便宜但粗略的方法放在前面，把昂贵但可靠的方法放在后面，能够避免把有限资源消耗在明显不合理的候选上。


筛选的第一个常见门槛是稳定性。无机晶体可参考形成能和凸包能量，合金和陶瓷还要考虑相分离、氧化和高温相变，玻璃则需要考虑玻璃形成能力、析晶倾向和熔制窗口。计算稳定不等于实验一定能够合成，热力学判断还要与动力学路径、前驱体和设备条件结合。


第二个门槛是多目标要求。AR 光波导玻璃不仅需要较高折射率，还要兼顾透过率、色散、玻璃化转变温度、热膨胀系数、密度和加工性能；microLED 发光材料除带隙外，还要考虑缺陷、载流子复合、热稳定和外延兼容性。相应地，候选排序不应只依据单一预测值，而应保留不同目标之间的折中关系。筛选漏斗中不同层级的作用见表 \ref{tab:screening-funnel-levels}。


\begin{formalTable}[htbp]
\caption{筛选漏斗中不同层级的作用}
\label{tab:screening-funnel-levels}
\begin{tabularx}{\textwidth}{@{}YYYY@{}}
\toprule
\textbf{筛选层级} & \textbf{常见工具} & \textbf{适合回答的问题} & \textbf{学习要点} \\
\midrule
规则层 & 元素价态、毒性、成本、工艺窗口 & 哪些候选明显不可取？ & 规则便宜但粗糙，适合第一轮排除 \\
代理模型层 & 回归/分类模型、图神经网络、集成模型 & 哪些候选可能性能较好？ & 模型用于排序，不等于最终证明 \\
物理计算层 & DFT、相图、分子动力学、热力学模型 & 候选是否稳定、机制是否可信？ & 精度更高但成本也更高 \\
实验层 & 合成、表征、服役测试 & 候选能否成为真实材料？ & 实验验证是材料发现的终点而非附属步骤 \\
\bottomrule
\end{tabularx}
\end{formalTable}

\subsubsection*{\textbf{1. 案例示范：高折射率玻璃候选筛选}}

以高折射率玻璃为例，首先把应用要求转化为可计算指标；然后限定各氧化物比例之和、原料安全和设备温度等边界；代理模型预测折射率、Tg、热膨胀系数和析晶倾向；物理或经验模型检查玻璃形成与黏度窗口；最后制备少量代表性配方并测量真实光学损耗。每一级都应记录排除理由，使候选清单可复核。

具体而言，第一层规则可先排除明显不合理的配方。例如，各氧化物比例之和必须为 100\%，单一组分含量不能超过玻璃形成的经验范围，含有明显不适合教学或工程应用的高风险元素时应直接排除。第二层可用代理模型预测折射率和热性质，把候选从几千个缩小到几十个。第三层再用更可靠的热力学判断或经验模型检查玻璃形成能力、析晶风险和黏度窗口。最后才选择少量配方进入熔制实验。

筛选漏斗的一个重要原则是“保留失败理由”。如果一个候选被排除，应记录它是因为折射率不够、Tg 太低、热膨胀系数过高、含有风险元素，还是因为预测不确定性太大。这样做有两个好处。第一，后续可复核筛选规则是否过于严格；第二，失败信息可反馈给下一轮模型或实验设计。如果只保留最后通过的候选，而删除中间淘汰信息，材料发现过程就会失去可追溯性。

在初学阶段，研究者可将筛选漏斗理解为一个逐层提问的过程：第一问，这个候选是否明显违反基本规则？第二问，它的目标性能是否接近要求？第三问，它是否可能稳定存在或可制备？第四问，它是否值得消耗实验资源？每一层都不是绝对正确的判断，而是在不同成本下尽量减少无效尝试。筛选过程中需要保留的判断信息见表 \ref{tab:screening-funnel-records}。

\begin{formalTable}[htbp]
\caption{筛选漏斗中需要保留的判断信息}
\label{tab:screening-funnel-records}
\footnotesize
\begin{tabularx}{\textwidth}{@{}YYYY@{}}
\toprule
\textbf{筛选问题} & \textbf{排除依据} & \textbf{保留信息} & \textbf{后续用途} \\
\midrule
是否违反基本规则 & 比例错误、风险元素、明显价态异常 & 排除原因和原始配方 & 修正规则或候选生成空间 \\
性能是否接近目标 & 折射率、带隙、硬度等预测值偏离 & 预测值和模型版本 & 比较不同代理模型 \\
是否可能稳定 & 凸包能、析晶倾向、相分离风险 & 稳定性指标和计算方法 & 决定是否进入高精度计算 \\
是否值得实验 & 成本、设备温度、样品数量限制 & 排序分数和实验批次 & 形成可执行实验计划 \\
\bottomrule
\end{tabularx}
\end{formalTable}

这一过程可借助 AI 辅助把筛选规则写成清单，但规则本身必须由材料任务决定。可用如下提示词把流程规范化。

\begin{samepage}
\begin{userbox}
我正在设计一个高折射率玻璃候选筛选漏斗。候选数据表中有氧化物比例、预测折射率、预测 Tg、热膨胀系数、是否含风险元素、预测析晶风险等字段。请帮我把筛选流程整理成逐层规则。

要求如下：
\begin{enumerate}
\item 先列出规则过滤层、代理模型层、物理/经验验证层和实验层；
\item 每一层说明主要排除什么候选；
\item 明确哪些排除理由需要记录；
\item 提醒候选排序不应只看折射率；
\item 不要编造真实配方数据。
\end{enumerate}
\end{userbox}

\begin{aibox}
可将筛选漏斗设计为四层：

\begin{enumerate}
\item \textbf{规则过滤层：}检查氧化物比例是否合理、是否含风险元素、熔制温度是否超出设备范围。被排除的候选应记录具体违反哪条规则。
\item \textbf{代理模型层：}预测折射率、Tg 和热膨胀系数。折射率高但 Tg 过低或热膨胀系数过高的候选不宜直接排在最前。
\item \textbf{物理/经验验证层：}进一步检查玻璃形成能力、析晶倾向、黏度窗口和配方稳定性。此处的目标是排除“预测性能好但难以形成玻璃”的候选。
\item \textbf{实验层：}选择少量代表性配方进行熔制、退火、折射率和透过率测试，并记录失败样品原因。
\end{enumerate}

筛选报告中应保留每一级候选数量、通过率和排除理由。这样可判断筛选规则是否过严，也便于后续主动学习利用失败信息。
\end{aibox}
\end{samepage}

筛选漏斗还有一个容易被忽视的问题：前一级筛选规则会影响后一级能看到什么候选。如果第一层规则过于严格，可能会把一些有潜力但不符合传统经验的候选提前排除；如果第一层规则过于宽松，又会把大量明显不可行的候选留给后续高成本计算。为此，筛选规则需要根据任务目标反复校准，而不是一次确定后不再修改。

在课堂练习中，可让研究者比较两种筛选策略。一种策略偏保守，优先排除风险候选，适合实验资源很有限的任务；另一种策略偏探索，保留更多边界候选，适合希望发现新体系的任务。两种策略没有绝对优劣，关键在于是否说明了筛选目的、资源条件和验证能力。这样可避免把筛选漏斗理解为机械流程，而是理解为材料决策工具。


以下给出一个最小候选排序代码示例。代码把目标带隙、凸包能量和元素风险写成显式评分规则，用于说明筛选假设如何转化为可复核的计算步骤。评分只负责初步排序，不宜代替稳定性复算和实验验证。


\begin{lstlisting}[language=Python,escapeinside={(*@}{@*)}]
import pandas as pd
# (*@\textnormal{gap 的单位为 eV，hull 的单位为 eV/atom}@*)
cands = pd.DataFrame({
    "id": ["mp-1", "mp-2", "mp-3", "mp-4"],
    "formula": ["M1O2", "M2O3", "M3S", "PbM4O"],
    "gap": [2.10, 1.55, 2.35, 2.00],
    "hull": [0.018, 0.006, 0.065, 0.012],
    "elements": ["M1,O", "M2,O", "M3,S", "Pb,M4,O"],
})
target_gap = 2.0
cands["gap_score"] = 1 - (cands["gap"] - target_gap).abs() / target_gap
cands["stability_score"] = (0.08 - cands["hull"]).clip(lower=0) / 0.08
cands["element_penalty"] = cands["elements"].str.contains("Pb|Cd|Hg").astype(int)
cands["priority"] = (
    0.45 * cands["gap_score"]
    + 0.45 * cands["stability_score"]
    - 0.50 * cands["element_penalty"]
)
cols = ["id", "formula", "gap", "hull", "priority"]
print(cands.sort_values("priority", ascending=False)[cols])
\end{lstlisting}

\manualfigcaption{\textbf{代码示例：}候选材料初步排序示例}

\subsection{主动学习与实验决策}

高通量筛选是在固定候选库中做过滤，主动学习则是在搜索过程中不断更新模型。研究者首先用少量实验或计算训练代理模型，模型同时给出候选性能和不确定性，再选择下一批最值得验证的候选；新结果返回后重新训练模型，形成序贯迭代。


\begin{figure}[H]
\centering
\BookFigureImage{images_11/4.png}
\caption{主动学习与生成设计闭环示意图}
\label{fig:active-learning-generative-loop}
\end{figure}
\FloatBarrier

图 \ref{fig:active-learning-generative-loop} 概括了主动学习与生成设计之间的闭环关系。贝叶斯优化是常见的主动学习方法，适合小样本和高成本实验。它通常包含代理模型和采集函数两个部分：代理模型近似真实实验或高精度计算，采集函数综合预测均值与不确定性，决定下一步优先验证哪些候选。其核心不是总选择当前分数最高的样品，而是在“利用”和“探索”之间取得平衡。


利用型候选靠近模型认为的高性能区域，有利于尽快得到改进材料；探索型候选位于模型不确定但物理上仍合理的区域，用于扩展知识边界；重复型或近重复型候选用于估计制备和测量噪声。真实实验批次通常同时包含这几类样品，而不是只给出一个单一的“最佳配方”。主动学习批次中的候选类型见表 \ref{tab:active-learning-batch-types}。


\begin{formalTable}[htbp]
\caption{主动学习批次中的候选类型}
\label{tab:active-learning-batch-types}
\begin{tabularx}{\textwidth}{@{}YYY@{}}
\toprule
\textbf{候选类型} & \textbf{选择依据} & \textbf{在实验轮次中的作用} \\
\midrule
利用型 & 预测性能高且约束满足 & 尽快逼近可用材料 \\
探索型 & 模型不确定但物理上可行 & 扩展模型对未知区域的认识 \\
边界型 & 靠近工艺或相稳定边界 & 识别可制备窗口 \\
重复型 & 已有样品或近似重复配方 & 估计实验噪声和批次误差 \\
\bottomrule
\end{tabularx}
\end{formalTable}

实验噪声和批次效应是主动学习能否落地的关键。同一成分可能因混料、升温速率、保温时间、气氛、坩埚污染或仪器校准而表现不同。如果模型把这些差异全部当成材料本征规律，就会错误更新搜索方向。这意味着，样品批次、设备、制备程序、测量条件和失败原因都应进入数据记录。


失败样品同样提供信息。某类成分若反复析晶、开裂或产生杂相，说明该区域可能超出可制备边界。例如，某轮玻璃熔制实验中如果发现 La\textsubscript{2}O\textsubscript{3} 含量高于 12\% 的配方全部析晶，下一轮模型排序就应把这一区域标记为高风险，并降低相邻候选的推荐优先级。主动学习不应只学习“哪些候选性能高”，还应学习“哪些区域无法稳定制备”。这使下一轮实验决策同时考虑性能、不确定性、工艺风险和已有负结果。


\subsubsection*{\textbf{1. 案例示范：AR 光波导玻璃主动学习}}

AR 光波导玻璃可作为连续案例。初始数据由已有氧化物配方及其折射率、Tg、热膨胀系数和析晶情况构成；代理模型预测候选性质和不确定性；每轮选择若干高性能候选、若干未知区域候选和少量重复样品；实验结果再用于修正模型和玻璃形成边界。这样，模型输出被转化为可执行的实验批次。

主动学习与普通筛选最大的不同，是它把“下一步做什么实验”作为核心问题。普通筛选常常一次性给出候选排序，而主动学习会根据已经完成的实验不断调整模型。第一轮实验后，模型可能发现某个组成区域折射率较高但容易析晶；第二轮就应减少该区域中高风险候选的比例，同时在相邻区域选择少量探索样品。这样，实验不是简单验证模型，而是在持续修正模型。

在真实实验中，一轮实验通常不应只选择预测性能最高的候选。如果全部选择利用型样品，模型可能很快陷入局部区域，对未知空间了解不足；如果全部选择探索型样品，短期内又可能得不到性能改进。为此，一个合理批次通常同时包含高性能候选、模型不确定候选、边界候选和重复样品。重复样品表面上浪费名额，但它可估计实验噪声和批次误差，对判断模型是否真的改进很重要。

主动学习还需要记录失败类型。对于玻璃配方来说，失败可能表现为不能熔清、冷却后析晶、样品开裂、气泡过多、折射率测试失败或热膨胀系数不满足要求。不同失败类型对应不同边界：不能熔清可能是工艺温度问题，析晶可能是玻璃形成能力不足，开裂可能与热膨胀系数或退火制度有关。把这些失败统一写成“失败”会损失大量材料信息。主动学习实验结果的记录方式见表 \ref{tab:active-learning-result-records}。

\begin{formalTable}[htbp]
\caption{主动学习实验结果的记录方式}
\label{tab:active-learning-result-records}
\footnotesize
\begin{tabularx}{\textwidth}{@{}YYYY@{}}
\toprule
\textbf{实验结果类型} & \textbf{记录内容} & \textbf{对模型的作用} & \textbf{下一轮决策} \\
\midrule
性能达标 & 折射率、透过率、Tg、热膨胀系数 & 提供高价值正样本 & 在附近细化搜索 \\
性能一般但可成玻 & 完整性质数据和配方 & 扩展模型覆盖范围 & 可作为训练样本保留 \\
析晶或开裂 & 失败现象、热处理条件 & 标记不可制备边界 & 减少相邻高风险候选 \\
重复样品 & 批次、设备、测试误差 & 估计噪声水平 & 修正不确定性判断 \\
\bottomrule
\end{tabularx}
\end{formalTable}

AI 可辅助把主动学习输出整理成实验批次说明，但提示词中必须要求它保留利用、探索、边界和重复样品的比例，并提醒不能自动执行模型给出的批次。

\begin{samepage}
\begin{userbox}
我正在做 AR 光波导玻璃配方的主动学习优化。模型已经给出候选配方的预测折射率、不确定性、Tg、热膨胀系数、距离已有样本的远近以及是否满足基本约束。请帮我整理下一轮实验批次选择原则。

要求如下：
\begin{enumerate}
\item 区分利用型、探索型、边界型和重复型样品；
\item 说明每类样品为什么需要进入实验批次；
\item 提醒哪些实验元数据必须记录；
\item 说明失败样品如何反馈到下一轮模型；
\item 不要直接把模型最高分样品全部选为下一轮实验。
\end{enumerate}
\end{userbox}

\begin{aibox}
下一轮实验批次可按四类候选组合：

\begin{enumerate}
\item \textbf{利用型样品：}预测折射率较高、Tg 和热膨胀系数满足约束，用于尽快获得性能改进。
\item \textbf{探索型样品：}预测不确定性较高但物理上仍可行，用于扩展模型对未知组成区域的认识。
\item \textbf{边界型样品：}靠近玻璃形成或工艺边界，用于判断哪些区域容易析晶、开裂或无法熔制。
\item \textbf{重复型样品：}选择已有样品或近似重复样品，用于估计实验噪声、批次差异和测试误差。
\end{enumerate}

每个样品都应记录配方、原料批次、熔制温度、保温时间、冷却制度、样品状态、测试设备和失败原因。失败样品不应删除，而应作为不可制备区域或高风险区域的信息反馈给下一轮模型。
\end{aibox}
\end{samepage}

以下代码示例保留了主动学习最核心的计算关系：预测值体现“利用”，不确定性体现“探索”，Tg 和热膨胀系数等惩罚项体现工艺约束。代码输出的是待人工复核的实验批次，而不是可自动执行的最终方案。


\begin{lstlisting}[language=Python,escapeinside={(*@}{@*)}]
import pandas as pd
cand = pd.DataFrame({
    "id": ["G01", "G02", "G03", "G04", "G05", "G06"],
    "n_pred": [1.86, 1.91, 1.88, 1.83, 1.94, 1.89],
    "sigma": [0.012, 0.010, 0.006, 0.030, 0.018, 0.026],
    "Tg": [565, 548, 570, 535, 505, 552],
    "cte": [8.1, 8.7, 8.4, 8.9, 9.8, 8.0],
    "dist": [0.25, 0.18, 0.35, 0.12, 0.08, 0.10],
    "ok": [True, True, True, True, False, True],
})
target_cte = 8.3
cand["penalty"] = (
    (cand["Tg"] < 520).astype(int) * 0.20
    + (cand["cte"] - target_cte).abs() * 0.03
    + (~cand["ok"]).astype(int) * 10
)
cand["score"] = cand["n_pred"] + 0.4 * cand["sigma"] - cand["penalty"]
feasible = cand[cand["ok"]].copy()
use = feasible.sort_values("n_pred", ascending=False).head(2)
unused = feasible.drop(use.index)
explore = unused.sort_values("sigma", ascending=False).head(2)
unused = unused.drop(explore.index)
boundary = unused.sort_values("dist").head(1)
# (*@\textnormal{保留 G03 作为重复样品，用于评估实验重复性}@*)
repeat = feasible[feasible["id"].eq("G03")].copy()
batch = pd.concat([
    use.assign(kind="use"),
    explore.assign(kind="explore"),
    boundary.assign(kind="boundary"),
    repeat.assign(kind="repeat"),
])
cols = ["id", "kind", "n_pred", "sigma", "Tg", "cte", "score"]
print(batch[cols])
\end{lstlisting}

\manualfigcaption{\textbf{代码示例：}主动学习实验批次评分示例}

\subsubsection*{\textbf{小结}}

本节说明候选筛选应按照由低成本到高成本、由粗判断到精验证的漏斗推进。
规则过滤、代理模型、物理计算和实验验证各有作用，每一级都应保留通过和排除理由。
主动学习把预测值、不确定性和实验反馈结合起来，失败样品和噪声记录同样是有价值的数据。

\subsection*{习题}

\begin{enumerate}[leftmargin=*]
\item 高通量筛选为什么常采用分层漏斗？
\item 比较规则过滤、代理模型、DFT 和实验验证的作用。
\item 主动学习中的“利用”和“探索”分别指什么？
\end{enumerate}

\clearpage
\section{生成设计与候选验证}


\begin{sectionkeypoints}
\item 生成设计改变的是候选产生方式，可在目标条件下提出新的组成、结构或配方，但不能取消材料约束和实验验证。
\item 材料表示决定生成模型能够处理的对象。晶体、玻璃、分子和聚合物需要不同表示方式，不宜简单套用同一种输入格式。
\item 条件生成需要把目标性能转化为可计算条件，例如带隙范围、形成能、元素体系、成本边界或供应链限制。
\item 生成候选必须经过有效性、稳定性、新颖性、可合成性和应用性等多层验证，不应只看生成数量。
\item 生成设计通常需要与筛选、主动学习、DFT 计算和实验验证串联使用，才能形成可信的材料发现流程。
\end{sectionkeypoints}

\subsection{条件生成与验证流程}

候选筛选只能在已经列出的空间中寻找材料。如果候选库没有包含某类新结构或新配方，筛选再快也无法发现它。生成设计尝试学习已有材料的分布，并在目标条件下提出新的组成、结构或配方。它改变的是候选产生方式，而不是取消物理约束和实验验证。


材料表示决定生成模型能够处理什么。分子可使用字符串或分子图，晶体需要同时表示元素种类、周期性晶格和原子坐标，玻璃与非晶材料常以成分和短程有序统计表示，聚合物还要考虑单体序列和拓扑。生成表示若不能表达周期性、价态和局域配位，模型容易产生形式正确但没有材料意义的候选。表中的 VAE\footnote{Variational Autoencoder，变分自编码器。} 指变分自编码器，GAN\footnote{Generative Adversarial Network，生成对抗网络。} 指生成对抗网络，自回归模型\footnote{Autoregressive Model，自回归模型，按序列顺序逐步生成数据。} 按步骤生成序列或结构片段，扩散模型\footnote{Diffusion Model，扩散模型，从噪声出发逐步去噪生成数据。} 从噪声中逐步生成结构，SMILES\footnote{Simplified Molecular Input Line Entry System，简化分子线性输入系统。} 是分子线性表示方法。不同生成模型路线与材料对象的对应关系见表 \ref{tab:generative-model-routes}。


\begin{formalTable}[htbp]
\caption{生成模型路线与材料对象}
\label{tab:generative-model-routes}
\begin{tabularx}{\textwidth}{@{}YYYY@{}}
\toprule
\textbf{模型路线} & \textbf{材料生成中的直觉} & \textbf{典型对象} & \textbf{需要特别注意} \\
\midrule
VAE & 把材料压缩到连续潜空间并在其中优化 & 分子、聚合物、部分晶体表示 & 潜空间是否有物理意义 \\
GAN & 生成器与判别器对抗，逼近真实样本分布 & 显微组织、分子、结构图像 & 训练稳定性与约束不足 \\
自回归模型 & 按序列或步骤逐步生成 & SMILES、合成路线、结构 token & 语法有效性和错误传播 \\
扩散模型 & 从噪声逐步去噪生成结构 & 晶体、分子三维结构 & 周期边界、对称性和坐标处理 \\
混合工作流 & 生成、筛选、DFT、实验闭环结合 & 无机晶体、玻璃、催化剂 & 验证流程比模型名称更重要 \\
\bottomrule
\end{tabularx}
\end{formalTable}

VAE、GAN、自回归模型和扩散模型可采用不同方式学习材料分布。VAE 更适合把分子或配方压缩到连续潜空间中再优化；GAN 常用于生成图像式显微组织或分子结构，但训练稳定性和约束控制需要特别检查；自回归模型适合 SMILES、合成路线或结构 token 这类有顺序的表示；扩散模型适合逐步生成分子或晶体三维结构，但必须处理周期性边界、坐标对称性和原子距离合理性。本章不展开各模型的数学细节，只保留共同逻辑：模型从已有材料中学习可行区域，目标条件用于引导生成方向，化学和结构约束用于排除无效样本，正向模型和高精度计算用于评价生成结果。


\subsubsection*{\textbf{1. 案例示范：发光材料条件生成}}

条件生成是逆向设计的关键。目标条件可以是带隙范围、形成能、体模量、磁性、元素体系或供应链限制。以发光材料为例，目标波长可转化为带隙范围，生成模型提出候选组成和结构，再由第\ref{chap:materials-cognition-prediction}章的带隙模型、DFT 和实验逐级验证。目标带隙只是入口，缺陷、掺杂、热稳定和器件兼容性仍需单独检查。

对于发光材料设计，条件生成可从两个层次理解。第一层是性能条件，例如目标发光波长、带隙范围、热稳定性和缺陷容忍度。第二层是材料边界，例如元素体系、晶体结构类型、掺杂元素、成本和安全要求。如果只给生成模型一个带隙条件，模型可能提出带隙合适但不稳定、含风险元素或难以合成的候选。因此，条件生成需要同时输入目标和边界。

发光材料还具有明显的“性质链条”。目标颜色对应发光波长，发光波长与带隙有关，但材料是否真正发光还取决于缺陷、能级、载流子复合和非辐射损失。换言之，带隙只是候选进入下一步验证的第一道门槛。一个候选材料即使带隙接近目标，如果存在深能级缺陷、热稳定性差或与器件结构不兼容，也不宜直接认为适合应用。

由此可见，生成设计更适合作为“候选扩展器”。它可帮助研究者跳出已有数据库或人工枚举空间，提出一些值得进一步检查的新组成或新结构。但候选是否有意义，要经过规则检查、正向预测、结构弛豫、高精度计算和实验验证。生成数量本身没有价值，真正有价值的是通过层层验证后仍然保留下来的候选。发光材料条件生成中的约束类型见表 \ref{tab:luminescent-material-constraints}。

\begin{formalTable}[htbp]
\caption{发光材料条件生成中的约束类型}
\label{tab:luminescent-material-constraints}
\footnotesize
\begin{tabularx}{\textwidth}{@{}YYYY@{}}
\toprule
\textbf{条件类型} & \textbf{发光材料中的含义} & \textbf{可能数据字段} & \textbf{验证方式} \\
\midrule
目标条件 & 希望获得的发光颜色或带隙范围 & 目标波长、带隙范围 & 光谱测试、带隙预测、DFT \\
结构条件 & 候选结构是否合理 & 空间群、原子坐标、局域配位 & 结构弛豫、距离检查、声子计算 \\
稳定条件 & 候选是否可能存在 & 形成能、凸包能、热稳定性 & 高精度计算、热分析、实验合成 \\
应用条件 & 是否适合器件环境 & 缺陷、热稳定、外延兼容 & 发光测试、器件测试、可靠性测试 \\
\bottomrule
\end{tabularx}
\end{formalTable}

在使用 AI 辅助生成设计时，提示词应强调“提出候选后的验证链条”，而不是让 AI 直接给出材料名称。示例提示词可这样组织。

\begin{samepage}
\begin{userbox}
我想设计一种目标发光波长在蓝光范围附近的无机发光材料。请帮我把这个目标转化为条件生成设计流程。

要求如下：
\begin{enumerate}
\item 说明目标波长如何转化为带隙或能级相关条件；
\item 说明生成模型可提出哪些类型的候选，例如组成、掺杂或晶体结构；
\item 列出生成后必须检查的化学、结构、稳定性和应用约束；
\item 说明为什么不应只根据带隙判断材料是否适合发光器件；
\item 不要直接虚构某个“最佳材料”。
\end{enumerate}
\end{userbox}

\begin{aibox}
可把任务拆成以下流程：

\begin{enumerate}
\item 目标条件设定：将目标发光波长转化为近似带隙范围，同时记录发光效率、热稳定性和器件兼容性等要求。
\item \textbf{候选生成：}生成模型可提出候选组成、掺杂元素、晶体结构或局域配位环境，但这些结果只是候选，不是最终材料。
\item \textbf{有效性检查：}检查化学式、价态、原子距离、晶格体积和结构格式是否合理，排除明显无意义候选。
\item 稳定性验证：用形成能、凸包能、结构弛豫和必要的声子或热稳定性分析判断候选是否可能存在。
\item \textbf{应用验证：}进一步检查缺陷、非辐射复合、热稳定、外延匹配和器件环境适配性。
\end{enumerate}

带隙合适只能说明候选可能满足发光颜色要求，不能说明它一定具有高发光效率或可制备性。生成模型的作用是扩展候选空间，最终判断仍然需要计算和实验验证。
\end{aibox}
\end{samepage}

条件生成还需要保留“负面条件”。很多时候，研究者不仅知道想要什么，也知道不想要什么。例如，不希望候选含有某些高风险元素，不希望晶体结构出现过短原子距离，不希望配方超出已知可制备范围，不希望候选与训练集中已有材料完全重复。这些负面条件如果不写入生成和筛选流程，模型可能不断提出形式上满足目标性能、但材料上明显不可接受的候选。

具体而言，条件生成的输入可分为三类：目标条件、边界条件和排除条件。目标条件推动模型朝目标性能靠近；边界条件保证候选处于可研究范围；排除条件提前去掉明显不可接受的结果。对于材料科学任务，三类条件同样重要。只写目标条件，容易得到“性能诱人但不可用”的候选；只写排除条件，又可能使搜索空间过窄，失去发现新材料的机会。


\subsection{候选验证与路线组合}

\subsubsection*{\textbf{1. 验证流程：从有效性到应用性}}

生成候选首先要检查有效性：化学式、价态、原子距离、晶格体积和周期性是否合理；其次检查稳定性：结构弛豫后是否保持、形成能和声子是否可接受；再次检查新颖性：候选是否只是训练集或数据库中的重复结构；最后检查可合成性和应用条件。任何一级不通过，都不应把候选直接称为材料发现。生成设计的主要验证维度见表 \ref{tab:generative-design-validation}。


\begin{formalTable}[htbp]
\caption{生成设计的验证维度}
\label{tab:generative-design-validation}
\begin{tabularx}{\textwidth}{@{}YYYY@{}}
\toprule
\textbf{验证指标} & \textbf{问题} & \textbf{常用判断} & \textbf{为什么重要} \\
\midrule
有效性 & 结构是否符合基本规则？ & 价态、距离、格式、晶格范围 & 避免无意义候选 \\
稳定性 & 候选是否可能存在？ & 形成能、凸包能、声子、动力学 & 减少纸上材料 \\
新颖性 & 是否只是训练集复述？ & 结构匹配、组成相似度、数据库检索 & 证明发现价值 \\
可合成性 & 实验能否做出来？ & 合成路线、温度窗口、前驱体、杂相 & 连接真实材料 \\
应用性 & 能否满足场景要求？ & 服役环境、成本、可靠性、规模化 & 决定工程价值 \\
\bottomrule
\end{tabularx}
\end{formalTable}

可靠的验证链条通常分为四层。第一层是格式、几何和重复性检查；第二层是快速正向模型打分；第三层是 DFT、分子动力学或热力学精算；第四层是实验合成、结构表征和性能测试。研究报告还应说明每一级候选的通过率，区分生成数量、有效数量、稳定数量和实验验证数量。


筛选、主动学习和生成设计可串联使用：生成模型扩展候选空间，规则和代理模型快速排除明显不合理候选，主动学习决定哪些候选最值得进入下一轮高精度计算或实验。最终判断仍然来自材料约束和验证证据，而不是模型名称。

可合成性是生成设计最容易被低估的环节。一个结构在计算上稳定，不代表实验上容易合成。实验可达性受到前驱体、温度、压力、气氛、反应路径、动力学障碍和杂相竞争影响。对于玻璃材料，候选配方还必须能熔制、澄清、冷却成玻璃并避免严重析晶；对于晶体材料，候选结构还要考虑相竞争和合成窗口。可合成性判断不足，是“纸上材料”大量出现的重要原因。生成候选在不同证据阶段的表述边界见表 \ref{tab:generated-candidate-evidence}。

\begin{formalTable}[htbp]
\caption{生成候选在不同证据阶段的表述边界}
\label{tab:generated-candidate-evidence}
\footnotesize
\begin{tabularx}{\textwidth}{@{}YYYY@{}}
\toprule
\textbf{候选阶段} & \textbf{可采用的表述} & \textbf{不能怎样表述} & \textbf{需要补充证据} \\
\midrule
模型生成 & 生成了一批候选结构或配方 & 发现了一批新材料 & 格式、价态和几何检查 \\
模型打分 & 获得高优先级候选 & 证明候选性能优异 & 正向模型误差和适用范围 \\
高精度计算 & 候选在计算上接近稳定 & 候选一定可合成 & 结构弛豫、凸包能、相竞争 \\
实验合成 & 候选已被制备并测试 & 可直接工程应用 & 服役、可靠性和规模化验证 \\
\bottomrule
\end{tabularx}
\end{formalTable}

这一部分也适合用 AI 辅助整理验证清单。关键是要求 AI 区分证据层级，避免把模型输出写成最终结论。

\begin{samepage}
\begin{userbox}
我用生成模型得到了一批候选无机晶体结构。请帮我设计一个候选验证清单，判断这些候选能否进入后续计算或实验。

要求如下：
\begin{enumerate}
\item 按有效性、稳定性、新颖性、可合成性和应用性分层；
\item 每一层说明需要检查哪些信息；
\item 区分“模型推荐候选”和“已验证材料”的表述；
\item 提醒如何记录每一级候选数量和通过率；
\item 不要把生成模型输出直接称为材料发现。
\end{enumerate}
\end{userbox}

\begin{aibox}
可建立五层验证清单：

\begin{enumerate}
\item \textbf{有效性：}检查化学式、价态、原子距离、晶格体积、周期性和文件格式是否合理。
\item 稳定性：进行结构弛豫，检查形成能、凸包能、声子稳定性或分子动力学稳定性。
\item \textbf{新颖性：}与训练集、材料数据库和文献记录比对，判断是否为已有材料或高度相似结构。
\item \textbf{可合成性：}分析前驱体、温度窗口、相竞争、动力学路径和实验设备条件。
\item \textbf{应用性：}结合目标场景检查性能、成本、安全、服役环境和长期可靠性。
\end{enumerate}

报告中应分别写出生成数量、有效数量、稳定数量、新颖数量、进入实验数量和实验验证数量。模型生成和模型打分只能说明候选值得进一步检查，不宜直接等同于材料发现。
\end{aibox}
\end{samepage}

在实际写作中，还应避免把验证流程写成一条单向直线。候选在任何一级验证失败，都可能需要返回前面的步骤。例如，结构弛豫后候选发生明显畸变，说明生成结构可能不合理；DFT 显示候选远离稳定区域，说明筛选规则需要加强稳定性约束；实验中反复出现杂相，说明可合成性判断不足。这些失败不是无效信息，而是帮助修正候选生成和筛选规则的重要反馈。

因此，候选验证更像是一个带反馈的流程。生成模型提出候选，验证流程排除或确认候选，失败原因再返回去修改约束、训练数据或生成条件。只有把反馈写清楚，生成设计才不是“模型输出一批材料”，而是“模型、计算和实验共同收敛到更可靠候选”的过程。

\subsubsection*{\textbf{小结}}

本节说明生成设计的作用是扩展候选空间，而不是直接给出最终材料。
条件生成需要同时设置目标条件、边界条件和排除条件，并选择适合材料对象的表示方式。
生成候选必须经过有效性、稳定性、新颖性、可合成性和应用性验证，失败结果也应回流流程。

\subsection*{习题}

\begin{enumerate}[leftmargin=*]
\item 生成候选为什么不能直接称为材料发现？
\item 为发光材料生成任务设计三类验证指标。
\end{enumerate}

\clearpage
\section{可信发现与路线选择}


\begin{sectionkeypoints}
\item 可信材料发现需要区分数据、模型、材料、验证和应用等不同证据层级，不应把模型推荐直接等同于真实材料发现。
\item 方法选择应由材料问题决定。候选库明确时优先筛选，实验昂贵时优先主动学习，候选库不足时再考虑生成模型。
\item 安全、环境、供应链、成本、可加工性和长期可靠性应作为并列约束保留，不能被压缩成无法解释的单一分数。
\item 第\ref{chap:materials-cognition-prediction}章的材料表示、性质预测、不确定性、机器学习势和可解释性，是第\ref{chap:materials-design-discovery}章候选筛选、生成验证和闭环优化的基础。
\item 全模块的完整逻辑是：定义目标、限定边界、提出候选、逐层验证、实验反馈，并在新数据基础上继续更新模型。
\end{sectionkeypoints}

\subsection{可靠评价与方法选择}

人工智能能够扩大材料搜索范围，也会放大数据偏差和模型外推风险。材料数据库中的计算参数、实验条件和样品质量并不完全一致，成功结果通常多于失败结果，热门体系的数据也远多于冷门体系。模型会继承这些差异，因此候选排名并不是客观真理，而是特定数据和模型边界下的判断。


材料性能还高度依赖结构、工艺和服役环境。同一化学式在不同缺陷、微结构和热处理条件下可能表现迥异，如果数据只记录成分和最终性能，模型就会把未记录的工艺差异当成噪声。逆向设计因此必须明确材料对象和数据适用范围，不应把材料简化为一个化学式或单一分数。


“纸上材料”指计算或模型预测较好，但缺少可合成性、规模化或应用验证的候选。真正的材料发现需要区分不同证据层级：模型预测提供线索，高精度计算提高物理可信度，实验合成证明材料可达，器件测试和长期稳定性才说明应用潜力。材料发现结果的分层评价见表 \ref{tab:materials-discovery-evaluation}。


\begin{formalTable}[htbp]
\caption{材料发现结果的分层评价}
\label{tab:materials-discovery-evaluation}
\begin{tabularx}{\textwidth}{@{}YYY@{}}
\toprule
\textbf{评价层面} & \textbf{核心判断} & \textbf{可靠表述示例} \\
\midrule
数据 & 来源、去重、划分是否清楚 & 在某数据库和划分方式下模型表现良好 \\
模型 & 是否有基线、不确定性和边界说明 & 模型推荐了若干高优先级候选 \\
材料 & 是否满足化学、结构和稳定性约束 & 候选在计算上接近稳定且结构合理 \\
验证 & 是否经过 DFT、实验或器件测试 & 某候选已合成并测得目标性质 \\
应用 & 是否考虑成本、安全和可靠性 & 候选具备进一步器件验证价值 \\
\bottomrule
\end{tabularx}
\end{formalTable}

可信评价的关键，是把“模型表现好”和“材料真的好”分开。模型表现好，通常指在某个数据集、某种划分方式和某个评价指标下，预测误差较小或排序效果较好。材料真的好，则意味着候选在真实制备、结构稳定、性能测试和服役环境中都能经受验证。二者之间存在距离。材料人工智能研究中常见的问题，是把模型层面的高分候选直接写成材料发现结果，忽略了验证证据不足。

评价结果时，还要说明模型适用范围。例如，模型是在无机晶体数据库上训练的，还是在某类玻璃配方上训练的；数据是计算数据、实验数据，还是二者混合；测试集划分是随机划分，还是按材料体系、时间或组成空间划分。如果训练集和测试集高度相似，模型指标可能看起来很好，但对真正未知材料的预测能力并不一定强。对于逆向设计来说，外推风险比普通插值预测更突出，因为候选往往位于已有数据边界附近。

不确定性也是可信评价的重要组成部分。一个候选预测性能高，但模型不确定性也高，说明它可能值得探索，但不宜直接排在确定性较高的候选之前。相反，一个候选预测性能略低但不确定性小、约束满足、可合成性强，也可能更适合作为近期实验对象。候选报告中应同时给出预测值、不确定性、约束是否满足和推荐理由。材料发现结果的可靠表述方式见表 \ref{tab:reliable-discovery-expression}。

\begin{formalTable}[htbp]
\caption{材料发现结果的可靠表述方式}
\label{tab:reliable-discovery-expression}
\footnotesize
\begin{tabularx}{\textwidth}{@{}YYYY@{}}
\toprule
\textbf{常见表述} & \textbf{风险} & \textbf{更稳妥的写法} & \textbf{还需补充} \\
\midrule
模型发现了某材料 & 夸大模型作用 & 模型推荐该候选进入验证 & 计算或实验验证 \\
该候选性能优异 & 忽略不确定性和边界 & 在当前模型和数据范围内预测值较高 & 不确定性、误差范围 \\
生成了大量新材料 & 混淆生成与发现 & 生成了若干待验证候选 & 去重、新颖性和稳定性检查 \\
候选可用于器件 & 跳过应用验证 & 候选具备进一步器件验证价值 & 器件测试和长期可靠性 \\
\bottomrule
\end{tabularx}
\end{formalTable}

候选报告还需要特别注意“数量表述”。生成模型可能一次输出几千个候选，但这并不意味着发现了几千种新材料。更准确的写法应区分：生成候选数量、格式有效数量、结构合理数量、计算稳定数量、去重后新颖数量、进入实验数量和实验验证成功数量。只有经过实验合成和性能测试的候选，才能接近“发现”的表述；只经过模型打分的候选，只能称为“推荐候选”或“高优先级候选”。

新颖性检查也属于可信评价的一部分。生成模型可能复述训练集中已经存在的材料，或者只对已有结构做很小改动。如果不做数据库检索和结构匹配，就可能把已有材料误写成新发现。对于晶体材料，可比较化学组成、空间群、晶格参数和结构相似度；对于玻璃或配方材料，可比较组分比例和已有专利、文献或实验记录。新颖性不是单靠模型输出判断的，而需要与数据库和文献记录对照。

\subsubsection*{\textbf{1. 路线选择：问题决定方法}}

方法选择应由问题决定，而不是由模型新旧决定。候选空间已经明确时，高通量筛选通常最直接；实验成本高且每轮样品有限时，主动学习更适合；需要突破已有候选库时，生成模型才体现优势。玻璃形成边界、相稳定、毒性和设备温度等硬约束，应尽可能在搜索前加入，而不是在模型给出结果后再补救。不同问题情形下的方法选择见表 \ref{tab:inverse-design-method-selection}。


\begin{formalTable}[htbp]
\caption{材料逆向设计的方法选择}
\label{tab:inverse-design-method-selection}
\begin{tabularx}{\textwidth}{@{}YYYY@{}}
\toprule
\textbf{问题情形} & \textbf{优先方法} & \textbf{理由} & \textbf{典型材料场景} \\
\midrule
候选库明确 & 高通量筛选 & 流程清楚、可解释、验证成本低 & 从数据库筛选稳定光电材料 \\
实验昂贵且小样本 & 贝叶斯优化/主动学习 & 每一轮实验都要最大化信息量 & 玻璃配方、合金热处理、催化条件 \\
候选库不足 & 生成模型 + 后筛选 & 需要提出库外候选 & 新晶体结构、功能分子、聚合物序列 \\
安全或法规约束强 & 规则约束 + 人工审核 & 先排除不可接受候选 & 含毒元素、稀缺元素或高风险反应 \\
\bottomrule
\end{tabularx}
\end{formalTable}

可靠评价还应考虑安全、环境和供应链。候选若含有有毒、稀缺或高风险元素，即使单项性能突出，也可能不适合真实应用。相应地，性能、成本、环境影响、可加工性和生命周期不应被压缩成一个无法解释的总分，而应作为并列约束保留在候选报告中。

方法选择还应考虑课程或项目阶段。对于入门学习，最合适的方法往往不是最新的生成模型，而是可解释、可复核的筛选流程。因为研究者需要先理解数据字段、模型输入、约束条件和验证证据之间的关系。对于研究生或真实科研项目，如果已有稳定的数据积累和实验反馈机制，再引入主动学习或生成模型才更合理。否则，复杂模型可能只是增加技术难度，并不能提高材料发现质量。

在写作时，可把方法选择写成一组判断问题：候选空间是否已经明确？每次实验成本是否很高？是否需要提出已有数据库之外的新候选？是否存在严格安全或工艺边界？是否有能力进行后续高精度计算和实验验证？这些问题比“选择哪个先进算法”更重要。

这一节中的方法选择和可信评价也可借助 AI 辅助整理，但提示词应要求它给出理由和证据边界，不应只推荐一个模型名称。

\begin{samepage}
\begin{userbox}
我正在写材料逆向设计方案，候选空间包括已有数据库中的无机晶体和少量可枚举组成。实验成本较高，每轮只能验证 5～10 个样品。请帮我判断优先使用高通量筛选、主动学习还是生成模型。

要求如下：
\begin{enumerate}
\item 根据候选空间、实验成本、数据量和验证能力进行判断；
\item 说明每种方法适合和不适合的情况；
\item 给出推荐路线，但必须说明理由；
\item 提醒结果报告中应如何避免夸大模型作用；
\item 不要只因为生成模型较新就优先推荐它。
\end{enumerate}
\end{userbox}

\begin{aibox}
在这个任务中，候选空间已有数据库和可枚举组成，说明不一定需要首先使用生成模型。更合理的路线可以是：

\begin{enumerate}
\item 第一阶段使用高通量筛选：先用规则、正向模型和稳定性判断从已有候选库中排除明显不合理候选。
\item \textbf{第二阶段使用主动学习：}由于每轮实验只能验证 5～10 个样品，应让模型同时考虑预测性能和不确定性，选择高价值实验批次。
\item 生成模型暂时作为补充：只有当已有候选库不能覆盖目标性能区域，或者需要提出库外结构时，再考虑生成模型。
\end{enumerate}

报告中应写成“模型推荐若干高优先级候选，并通过后续计算或实验验证”，而不是“模型发现了最终材料”。如果没有实验或高精度计算支持，应明确说明候选仍处于待验证阶段。
\end{aibox}
\end{samepage}

\subsection{能力复用与闭环迭代}

\subsubsection*{\textbf{1. 两章闭环：正向模型服务逆向设计}}

两章之间形成正反闭环。第\ref{chap:materials-cognition-prediction}章把材料数据转化为描述符或结构图，训练正向模型并解释模型边界；第\ref{chap:materials-design-discovery}章把这些模型作为候选筛选和生成验证的评价器，再用新计算和实验结果更新数据。没有正向预测，逆向搜索缺少可靠评价；没有逆向设计，正向模型容易停留在被动打分。第\ref{chap:materials-cognition-prediction}章能力在第\ref{chap:materials-design-discovery}章中的复用关系见表 \ref{tab:chapter-one-capability-reuse}。


\begin{formalTable}[htbp]
\caption{第\ref{chap:materials-cognition-prediction}章能力在第\ref{chap:materials-design-discovery}章中的复用}
\label{tab:chapter-one-capability-reuse}
\begin{tabularx}{\textwidth}{@{}YYY@{}}
\toprule
\textbf{第\ref{chap:materials-cognition-prediction}章能力} & \textbf{第\ref{chap:materials-design-discovery}章中的用途} & \textbf{材料例子} \\
\midrule
性质预测 & 作为筛选和生成后的快速打分器 & 带隙筛选 microLED 候选 \\
材料表示 & 决定候选能否被模型有效评价 & 晶体图、玻璃成分、分子图 \\
不确定性量化 & 指导主动学习下一步实验 & 高折射率玻璃小样本优化 \\
机器学习势 & 验证候选结构和高温行为 & 超高温陶瓷和高熵合金 \\
可解释性 & 把模型推荐转化为材料规律 & 形成能力、熔点或缺陷趋势分析 \\
\bottomrule
\end{tabularx}
\end{formalTable}

完整逻辑可概括为：目标性质定义问题，材料知识限定边界，筛选、主动学习或生成模型提出候选，正向模型和物理计算逐层评价，实验结果决定最终结论并反馈到数据集。人工智能的价值在于提高搜索与决策效率，材料科学的责任则是保证问题、约束和证据真实可靠。

闭环迭代要落实到可操作的数据管理上。每一轮推荐都应记录模型版本、训练数据版本、候选来源、预测值、不确定性、约束检查结果和推荐理由；每一轮验证后，还要把成功样品、失败样品、测试条件和失败原因按同一格式回写数据表。这样，下一轮模型更新才知道哪些变化来自材料本身，哪些变化来自数据版本、实验批次或测试条件。

例如，在高折射率玻璃任务中，第\ref{chap:materials-cognition-prediction}章中讨论的数据整理和描述符构建可帮助建立折射率、Tg 和热膨胀系数预测模型；第\ref{chap:materials-design-discovery}章则利用这些模型筛选候选配方，并通过主动学习选择下一轮熔制实验。在高熵合金涂层任务中，第\ref{chap:materials-cognition-prediction}章的硬度预测、图神经网络和机器学习势可提供性能估计与结构理解；第\ref{chap:materials-design-discovery}章则进一步用这些模型决定哪些成分或硬质相值得验证。这样的闭环，才是人工智能真正进入材料发现流程的方式。

闭环迭代还要求新数据及时回流。实验失败、性能一般、模型预测错误的样品，都应被记录下来。只把成功样品加入数据库，会使模型越来越偏向成功区域，无法学习失败边界；只记录最终性能而不记录工艺条件，也会使模型无法区分成分效应和工艺效应。闭环发现不仅是“模型推荐--实验验证”，还包括“实验记录--数据更新--模型修正”。材料智能发现闭环中的数据回流内容见表 \ref{tab:closed-loop-data-feedback}。

\begin{formalTable}[htbp]
\caption{材料智能发现闭环中的数据回流}
\label{tab:closed-loop-data-feedback}
\footnotesize
\begin{tabularx}{\textwidth}{@{}YYYY@{}}
\toprule
\textbf{闭环环节} & \textbf{主要任务} & \textbf{需要记录的信息} & \textbf{常见问题} \\
\midrule
模型推荐 & 根据目标和约束提出候选 & 预测值、不确定性、推荐理由 & 只给排名，不说明边界 \\
计算验证 & 用更高精度方法复核候选 & 结构、能量、稳定性、失败原因 & 只保留通过候选 \\
实验验证 & 合成、表征和性能测试 & 工艺、样品状态、测试条件 & 缺少失败样品记录 \\
数据回流 & 更新数据集并重训模型 & 新数据、元数据、版本信息 & 数据来源不可追溯 \\
\bottomrule
\end{tabularx}
\end{formalTable}

AI 可辅助整理闭环报告结构，尤其是检查是否遗漏数据来源、验证证据和失败记录。但同样需要明确：AI 只能帮助组织报告，不宜替代实验和模型审查。

\begin{samepage}
\begin{userbox}
我想把一个材料智能发现项目写成闭环流程报告。项目包括正向模型预测、候选筛选、实验验证和数据回流。请帮我设计报告结构。

要求如下：
\begin{enumerate}
\item 说明第\ref{chap:materials-cognition-prediction}章中的材料表示、性能预测和模型解释如何服务第\ref{chap:materials-design-discovery}章的逆向设计；
\item 按模型推荐、计算验证、实验验证和数据回流四个环节组织；
\item 提醒每个环节需要记录哪些信息；
\item 特别说明失败样品和预测错误样品为什么也要记录；
\item 不要把闭环写成模型自动发现材料。
\end{enumerate}
\end{userbox}

\begin{aibox}
闭环流程报告可包括四部分：

\begin{enumerate}
\item \textbf{模型推荐：}说明使用了哪些材料表示和正向预测模型，给出候选的预测值、不确定性和推荐理由。
\item \textbf{计算验证：}记录结构弛豫、稳定性判断、形成能或其他高精度计算结果，并保留未通过候选的原因。
\item \textbf{实验验证：}记录制备工艺、样品状态、表征结果、性能测试条件和失败现象，不应只记录成功样品。
\item \textbf{数据回流：}说明新数据如何加入数据库，如何更新训练集，模型是否重新训练，以及新模型与旧模型相比是否改进。
\end{enumerate}

失败样品和预测错误样品同样重要，因为它们可帮助模型学习不可制备区域、实验噪声和模型外推边界。材料智能发现不是模型自动给出答案，而是模型、计算、实验和数据更新共同形成的迭代过程。
\end{aibox}
\end{samepage}

如果把两章放在一个完整模块中理解，第\ref{chap:materials-cognition-prediction}章更像是“建立评价能力”，第\ref{chap:materials-design-discovery}章更像是“利用评价能力做决策”。没有第\ref{chap:materials-cognition-prediction}章的数据整理、表示学习和模型评估，第\ref{chap:materials-design-discovery}章的筛选和生成就缺少可靠判断依据；没有第\ref{chap:materials-design-discovery}章的候选选择和实验反馈，第\ref{chap:materials-cognition-prediction}章的模型又容易停留在已有数据上的被动预测。二者结合起来，才能形成从材料认知到材料发现的完整教学闭环。

最后一节不只是总结第\ref{chap:materials-design-discovery}章，也是在说明整个模块的逻辑边界：人工智能可帮助研究者更快地整理材料数据、构建候选、预测性能、安排验证和更新模型，但不应跳过材料科学中的基本问题。材料是否真实存在，能否制备出来，是否满足应用场景，仍然要靠实验、计算和工程约束共同判断。

\subsubsection*{\textbf{小结}}

本节强调可信材料发现需要区分模型推荐、计算证据、实验验证和应用潜力。
方法选择应由问题决定，候选库明确时优先筛选，实验昂贵时优先主动学习，候选不足时再考虑生成。
第\ref{chap:materials-cognition-prediction}章建立评价能力，第\ref{chap:materials-design-discovery}章利用评价能力做决策，数据回流使整个模块形成闭环。

\subsection*{习题}

\begin{enumerate}[leftmargin=*]
\item 数据回流为什么是材料智能发现闭环的关键？
\item 如何根据候选库、实验成本和数据量选择材料发现路线？
\end{enumerate}

\section*{本章小结}

本章讨论了材料设计与智能发现的基本思路。逆向设计首先要求将应用需求转化为可计算目标、变量和约束；候选筛选通过规则、代理模型、高精度计算和实验验证逐层缩小范围；主动学习利用预测值和不确定性安排下一批实验或计算任务；生成设计能够提出新的结构或配方，但必须接受有效性、稳定性、新颖性、可合成性和应用条件等多层验证。全章强调，智能发现不是单次模型输出，而是目标定义、候选构建、分层验证和数据回流共同组成的闭环过程。
